\documentclass[10pt,a4paper]{article}
\usepackage[utf8]{inputenc}
\usepackage[T1]{fontenc}
\usepackage{lmodern}
\usepackage[a4paper, margin=1.8cm, top=2.2cm, bottom=2.2cm, headheight=16pt, footskip=14pt]{geometry}
\usepackage{xcolor}
\usepackage{graphicx}
\usepackage{tikz}
\usepackage{amsmath,amssymb}
\usepackage{pifont}
\usepackage{fancyhdr}
\usepackage{titlesec}
\usepackage{microtype}
\usepackage{longtable}
\usepackage{booktabs}
\usepackage{array}
\usepackage{tcolorbox}
\usepackage{enumitem}
\usepackage{hyperref}

% Tight options list for questions
\newlist{qoptions}{enumerate}{1}
\setlist[qoptions]{label=\textbf{\Alph*.}, leftmargin=2em, itemsep=1.5pt, topsep=2.5pt, parsep=0pt}

% Color Definitions
\definecolor{hblue}{HTML}{1E3A8A}      % Navy Blue
\definecolor{hred}{HTML}{DC2626}       % Huawei Red accent
\definecolor{singlecol}{HTML}{0D9488}   % Teal for Single Choice
\definecolor{multicol}{HTML}{6D28D9}    % Purple for Multi Choice
\definecolor{tfcol}{HTML}{D97706}       % Amber for True/False
\definecolor{ansgreen}{HTML}{047857}    % Forest Green
\definecolor{ansbg}{HTML}{F0FDF4}       % Light Green Background
\definecolor{qbg}{HTML}{F8FAFC}         % Light Slate Background
\definecolor{qborder}{HTML}{CBD5E1}     % Border Slate
\definecolor{darkslate}{HTML}{1E293B}   % Slate 800

% Hyperref Setup
\hypersetup{
    colorlinks=true,
    linkcolor=hblue,
    citecolor=hblue,
    urlcolor=hblue,
    pdfauthor={Huawei ICT Academy / HCIA-Security Preparation},
    pdftitle={HCIA-Security Comprehensive Exam Question Bank},
    pdfsubject={HCIA-Security V3.0 Certification Exam Preparation},
    pdfkeywords={HCIA, Security, Huawei, Firewall, NAT, VRRP, IPS, PKI, IPsec, VPN}
}

% Header & Footer Configuration
\pagestyle{fancy}
\fancyhf{}
\fancyhead[L]{\small\color{darkslate!70}\textbf{HCIA-Security Exam Bank}}
\fancyhead[R]{\small\color{darkslate!70}\nouppercase{\leftmark}}
\fancyfoot[L]{\small\color{darkslate!50}Huawei Certified ICT Associate - Security}
\fancyfoot[R]{\small\color{darkslate!80}\textbf{Page \thepage}}
\renewcommand{\headrulewidth}{0.6pt}
\renewcommand{\footrulewidth}{0.4pt}

% Section Titles Styling
\titleformat{\part}[display]
  {\normalfont\huge\bfseries\centering\color{hblue}}
  {\Large\MakeUppercase{\partname} \thepart}
  {1ex}
  {\titlerule\vspace{1ex}\Huge}
  [\vspace{1ex}\titlerule]

\titleformat{\section}
  {\normalfont\Large\bfseries\color{hblue}}
  {\thesection}{1em}{}
  [\vspace{0.2ex}{\color{hblue}\titlerule[1.2pt]}]

\titleformat{\subsection}
  {\normalfont\large\bfseries\color{darkslate}}
  {\thesubsection}{0.8em}{}

% Custom Question Box
\newtcolorbox{qcard}[2]{
    colback=qbg,
    colframe=#2,
    arc=1.8mm,
    boxrule=0.7pt,
    leftrule=3.8pt,
    left=7pt,
    right=7pt,
    top=5pt,
    bottom=5pt,
    title={\textbf{Question #1}},
    coltitle=white,
    colbacktitle=#2,
    before=\par\vspace{0.6em}\noindent,
    after=\par\vspace{0.1em}
}

% Custom Explanation Box
\newtcolorbox{explainbox}[1]{
    colback=ansbg,
    colframe=ansgreen,
    arc=1.5mm,
    boxrule=0.5pt,
    leftrule=3.2pt,
    left=7pt,
    right=7pt,
    top=4pt,
    bottom=4pt,
    title={\small\textbf{Correct Answer: #1}},
    coltitle=white,
    colbacktitle=ansgreen,
    before=\par\vspace{0.1em}\noindent,
    after=\par\vspace{0.6em}
}

\begin{document}

\begin{titlepage}
\centering
\vspace*{2cm}
{\Huge\bfseries\color{hblue} HCIA-Security Comprehensive Exam Question Bank\\[0.8cm]}
{\Large\color{darkslate} V3.0 Certification Preparation\\[0.4cm]}
\vspace{1.5cm}
\begin{tcolorbox}[colback=qbg,colframe=hblue,arc=3mm,boxrule=1pt,width=0.85\textwidth]
\centering
\textbf{Total Questions:} 1070\\[0.3cm]
\textbf{Question Types:} Single Choice, Multiple Choice, True/False\\[0.3cm]
\textbf{Coverage:} All HCIA-Security V3.0 Exam Topics
\end{tcolorbox}
\vfill
{\small Generated for HCIA-Security exam preparation}
\end{titlepage}

\tableofcontents
\newpage

\part{Single-Answer Questions}
\section{Network Security Concepts}
\begin{questioncard}{1}{singlecol}
\qtypebadge{singlecol}{Single Choice} \hfill \textcolor{darkslate!70}{\small Topic: Network Security Concepts}

In a broad sense, network security refers to:

\begin{optlist}
\item Physical security of devices only
\item Cybersecurity / cyberspace security at a national level
\item Only firewall deployment
\item Only antivirus protection
\end{optlist}
\begin{explainbox}{}
\textbf{Correct Answer: B}\\[0.3em]
\textbf{Concept:} Broadly, network security is synonymous with cybersecurity, covering national laws, regulations, and processes to protect cyberspace.
\end{explainbox}
\end{questioncard}

\begin{questioncard}{2}{singlecol}
\qtypebadge{singlecol}{Single Choice} \hfill \textcolor{darkslate!70}{\small Topic: Network Security Concepts}

In a narrow sense, network security focuses on:

\begin{optlist}
\item National cyberspace policies only
\item Vendor devices and solutions that ensure secure running of enterprise networks
\item Only social engineering defense
\item Only data-center cooling
\end{optlist}
\begin{explainbox}{}
\textbf{Correct Answer: B}\\[0.3em]
\textbf{Concept:} In a narrow sense, network security refers to measures such as firewalls, IPS, VPNs, and policies that protect enterprise networks.
\end{explainbox}
\end{questioncard}

\begin{questioncard}{3}{singlecol}
\qtypebadge{singlecol}{Single Choice} \hfill \textcolor{darkslate!70}{\small Topic: Network Security Concepts}

Which period of security history mainly used stream ciphers and physical protection?

\begin{optlist}
\item Information security period
\item Communication security period
\item Cyberspace security period
\item Information assurance period
\end{optlist}
\begin{explainbox}{}
\textbf{Correct Answer: B}\\[0.3em]
\textbf{Concept:} The communication security period (1940s) relied on physical security and cryptographic communication, mainly stream ciphers.
\end{explainbox}
\end{questioncard}

\begin{questioncard}{4}{singlecol}
\qtypebadge{singlecol}{Single Choice} \hfill \textcolor{darkslate!70}{\small Topic: Network Security Concepts}

During the information security period, which three goals were the main focus?

\begin{optlist}
\item Confidentiality, integrity, availability
\item Speed, bandwidth, latency
\item Power, cooling, cabling
\item Routing, switching, NAT
\end{optlist}
\begin{explainbox}{}
\textbf{Correct Answer: A}\\[0.3em]
\textbf{Concept:} The classic CIA triad — confidentiality, integrity, and availability — became the core objectives of information security.
\end{explainbox}
\end{questioncard}

\begin{questioncard}{5}{singlecol}
\qtypebadge{singlecol}{Single Choice} \hfill \textcolor{darkslate!70}{\small Topic: Network Security Concepts}

Which two properties were added to the CIA triad during the information assurance period?

\begin{optlist}
\item Controllability and non-repudiation
\item Speed and redundancy
\item Encryption and hashing
\item NAT and VLAN
\end{optlist}
\begin{explainbox}{}
\textbf{Correct Answer: A}\\[0.3em]
\textbf{Concept:} In the information assurance period, controllability (monitoring and control) and non-repudiation (proof of actions) joined the CIA triad.
\end{explainbox}
\end{questioncard}

\begin{questioncard}{6}{singlecol}
\qtypebadge{singlecol}{Single Choice} \hfill \textcolor{darkslate!70}{\small Topic: Network Security Concepts}

The cyberspace security period aims to protect:

\begin{optlist}
\item Only web servers
\item Only user passwords
\item Facilities, data, users, and operations across cyberspace
\item Only email gateways
\end{optlist}
\begin{explainbox}{}
\textbf{Correct Answer: C}\\[0.3em]
\textbf{Concept:} Cyberspace security expands protection to the entire digital ecosystem: facilities, data, users, and operations.
\end{explainbox}
\end{questioncard}

\begin{questioncard}{7}{singlecol}
\qtypebadge{singlecol}{Single Choice} \hfill \textcolor{darkslate!70}{\small Topic: Network Security Concepts}

Which term describes the practice of protecting networks and data from unauthorized access?

\begin{optlist}
\item Cybersecurity
\item Marketing
\item Accounting
\item Logistics
\end{optlist}
\begin{explainbox}{}
\textbf{Correct Answer: A}\\[0.3em]
\textbf{Concept:} Cybersecurity protects networks, devices, and data from unauthorized access or attack.
\end{explainbox}
\end{questioncard}

\begin{questioncard}{8}{singlecol}
\qtypebadge{singlecol}{Single Choice} \hfill \textcolor{darkslate!70}{\small Topic: Network Security Concepts}

Which of the following best defines a vulnerability?

\begin{optlist}
\item A weakness that can be exploited
\item A security patch
\item A firewall rule
\item A network diagram
\end{optlist}
\begin{explainbox}{}
\textbf{Correct Answer: A}\\[0.3em]
\textbf{Concept:} A vulnerability is a weakness in a system that a threat actor can exploit.
\end{explainbox}
\end{questioncard}

\begin{questioncard}{9}{singlecol}
\qtypebadge{singlecol}{Single Choice} \hfill \textcolor{darkslate!70}{\small Topic: Network Security Concepts}

What is an exploit?

\begin{optlist}
\item A software patch
\item A method to take advantage of a vulnerability
\item A type of firewall
\item A routing protocol
\end{optlist}
\begin{explainbox}{}
\textbf{Correct Answer: B}\\[0.3em]
\textbf{Concept:} An exploit is a technique or code that leverages a vulnerability.
\end{explainbox}
\end{questioncard}

\begin{questioncard}{10}{singlecol}
\qtypebadge{singlecol}{Single Choice} \hfill \textcolor{darkslate!70}{\small Topic: Network Security Concepts}

Which term describes a potential danger to an asset?

\begin{optlist}
\item Risk
\item Threat
\item Vulnerability
\item Control
\end{optlist}
\begin{explainbox}{}
\textbf{Correct Answer: B}\\[0.3em]
\textbf{Concept:} A threat is any potential danger that could exploit a vulnerability.
\end{explainbox}
\end{questioncard}

\begin{questioncard}{11}{singlecol}
\qtypebadge{singlecol}{Single Choice} \hfill \textcolor{darkslate!70}{\small Topic: Network Security Concepts}

Which term describes the likelihood and impact of a threat exploiting a vulnerability?

\begin{optlist}
\item Threat
\item Risk
\item Asset
\item Control
\end{optlist}
\begin{explainbox}{}
\textbf{Correct Answer: B}\\[0.3em]
\textbf{Concept:} Risk measures the probability and impact of a security incident.
\end{explainbox}
\end{questioncard}

\begin{questioncard}{12}{singlecol}
\qtypebadge{singlecol}{Single Choice} \hfill \textcolor{darkslate!70}{\small Topic: Network Security Concepts}

Future security architecture is moving toward:

\begin{optlist}
\item Single perimeter defense
\item Zero-trust and dynamic trust models
\item Open anonymous access
\item No logging
\end{optlist}
\begin{explainbox}{}
\textbf{Correct Answer: B}\\[0.3em]
\textbf{Concept:} Zero trust assumes no implicit trust and verifies every access request dynamically.
\end{explainbox}
\end{questioncard}

\begin{questioncard}{13}{singlecol}
\qtypebadge{singlecol}{Single Choice} \hfill \textcolor{darkslate!70}{\small Topic: Network Security Concepts}

Which technology is increasingly used to detect unknown threats?

\begin{optlist}
\item Signature-only antivirus
\item Artificial intelligence and machine learning
\item Hub-based networks
\item Plaintext protocols
\end{optlist}
\begin{explainbox}{}
\textbf{Correct Answer: B}\\[0.3em]
\textbf{Concept:} AI/ML helps analyze behavior and detect anomalies and unknown threats.
\end{explainbox}
\end{questioncard}

\begin{questioncard}{14}{singlecol}
\qtypebadge{singlecol}{Single Choice} \hfill \textcolor{darkslate!70}{\small Topic: Network Security Concepts}

What does confidentiality ensure?

\begin{optlist}
\item Only speed
\item Information is accessible only to authorized users
\item Unlimited access
\item Data duplication
\end{optlist}
\begin{explainbox}{}
\textbf{Correct Answer: B}\\[0.3em]
\textbf{Concept:} Confidentiality ensures that information is disclosed only to authorized individuals.
\end{explainbox}
\end{questioncard}

\begin{questioncard}{15}{singlecol}
\qtypebadge{singlecol}{Single Choice} \hfill \textcolor{darkslate!70}{\small Topic: Network Security Concepts}

What does integrity ensure?

\begin{optlist}
\item Fast transmission
\item Information is not tampered with during transmission/storage
\item Free access
\item Unencrypted storage
\end{optlist}
\begin{explainbox}{}
\textbf{Correct Answer: B}\\[0.3em]
\textbf{Concept:} Integrity ensures data is not altered or destroyed in an unauthorized manner.
\end{explainbox}
\end{questioncard}

\begin{questioncard}{16}{singlecol}
\qtypebadge{singlecol}{Single Choice} \hfill \textcolor{darkslate!70}{\small Topic: Network Security Concepts}

What does availability ensure?

\begin{optlist}
\item Authorized users can access information when needed
\item All users can access
\item Systems are offline
\item Data is hidden
\end{optlist}
\begin{explainbox}{}
\textbf{Correct Answer: A}\\[0.3em]
\textbf{Concept:} Availability guarantees reliable and timely access to information for authorized users.
\end{explainbox}
\end{questioncard}

\begin{questioncard}{17}{singlecol}
\qtypebadge{singlecol}{Single Choice} \hfill \textcolor{darkslate!70}{\small Topic: Network Security Concepts}

What does controllability mean in information security?

\begin{optlist}
\item No monitoring
\item Ability to monitor and control information/systems
\item Unlimited downloads
\item No access controls
\end{optlist}
\begin{explainbox}{}
\textbf{Correct Answer: B}\\[0.3em]
\textbf{Concept:} Controllability enables organizations to monitor, manage, and protect information and systems.
\end{explainbox}
\end{questioncard}

\begin{questioncard}{18}{singlecol}
\qtypebadge{singlecol}{Single Choice} \hfill \textcolor{darkslate!70}{\small Topic: Network Security Concepts}

What does non-repudiation prevent?

\begin{optlist}
\item Fast data transfer
\item Senders or receivers from denying their actions
\item Encryption
\item Backups
\end{optlist}
\begin{explainbox}{}
\textbf{Correct Answer: B}\\[0.3em]
\textbf{Concept:} Non-repudiation provides proof of origin and delivery so parties cannot deny participation.
\end{explainbox}
\end{questioncard}

\begin{questioncard}{19}{singlecol}
\qtypebadge{singlecol}{Single Choice} \hfill \textcolor{darkslate!70}{\small Topic: Network Security Concepts}

The broad definition of network security is closest to:

\begin{optlist}
\item Cybersecurity
\item Firewall configuration
\item Antivirus deployment
\item Physical security
\end{optlist}
\begin{explainbox}{}
\textbf{Correct Answer: A}\\[0.3em]
\textbf{Concept:} Broadly, network security equates to cybersecurity covering national and societal cyberspace protection.
\end{explainbox}
\end{questioncard}

\begin{questioncard}{20}{singlecol}
\qtypebadge{singlecol}{Single Choice} \hfill \textcolor{darkslate!70}{\small Topic: Network Security Concepts}

The narrow definition of network security focuses on:

\begin{optlist}
\item National laws only
\item Enterprise network protection solutions
\item Social media policies
\item Hardware warranties
\end{optlist}
\begin{explainbox}{}
\textbf{Correct Answer: B}\\[0.3em]
\textbf{Concept:} Narrowly, network security means vendor solutions and measures that keep enterprise networks secure.
\end{explainbox}
\end{questioncard}

\begin{questioncard}{21}{singlecol}
\qtypebadge{singlecol}{Single Choice} \hfill \textcolor{darkslate!70}{\small Topic: Network Security Concepts}

During which period did communication cipher security emerge?

\begin{optlist}
\item Cyberspace security period
\item Communication security period
\item Information assurance period
\item Cloud security period
\end{optlist}
\begin{explainbox}{}
\textbf{Correct Answer: B}\\[0.3em]
\textbf{Concept:} Stream ciphers and communication confidentiality appeared during the communication security period.
\end{explainbox}
\end{questioncard}

\begin{questioncard}{22}{singlecol}
\qtypebadge{singlecol}{Single Choice} \hfill \textcolor{darkslate!70}{\small Topic: Network Security Concepts}

The information security period added which goal to the earlier communication security focus?

\begin{optlist}
\item Availability
\item Confidentiality, integrity, and availability
\item Non-repudiation only
\item Controllability only
\end{optlist}
\begin{explainbox}{}
\textbf{Correct Answer: B}\\[0.3em]
\textbf{Concept:} The information security period formalized confidentiality, integrity, and availability as core goals.
\end{explainbox}
\end{questioncard}

\begin{questioncard}{23}{singlecol}
\qtypebadge{singlecol}{Single Choice} \hfill \textcolor{darkslate!70}{\small Topic: Network Security Concepts}

Controllability and non-repudiation became prominent during:

\begin{optlist}
\item Communication security period
\item Information security period
\item Information assurance period
\item Cloud computing period
\end{optlist}
\begin{explainbox}{}
\textbf{Correct Answer: C}\\[0.3em]
\textbf{Concept:} The information assurance period expanded security to include monitoring and proof of actions.
\end{explainbox}
\end{questioncard}

\begin{questioncard}{24}{singlecol}
\qtypebadge{singlecol}{Single Choice} \hfill \textcolor{darkslate!70}{\small Topic: Network Security Concepts}

Which concept is described as: is synonymous with national cyberspace security?

\begin{optlist}
\item Cyberspace security
\item Broad network security
\item Non-repudiation
\item Controllability
\end{optlist}
\begin{explainbox}{}
\textbf{Correct Answer: B}\\[0.3em]
\textbf{Concept:} Broad network security refers to the concept where is synonymous with national cyberspace security.
\end{explainbox}
\end{questioncard}

\begin{questioncard}{25}{singlecol}
\qtypebadge{singlecol}{Single Choice} \hfill \textcolor{darkslate!70}{\small Topic: Network Security Concepts}

Which concept is described as: focuses on enterprise network protection solutions?

\begin{optlist}
\item Integrity
\item AI and machine learning
\item Narrow network security
\item Confidentiality
\end{optlist}
\begin{explainbox}{}
\textbf{Correct Answer: C}\\[0.3em]
\textbf{Concept:} Narrow network security refers to the concept where focuses on enterprise network protection solutions.
\end{explainbox}
\end{questioncard}

\begin{questioncard}{26}{singlecol}
\qtypebadge{singlecol}{Single Choice} \hfill \textcolor{darkslate!70}{\small Topic: Network Security Concepts}

Which concept is described as: consists of confidentiality, integrity, and availability?

\begin{optlist}
\item Integrity
\item Vulnerability
\item Narrow network security
\item CIA triad
\end{optlist}
\begin{explainbox}{}
\textbf{Correct Answer: D}\\[0.3em]
\textbf{Concept:} CIA triad refers to the concept where consists of confidentiality, integrity, and availability.
\end{explainbox}
\end{questioncard}

\begin{questioncard}{27}{singlecol}
\qtypebadge{singlecol}{Single Choice} \hfill \textcolor{darkslate!70}{\small Topic: Network Security Concepts}

Which concept is described as: added controllability and non-repudiation?

\begin{optlist}
\item Information assurance period
\item AI and machine learning
\item Non-repudiation
\item Threat
\end{optlist}
\begin{explainbox}{}
\textbf{Correct Answer: A}\\[0.3em]
\textbf{Concept:} Information assurance period refers to the concept where added controllability and non-repudiation.
\end{explainbox}
\end{questioncard}

\begin{questioncard}{28}{singlecol}
\qtypebadge{singlecol}{Single Choice} \hfill \textcolor{darkslate!70}{\small Topic: Network Security Concepts}

Which concept is described as: protects facilities, data, users, and operations?

\begin{optlist}
\item CIA triad
\item Availability
\item Cyberspace security
\item Vulnerability
\end{optlist}
\begin{explainbox}{}
\textbf{Correct Answer: C}\\[0.3em]
\textbf{Concept:} Cyberspace security refers to the concept where protects facilities, data, users, and operations.
\end{explainbox}
\end{questioncard}

\begin{questioncard}{29}{singlecol}
\qtypebadge{singlecol}{Single Choice} \hfill \textcolor{darkslate!70}{\small Topic: Network Security Concepts}

Which concept is described as: ensures information is accessible only to authorized users?

\begin{optlist}
\item Availability
\item Risk
\item Non-repudiation
\item Confidentiality
\end{optlist}
\begin{explainbox}{}
\textbf{Correct Answer: D}\\[0.3em]
\textbf{Concept:} Confidentiality refers to the concept where ensures information is accessible only to authorized users.
\end{explainbox}
\end{questioncard}

\begin{questioncard}{30}{singlecol}
\qtypebadge{singlecol}{Single Choice} \hfill \textcolor{darkslate!70}{\small Topic: Network Security Concepts}

Which concept is described as: ensures data is not tampered with during transmission or storage?

\begin{optlist}
\item Availability
\item Broad network security
\item Future security trend
\item Integrity
\end{optlist}
\begin{explainbox}{}
\textbf{Correct Answer: D}\\[0.3em]
\textbf{Concept:} Integrity refers to the concept where ensures data is not tampered with during transmission or storage.
\end{explainbox}
\end{questioncard}

\begin{questioncard}{31}{singlecol}
\qtypebadge{singlecol}{Single Choice} \hfill \textcolor{darkslate!70}{\small Topic: Network Security Concepts}

Which concept is described as: ensures authorized users can access information when needed?

\begin{optlist}
\item Integrity
\item Exploit
\item AI and machine learning
\item Availability
\end{optlist}
\begin{explainbox}{}
\textbf{Correct Answer: D}\\[0.3em]
\textbf{Concept:} Availability refers to the concept where ensures authorized users can access information when needed.
\end{explainbox}
\end{questioncard}

\begin{questioncard}{32}{singlecol}
\qtypebadge{singlecol}{Single Choice} \hfill \textcolor{darkslate!70}{\small Topic: Network Security Concepts}

Which concept is described as: enables monitoring and control of information and systems?

\begin{optlist}
\item Non-repudiation
\item Broad network security
\item Controllability
\item Narrow network security
\end{optlist}
\begin{explainbox}{}
\textbf{Correct Answer: C}\\[0.3em]
\textbf{Concept:} Controllability refers to the concept where enables monitoring and control of information and systems.
\end{explainbox}
\end{questioncard}

\begin{questioncard}{33}{singlecol}
\qtypebadge{singlecol}{Single Choice} \hfill \textcolor{darkslate!70}{\small Topic: Network Security Concepts}

Which concept is described as: prevents senders or receivers from denying actions?

\begin{optlist}
\item Threat
\item Vulnerability
\item Non-repudiation
\item Controllability
\end{optlist}
\begin{explainbox}{}
\textbf{Correct Answer: C}\\[0.3em]
\textbf{Concept:} Non-repudiation refers to the concept where prevents senders or receivers from denying actions.
\end{explainbox}
\end{questioncard}

\begin{questioncard}{34}{singlecol}
\qtypebadge{singlecol}{Single Choice} \hfill \textcolor{darkslate!70}{\small Topic: Network Security Concepts}

Which concept is described as: is the combination of threat likelihood and impact?

\begin{optlist}
\item Exploit
\item Risk
\item Cyberspace security
\item Information assurance period
\end{optlist}
\begin{explainbox}{}
\textbf{Correct Answer: B}\\[0.3em]
\textbf{Concept:} Risk refers to the concept where is the combination of threat likelihood and impact.
\end{explainbox}
\end{questioncard}

\begin{questioncard}{35}{singlecol}
\qtypebadge{singlecol}{Single Choice} \hfill \textcolor{darkslate!70}{\small Topic: Network Security Concepts}

Which concept is described as: is a weakness that can be exploited by threats?

\begin{optlist}
\item Controllability
\item Risk
\item Confidentiality
\item Vulnerability
\end{optlist}
\begin{explainbox}{}
\textbf{Correct Answer: D}\\[0.3em]
\textbf{Concept:} Vulnerability refers to the concept where is a weakness that can be exploited by threats.
\end{explainbox}
\end{questioncard}

\begin{questioncard}{36}{singlecol}
\qtypebadge{singlecol}{Single Choice} \hfill \textcolor{darkslate!70}{\small Topic: Network Security Concepts}

Which concept is described as: is a potential danger to an asset?

\begin{optlist}
\item Confidentiality
\item Vulnerability
\item Threat
\item Exploit
\end{optlist}
\begin{explainbox}{}
\textbf{Correct Answer: C}\\[0.3em]
\textbf{Concept:} Threat refers to the concept where is a potential danger to an asset.
\end{explainbox}
\end{questioncard}

\begin{questioncard}{37}{singlecol}
\qtypebadge{singlecol}{Single Choice} \hfill \textcolor{darkslate!70}{\small Topic: Network Security Concepts}

Which concept is described as: is a method that takes advantage of a vulnerability?

\begin{optlist}
\item Non-repudiation
\item Confidentiality
\item Controllability
\item Exploit
\end{optlist}
\begin{explainbox}{}
\textbf{Correct Answer: D}\\[0.3em]
\textbf{Concept:} Exploit refers to the concept where is a method that takes advantage of a vulnerability.
\end{explainbox}
\end{questioncard}

\begin{questioncard}{38}{singlecol}
\qtypebadge{singlecol}{Single Choice} \hfill \textcolor{darkslate!70}{\small Topic: Network Security Concepts}

Which concept is described as: includes zero-trust architecture?

\begin{optlist}
\item Vulnerability
\item Future security trend
\item Broad network security
\item Integrity
\end{optlist}
\begin{explainbox}{}
\textbf{Correct Answer: B}\\[0.3em]
\textbf{Concept:} Future security trend refers to the concept where includes zero-trust architecture.
\end{explainbox}
\end{questioncard}

\begin{questioncard}{39}{singlecol}
\qtypebadge{singlecol}{Single Choice} \hfill \textcolor{darkslate!70}{\small Topic: Network Security Concepts}

Which concept is described as: help detect unknown threats and anomalies?

\begin{optlist}
\item AI and machine learning
\item Availability
\item Broad network security
\item Controllability
\end{optlist}
\begin{explainbox}{}
\textbf{Correct Answer: A}\\[0.3em]
\textbf{Concept:} AI and machine learning refers to the concept where help detect unknown threats and anomalies.
\end{explainbox}
\end{questioncard}

\section{Network Security Standards}
\begin{questioncard}{1}{singlecol}
\qtypebadge{singlecol}{Single Choice} \hfill \textcolor{darkslate!70}{\small Topic: Network Security Standards}

ISO 27001 specifies requirements for:

\begin{optlist}
\item An information security management system (ISMS)
\item A firewall hardware model
\item A routing protocol
\item A programming language
\end{optlist}
\begin{explainbox}{}
\textbf{Correct Answer: A}\\[0.3em]
\textbf{Concept:} ISO/IEC 27001 is the international standard for establishing, implementing, maintaining, and continually improving an ISMS.
\end{explainbox}
\end{questioncard}

\begin{questioncard}{2}{singlecol}
\qtypebadge{singlecol}{Single Choice} \hfill \textcolor{darkslate!70}{\small Topic: Network Security Standards}

ISO 27002 is best described as:

\begin{optlist}
\item A certification standard
\item A code of practice with security controls
\item A firewall configuration guide
\item A wireless protocol
\end{optlist}
\begin{explainbox}{}
\textbf{Correct Answer: B}\\[0.3em]
\textbf{Concept:} ISO/IEC 27002 provides best-practice security controls, while ISO 27001 is the certifiable requirements standard.
\end{explainbox}
\end{questioncard}

\begin{questioncard}{3}{singlecol}
\qtypebadge{singlecol}{Single Choice} \hfill \textcolor{darkslate!70}{\small Topic: Network Security Standards}

Cybersecurity Classified Protection 2.0 is issued by:

\begin{optlist}
\item Huawei
\item China (as a national standard/policy)
\item IEEE
\item IETF
\end{optlist}
\begin{explainbox}{}
\textbf{Correct Answer: B}\\[0.3em]
\textbf{Concept:} Classified Protection 2.0 is a Chinese national cybersecurity baseline and compliance framework.
\end{explainbox}
\end{questioncard}

\begin{questioncard}{4}{singlecol}
\qtypebadge{singlecol}{Single Choice} \hfill \textcolor{darkslate!70}{\small Topic: Network Security Standards}

Classified Protection 2.0 uses a protection-level scale from:

\begin{optlist}
\item Level 1 to Level 5
\item Level 0 to Level 4
\item Level A to Level E
\item Level I to Level V
\end{optlist}
\begin{explainbox}{}
\textbf{Correct Answer: A}\\[0.3em]
\textbf{Concept:} The Chinese classified protection system defines five levels, with Level 5 being the highest.
\end{explainbox}
\end{questioncard}

\begin{questioncard}{5}{singlecol}
\qtypebadge{singlecol}{Single Choice} \hfill \textcolor{darkslate!70}{\small Topic: Network Security Standards}

Which of the following is NOT a typical Classified Protection 2.0 lifecycle phase?

\begin{optlist}
\item Classification
\item Filing
\item Construction and rectification
\item Marketing and sales
\end{optlist}
\begin{explainbox}{}
\textbf{Correct Answer: D}\\[0.3em]
\textbf{Concept:} The Classified Protection lifecycle includes classification, filing, construction/rectification, level assessment, and supervision/inspection.
\end{explainbox}
\end{questioncard}

\begin{questioncard}{6}{singlecol}
\qtypebadge{singlecol}{Single Choice} \hfill \textcolor{darkslate!70}{\small Topic: Network Security Standards}

Which ISO standard family addresses information security management?

\begin{optlist}
\item ISO 9000
\item ISO 27000
\item ISO 14000
\item ISO 20000
\end{optlist}
\begin{explainbox}{}
\textbf{Correct Answer: B}\\[0.3em]
\textbf{Concept:} The ISO/IEC 27000 family covers information security management systems and controls.
\end{explainbox}
\end{questioncard}

\begin{questioncard}{7}{singlecol}
\qtypebadge{singlecol}{Single Choice} \hfill \textcolor{darkslate!70}{\small Topic: Network Security Standards}

Which Chinese regulation is the basis for Classified Protection?

\begin{optlist}
\item Cybersecurity Law
\item GDPR
\item HIPAA
\item PCI DSS
\end{optlist}
\begin{explainbox}{}
\textbf{Correct Answer: A}\\[0.3em]
\textbf{Concept:} China Cybersecurity Law underpins the Classified Protection compliance framework.
\end{explainbox}
\end{questioncard}

\begin{questioncard}{8}{singlecol}
\qtypebadge{singlecol}{Single Choice} \hfill \textcolor{darkslate!70}{\small Topic: Network Security Standards}

ISO 27001 uses which methodology for risk management?

\begin{optlist}
\item Ignore risks
\item Plan-Do-Check-Act (PDCA)
\item Random patching
\item No reviews
\end{optlist}
\begin{explainbox}{}
\textbf{Correct Answer: B}\\[0.3em]
\textbf{Concept:} ISO 27001 follows the PDCA cycle for continuous improvement of the ISMS.
\end{explainbox}
\end{questioncard}

\begin{questioncard}{9}{singlecol}
\qtypebadge{singlecol}{Single Choice} \hfill \textcolor{darkslate!70}{\small Topic: Network Security Standards}

Which document in ISO 27000 family provides implementation guidance for controls?

\begin{optlist}
\item ISO 27001
\item ISO 27002
\item ISO 27003
\item ISO 27005
\end{optlist}
\begin{explainbox}{}
\textbf{Correct Answer: B}\\[0.3em]
\textbf{Concept:} ISO 27002 lists security controls and implementation guidance.
\end{explainbox}
\end{questioncard}

\begin{questioncard}{10}{singlecol}
\qtypebadge{singlecol}{Single Choice} \hfill \textcolor{darkslate!70}{\small Topic: Network Security Standards}

Which is a key principle of Classified Protection 2.0?

\begin{optlist}
\item Security synchronization with construction
\item Security as an afterthought
\item No auditing
\item Anonymous access
\end{optlist}
\begin{explainbox}{}
\textbf{Correct Answer: A}\\[0.3em]
\textbf{Concept:} Classified Protection requires security to be planned, constructed, and used simultaneously with information systems.
\end{explainbox}
\end{questioncard}

\begin{questioncard}{11}{singlecol}
\qtypebadge{singlecol}{Single Choice} \hfill \textcolor{darkslate!70}{\small Topic: Network Security Standards}

In Classified Protection 2.0, the security protection level is preliminarily determined by the network operator, but when the level reaches which level or above, expert review and approval are required?

\begin{optlist}
\item Level 1
\item Level 2
\item Level 3
\item Level 4
\end{optlist}
\begin{explainbox}{}
\textbf{Correct Answer: B}\\[0.3em]
\textbf{Concept:} For protection Level 2 and above, the preliminary determination must be reviewed and approved; Level 1 can be determined by the operator.
\end{explainbox}
\end{questioncard}

\begin{questioncard}{12}{singlecol}
\qtypebadge{singlecol}{Single Choice} \hfill \textcolor{darkslate!70}{\small Topic: Network Security Standards}

Which standard is used to certify an ISMS?

\begin{optlist}
\item ISO 27001
\item ISO 27002
\item ISO 9001
\item ISO 14001
\end{optlist}
\begin{explainbox}{}
\textbf{Correct Answer: A}\\[0.3em]
\textbf{Concept:} ISO 27001 contains the certifiable requirements for an ISMS.
\end{explainbox}
\end{questioncard}

\begin{questioncard}{13}{singlecol}
\qtypebadge{singlecol}{Single Choice} \hfill \textcolor{darkslate!70}{\small Topic: Network Security Standards}

ISO 27002 is primarily a:

\begin{optlist}
\item Certification standard
\item Code of practice for security controls
\item Risk assessment methodology
\item Networking protocol
\end{optlist}
\begin{explainbox}{}
\textbf{Correct Answer: B}\\[0.3em]
\textbf{Concept:} ISO 27002 provides best-practice controls but is not itself a certification standard.
\end{explainbox}
\end{questioncard}

\begin{questioncard}{14}{singlecol}
\qtypebadge{singlecol}{Single Choice} \hfill \textcolor{darkslate!70}{\small Topic: Network Security Standards}

Classified Protection 2.0 compliance is primarily relevant in:

\begin{optlist}
\item China
\item United States only
\item European Union only
\item Every country equally
\end{optlist}
\begin{explainbox}{}
\textbf{Correct Answer: A}\\[0.3em]
\textbf{Concept:} Classified Protection 2.0 is a Chinese national cybersecurity compliance framework.
\end{explainbox}
\end{questioncard}

\begin{questioncard}{15}{singlecol}
\qtypebadge{singlecol}{Single Choice} \hfill \textcolor{darkslate!70}{\small Topic: Network Security Standards}

Which is the highest protection level in Classified Protection 2.0?

\begin{optlist}
\item Level 1
\item Level 3
\item Level 5
\item Level 10
\end{optlist}
\begin{explainbox}{}
\textbf{Correct Answer: C}\\[0.3em]
\textbf{Concept:} Classified Protection defines five levels, with Level 5 being the most stringent.
\end{explainbox}
\end{questioncard}

\begin{questioncard}{16}{singlecol}
\qtypebadge{singlecol}{Single Choice} \hfill \textcolor{darkslate!70}{\small Topic: Network Security Standards}

Expert review is required for Classified Protection levels starting at:

\begin{optlist}
\item Level 1
\item Level 2
\item Level 3
\item Level 4
\end{optlist}
\begin{explainbox}{}
\textbf{Correct Answer: B}\\[0.3em]
\textbf{Concept:} For Level 2 and above, expert review and approval are required.
\end{explainbox}
\end{questioncard}

\begin{questioncard}{17}{singlecol}
\qtypebadge{singlecol}{Single Choice} \hfill \textcolor{darkslate!70}{\small Topic: Network Security Standards}

Which concept is described as: specifies certifiable ISMS requirements?

\begin{optlist}
\item Classified Protection 2.0
\item ISO 27001:2022
\item ISO 27002
\item ISO 27001
\end{optlist}
\begin{explainbox}{}
\textbf{Correct Answer: D}\\[0.3em]
\textbf{Concept:} ISO 27001 refers to the concept where specifies certifiable ISMS requirements.
\end{explainbox}
\end{questioncard}

\begin{questioncard}{18}{singlecol}
\qtypebadge{singlecol}{Single Choice} \hfill \textcolor{darkslate!70}{\small Topic: Network Security Standards}

Which concept is described as: provides a code of practice for security controls?

\begin{optlist}
\item PDCA
\item Level 2 and above
\item Classified Protection 2.0
\item ISO 27002
\end{optlist}
\begin{explainbox}{}
\textbf{Correct Answer: D}\\[0.3em]
\textbf{Concept:} ISO 27002 refers to the concept where provides a code of practice for security controls.
\end{explainbox}
\end{questioncard}

\begin{questioncard}{19}{singlecol}
\qtypebadge{singlecol}{Single Choice} \hfill \textcolor{darkslate!70}{\small Topic: Network Security Standards}

Which concept is described as: is the continuous improvement methodology used in ISO 27001?

\begin{optlist}
\item Classified Protection lifecycle
\item Classified Protection 2.0
\item PDCA
\item Security synchronization
\end{optlist}
\begin{explainbox}{}
\textbf{Correct Answer: C}\\[0.3em]
\textbf{Concept:} PDCA refers to the concept where is the continuous improvement methodology used in ISO 27001.
\end{explainbox}
\end{questioncard}

\begin{questioncard}{20}{singlecol}
\qtypebadge{singlecol}{Single Choice} \hfill \textcolor{darkslate!70}{\small Topic: Network Security Standards}

Which concept is described as: is a Chinese national cybersecurity framework?

\begin{optlist}
\item ISO 27002
\item ISO 27001
\item Classified Protection levels
\item Classified Protection 2.0
\end{optlist}
\begin{explainbox}{}
\textbf{Correct Answer: D}\\[0.3em]
\textbf{Concept:} Classified Protection 2.0 refers to the concept where is a Chinese national cybersecurity framework.
\end{explainbox}
\end{questioncard}

\begin{questioncard}{21}{singlecol}
\qtypebadge{singlecol}{Single Choice} \hfill \textcolor{darkslate!70}{\small Topic: Network Security Standards}

Which concept is described as: range from Level 1 to Level 5?

\begin{optlist}
\item Classified Protection levels
\item Classified Protection lifecycle
\item ISO 27001
\item PDCA
\end{optlist}
\begin{explainbox}{}
\textbf{Correct Answer: A}\\[0.3em]
\textbf{Concept:} Classified Protection levels refers to the concept where range from Level 1 to Level 5.
\end{explainbox}
\end{questioncard}

\begin{questioncard}{22}{singlecol}
\qtypebadge{singlecol}{Single Choice} \hfill \textcolor{darkslate!70}{\small Topic: Network Security Standards}

Which concept is described as: require expert review and approval?

\begin{optlist}
\item Classified Protection 2.0
\item ISO 27002
\item Classified Protection lifecycle
\item Level 2 and above
\end{optlist}
\begin{explainbox}{}
\textbf{Correct Answer: D}\\[0.3em]
\textbf{Concept:} Level 2 and above refers to the concept where require expert review and approval.
\end{explainbox}
\end{questioncard}

\begin{questioncard}{23}{singlecol}
\qtypebadge{singlecol}{Single Choice} \hfill \textcolor{darkslate!70}{\small Topic: Network Security Standards}

Which concept is described as: includes classification, filing, construction, assessment, and supervision?

\begin{optlist}
\item Classified Protection lifecycle
\item Classified Protection levels
\item PDCA
\item ISO 27001
\end{optlist}
\begin{explainbox}{}
\textbf{Correct Answer: A}\\[0.3em]
\textbf{Concept:} Classified Protection lifecycle refers to the concept where includes classification, filing, construction, assessment, and supervision.
\end{explainbox}
\end{questioncard}

\begin{questioncard}{24}{singlecol}
\qtypebadge{singlecol}{Single Choice} \hfill \textcolor{darkslate!70}{\small Topic: Network Security Standards}

Which concept is described as: means planning security with information system construction?

\begin{optlist}
\item Classified Protection levels
\item PDCA
\item Security synchronization
\item Classified Protection 2.0
\end{optlist}
\begin{explainbox}{}
\textbf{Correct Answer: C}\\[0.3em]
\textbf{Concept:} Security synchronization refers to the concept where means planning security with information system construction.
\end{explainbox}
\end{questioncard}

\begin{questioncard}{25}{singlecol}
\qtypebadge{singlecol}{Single Choice} \hfill \textcolor{darkslate!70}{\small Topic: Network Security Standards}

Which concept is described as: organizes controls into four themes?

\begin{optlist}
\item ISO 27002
\item ISO 27001:2022
\item Classified Protection levels
\item Cybersecurity Law
\end{optlist}
\begin{explainbox}{}
\textbf{Correct Answer: B}\\[0.3em]
\textbf{Concept:} ISO 27001:2022 refers to the concept where organizes controls into four themes.
\end{explainbox}
\end{questioncard}

\begin{questioncard}{26}{singlecol}
\qtypebadge{singlecol}{Single Choice} \hfill \textcolor{darkslate!70}{\small Topic: Network Security Standards}

Which concept is described as: underpins Classified Protection in China?

\begin{optlist}
\item Cybersecurity Law
\item ISO 27001:2022
\item ISO 27002
\item Classified Protection lifecycle
\end{optlist}
\begin{explainbox}{}
\textbf{Correct Answer: A}\\[0.3em]
\textbf{Concept:} Cybersecurity Law refers to the concept where underpins Classified Protection in China.
\end{explainbox}
\end{questioncard}

\section{Network Fundamentals}
\begin{questioncard}{1}{singlecol}
\qtypebadge{singlecol}{Single Choice} \hfill \textcolor{darkslate!70}{\small Topic: Network Fundamentals}

What is the default subnet mask for a Class C IPv4 address?

\begin{optlist}
\item 255.0.0.0
\item 255.255.0.0
\item 255.255.255.0
\item 255.255.255.255
\end{optlist}
\begin{explainbox}{}
\textbf{Correct Answer: C}\\[0.3em]
\textbf{Concept:} Class C networks use a default /24 subnet mask: 255.255.255.0.
\end{explainbox}
\end{questioncard}

\begin{questioncard}{2}{singlecol}
\qtypebadge{singlecol}{Single Choice} \hfill \textcolor{darkslate!70}{\small Topic: Network Fundamentals}

Which protocol is used to automatically assign IPv6 addresses using MAC addresses?

\begin{optlist}
\item DHCPv6
\item SLAAC
\item ARP
\item DNS
\end{optlist}
\begin{explainbox}{}
\textbf{Correct Answer: B}\\[0.3em]
\textbf{Concept:} SLAAC can generate IPv6 addresses from MAC addresses.
\end{explainbox}
\end{questioncard}

\begin{questioncard}{3}{singlecol}
\qtypebadge{singlecol}{Single Choice} \hfill \textcolor{darkslate!70}{\small Topic: Network Fundamentals}

Which layer-2 protocol prevents loops in redundant Ethernet topologies?

\begin{optlist}
\item VTP
\item STP
\item RSTP
\item Both STP and RSTP
\end{optlist}
\begin{explainbox}{}
\textbf{Correct Answer: D}\\[0.3em]
\textbf{Concept:} Spanning Tree Protocol and Rapid STP prevent Layer 2 loops.
\end{explainbox}
\end{questioncard}

\begin{questioncard}{4}{singlecol}
\qtypebadge{singlecol}{Single Choice} \hfill \textcolor{darkslate!70}{\small Topic: Network Fundamentals}

What does VLAN stand for?

\begin{optlist}
\item Virtual Local Area Network
\item Virtual Large Area Network
\item Verified Local Access Node
\item Virtual Link Aggregation Network
\end{optlist}
\begin{explainbox}{}
\textbf{Correct Answer: A}\\[0.3em]
\textbf{Concept:} VLANs logically segment a switch into separate broadcast domains.
\end{explainbox}
\end{questioncard}

\begin{questioncard}{5}{singlecol}
\qtypebadge{singlecol}{Single Choice} \hfill \textcolor{darkslate!70}{\small Topic: Network Fundamentals}

Which device interconnects VLANs at Layer 3?

\begin{optlist}
\item Layer 2 switch
\item Router or Layer 3 switch
\item Hub
\item Bridge
\end{optlist}
\begin{explainbox}{}
\textbf{Correct Answer: B}\\[0.3em]
\textbf{Concept:} Routers or Layer 3 switches perform inter-VLAN routing.
\end{explainbox}
\end{questioncard}

\begin{questioncard}{6}{singlecol}
\qtypebadge{singlecol}{Single Choice} \hfill \textcolor{darkslate!70}{\small Topic: Network Fundamentals}

Which OSI layer is responsible for reliable data transfer between hosts?

\begin{optlist}
\item Network
\item Transport
\item Data link
\item Physical
\end{optlist}
\begin{explainbox}{}
\textbf{Correct Answer: B}\\[0.3em]
\textbf{Concept:} The transport layer manages end-to-end reliability, flow control, and error recovery.
\end{explainbox}
\end{questioncard}

\begin{questioncard}{7}{singlecol}
\qtypebadge{singlecol}{Single Choice} \hfill \textcolor{darkslate!70}{\small Topic: Network Fundamentals}

In the TCP/IP model, which layer corresponds to the OSI network layer?

\begin{optlist}
\item Internet layer
\item Network interface layer
\item Transport layer
\item Application layer
\end{optlist}
\begin{explainbox}{}
\textbf{Correct Answer: A}\\[0.3em]
\textbf{Concept:} The TCP/IP internet layer maps to the OSI network layer and handles IP routing.
\end{explainbox}
\end{questioncard}

\begin{questioncard}{8}{singlecol}
\qtypebadge{singlecol}{Single Choice} \hfill \textcolor{darkslate!70}{\small Topic: Network Fundamentals}

Which layer handles encryption and data format conversion?

\begin{optlist}
\item Application
\item Presentation
\item Session
\item Transport
\end{optlist}
\begin{explainbox}{}
\textbf{Correct Answer: B}\\[0.3em]
\textbf{Concept:} The presentation layer handles syntax, encryption, compression, and character encoding.
\end{explainbox}
\end{questioncard}

\begin{questioncard}{9}{singlecol}
\qtypebadge{singlecol}{Single Choice} \hfill \textcolor{darkslate!70}{\small Topic: Network Fundamentals}

Which layer establishes, manages, and terminates application dialogs?

\begin{optlist}
\item Application
\item Presentation
\item Session
\item Transport
\end{optlist}
\begin{explainbox}{}
\textbf{Correct Answer: C}\\[0.3em]
\textbf{Concept:} The session layer controls dialog establishment, maintenance, and termination.
\end{explainbox}
\end{questioncard}

\begin{questioncard}{10}{singlecol}
\qtypebadge{singlecol}{Single Choice} \hfill \textcolor{darkslate!70}{\small Topic: Network Fundamentals}

Which protocol is used to send email between servers?

\begin{optlist}
\item HTTP
\item SMTP
\item POP3
\item IMAP
\end{optlist}
\begin{explainbox}{}
\textbf{Correct Answer: B}\\[0.3em]
\textbf{Concept:} SMTP is used to send and relay email.
\end{explainbox}
\end{questioncard}

\begin{questioncard}{11}{singlecol}
\qtypebadge{singlecol}{Single Choice} \hfill \textcolor{darkslate!70}{\small Topic: Network Fundamentals}

Which protocol is used by clients to retrieve email from a server while keeping messages on the server?

\begin{optlist}
\item POP3
\item IMAP
\item SMTP
\item SNMP
\end{optlist}
\begin{explainbox}{}
\textbf{Correct Answer: B}\\[0.3em]
\textbf{Concept:} IMAP allows clients to access and manage email while messages remain on the server.
\end{explainbox}
\end{questioncard}

\begin{questioncard}{12}{singlecol}
\qtypebadge{singlecol}{Single Choice} \hfill \textcolor{darkslate!70}{\small Topic: Network Fundamentals}

Which protocol dynamically assigns IP addresses to hosts?

\begin{optlist}
\item DNS
\item DHCP
\item ARP
\item FTP
\end{optlist}
\begin{explainbox}{}
\textbf{Correct Answer: B}\\[0.3em]
\textbf{Concept:} DHCP dynamically assigns IP addresses and network configuration parameters.
\end{explainbox}
\end{questioncard}

\begin{questioncard}{13}{singlecol}
\qtypebadge{singlecol}{Single Choice} \hfill \textcolor{darkslate!70}{\small Topic: Network Fundamentals}

Which protocol maps IP addresses to MAC addresses on a local network?

\begin{optlist}
\item DNS
\item ARP
\item DHCP
\item ICMP
\end{optlist}
\begin{explainbox}{}
\textbf{Correct Answer: B}\\[0.3em]
\textbf{Concept:} ARP resolves IP addresses to MAC addresses.
\end{explainbox}
\end{questioncard}

\begin{questioncard}{14}{singlecol}
\qtypebadge{singlecol}{Single Choice} \hfill \textcolor{darkslate!70}{\small Topic: Network Fundamentals}

Which protocol is used for network management and monitoring?

\begin{optlist}
\item SNMP
\item SMTP
\item FTP
\item Telnet
\end{optlist}
\begin{explainbox}{}
\textbf{Correct Answer: A}\\[0.3em]
\textbf{Concept:} SNMP manages and monitors network devices.
\end{explainbox}
\end{questioncard}

\begin{questioncard}{15}{singlecol}
\qtypebadge{singlecol}{Single Choice} \hfill \textcolor{darkslate!70}{\small Topic: Network Fundamentals}

What does ICMP primarily provide?

\begin{optlist}
\item File transfer
\item Error reporting and diagnostics
\item Email delivery
\item Web browsing
\end{optlist}
\begin{explainbox}{}
\textbf{Correct Answer: B}\\[0.3em]
\textbf{Concept:} ICMP reports errors and provides diagnostic functions such as ping and traceroute.
\end{explainbox}
\end{questioncard}

\begin{questioncard}{16}{singlecol}
\qtypebadge{singlecol}{Single Choice} \hfill \textcolor{darkslate!70}{\small Topic: Network Fundamentals}

Which port does HTTP use by default?

\begin{optlist}
\item 20
\item 21
\item 80
\item 443
\end{optlist}
\begin{explainbox}{}
\textbf{Correct Answer: C}\\[0.3em]
\textbf{Concept:} HTTP uses TCP port 80 by default.
\end{explainbox}
\end{questioncard}

\begin{questioncard}{17}{singlecol}
\qtypebadge{singlecol}{Single Choice} \hfill \textcolor{darkslate!70}{\small Topic: Network Fundamentals}

Which port does HTTPS use by default?

\begin{optlist}
\item 80
\item 443
\item 22
\item 25
\end{optlist}
\begin{explainbox}{}
\textbf{Correct Answer: B}\\[0.3em]
\textbf{Concept:} HTTPS uses TCP port 443 by default, providing encrypted web traffic.
\end{explainbox}
\end{questioncard}

\begin{questioncard}{18}{singlecol}
\qtypebadge{singlecol}{Single Choice} \hfill \textcolor{darkslate!70}{\small Topic: Network Fundamentals}

Which port does SSH use by default?

\begin{optlist}
\item 21
\item 22
\item 23
\item 25
\end{optlist}
\begin{explainbox}{}
\textbf{Correct Answer: B}\\[0.3em]
\textbf{Concept:} SSH uses TCP port 22 for secure remote access.
\end{explainbox}
\end{questioncard}

\begin{questioncard}{19}{singlecol}
\qtypebadge{singlecol}{Single Choice} \hfill \textcolor{darkslate!70}{\small Topic: Network Fundamentals}

Telnet uses which port and is considered insecure because it transmits data in plaintext?

\begin{optlist}
\item 21
\item 22
\item 23
\item 25
\end{optlist}
\begin{explainbox}{}
\textbf{Correct Answer: C}\\[0.3em]
\textbf{Concept:} Telnet uses TCP port 23 and sends data, including credentials, in plaintext.
\end{explainbox}
\end{questioncard}

\begin{questioncard}{20}{singlecol}
\qtypebadge{singlecol}{Single Choice} \hfill \textcolor{darkslate!70}{\small Topic: Network Fundamentals}

Which address is a Layer 2 address?

\begin{optlist}
\item IP address
\item MAC address
\item Port number
\item Domain name
\end{optlist}
\begin{explainbox}{}
\textbf{Correct Answer: B}\\[0.3em]
\textbf{Concept:} A MAC address is a hardware address used at the data link layer.
\end{explainbox}
\end{questioncard}

\begin{questioncard}{21}{singlecol}
\qtypebadge{singlecol}{Single Choice} \hfill \textcolor{darkslate!70}{\small Topic: Network Fundamentals}

A router forwards packets based on:

\begin{optlist}
\item MAC address
\item IP address
\item Port number
\item VLAN ID
\end{optlist}
\begin{explainbox}{}
\textbf{Correct Answer: B}\\[0.3em]
\textbf{Concept:} Routers use IP addresses to make forwarding decisions.
\end{explainbox}
\end{questioncard}

\begin{questioncard}{22}{singlecol}
\qtypebadge{singlecol}{Single Choice} \hfill \textcolor{darkslate!70}{\small Topic: Network Fundamentals}

What is a subnet mask used for?

\begin{optlist}
\item MAC address
\item IP address class
\item Network and host portions of an IP address
\item DNS server
\end{optlist}
\begin{explainbox}{}
\textbf{Correct Answer: C}\\[0.3em]
\textbf{Concept:} A subnet mask divides an IP address into network and host portions.
\end{explainbox}
\end{questioncard}

\begin{questioncard}{23}{singlecol}
\qtypebadge{singlecol}{Single Choice} \hfill \textcolor{darkslate!70}{\small Topic: Network Fundamentals}

The OSI model layer closest to the physical medium is:

\begin{optlist}
\item Data link
\item Physical
\item Network
\item Transport
\end{optlist}
\begin{explainbox}{}
\textbf{Correct Answer: B}\\[0.3em]
\textbf{Concept:} The physical layer deals with bits, cables, and signals.
\end{explainbox}
\end{questioncard}

\begin{questioncard}{24}{singlecol}
\qtypebadge{singlecol}{Single Choice} \hfill \textcolor{darkslate!70}{\small Topic: Network Fundamentals}

TCP is considered reliable because it uses:

\begin{optlist}
\item Sequence numbers and acknowledgments
\item Broadcasting
\item No connection setup
\item Fixed ports only
\end{optlist}
\begin{explainbox}{}
\textbf{Correct Answer: A}\\[0.3em]
\textbf{Concept:} TCP uses sequence numbers, acknowledgments, and retransmissions for reliability.
\end{explainbox}
\end{questioncard}

\begin{questioncard}{25}{singlecol}
\qtypebadge{singlecol}{Single Choice} \hfill \textcolor{darkslate!70}{\small Topic: Network Fundamentals}

UDP is preferred for:

\begin{optlist}
\item Reliable file transfer
\item Real-time streaming and low-latency applications
\item Email delivery
\item Web browsing requiring guaranteed delivery
\end{optlist}
\begin{explainbox}{}
\textbf{Correct Answer: B}\\[0.3em]
\textbf{Concept:} UDP low overhead and lack of retransmission make it suitable for real-time traffic.
\end{explainbox}
\end{questioncard}

\begin{questioncard}{26}{singlecol}
\qtypebadge{singlecol}{Single Choice} \hfill \textcolor{darkslate!70}{\small Topic: Network Fundamentals}

DNS primarily uses:

\begin{optlist}
\item TCP port 53
\item UDP port 53
\item TCP port 80
\item UDP port 161
\end{optlist}
\begin{explainbox}{}
\textbf{Correct Answer: B}\\[0.3em]
\textbf{Concept:} DNS queries primarily use UDP port 53.
\end{explainbox}
\end{questioncard}

\begin{questioncard}{27}{singlecol}
\qtypebadge{singlecol}{Single Choice} \hfill \textcolor{darkslate!70}{\small Topic: Network Fundamentals}

A Layer 3 switch combines functions of a:

\begin{optlist}
\item Hub and repeater
\item Switch and router
\item Firewall and IDS
\item Modem and gateway
\end{optlist}
\begin{explainbox}{}
\textbf{Correct Answer: B}\\[0.3em]
\textbf{Concept:} Layer 3 switches perform both Layer 2 switching and Layer 3 routing.
\end{explainbox}
\end{questioncard}

\begin{questioncard}{28}{singlecol}
\qtypebadge{singlecol}{Single Choice} \hfill \textcolor{darkslate!70}{\small Topic: Network Fundamentals}

Which concept is described as: has seven layers?

\begin{optlist}
\item Hub
\item OSI model
\item SSH
\item TCP/IP equivalent model
\end{optlist}
\begin{explainbox}{}
\textbf{Correct Answer: B}\\[0.3em]
\textbf{Concept:} OSI model refers to the concept where has seven layers.
\end{explainbox}
\end{questioncard}

\begin{questioncard}{29}{singlecol}
\qtypebadge{singlecol}{Single Choice} \hfill \textcolor{darkslate!70}{\small Topic: Network Fundamentals}

Which concept is described as: has four layers?

\begin{optlist}
\item Physical layer
\item TCP/IP standard model
\item RFC 1918
\item SMTP
\end{optlist}
\begin{explainbox}{}
\textbf{Correct Answer: B}\\[0.3em]
\textbf{Concept:} TCP/IP standard model refers to the concept where has four layers.
\end{explainbox}
\end{questioncard}

\begin{questioncard}{30}{singlecol}
\qtypebadge{singlecol}{Single Choice} \hfill \textcolor{darkslate!70}{\small Topic: Network Fundamentals}

Which concept is described as: has five layers?

\begin{optlist}
\item MAC address
\item HTTP
\item TCP/IP equivalent model
\item STP
\end{optlist}
\begin{explainbox}{}
\textbf{Correct Answer: C}\\[0.3em]
\textbf{Concept:} TCP/IP equivalent model refers to the concept where has five layers.
\end{explainbox}
\end{questioncard}

\begin{questioncard}{31}{singlecol}
\qtypebadge{singlecol}{Single Choice} \hfill \textcolor{darkslate!70}{\small Topic: Network Fundamentals}

Which concept is described as: transmits raw bits over a medium?

\begin{optlist}
\item Data link layer
\item SMTP
\item Physical layer
\item MAC address
\end{optlist}
\begin{explainbox}{}
\textbf{Correct Answer: C}\\[0.3em]
\textbf{Concept:} Physical layer refers to the concept where transmits raw bits over a medium.
\end{explainbox}
\end{questioncard}

\begin{questioncard}{32}{singlecol}
\qtypebadge{singlecol}{Single Choice} \hfill \textcolor{darkslate!70}{\small Topic: Network Fundamentals}

Which concept is described as: handles framing and MAC addressing?

\begin{optlist}
\item TCP/IP equivalent model
\item Switch
\item Data link layer
\item Application layer
\end{optlist}
\begin{explainbox}{}
\textbf{Correct Answer: C}\\[0.3em]
\textbf{Concept:} Data link layer refers to the concept where handles framing and MAC addressing.
\end{explainbox}
\end{questioncard}

\begin{questioncard}{33}{singlecol}
\qtypebadge{singlecol}{Single Choice} \hfill \textcolor{darkslate!70}{\small Topic: Network Fundamentals}

Which concept is described as: handles logical addressing and routing?

\begin{optlist}
\item ICMP
\item Network layer
\item Transport layer
\item Application layer
\end{optlist}
\begin{explainbox}{}
\textbf{Correct Answer: B}\\[0.3em]
\textbf{Concept:} Network layer refers to the concept where handles logical addressing and routing.
\end{explainbox}
\end{questioncard}

\begin{questioncard}{34}{singlecol}
\qtypebadge{singlecol}{Single Choice} \hfill \textcolor{darkslate!70}{\small Topic: Network Fundamentals}

Which concept is described as: provides end-to-end communication and reliability?

\begin{optlist}
\item SMTP
\item Transport layer
\item Application layer
\item Hub
\end{optlist}
\begin{explainbox}{}
\textbf{Correct Answer: B}\\[0.3em]
\textbf{Concept:} Transport layer refers to the concept where provides end-to-end communication and reliability.
\end{explainbox}
\end{questioncard}

\begin{questioncard}{35}{singlecol}
\qtypebadge{singlecol}{Single Choice} \hfill \textcolor{darkslate!70}{\small Topic: Network Fundamentals}

Which concept is described as: manages application dialogs?

\begin{optlist}
\item MAC address
\item ARP
\item Session layer
\item DNS
\end{optlist}
\begin{explainbox}{}
\textbf{Correct Answer: C}\\[0.3em]
\textbf{Concept:} Session layer refers to the concept where manages application dialogs.
\end{explainbox}
\end{questioncard}

\begin{questioncard}{36}{singlecol}
\qtypebadge{singlecol}{Single Choice} \hfill \textcolor{darkslate!70}{\small Topic: Network Fundamentals}

Which concept is described as: handles encryption and data format conversion?

\begin{optlist}
\item ICMP
\item HTTPS
\item Presentation layer
\item Hub
\end{optlist}
\begin{explainbox}{}
\textbf{Correct Answer: C}\\[0.3em]
\textbf{Concept:} Presentation layer refers to the concept where handles encryption and data format conversion.
\end{explainbox}
\end{questioncard}

\begin{questioncard}{37}{singlecol}
\qtypebadge{singlecol}{Single Choice} \hfill \textcolor{darkslate!70}{\small Topic: Network Fundamentals}

Which concept is described as: provides network services to applications?

\begin{optlist}
\item SMTP
\item Application layer
\item FTP
\item UDP
\end{optlist}
\begin{explainbox}{}
\textbf{Correct Answer: B}\\[0.3em]
\textbf{Concept:} Application layer refers to the concept where provides network services to applications.
\end{explainbox}
\end{questioncard}

\begin{questioncard}{38}{singlecol}
\qtypebadge{singlecol}{Single Choice} \hfill \textcolor{darkslate!70}{\small Topic: Network Fundamentals}

Which concept is described as: is connection-oriented and reliable?

\begin{optlist}
\item Switch
\item TCP
\item UDP
\item DNS
\end{optlist}
\begin{explainbox}{}
\textbf{Correct Answer: B}\\[0.3em]
\textbf{Concept:} TCP refers to the concept where is connection-oriented and reliable.
\end{explainbox}
\end{questioncard}

\begin{questioncard}{39}{singlecol}
\qtypebadge{singlecol}{Single Choice} \hfill \textcolor{darkslate!70}{\small Topic: Network Fundamentals}

Which concept is described as: is connectionless and has low overhead?

\begin{optlist}
\item Switch
\item UDP
\item STP
\item SSH
\end{optlist}
\begin{explainbox}{}
\textbf{Correct Answer: B}\\[0.3em]
\textbf{Concept:} UDP refers to the concept where is connectionless and has low overhead.
\end{explainbox}
\end{questioncard}

\begin{questioncard}{40}{singlecol}
\qtypebadge{singlecol}{Single Choice} \hfill \textcolor{darkslate!70}{\small Topic: Network Fundamentals}

Which concept is described as: uses TCP ports 20 and 21?

\begin{optlist}
\item HTTPS
\item Presentation layer
\item OSI model
\item FTP
\end{optlist}
\begin{explainbox}{}
\textbf{Correct Answer: D}\\[0.3em]
\textbf{Concept:} FTP refers to the concept where uses TCP ports 20 and 21.
\end{explainbox}
\end{questioncard}

\begin{questioncard}{41}{singlecol}
\qtypebadge{singlecol}{Single Choice} \hfill \textcolor{darkslate!70}{\small Topic: Network Fundamentals}

Which concept is described as: uses TCP port 80?

\begin{optlist}
\item ARP
\item ICMP
\item FTP
\item HTTP
\end{optlist}
\begin{explainbox}{}
\textbf{Correct Answer: D}\\[0.3em]
\textbf{Concept:} HTTP refers to the concept where uses TCP port 80.
\end{explainbox}
\end{questioncard}

\begin{questioncard}{42}{singlecol}
\qtypebadge{singlecol}{Single Choice} \hfill \textcolor{darkslate!70}{\small Topic: Network Fundamentals}

Which concept is described as: uses TCP port 443?

\begin{optlist}
\item STP
\item VLAN
\item SSH
\item HTTPS
\end{optlist}
\begin{explainbox}{}
\textbf{Correct Answer: D}\\[0.3em]
\textbf{Concept:} HTTPS refers to the concept where uses TCP port 443.
\end{explainbox}
\end{questioncard}

\begin{questioncard}{43}{singlecol}
\qtypebadge{singlecol}{Single Choice} \hfill \textcolor{darkslate!70}{\small Topic: Network Fundamentals}

Which concept is described as: uses TCP port 22?

\begin{optlist}
\item FTP
\item Network layer
\item SSH
\item DHCP
\end{optlist}
\begin{explainbox}{}
\textbf{Correct Answer: C}\\[0.3em]
\textbf{Concept:} SSH refers to the concept where uses TCP port 22.
\end{explainbox}
\end{questioncard}

\begin{questioncard}{44}{singlecol}
\qtypebadge{singlecol}{Single Choice} \hfill \textcolor{darkslate!70}{\small Topic: Network Fundamentals}

Which concept is described as: uses TCP port 23 and is insecure?

\begin{optlist}
\item Telnet
\item Application layer
\item SMTP
\item IPv4 address
\end{optlist}
\begin{explainbox}{}
\textbf{Correct Answer: A}\\[0.3em]
\textbf{Concept:} Telnet refers to the concept where uses TCP port 23 and is insecure.
\end{explainbox}
\end{questioncard}

\begin{questioncard}{45}{singlecol}
\qtypebadge{singlecol}{Single Choice} \hfill \textcolor{darkslate!70}{\small Topic: Network Fundamentals}

Which concept is described as: primarily uses UDP port 53?

\begin{optlist}
\item Switch
\item DNS
\item Router
\item SMTP
\end{optlist}
\begin{explainbox}{}
\textbf{Correct Answer: B}\\[0.3em]
\textbf{Concept:} DNS refers to the concept where primarily uses UDP port 53.
\end{explainbox}
\end{questioncard}

\begin{questioncard}{46}{singlecol}
\qtypebadge{singlecol}{Single Choice} \hfill \textcolor{darkslate!70}{\small Topic: Network Fundamentals}

Which concept is described as: uses TCP port 25?

\begin{optlist}
\item TCP/IP equivalent model
\item UDP
\item Presentation layer
\item SMTP
\end{optlist}
\begin{explainbox}{}
\textbf{Correct Answer: D}\\[0.3em]
\textbf{Concept:} SMTP refers to the concept where uses TCP port 25.
\end{explainbox}
\end{questioncard}

\begin{questioncard}{47}{singlecol}
\qtypebadge{singlecol}{Single Choice} \hfill \textcolor{darkslate!70}{\small Topic: Network Fundamentals}

Which concept is described as: dynamically assigns IP addresses?

\begin{optlist}
\item RFC 1918
\item FTP
\item VLAN
\item DHCP
\end{optlist}
\begin{explainbox}{}
\textbf{Correct Answer: D}\\[0.3em]
\textbf{Concept:} DHCP refers to the concept where dynamically assigns IP addresses.
\end{explainbox}
\end{questioncard}

\begin{questioncard}{48}{singlecol}
\qtypebadge{singlecol}{Single Choice} \hfill \textcolor{darkslate!70}{\small Topic: Network Fundamentals}

Which concept is described as: maps IP addresses to MAC addresses?

\begin{optlist}
\item SSH
\item Presentation layer
\item MAC address
\item ARP
\end{optlist}
\begin{explainbox}{}
\textbf{Correct Answer: D}\\[0.3em]
\textbf{Concept:} ARP refers to the concept where maps IP addresses to MAC addresses.
\end{explainbox}
\end{questioncard}

\begin{questioncard}{49}{singlecol}
\qtypebadge{singlecol}{Single Choice} \hfill \textcolor{darkslate!70}{\small Topic: Network Fundamentals}

Which concept is described as: provides error reporting and diagnostics?

\begin{optlist}
\item ICMP
\item STP
\item ARP
\item Session layer
\end{optlist}
\begin{explainbox}{}
\textbf{Correct Answer: A}\\[0.3em]
\textbf{Concept:} ICMP refers to the concept where provides error reporting and diagnostics.
\end{explainbox}
\end{questioncard}

\begin{questioncard}{50}{singlecol}
\qtypebadge{singlecol}{Single Choice} \hfill \textcolor{darkslate!70}{\small Topic: Network Fundamentals}

Which concept is described as: is 32 bits long?

\begin{optlist}
\item IPv4 address
\item OSI model
\item MAC address
\item Router
\end{optlist}
\begin{explainbox}{}
\textbf{Correct Answer: A}\\[0.3em]
\textbf{Concept:} IPv4 address refers to the concept where is 32 bits long.
\end{explainbox}
\end{questioncard}

\begin{questioncard}{51}{singlecol}
\qtypebadge{singlecol}{Single Choice} \hfill \textcolor{darkslate!70}{\small Topic: Network Fundamentals}

Which concept is described as: is 128 bits long?

\begin{optlist}
\item RFC 1918
\item IPv6 address
\item Telnet
\item STP
\end{optlist}
\begin{explainbox}{}
\textbf{Correct Answer: B}\\[0.3em]
\textbf{Concept:} IPv6 address refers to the concept where is 128 bits long.
\end{explainbox}
\end{questioncard}

\begin{questioncard}{52}{singlecol}
\qtypebadge{singlecol}{Single Choice} \hfill \textcolor{darkslate!70}{\small Topic: Network Fundamentals}

Which concept is described as: is a Layer 2 hardware address?

\begin{optlist}
\item VLAN
\item MAC address
\item Network layer
\item OSI model
\end{optlist}
\begin{explainbox}{}
\textbf{Correct Answer: B}\\[0.3em]
\textbf{Concept:} MAC address refers to the concept where is a Layer 2 hardware address.
\end{explainbox}
\end{questioncard}

\begin{questioncard}{53}{singlecol}
\qtypebadge{singlecol}{Single Choice} \hfill \textcolor{darkslate!70}{\small Topic: Network Fundamentals}

Which concept is described as: operates at Layer 2 and forwards frames by MAC?

\begin{optlist}
\item Switch
\item STP
\item IPv6 address
\item Transport layer
\end{optlist}
\begin{explainbox}{}
\textbf{Correct Answer: A}\\[0.3em]
\textbf{Concept:} Switch refers to the concept where operates at Layer 2 and forwards frames by MAC.
\end{explainbox}
\end{questioncard}

\begin{questioncard}{54}{singlecol}
\qtypebadge{singlecol}{Single Choice} \hfill \textcolor{darkslate!70}{\small Topic: Network Fundamentals}

Which concept is described as: operates at Layer 3 and forwards packets by IP?

\begin{optlist}
\item FTP
\item VLAN
\item SSH
\item Router
\end{optlist}
\begin{explainbox}{}
\textbf{Correct Answer: D}\\[0.3em]
\textbf{Concept:} Router refers to the concept where operates at Layer 3 and forwards packets by IP.
\end{explainbox}
\end{questioncard}

\begin{questioncard}{55}{singlecol}
\qtypebadge{singlecol}{Single Choice} \hfill \textcolor{darkslate!70}{\small Topic: Network Fundamentals}

Which concept is described as: operates at Layer 1 and repeats signals?

\begin{optlist}
\item DHCP
\item TCP
\item Hub
\item DNS
\end{optlist}
\begin{explainbox}{}
\textbf{Correct Answer: C}\\[0.3em]
\textbf{Concept:} Hub refers to the concept where operates at Layer 1 and repeats signals.
\end{explainbox}
\end{questioncard}

\begin{questioncard}{56}{singlecol}
\qtypebadge{singlecol}{Single Choice} \hfill \textcolor{darkslate!70}{\small Topic: Network Fundamentals}

Which concept is described as: logically segments a network into broadcast domains?

\begin{optlist}
\item UDP
\item SSH
\item VLAN
\item DNS
\end{optlist}
\begin{explainbox}{}
\textbf{Correct Answer: C}\\[0.3em]
\textbf{Concept:} VLAN refers to the concept where logically segments a network into broadcast domains.
\end{explainbox}
\end{questioncard}

\begin{questioncard}{57}{singlecol}
\qtypebadge{singlecol}{Single Choice} \hfill \textcolor{darkslate!70}{\small Topic: Network Fundamentals}

Which concept is described as: prevents Layer 2 loops?

\begin{optlist}
\item STP
\item SMTP
\item VLAN
\item MAC address
\end{optlist}
\begin{explainbox}{}
\textbf{Correct Answer: A}\\[0.3em]
\textbf{Concept:} STP refers to the concept where prevents Layer 2 loops.
\end{explainbox}
\end{questioncard}

\begin{questioncard}{58}{singlecol}
\qtypebadge{singlecol}{Single Choice} \hfill \textcolor{darkslate!70}{\small Topic: Network Fundamentals}

Which concept is described as: defines private IPv4 address ranges?

\begin{optlist}
\item RFC 1918
\item TCP/IP standard model
\item DHCP
\item Physical layer
\end{optlist}
\begin{explainbox}{}
\textbf{Correct Answer: A}\\[0.3em]
\textbf{Concept:} RFC 1918 refers to the concept where defines private IPv4 address ranges.
\end{explainbox}
\end{questioncard}

\section{OSI and TCP/IP Models}
\begin{questioncard}{1}{singlecol}
\qtypebadge{singlecol}{Single Choice} \hfill \textcolor{darkslate!70}{\small Topic: OSI and TCP/IP Models}

How many layers does the OSI reference model have?

\begin{optlist}
\item 4
\item 5
\item 7
\item 9
\end{optlist}
\begin{explainbox}{}
\textbf{Correct Answer: C}\\[0.3em]
\textbf{Concept:} The OSI model has seven layers: physical, data link, network, transport, session, presentation, and application.
\end{explainbox}
\end{questioncard}

\begin{questioncard}{2}{singlecol}
\qtypebadge{singlecol}{Single Choice} \hfill \textcolor{darkslate!70}{\small Topic: OSI and TCP/IP Models}

How many layers does the TCP/IP standard model have?

\begin{optlist}
\item 4
\item 5
\item 6
\item 7
\end{optlist}
\begin{explainbox}{}
\textbf{Correct Answer: A}\\[0.3em]
\textbf{Concept:} The TCP/IP standard model has four layers: network interface, internet, transport, and application.
\end{explainbox}
\end{questioncard}

\begin{questioncard}{3}{singlecol}
\qtypebadge{singlecol}{Single Choice} \hfill \textcolor{darkslate!70}{\small Topic: OSI and TCP/IP Models}

Which layer is responsible for end-to-end communication and flow control?

\begin{optlist}
\item Network layer
\item Transport layer
\item Data link layer
\item Physical layer
\end{optlist}
\begin{explainbox}{}
\textbf{Correct Answer: B}\\[0.3em]
\textbf{Concept:} The transport layer provides end-to-end communication, reliability, and flow control (e.g., TCP).
\end{explainbox}
\end{questioncard}

\begin{questioncard}{4}{singlecol}
\qtypebadge{singlecol}{Single Choice} \hfill \textcolor{darkslate!70}{\small Topic: OSI and TCP/IP Models}

Which layer handles logical addressing and routing?

\begin{optlist}
\item Transport layer
\item Network layer
\item Session layer
\item Presentation layer
\end{optlist}
\begin{explainbox}{}
\textbf{Correct Answer: B}\\[0.3em]
\textbf{Concept:} The network layer uses logical addresses (IP) and routes packets between networks.
\end{explainbox}
\end{questioncard}

\begin{questioncard}{5}{singlecol}
\qtypebadge{singlecol}{Single Choice} \hfill \textcolor{darkslate!70}{\small Topic: OSI and TCP/IP Models}

Which layer is responsible for framing, MAC addressing, and error detection on a local link?

\begin{optlist}
\item Network layer
\item Transport layer
\item Data link layer
\item Application layer
\end{optlist}
\begin{explainbox}{}
\textbf{Correct Answer: C}\\[0.3em]
\textbf{Concept:} The data link layer handles framing, MAC addresses, and local link error detection.
\end{explainbox}
\end{questioncard}

\section{Network Protocols}
\begin{questioncard}{1}{singlecol}
\qtypebadge{singlecol}{Single Choice} \hfill \textcolor{darkslate!70}{\small Topic: Network Protocols}

Which protocol uses ports 20 and 21 for file transfer?

\begin{optlist}
\item HTTP
\item FTP
\item SMTP
\item DNS
\end{optlist}
\begin{explainbox}{}
\textbf{Correct Answer: B}\\[0.3em]
\textbf{Concept:} FTP uses TCP port 21 for control and port 20 for data transfer.
\end{explainbox}
\end{questioncard}

\begin{questioncard}{2}{singlecol}
\qtypebadge{singlecol}{Single Choice} \hfill \textcolor{darkslate!70}{\small Topic: Network Protocols}

Which protocol resolves domain names to IP addresses?

\begin{optlist}
\item DHCP
\item DNS
\item SNMP
\item NTP
\end{optlist}
\begin{explainbox}{}
\textbf{Correct Answer: B}\\[0.3em]
\textbf{Concept:} DNS (Domain Name System) maps human-readable domain names to IP addresses.
\end{explainbox}
\end{questioncard}

\begin{questioncard}{3}{singlecol}
\qtypebadge{singlecol}{Single Choice} \hfill \textcolor{darkslate!70}{\small Topic: Network Protocols}

Which transport protocol is connection-oriented and reliable?

\begin{optlist}
\item UDP
\item TCP
\item ICMP
\item IP
\end{optlist}
\begin{explainbox}{}
\textbf{Correct Answer: B}\\[0.3em]
\textbf{Concept:} TCP provides connection-oriented, reliable, ordered delivery with flow and congestion control.
\end{explainbox}
\end{questioncard}

\begin{questioncard}{4}{singlecol}
\qtypebadge{singlecol}{Single Choice} \hfill \textcolor{darkslate!70}{\small Topic: Network Protocols}

Which transport protocol is connectionless and has low overhead?

\begin{optlist}
\item TCP
\item UDP
\item HTTP
\item FTP
\end{optlist}
\begin{explainbox}{}
\textbf{Correct Answer: B}\\[0.3em]
\textbf{Concept:} UDP is connectionless, offers low latency and overhead, but does not guarantee delivery.
\end{explainbox}
\end{questioncard}

\section{Network Devices}
\begin{questioncard}{1}{singlecol}
\qtypebadge{singlecol}{Single Choice} \hfill \textcolor{darkslate!70}{\small Topic: Network Devices}

Which device operates at Layer 2 and forwards frames based on MAC addresses?

\begin{optlist}
\item Router
\item Switch
\item Hub
\item Firewall
\end{optlist}
\begin{explainbox}{}
\textbf{Correct Answer: B}\\[0.3em]
\textbf{Concept:} A switch operates at the data link layer and forwards frames using MAC address tables.
\end{explainbox}
\end{questioncard}

\begin{questioncard}{2}{singlecol}
\qtypebadge{singlecol}{Single Choice} \hfill \textcolor{darkslate!70}{\small Topic: Network Devices}

Which device operates at Layer 3 and forwards packets based on IP addresses?

\begin{optlist}
\item Switch
\item Router
\item Hub
\item Bridge
\end{optlist}
\begin{explainbox}{}
\textbf{Correct Answer: B}\\[0.3em]
\textbf{Concept:} A router operates at the network layer and makes forwarding decisions based on IP routing tables.
\end{explainbox}
\end{questioncard}

\begin{questioncard}{3}{singlecol}
\qtypebadge{singlecol}{Single Choice} \hfill \textcolor{darkslate!70}{\small Topic: Network Devices}

Which device simply repeats signals to all connected ports?

\begin{optlist}
\item Switch
\item Router
\item Hub
\item Firewall
\end{optlist}
\begin{explainbox}{}
\textbf{Correct Answer: C}\\[0.3em]
\textbf{Concept:} A hub is a Layer 1 device that broadcasts signals to all connected devices.
\end{explainbox}
\end{questioncard}

\section{Enterprise Security Threats}
\begin{questioncard}{1}{singlecol}
\qtypebadge{singlecol}{Single Choice} \hfill \textcolor{darkslate!70}{\small Topic: Enterprise Security Threats}

Which attack sends an overwhelming amount of traffic from many sources?

\begin{optlist}
\item DoS
\item DDoS
\item Phishing
\item Spoofing
\end{optlist}
\begin{explainbox}{}
\textbf{Correct Answer: B}\\[0.3em]
\textbf{Concept:} DDoS uses multiple distributed sources to flood a target.
\end{explainbox}
\end{questioncard}

\begin{questioncard}{2}{singlecol}
\qtypebadge{singlecol}{Single Choice} \hfill \textcolor{darkslate!70}{\small Topic: Enterprise Security Threats}

Which attack falsifies the source address of packets?

\begin{optlist}
\item Phishing
\item Spoofing
\item Sniffing
\item Spamming
\end{optlist}
\begin{explainbox}{}
\textbf{Correct Answer: B}\\[0.3em]
\textbf{Concept:} Spoofing impersonates a trusted source by falsifying addresses or identities.
\end{explainbox}
\end{questioncard}

\begin{questioncard}{3}{singlecol}
\qtypebadge{singlecol}{Single Choice} \hfill \textcolor{darkslate!70}{\small Topic: Enterprise Security Threats}

Which type of malware replicates itself across networks without user action?

\begin{optlist}
\item Virus
\item Worm
\item Trojan
\item Spyware
\end{optlist}
\begin{explainbox}{}
\textbf{Correct Answer: B}\\[0.3em]
\textbf{Concept:} Worms can self-replicate and spread across networks, often without user interaction.
\end{explainbox}
\end{questioncard}

\begin{questioncard}{4}{singlecol}
\qtypebadge{singlecol}{Single Choice} \hfill \textcolor{darkslate!70}{\small Topic: Enterprise Security Threats}

Which malware disguises itself as legitimate software?

\begin{optlist}
\item Virus
\item Worm
\item Trojan
\item Adware
\end{optlist}
\begin{explainbox}{}
\textbf{Correct Answer: C}\\[0.3em]
\textbf{Concept:} A Trojan horse masquerades as legitimate software to trick users into executing it.
\end{explainbox}
\end{questioncard}

\begin{questioncard}{5}{singlecol}
\qtypebadge{singlecol}{Single Choice} \hfill \textcolor{darkslate!70}{\small Topic: Enterprise Security Threats}

Ransomware primarily does what?

\begin{optlist}
\item Steals credentials
\item Encrypts data and demands payment
\item Deletes logs
\item Monitors traffic
\end{optlist}
\begin{explainbox}{}
\textbf{Correct Answer: B}\\[0.3em]
\textbf{Concept:} Ransomware encrypts files and demands a ransom for decryption.
\end{explainbox}
\end{questioncard}

\begin{questioncard}{6}{singlecol}
\qtypebadge{singlecol}{Single Choice} \hfill \textcolor{darkslate!70}{\small Topic: Enterprise Security Threats}

Which attack intercepts communications between two parties without their knowledge?

\begin{optlist}
\item Man-in-the-middle
\item DoS
\item Phishing
\item Tailgating
\end{optlist}
\begin{explainbox}{}
\textbf{Correct Answer: A}\\[0.3em]
\textbf{Concept:} A man-in-the-middle attack intercepts and possibly alters communication between parties.
\end{explainbox}
\end{questioncard}

\begin{questioncard}{7}{singlecol}
\qtypebadge{singlecol}{Single Choice} \hfill \textcolor{darkslate!70}{\small Topic: Enterprise Security Threats}

Which protocol can help prevent man-in-the-middle attacks by verifying server identity?

\begin{optlist}
\item HTTP
\item HTTPS
\item FTP
\item Telnet
\end{optlist}
\begin{explainbox}{}
\textbf{Correct Answer: B}\\[0.3em]
\textbf{Concept:} HTTPS uses TLS/SSL to encrypt traffic and authenticate the server.
\end{explainbox}
\end{questioncard}

\begin{questioncard}{8}{singlecol}
\qtypebadge{singlecol}{Single Choice} \hfill \textcolor{darkslate!70}{\small Topic: Enterprise Security Threats}

Which wireless security protocol is more secure than WEP?

\begin{optlist}
\item WPA2/WPA3
\item WEP
\item WPA (TKIP only)
\item Open network
\end{optlist}
\begin{explainbox}{}
\textbf{Correct Answer: A}\\[0.3em]
\textbf{Concept:} WPA2 and WPA3 provide stronger encryption and authentication than WEP.
\end{explainbox}
\end{questioncard}

\begin{questioncard}{9}{singlecol}
\qtypebadge{singlecol}{Single Choice} \hfill \textcolor{darkslate!70}{\small Topic: Enterprise Security Threats}

What is a DMZ typically used for?

\begin{optlist}
\item Hosting internal databases only
\item Hosting public-facing servers
\item User workstations
\item Printer sharing
\end{optlist}
\begin{explainbox}{}
\textbf{Correct Answer: B}\\[0.3em]
\textbf{Concept:} A DMZ hosts public-facing services between the Internet and internal network.
\end{explainbox}
\end{questioncard}

\begin{questioncard}{10}{singlecol}
\qtypebadge{singlecol}{Single Choice} \hfill \textcolor{darkslate!70}{\small Topic: Enterprise Security Threats}

Which security zone usually has the highest priority in Huawei firewalls?

\begin{optlist}
\item Untrust
\item Trust
\item DMZ
\item Local
\end{optlist}
\begin{explainbox}{}
\textbf{Correct Answer: D}\\[0.3em]
\textbf{Concept:} The Local zone (firewall itself) has priority 100, the highest default priority.
\end{explainbox}
\end{questioncard}

\begin{questioncard}{11}{singlecol}
\qtypebadge{singlecol}{Single Choice} \hfill \textcolor{darkslate!70}{\small Topic: Enterprise Security Threats}

Which practice limits the damage from compromised accounts?

\begin{optlist}
\item Least privilege
\item Sharing admin passwords
\item Disabling logs
\item Open shares
\end{optlist}
\begin{explainbox}{}
\textbf{Correct Answer: A}\\[0.3em]
\textbf{Concept:} Least privilege gives users only the access they need, limiting blast radius.
\end{explainbox}
\end{questioncard}

\begin{questioncard}{12}{singlecol}
\qtypebadge{singlecol}{Single Choice} \hfill \textcolor{darkslate!70}{\small Topic: Enterprise Security Threats}

What is the purpose of host-based intrusion prevention?

\begin{optlist}
\item Route packets
\item Detect and block malicious host activity
\item Assign IP addresses
\item Synchronize time
\end{optlist}
\begin{explainbox}{}
\textbf{Correct Answer: B}\\[0.3em]
\textbf{Concept:} Host-based IPS monitors and blocks suspicious activity on an endpoint.
\end{explainbox}
\end{questioncard}

\begin{questioncard}{13}{singlecol}
\qtypebadge{singlecol}{Single Choice} \hfill \textcolor{darkslate!70}{\small Topic: Enterprise Security Threats}

Which technology collects and analyzes security logs centrally?

\begin{optlist}
\item SIEM
\item SNMP
\item DHCP
\item ARP
\end{optlist}
\begin{explainbox}{}
\textbf{Correct Answer: A}\\[0.3em]
\textbf{Concept:} SIEM centralizes log collection and analysis.
\end{explainbox}
\end{questioncard}

\begin{questioncard}{14}{singlecol}
\qtypebadge{singlecol}{Single Choice} \hfill \textcolor{darkslate!70}{\small Topic: Enterprise Security Threats}

What is a key benefit of centralized security management?

\begin{optlist}
\item Reduced visibility
\item Unified policy and monitoring
\item More complex passwords
\item Faster malware spread
\end{optlist}
\begin{explainbox}{}
\textbf{Correct Answer: B}\\[0.3em]
\textbf{Concept:} Centralized management provides unified visibility, policy control, and monitoring.
\end{explainbox}
\end{questioncard}

\begin{questioncard}{15}{singlecol}
\qtypebadge{singlecol}{Single Choice} \hfill \textcolor{darkslate!70}{\small Topic: Enterprise Security Threats}

What is an insider threat?

\begin{optlist}
\item An attack from outside the network
\item A security risk originating from within the organization
\item A hardware vendor
\item A natural disaster
\end{optlist}
\begin{explainbox}{}
\textbf{Correct Answer: B}\\[0.3em]
\textbf{Concept:} Insider threats come from employees, contractors, or partners with internal access.
\end{explainbox}
\end{questioncard}

\begin{questioncard}{16}{singlecol}
\qtypebadge{singlecol}{Single Choice} \hfill \textcolor{darkslate!70}{\small Topic: Enterprise Security Threats}

Which attack sends oversized ICMP packets to crash a target?

\begin{optlist}
\item Ping of death
\item Ping sweep
\item ARP spoofing
\item DNS hijacking
\end{optlist}
\begin{explainbox}{}
\textbf{Correct Answer: A}\\[0.3em]
\textbf{Concept:} Ping of death sends ICMP packets larger than the maximum allowed size.
\end{explainbox}
\end{questioncard}

\begin{questioncard}{17}{singlecol}
\qtypebadge{singlecol}{Single Choice} \hfill \textcolor{darkslate!70}{\small Topic: Enterprise Security Threats}

Which attack uses broadcast pings to amplify traffic?

\begin{optlist}
\item Smurf attack
\item SYN flood
\item Fraggle
\item Session hijacking
\end{optlist}
\begin{explainbox}{}
\textbf{Correct Answer: A}\\[0.3em]
\textbf{Concept:} Smurf attacks send ICMP echo requests to broadcast addresses with a spoofed source.
\end{explainbox}
\end{questioncard}

\begin{questioncard}{18}{singlecol}
\qtypebadge{singlecol}{Single Choice} \hfill \textcolor{darkslate!70}{\small Topic: Enterprise Security Threats}

Which attack intercepts ARP messages to redirect traffic?

\begin{optlist}
\item ARP spoofing
\item DNS spoofing
\item IP spoofing
\item MAC flooding
\end{optlist}
\begin{explainbox}{}
\textbf{Correct Answer: A}\\[0.3em]
\textbf{Concept:} ARP spoofing sends forged ARP replies to associate an attacker MAC with another IP.
\end{explainbox}
\end{questioncard}

\begin{questioncard}{19}{singlecol}
\qtypebadge{singlecol}{Single Choice} \hfill \textcolor{darkslate!70}{\small Topic: Enterprise Security Threats}

A SYN flood is a type of:

\begin{optlist}
\item Denial of Service attack
\item Phishing attack
\item Data exfiltration
\item Physical attack
\end{optlist}
\begin{explainbox}{}
\textbf{Correct Answer: A}\\[0.3em]
\textbf{Concept:} SYN floods exhaust server connection resources with half-open TCP connections.
\end{explainbox}
\end{questioncard}

\begin{questioncard}{20}{singlecol}
\qtypebadge{singlecol}{Single Choice} \hfill \textcolor{darkslate!70}{\small Topic: Enterprise Security Threats}

ARP spoofing is used to:

\begin{optlist}
\item Intercept traffic on a local network
\item Encrypt files
\item Speed up routing
\item Authenticate users
\end{optlist}
\begin{explainbox}{}
\textbf{Correct Answer: A}\\[0.3em]
\textbf{Concept:} Attackers use forged ARP messages to redirect local traffic through themselves.
\end{explainbox}
\end{questioncard}

\begin{questioncard}{21}{singlecol}
\qtypebadge{singlecol}{Single Choice} \hfill \textcolor{darkslate!70}{\small Topic: Enterprise Security Threats}

The principle of least privilege aims to:

\begin{optlist}
\item Give everyone admin rights
\item Limit user access to the minimum necessary
\item Disable all logging
\item Remove firewalls
\end{optlist}
\begin{explainbox}{}
\textbf{Correct Answer: B}\\[0.3em]
\textbf{Concept:} Least privilege reduces risk by restricting access to only what is needed.
\end{explainbox}
\end{questioncard}

\begin{questioncard}{22}{singlecol}
\qtypebadge{singlecol}{Single Choice} \hfill \textcolor{darkslate!70}{\small Topic: Enterprise Security Threats}

A SIEM system is primarily used for:

\begin{optlist}
\item Centralized log analysis and threat detection
\item Assigning IP addresses
\item Switching frames
\item Manufacturing cables
\end{optlist}
\begin{explainbox}{}
\textbf{Correct Answer: A}\\[0.3em]
\textbf{Concept:} SIEM aggregates and analyzes security events and logs.
\end{explainbox}
\end{questioncard}

\begin{questioncard}{23}{singlecol}
\qtypebadge{singlecol}{Single Choice} \hfill \textcolor{darkslate!70}{\small Topic: Enterprise Security Threats}

Which concept is described as: monitors traffic without altering it?

\begin{optlist}
\item SIEM
\item Passive attack
\item Worm
\item SYN flood
\end{optlist}
\begin{explainbox}{}
\textbf{Correct Answer: B}\\[0.3em]
\textbf{Concept:} Passive attack refers to the concept where monitors traffic without altering it.
\end{explainbox}
\end{questioncard}

\begin{questioncard}{24}{singlecol}
\qtypebadge{singlecol}{Single Choice} \hfill \textcolor{darkslate!70}{\small Topic: Enterprise Security Threats}

Which concept is described as: modifies or disrupts systems or data?

\begin{optlist}
\item Ransomware
\item Passive attack
\item Patch management
\item Active attack
\end{optlist}
\begin{explainbox}{}
\textbf{Correct Answer: D}\\[0.3em]
\textbf{Concept:} Active attack refers to the concept where modifies or disrupts systems or data.
\end{explainbox}
\end{questioncard}

\begin{questioncard}{25}{singlecol}
\qtypebadge{singlecol}{Single Choice} \hfill \textcolor{darkslate!70}{\small Topic: Enterprise Security Threats}

Which concept is described as: includes viruses, worms, Trojans, and ransomware?

\begin{optlist}
\item Smurf attack
\item Passive attack
\item Malware
\item Worm
\end{optlist}
\begin{explainbox}{}
\textbf{Correct Answer: C}\\[0.3em]
\textbf{Concept:} Malware refers to the concept where includes viruses, worms, Trojans, and ransomware.
\end{explainbox}
\end{questioncard}

\begin{questioncard}{26}{singlecol}
\qtypebadge{singlecol}{Single Choice} \hfill \textcolor{darkslate!70}{\small Topic: Enterprise Security Threats}

Which concept is described as: requires a host program to spread?

\begin{optlist}
\item Worm
\item DDoS attack
\item Virus
\item Insider threat
\end{optlist}
\begin{explainbox}{}
\textbf{Correct Answer: C}\\[0.3em]
\textbf{Concept:} Virus refers to the concept where requires a host program to spread.
\end{explainbox}
\end{questioncard}

\begin{questioncard}{27}{singlecol}
\qtypebadge{singlecol}{Single Choice} \hfill \textcolor{darkslate!70}{\small Topic: Enterprise Security Threats}

Which concept is described as: can self-replicate across networks?

\begin{optlist}
\item Worm
\item Ping of death
\item Virus
\item Passive attack
\end{optlist}
\begin{explainbox}{}
\textbf{Correct Answer: A}\\[0.3em]
\textbf{Concept:} Worm refers to the concept where can self-replicate across networks.
\end{explainbox}
\end{questioncard}

\begin{questioncard}{28}{singlecol}
\qtypebadge{singlecol}{Single Choice} \hfill \textcolor{darkslate!70}{\small Topic: Enterprise Security Threats}

Which concept is described as: disguises itself as legitimate software?

\begin{optlist}
\item Trojan
\item Active attack
\item Worm
\item Man-in-the-middle attack
\end{optlist}
\begin{explainbox}{}
\textbf{Correct Answer: A}\\[0.3em]
\textbf{Concept:} Trojan refers to the concept where disguises itself as legitimate software.
\end{explainbox}
\end{questioncard}

\begin{questioncard}{29}{singlecol}
\qtypebadge{singlecol}{Single Choice} \hfill \textcolor{darkslate!70}{\small Topic: Enterprise Security Threats}

Which concept is described as: encrypts data and demands payment?

\begin{optlist}
\item SIEM
\item Smurf attack
\item Social engineering
\item Ransomware
\end{optlist}
\begin{explainbox}{}
\textbf{Correct Answer: D}\\[0.3em]
\textbf{Concept:} Ransomware refers to the concept where encrypts data and demands payment.
\end{explainbox}
\end{questioncard}

\begin{questioncard}{30}{singlecol}
\qtypebadge{singlecol}{Single Choice} \hfill \textcolor{darkslate!70}{\small Topic: Enterprise Security Threats}

Which concept is described as: uses deception to steal information?

\begin{optlist}
\item Worm
\item Patch management
\item Phishing
\item Virus
\end{optlist}
\begin{explainbox}{}
\textbf{Correct Answer: C}\\[0.3em]
\textbf{Concept:} Phishing refers to the concept where uses deception to steal information.
\end{explainbox}
\end{questioncard}

\begin{questioncard}{31}{singlecol}
\qtypebadge{singlecol}{Single Choice} \hfill \textcolor{darkslate!70}{\small Topic: Enterprise Security Threats}

Which concept is described as: manipulates people into breaking security rules?

\begin{optlist}
\item Patch management
\item Social engineering
\item Worm
\item ARP spoofing
\end{optlist}
\begin{explainbox}{}
\textbf{Correct Answer: B}\\[0.3em]
\textbf{Concept:} Social engineering refers to the concept where manipulates people into breaking security rules.
\end{explainbox}
\end{questioncard}

\begin{questioncard}{32}{singlecol}
\qtypebadge{singlecol}{Single Choice} \hfill \textcolor{darkslate!70}{\small Topic: Enterprise Security Threats}

Which concept is described as: aims to make a service unavailable?

\begin{optlist}
\item SIEM
\item Worm
\item Ransomware
\item DoS attack
\end{optlist}
\begin{explainbox}{}
\textbf{Correct Answer: D}\\[0.3em]
\textbf{Concept:} DoS attack refers to the concept where aims to make a service unavailable.
\end{explainbox}
\end{questioncard}

\begin{questioncard}{33}{singlecol}
\qtypebadge{singlecol}{Single Choice} \hfill \textcolor{darkslate!70}{\small Topic: Enterprise Security Threats}

Which concept is described as: uses many distributed sources to flood a target?

\begin{optlist}
\item Virus
\item DDoS attack
\item ARP spoofing
\item Man-in-the-middle attack
\end{optlist}
\begin{explainbox}{}
\textbf{Correct Answer: B}\\[0.3em]
\textbf{Concept:} DDoS attack refers to the concept where uses many distributed sources to flood a target.
\end{explainbox}
\end{questioncard}

\begin{questioncard}{34}{singlecol}
\qtypebadge{singlecol}{Single Choice} \hfill \textcolor{darkslate!70}{\small Topic: Enterprise Security Threats}

Which concept is described as: intercepts communication between two parties?

\begin{optlist}
\item ARP spoofing
\item Patch management
\item Man-in-the-middle attack
\item Social engineering
\end{optlist}
\begin{explainbox}{}
\textbf{Correct Answer: C}\\[0.3em]
\textbf{Concept:} Man-in-the-middle attack refers to the concept where intercepts communication between two parties.
\end{explainbox}
\end{questioncard}

\begin{questioncard}{35}{singlecol}
\qtypebadge{singlecol}{Single Choice} \hfill \textcolor{darkslate!70}{\small Topic: Enterprise Security Threats}

Which concept is described as: forges ARP replies to redirect traffic?

\begin{optlist}
\item ARP spoofing
\item Least privilege
\item DDoS attack
\item Malware
\end{optlist}
\begin{explainbox}{}
\textbf{Correct Answer: A}\\[0.3em]
\textbf{Concept:} ARP spoofing refers to the concept where forges ARP replies to redirect traffic.
\end{explainbox}
\end{questioncard}

\begin{questioncard}{36}{singlecol}
\qtypebadge{singlecol}{Single Choice} \hfill \textcolor{darkslate!70}{\small Topic: Enterprise Security Threats}

Which concept is described as: sends many half-open TCP connections?

\begin{optlist}
\item Ransomware
\item DMZ
\item Trojan
\item SYN flood
\end{optlist}
\begin{explainbox}{}
\textbf{Correct Answer: D}\\[0.3em]
\textbf{Concept:} SYN flood refers to the concept where sends many half-open TCP connections.
\end{explainbox}
\end{questioncard}

\begin{questioncard}{37}{singlecol}
\qtypebadge{singlecol}{Single Choice} \hfill \textcolor{darkslate!70}{\small Topic: Enterprise Security Threats}

Which concept is described as: sends oversized ICMP packets?

\begin{optlist}
\item Ping of death
\item Phishing
\item SIEM
\item Malware
\end{optlist}
\begin{explainbox}{}
\textbf{Correct Answer: A}\\[0.3em]
\textbf{Concept:} Ping of death refers to the concept where sends oversized ICMP packets.
\end{explainbox}
\end{questioncard}

\begin{questioncard}{38}{singlecol}
\qtypebadge{singlecol}{Single Choice} \hfill \textcolor{darkslate!70}{\small Topic: Enterprise Security Threats}

Which concept is described as: uses broadcast pings to amplify traffic?

\begin{optlist}
\item Social engineering
\item ARP spoofing
\item Trojan
\item Smurf attack
\end{optlist}
\begin{explainbox}{}
\textbf{Correct Answer: D}\\[0.3em]
\textbf{Concept:} Smurf attack refers to the concept where uses broadcast pings to amplify traffic.
\end{explainbox}
\end{questioncard}

\begin{questioncard}{39}{singlecol}
\qtypebadge{singlecol}{Single Choice} \hfill \textcolor{darkslate!70}{\small Topic: Enterprise Security Threats}

Which concept is described as: limits user access to the minimum necessary?

\begin{optlist}
\item Passive attack
\item Social engineering
\item Worm
\item Least privilege
\end{optlist}
\begin{explainbox}{}
\textbf{Correct Answer: D}\\[0.3em]
\textbf{Concept:} Least privilege refers to the concept where limits user access to the minimum necessary.
\end{explainbox}
\end{questioncard}

\begin{questioncard}{40}{singlecol}
\qtypebadge{singlecol}{Single Choice} \hfill \textcolor{darkslate!70}{\small Topic: Enterprise Security Threats}

Which concept is described as: eliminates known vulnerabilities?

\begin{optlist}
\item Phishing
\item SYN flood
\item Patch management
\item Trojan
\end{optlist}
\begin{explainbox}{}
\textbf{Correct Answer: C}\\[0.3em]
\textbf{Concept:} Patch management refers to the concept where eliminates known vulnerabilities.
\end{explainbox}
\end{questioncard}

\begin{questioncard}{41}{singlecol}
\qtypebadge{singlecol}{Single Choice} \hfill \textcolor{darkslate!70}{\small Topic: Enterprise Security Threats}

Which concept is described as: centralizes log collection and analysis?

\begin{optlist}
\item Malware
\item Active attack
\item Phishing
\item SIEM
\end{optlist}
\begin{explainbox}{}
\textbf{Correct Answer: D}\\[0.3em]
\textbf{Concept:} SIEM refers to the concept where centralizes log collection and analysis.
\end{explainbox}
\end{questioncard}

\begin{questioncard}{42}{singlecol}
\qtypebadge{singlecol}{Single Choice} \hfill \textcolor{darkslate!70}{\small Topic: Enterprise Security Threats}

Which concept is described as: hosts public-facing servers between Internet and internal network?

\begin{optlist}
\item DMZ
\item DDoS attack
\item DoS attack
\item Patch management
\end{optlist}
\begin{explainbox}{}
\textbf{Correct Answer: A}\\[0.3em]
\textbf{Concept:} DMZ refers to the concept where hosts public-facing servers between Internet and internal network.
\end{explainbox}
\end{questioncard}

\begin{questioncard}{43}{singlecol}
\qtypebadge{singlecol}{Single Choice} \hfill \textcolor{darkslate!70}{\small Topic: Enterprise Security Threats}

Which concept is described as: originates from within the organization?

\begin{optlist}
\item Insider threat
\item Worm
\item Man-in-the-middle attack
\item Active attack
\end{optlist}
\begin{explainbox}{}
\textbf{Correct Answer: A}\\[0.3em]
\textbf{Concept:} Insider threat refers to the concept where originates from within the organization.
\end{explainbox}
\end{questioncard}

\section{Communication Network Security}
\begin{questioncard}{1}{singlecol}
\qtypebadge{singlecol}{Single Choice} \hfill \textcolor{darkslate!70}{\small Topic: Communication Network Security}

Which threat involves intercepting data traveling across a network?

\begin{optlist}
\item Eavesdropping
\item DoS
\item Phishing
\item Tailgating
\end{optlist}
\begin{explainbox}{}
\textbf{Correct Answer: A}\\[0.3em]
\textbf{Concept:} Eavesdropping intercepts network communications to steal information.
\end{explainbox}
\end{questioncard}

\section{Zone Border Security}
\begin{questioncard}{1}{singlecol}
\qtypebadge{singlecol}{Single Choice} \hfill \textcolor{darkslate!70}{\small Topic: Zone Border Security}

What is the main purpose of a security zone?

\begin{optlist}
\item Increase Internet speed
\item Group interfaces with similar security levels
\item Replace antivirus
\item Disable routing
\end{optlist}
\begin{explainbox}{}
\textbf{Correct Answer: B}\\[0.3em]
\textbf{Concept:} Security zones group network interfaces with similar trust levels to enforce inter-zone policies.
\end{explainbox}
\end{questioncard}

\begin{questioncard}{2}{singlecol}
\qtypebadge{singlecol}{Single Choice} \hfill \textcolor{darkslate!70}{\small Topic: Zone Border Security}

Which device is typically deployed at a zone border to enforce access control?

\begin{optlist}
\item Hub
\item Firewall
\item Repeater
\item Bridge
\end{optlist}
\begin{explainbox}{}
\textbf{Correct Answer: B}\\[0.3em]
\textbf{Concept:} Firewalls are deployed at network boundaries to enforce security policies between zones.
\end{explainbox}
\end{questioncard}

\section{Computing Environment Security}
\begin{questioncard}{1}{singlecol}
\qtypebadge{singlecol}{Single Choice} \hfill \textcolor{darkslate!70}{\small Topic: Computing Environment Security}

Which control helps protect endpoints from malware?

\begin{optlist}
\item Antivirus/EDR
\item Physical locks only
\item Hub replacement
\item Disabling DNS
\end{optlist}
\begin{explainbox}{}
\textbf{Correct Answer: A}\\[0.3em]
\textbf{Concept:} Antivirus and endpoint detection and response (EDR) protect hosts from malware and suspicious behavior.
\end{explainbox}
\end{questioncard}

\section{Firewall Basics}
\begin{questioncard}{1}{singlecol}
\qtypebadge{singlecol}{Single Choice} \hfill \textcolor{darkslate!70}{\small Topic: Firewall Basics}

Which firewall feature inspects application-layer data to detect attacks?

\begin{optlist}
\item Packet filtering
\item Application inspection
\item Hub forwarding
\item NAT
\end{optlist}
\begin{explainbox}{}
\textbf{Correct Answer: B}\\[0.3em]
\textbf{Concept:} Application inspection examines payload content for malicious patterns.
\end{explainbox}
\end{questioncard}

\begin{questioncard}{2}{singlecol}
\qtypebadge{singlecol}{Single Choice} \hfill \textcolor{darkslate!70}{\small Topic: Firewall Basics}

Which Huawei firewall command mode is used for system-level configuration?

\begin{optlist}
\item User view
\item System view
\item Interface view
\item Policy view
\end{optlist}
\begin{explainbox}{}
\textbf{Correct Answer: B}\\[0.3em]
\textbf{Concept:} System view allows global-level configuration on Huawei devices.
\end{explainbox}
\end{questioncard}

\begin{questioncard}{3}{singlecol}
\qtypebadge{singlecol}{Single Choice} \hfill \textcolor{darkslate!70}{\small Topic: Firewall Basics}

Which view is entered after the system-view command?

\begin{optlist}
\item User view
\item System view
\item Interface view
\item None
\end{optlist}
\begin{explainbox}{}
\textbf{Correct Answer: B}\\[0.3em]
\textbf{Concept:} The system-view command enters system view from user view.
\end{explainbox}
\end{questioncard}

\begin{questioncard}{4}{singlecol}
\qtypebadge{singlecol}{Single Choice} \hfill \textcolor{darkslate!70}{\small Topic: Firewall Basics}

What is the purpose of a security profile in NGFW?

\begin{optlist}
\item Physical installation
\item Applying advanced security features such as IPS/AV/URL filtering
\item Layer 1 forwarding
\item DHCP relay
\end{optlist}
\begin{explainbox}{}
\textbf{Correct Answer: B}\\[0.3em]
\textbf{Concept:} Security profiles group advanced inspection features applied to policies.
\end{explainbox}
\end{questioncard}

\begin{questioncard}{5}{singlecol}
\qtypebadge{singlecol}{Single Choice} \hfill \textcolor{darkslate!70}{\small Topic: Firewall Basics}

Which firewall generation combines traditional firewall functions with application awareness and intrusion prevention?

\begin{optlist}
\item Packet-filtering firewall
\item Next-generation firewall
\item Stateful firewall
\item Proxy firewall
\end{optlist}
\begin{explainbox}{}
\textbf{Correct Answer: B}\\[0.3em]
\textbf{Concept:} Next-generation firewalls integrate application awareness, IPS, antivirus, and user identity.
\end{explainbox}
\end{questioncard}

\begin{questioncard}{6}{singlecol}
\qtypebadge{singlecol}{Single Choice} \hfill \textcolor{darkslate!70}{\small Topic: Firewall Basics}

What is the purpose of a default security policy deny action?

\begin{optlist}
\item Improve performance
\item Ensure least-privilege access
\item Enable logging
\item Allow management
\end{optlist}
\begin{explainbox}{}
\textbf{Correct Answer: B}\\[0.3em]
\textbf{Concept:} Default deny enforces least privilege by blocking unspecified traffic.
\end{explainbox}
\end{questioncard}

\begin{questioncard}{7}{singlecol}
\qtypebadge{singlecol}{Single Choice} \hfill \textcolor{darkslate!70}{\small Topic: Firewall Basics}

Which zone would typically contain public web servers?

\begin{optlist}
\item Trust
\item Untrust
\item DMZ
\item Local
\end{optlist}
\begin{explainbox}{}
\textbf{Correct Answer: C}\\[0.3em]
\textbf{Concept:} Public-facing servers are placed in the DMZ for isolation.
\end{explainbox}
\end{questioncard}

\begin{questioncard}{8}{singlecol}
\qtypebadge{singlecol}{Single Choice} \hfill \textcolor{darkslate!70}{\small Topic: Firewall Basics}

What does policy hit counting help with?

\begin{optlist}
\item Routing optimization
\item Verifying policy usage and troubleshooting
\item MAC learning
\item DHCP assignment
\end{optlist}
\begin{explainbox}{}
\textbf{Correct Answer: B}\\[0.3em]
\textbf{Concept:} Hit counts show which policies are triggered, aiding auditing and troubleshooting.
\end{explainbox}
\end{questioncard}

\begin{questioncard}{9}{singlecol}
\qtypebadge{singlecol}{Single Choice} \hfill \textcolor{darkslate!70}{\small Topic: Firewall Basics}

Which of the following is NOT a firewall deployment mode?

\begin{optlist}
\item Routing mode
\item Transparent mode
\item Hybrid mode
\item Broadcast mode
\end{optlist}
\begin{explainbox}{}
\textbf{Correct Answer: D}\\[0.3em]
\textbf{Concept:} Firewalls deploy in routing, transparent (bridge), or hybrid modes; broadcast mode is not standard.
\end{explainbox}
\end{questioncard}

\begin{questioncard}{10}{singlecol}
\qtypebadge{singlecol}{Single Choice} \hfill \textcolor{darkslate!70}{\small Topic: Firewall Basics}

What happens when traffic moves from a higher-priority zone to a lower-priority zone?

\begin{optlist}
\item Inbound traffic
\item Outbound traffic
\item No traffic allowed
\item Always dropped
\end{optlist}
\begin{explainbox}{}
\textbf{Correct Answer: B}\\[0.3em]
\textbf{Concept:} Traffic from a higher-priority zone to a lower-priority zone is considered outbound.
\end{explainbox}
\end{questioncard}

\begin{questioncard}{11}{singlecol}
\qtypebadge{singlecol}{Single Choice} \hfill \textcolor{darkslate!70}{\small Topic: Firewall Basics}

What is the priority range for custom security zones on Huawei firewalls?

\begin{optlist}
\item 0-50
\item 1-100
\item 1-99 excluding defaults
\item 0-100
\end{optlist}
\begin{explainbox}{}
\textbf{Correct Answer: C}\\[0.3em]
\textbf{Concept:} Custom zone priorities range from 1 to 99, avoiding reserved default values.
\end{explainbox}
\end{questioncard}

\begin{questioncard}{12}{singlecol}
\qtypebadge{singlecol}{Single Choice} \hfill \textcolor{darkslate!70}{\small Topic: Firewall Basics}

Which action in a security policy can log permitted or denied traffic?

\begin{optlist}
\item Permit
\item Deny
\item Both permit and deny
\item Neither
\end{optlist}
\begin{explainbox}{}
\textbf{Correct Answer: C}\\[0.3em]
\textbf{Concept:} Security policies can enable logging for both permit and deny actions.
\end{explainbox}
\end{questioncard}

\begin{questioncard}{13}{singlecol}
\qtypebadge{singlecol}{Single Choice} \hfill \textcolor{darkslate!70}{\small Topic: Firewall Basics}

Which security policy action is typically applied to suspicious traffic for further inspection?

\begin{optlist}
\item Permit
\item Deny
\item Block
\item Redirect
\end{optlist}
\begin{explainbox}{}
\textbf{Correct Answer: B}\\[0.3em]
\textbf{Concept:} Suspicious traffic is commonly denied or blocked by policy.
\end{explainbox}
\end{questioncard}

\begin{questioncard}{14}{singlecol}
\qtypebadge{singlecol}{Single Choice} \hfill \textcolor{darkslate!70}{\small Topic: Firewall Basics}

When a return packet arrives, what does a stateful firewall check first?

\begin{optlist}
\item Routing table
\item Session table
\item MAC table
\item ARP table
\end{optlist}
\begin{explainbox}{}
\textbf{Correct Answer: B}\\[0.3em]
\textbf{Concept:} A stateful firewall first checks the session table to find a matching connection.
\end{explainbox}
\end{questioncard}

\begin{questioncard}{15}{singlecol}
\qtypebadge{singlecol}{Single Choice} \hfill \textcolor{darkslate!70}{\small Topic: Firewall Basics}

What happens if a return packet does not match any session table entry?

\begin{optlist}
\item Forwarded normally
\item Dropped or inspected against policies
\item Always permitted
\item Sent to DMZ
\end{optlist}
\begin{explainbox}{}
\textbf{Correct Answer: B}\\[0.3em]
\textbf{Concept:} Without a matching session, the firewall treats the packet as a new connection and applies policies.
\end{explainbox}
\end{questioncard}

\begin{questioncard}{16}{singlecol}
\qtypebadge{singlecol}{Single Choice} \hfill \textcolor{darkslate!70}{\small Topic: Firewall Basics}

Which protocol commonly requires ASPF support due to dynamic port negotiation?

\begin{optlist}
\item HTTP
\item FTP
\item DNS
\item NTP
\end{optlist}
\begin{explainbox}{}
\textbf{Correct Answer: B}\\[0.3em]
\textbf{Concept:} FTP negotiates dynamic data ports in its control channel, requiring ASPF.
\end{explainbox}
\end{questioncard}

\begin{questioncard}{17}{singlecol}
\qtypebadge{singlecol}{Single Choice} \hfill \textcolor{darkslate!70}{\small Topic: Firewall Basics}

ASPF primarily inspects traffic at which layer?

\begin{optlist}
\item Network layer only
\item Application layer
\item Physical layer
\item Data link layer
\end{optlist}
\begin{explainbox}{}
\textbf{Correct Answer: B}\\[0.3em]
\textbf{Concept:} ASPF performs application-layer inspection to identify negotiated connections.
\end{explainbox}
\end{questioncard}

\begin{questioncard}{18}{singlecol}
\qtypebadge{singlecol}{Single Choice} \hfill \textcolor{darkslate!70}{\small Topic: Firewall Basics}

In Huawei firewalls, the default priority of Trust is:

\begin{optlist}
\item 5
\item 50
\item 85
\item 100
\end{optlist}
\begin{explainbox}{}
\textbf{Correct Answer: C}\\[0.3em]
\textbf{Concept:} Trust zone default priority is 85.
\end{explainbox}
\end{questioncard}

\begin{questioncard}{19}{singlecol}
\qtypebadge{singlecol}{Single Choice} \hfill \textcolor{darkslate!70}{\small Topic: Firewall Basics}

Traffic from a low-priority zone to a high-priority zone is considered:

\begin{optlist}
\item Outbound
\item Inbound
\item Blocked by default
\item Always permitted
\end{optlist}
\begin{explainbox}{}
\textbf{Correct Answer: B}\\[0.3em]
\textbf{Concept:} Traffic entering a higher-priority zone from a lower-priority zone is inbound.
\end{explainbox}
\end{questioncard}

\begin{questioncard}{20}{singlecol}
\qtypebadge{singlecol}{Single Choice} \hfill \textcolor{darkslate!70}{\small Topic: Firewall Basics}

The default security policy action is:

\begin{optlist}
\item Permit
\item Deny
\item Log
\item Redirect
\end{optlist}
\begin{explainbox}{}
\textbf{Correct Answer: B}\\[0.3em]
\textbf{Concept:} Huawei firewalls deny traffic by default unless explicitly permitted.
\end{explainbox}
\end{questioncard}

\begin{questioncard}{21}{singlecol}
\qtypebadge{singlecol}{Single Choice} \hfill \textcolor{darkslate!70}{\small Topic: Firewall Basics}

Stateful inspection relies on:

\begin{optlist}
\item Session tables
\item MAC tables only
\item DNS records
\item DHCP leases
\end{optlist}
\begin{explainbox}{}
\textbf{Correct Answer: A}\\[0.3em]
\textbf{Concept:} Session tables track connection state for stateful inspection.
\end{explainbox}
\end{questioncard}

\begin{questioncard}{22}{singlecol}
\qtypebadge{singlecol}{Single Choice} \hfill \textcolor{darkslate!70}{\small Topic: Firewall Basics}

ASPF helps firewalls handle:

\begin{optlist}
\item Multi-channel protocols
\item Static routing
\item MAC learning
\item DHCP relay
\end{optlist}
\begin{explainbox}{}
\textbf{Correct Answer: A}\\[0.3em]
\textbf{Concept:} ASPF inspects protocols that negotiate dynamic secondary channels.
\end{explainbox}
\end{questioncard}

\begin{questioncard}{23}{singlecol}
\qtypebadge{singlecol}{Single Choice} \hfill \textcolor{darkslate!70}{\small Topic: Firewall Basics}

Which concept is described as: is 85 by default?

\begin{optlist}
\item Stateful firewall
\item Hybrid mode
\item Untrust zone priority
\item Trust zone priority
\end{optlist}
\begin{explainbox}{}
\textbf{Correct Answer: D}\\[0.3em]
\textbf{Concept:} Trust zone priority refers to the concept where is 85 by default.
\end{explainbox}
\end{questioncard}

\begin{questioncard}{24}{singlecol}
\qtypebadge{singlecol}{Single Choice} \hfill \textcolor{darkslate!70}{\small Topic: Firewall Basics}

Which concept is described as: is 50 by default?

\begin{optlist}
\item DMZ zone priority
\item Trust zone priority
\item Security zones
\item Policy hit count
\end{optlist}
\begin{explainbox}{}
\textbf{Correct Answer: A}\\[0.3em]
\textbf{Concept:} DMZ zone priority refers to the concept where is 50 by default.
\end{explainbox}
\end{questioncard}

\begin{questioncard}{25}{singlecol}
\qtypebadge{singlecol}{Single Choice} \hfill \textcolor{darkslate!70}{\small Topic: Firewall Basics}

Which concept is described as: is 5 by default?

\begin{optlist}
\item Next-generation firewall
\item Stateful firewall
\item Security zones
\item Untrust zone priority
\end{optlist}
\begin{explainbox}{}
\textbf{Correct Answer: D}\\[0.3em]
\textbf{Concept:} Untrust zone priority refers to the concept where is 5 by default.
\end{explainbox}
\end{questioncard}

\begin{questioncard}{26}{singlecol}
\qtypebadge{singlecol}{Single Choice} \hfill \textcolor{darkslate!70}{\small Topic: Firewall Basics}

Which concept is described as: is 100 by default?

\begin{optlist}
\item Transparent mode
\item Untrust zone priority
\item Default security policy action
\item Local zone priority
\end{optlist}
\begin{explainbox}{}
\textbf{Correct Answer: D}\\[0.3em]
\textbf{Concept:} Local zone priority refers to the concept where is 100 by default.
\end{explainbox}
\end{questioncard}

\begin{questioncard}{27}{singlecol}
\qtypebadge{singlecol}{Single Choice} \hfill \textcolor{darkslate!70}{\small Topic: Firewall Basics}

Which concept is described as: is deny?

\begin{optlist}
\item Intra-zone traffic
\item Routing mode
\item Default security policy action
\item Hybrid mode
\end{optlist}
\begin{explainbox}{}
\textbf{Correct Answer: C}\\[0.3em]
\textbf{Concept:} Default security policy action refers to the concept where is deny.
\end{explainbox}
\end{questioncard}

\begin{questioncard}{28}{singlecol}
\qtypebadge{singlecol}{Single Choice} \hfill \textcolor{darkslate!70}{\small Topic: Firewall Basics}

Which concept is described as: tracks connections in a session table?

\begin{optlist}
\item Next-generation firewall
\item Security policy logging
\item Untrust zone priority
\item Stateful firewall
\end{optlist}
\begin{explainbox}{}
\textbf{Correct Answer: D}\\[0.3em]
\textbf{Concept:} Stateful firewall refers to the concept where tracks connections in a session table.
\end{explainbox}
\end{questioncard}

\begin{questioncard}{29}{singlecol}
\qtypebadge{singlecol}{Single Choice} \hfill \textcolor{darkslate!70}{\small Topic: Firewall Basics}

Which concept is described as: inspects packets individually without state?

\begin{optlist}
\item Packet-filtering firewall
\item Hybrid mode
\item DMZ zone priority
\item Untrust zone priority
\end{optlist}
\begin{explainbox}{}
\textbf{Correct Answer: A}\\[0.3em]
\textbf{Concept:} Packet-filtering firewall refers to the concept where inspects packets individually without state.
\end{explainbox}
\end{questioncard}

\begin{questioncard}{30}{singlecol}
\qtypebadge{singlecol}{Single Choice} \hfill \textcolor{darkslate!70}{\small Topic: Firewall Basics}

Which concept is described as: inspects multi-channel protocols like FTP?

\begin{optlist}
\item Intra-zone traffic
\item ASPF
\item Routing mode
\item Local zone priority
\end{optlist}
\begin{explainbox}{}
\textbf{Correct Answer: B}\\[0.3em]
\textbf{Concept:} ASPF refers to the concept where inspects multi-channel protocols like FTP.
\end{explainbox}
\end{questioncard}

\begin{questioncard}{31}{singlecol}
\qtypebadge{singlecol}{Single Choice} \hfill \textcolor{darkslate!70}{\small Topic: Firewall Basics}

Which concept is described as: includes application awareness and IPS?

\begin{optlist}
\item Hybrid mode
\item Security policy logging
\item Next-generation firewall
\item Stateful firewall
\end{optlist}
\begin{explainbox}{}
\textbf{Correct Answer: C}\\[0.3em]
\textbf{Concept:} Next-generation firewall refers to the concept where includes application awareness and IPS.
\end{explainbox}
\end{questioncard}

\begin{questioncard}{32}{singlecol}
\qtypebadge{singlecol}{Single Choice} \hfill \textcolor{darkslate!70}{\small Topic: Firewall Basics}

Which concept is described as: group interfaces with similar trust levels?

\begin{optlist}
\item Next-generation firewall
\item Security zones
\item Default security policy action
\item Untrust zone priority
\end{optlist}
\begin{explainbox}{}
\textbf{Correct Answer: B}\\[0.3em]
\textbf{Concept:} Security zones refers to the concept where group interfaces with similar trust levels.
\end{explainbox}
\end{questioncard}

\begin{questioncard}{33}{singlecol}
\qtypebadge{singlecol}{Single Choice} \hfill \textcolor{darkslate!70}{\small Topic: Firewall Basics}

Which concept is described as: can be controlled by policies?

\begin{optlist}
\item Security policy logging
\item Intra-zone traffic
\item Untrust zone priority
\item Next-generation firewall
\end{optlist}
\begin{explainbox}{}
\textbf{Correct Answer: B}\\[0.3em]
\textbf{Concept:} Intra-zone traffic refers to the concept where can be controlled by policies.
\end{explainbox}
\end{questioncard}

\begin{questioncard}{34}{singlecol}
\qtypebadge{singlecol}{Single Choice} \hfill \textcolor{darkslate!70}{\small Topic: Firewall Basics}

Which concept is described as: helps verify which rules are used?

\begin{optlist}
\item Next-generation firewall
\item Policy hit count
\item Intra-zone traffic
\item Routing mode
\end{optlist}
\begin{explainbox}{}
\textbf{Correct Answer: B}\\[0.3em]
\textbf{Concept:} Policy hit count refers to the concept where helps verify which rules are used.
\end{explainbox}
\end{questioncard}

\begin{questioncard}{35}{singlecol}
\qtypebadge{singlecol}{Single Choice} \hfill \textcolor{darkslate!70}{\small Topic: Firewall Basics}

Which concept is described as: can be enabled for permit and deny actions?

\begin{optlist}
\item Security policy logging
\item Policy hit count
\item DMZ zone priority
\item Packet-filtering firewall
\end{optlist}
\begin{explainbox}{}
\textbf{Correct Answer: A}\\[0.3em]
\textbf{Concept:} Security policy logging refers to the concept where can be enabled for permit and deny actions.
\end{explainbox}
\end{questioncard}

\begin{questioncard}{36}{singlecol}
\qtypebadge{singlecol}{Single Choice} \hfill \textcolor{darkslate!70}{\small Topic: Firewall Basics}

Which concept is described as: is a firewall deployment mode?

\begin{optlist}
\item Trust zone priority
\item Routing mode
\item DMZ zone priority
\item Untrust zone priority
\end{optlist}
\begin{explainbox}{}
\textbf{Correct Answer: B}\\[0.3em]
\textbf{Concept:} Routing mode refers to the concept where is a firewall deployment mode.
\end{explainbox}
\end{questioncard}

\begin{questioncard}{37}{singlecol}
\qtypebadge{singlecol}{Single Choice} \hfill \textcolor{darkslate!70}{\small Topic: Firewall Basics}

Which concept is described as: is a firewall deployment mode?

\begin{optlist}
\item Intra-zone traffic
\item ASPF
\item Transparent mode
\item Next-generation firewall
\end{optlist}
\begin{explainbox}{}
\textbf{Correct Answer: C}\\[0.3em]
\textbf{Concept:} Transparent mode refers to the concept where is a firewall deployment mode.
\end{explainbox}
\end{questioncard}

\begin{questioncard}{38}{singlecol}
\qtypebadge{singlecol}{Single Choice} \hfill \textcolor{darkslate!70}{\small Topic: Firewall Basics}

Which concept is described as: is a firewall deployment mode?

\begin{optlist}
\item Hybrid mode
\item DMZ zone priority
\item ASPF
\item Default security policy action
\end{optlist}
\begin{explainbox}{}
\textbf{Correct Answer: A}\\[0.3em]
\textbf{Concept:} Hybrid mode refers to the concept where is a firewall deployment mode.
\end{explainbox}
\end{questioncard}

\section{Security Zones}
\begin{questioncard}{1}{singlecol}
\qtypebadge{singlecol}{Single Choice} \hfill \textcolor{darkslate!70}{\small Topic: Security Zones}

What is the default priority of the Trust zone on a Huawei firewall?

\begin{optlist}
\item 5
\item 50
\item 85
\item 100
\end{optlist}
\begin{explainbox}{}
\textbf{Correct Answer: C}\\[0.3em]
\textbf{Concept:} The Trust zone has a default priority of 85.
\end{explainbox}
\end{questioncard}

\begin{questioncard}{2}{singlecol}
\qtypebadge{singlecol}{Single Choice} \hfill \textcolor{darkslate!70}{\small Topic: Security Zones}

What is the default priority of the DMZ zone on a Huawei firewall?

\begin{optlist}
\item 5
\item 50
\item 85
\item 100
\end{optlist}
\begin{explainbox}{}
\textbf{Correct Answer: B}\\[0.3em]
\textbf{Concept:} The DMZ zone has a default priority of 50.
\end{explainbox}
\end{questioncard}

\begin{questioncard}{3}{singlecol}
\qtypebadge{singlecol}{Single Choice} \hfill \textcolor{darkslate!70}{\small Topic: Security Zones}

What is the default priority of the Untrust zone on a Huawei firewall?

\begin{optlist}
\item 5
\item 50
\item 85
\item 100
\end{optlist}
\begin{explainbox}{}
\textbf{Correct Answer: A}\\[0.3em]
\textbf{Concept:} The Untrust zone has a default priority of 5.
\end{explainbox}
\end{questioncard}

\begin{questioncard}{4}{singlecol}
\qtypebadge{singlecol}{Single Choice} \hfill \textcolor{darkslate!70}{\small Topic: Security Zones}

What is the default priority of the Local zone on a Huawei firewall?

\begin{optlist}
\item 5
\item 50
\item 85
\item 100
\end{optlist}
\begin{explainbox}{}
\textbf{Correct Answer: D}\\[0.3em]
\textbf{Concept:} The Local zone represents the firewall itself and has the highest default priority of 100.
\end{explainbox}
\end{questioncard}

\section{Security Policies}
\begin{questioncard}{1}{singlecol}
\qtypebadge{singlecol}{Single Choice} \hfill \textcolor{darkslate!70}{\small Topic: Security Policies}

What is the default action of a Huawei firewall security policy if no rule matches?

\begin{optlist}
\item Permit
\item Deny
\item Log only
\item Redirect
\end{optlist}
\begin{explainbox}{}
\textbf{Correct Answer: B}\\[0.3em]
\textbf{Concept:} By default, unmatched traffic is denied.
\end{explainbox}
\end{questioncard}

\begin{questioncard}{2}{singlecol}
\qtypebadge{singlecol}{Single Choice} \hfill \textcolor{darkslate!70}{\small Topic: Security Policies}

Which components can be used as matching conditions in a security policy? (single best answer)

\begin{optlist}
\item Source/destination zone, address, service, application, user, time
\item Only source IP
\item Only destination port
\item Only protocol
\end{optlist}
\begin{explainbox}{}
\textbf{Correct Answer: A}\\[0.3em]
\textbf{Concept:} Security policies can match on zones, addresses, users, services, applications, and time.
\end{explainbox}
\end{questioncard}

\section{Stateful Inspection}
\begin{questioncard}{1}{singlecol}
\qtypebadge{singlecol}{Single Choice} \hfill \textcolor{darkslate!70}{\small Topic: Stateful Inspection}

What is a session table used for in a stateful firewall?

\begin{optlist}
\item Storing user passwords
\item Tracking established connections
\item Routing packets
\item Encrypting traffic
\end{optlist}
\begin{explainbox}{}
\textbf{Correct Answer: B}\\[0.3em]
\textbf{Concept:} The session table tracks connection state so return traffic can be permitted without a new policy lookup.
\end{explainbox}
\end{questioncard}

\begin{questioncard}{2}{singlecol}
\qtypebadge{singlecol}{Single Choice} \hfill \textcolor{darkslate!70}{\small Topic: Stateful Inspection}

Which type of firewall inspects packets individually without connection state?

\begin{optlist}
\item Stateful firewall
\item Packet-filtering firewall
\item Proxy firewall
\item Next-generation firewall
\end{optlist}
\begin{explainbox}{}
\textbf{Correct Answer: B}\\[0.3em]
\textbf{Concept:} Packet-filtering firewalls inspect each packet independently based on header information.
\end{explainbox}
\end{questioncard}

\section{ASPF}
\begin{questioncard}{1}{singlecol}
\qtypebadge{singlecol}{Single Choice} \hfill \textcolor{darkslate!70}{\small Topic: ASPF}

What does ASPF stand for?

\begin{optlist}
\item Advanced Stateful Packet Filtering
\item Application Specific Packet Filtering
\item Automatic Security Policy Forwarding
\item Advanced Stateful Policy Firewall
\end{optlist}
\begin{explainbox}{}
\textbf{Correct Answer: B}\\[0.3em]
\textbf{Concept:} ASPF is Application Specific Packet Filtering, used to inspect multi-channel protocols.
\end{explainbox}
\end{questioncard}

\begin{questioncard}{2}{singlecol}
\qtypebadge{singlecol}{Single Choice} \hfill \textcolor{darkslate!70}{\small Topic: ASPF}

Why is ASPF needed for protocols like FTP?

\begin{optlist}
\item FTP only uses one port
\item FTP opens dynamic data channels
\item FTP encrypts all traffic
\item FTP does not use TCP
\end{optlist}
\begin{explainbox}{}
\textbf{Correct Answer: B}\\[0.3em]
\textbf{Concept:} FTP opens a dynamic data channel, and ASPF inspects the control channel to permit the negotiated data channel.
\end{explainbox}
\end{questioncard}

\section{NAT}
\begin{questioncard}{1}{singlecol}
\qtypebadge{singlecol}{Single Choice} \hfill \textcolor{darkslate!70}{\small Topic: NAT}

Which NAT type is also called many-to-one NAT?

\begin{optlist}
\item Static NAT
\item Dynamic NAT
\item NAPT
\item NAT Server
\end{optlist}
\begin{explainbox}{}
\textbf{Correct Answer: C}\\[0.3em]
\textbf{Concept:} NAPT maps multiple private addresses to one public IP using port numbers.
\end{explainbox}
\end{questioncard}

\begin{questioncard}{2}{singlecol}
\qtypebadge{singlecol}{Single Choice} \hfill \textcolor{darkslate!70}{\small Topic: NAT}

Which NAT type preserves a one-to-one mapping between public and private addresses?

\begin{optlist}
\item Static NAT
\item NAPT
\item Easy-IP
\item Dynamic PAT
\end{optlist}
\begin{explainbox}{}
\textbf{Correct Answer: A}\\[0.3em]
\textbf{Concept:} Static NAT maps one public IP to one private IP consistently.
\end{explainbox}
\end{questioncard}

\begin{questioncard}{3}{singlecol}
\qtypebadge{singlecol}{Single Choice} \hfill \textcolor{darkslate!70}{\small Topic: NAT}

Which NAT type uses the outbound interface IP as the public address?

\begin{optlist}
\item NAPT
\item NAT No-PAT
\item Easy-IP
\item Destination NAT
\end{optlist}
\begin{explainbox}{}
\textbf{Correct Answer: C}\\[0.3em]
\textbf{Concept:} Easy-IP uses the firewall egress interface IP for translation.
\end{explainbox}
\end{questioncard}

\begin{questioncard}{4}{singlecol}
\qtypebadge{singlecol}{Single Choice} \hfill \textcolor{darkslate!70}{\small Topic: NAT}

Which source NAT mode maps one private address to one public address without port translation?

\begin{optlist}
\item NAPT
\item NAT No-PAT
\item Easy-IP
\item NAT Server
\end{optlist}
\begin{explainbox}{}
\textbf{Correct Answer: B}\\[0.3em]
\textbf{Concept:} NAT No-PAT performs one-to-one address translation without changing port numbers.
\end{explainbox}
\end{questioncard}

\begin{questioncard}{5}{singlecol}
\qtypebadge{singlecol}{Single Choice} \hfill \textcolor{darkslate!70}{\small Topic: NAT}

When many internal hosts access the Internet using one public IP, which NAT is used?

\begin{optlist}
\item NAT No-PAT
\item NAPT
\item NAT Server
\item Static NAT
\end{optlist}
\begin{explainbox}{}
\textbf{Correct Answer: B}\\[0.3em]
\textbf{Concept:} NAPT allows many-to-one address translation using port numbers.
\end{explainbox}
\end{questioncard}

\begin{questioncard}{6}{singlecol}
\qtypebadge{singlecol}{Single Choice} \hfill \textcolor{darkslate!70}{\small Topic: NAT}

Destination NAT is also known as:

\begin{optlist}
\item Static NAT
\item Port forwarding / virtual IP
\item PAT
\item Proxy ARP
\end{optlist}
\begin{explainbox}{}
\textbf{Correct Answer: B}\\[0.3em]
\textbf{Concept:} Destination NAT maps external destinations to internal hosts, commonly called port forwarding or virtual IP.
\end{explainbox}
\end{questioncard}

\begin{questioncard}{7}{singlecol}
\qtypebadge{singlecol}{Single Choice} \hfill \textcolor{darkslate!70}{\small Topic: NAT}

Which NAT type maps a public IP/port to a private server IP/port?

\begin{optlist}
\item Source NAT
\item NAT Server
\item NAPT
\item Bidirectional NAT
\end{optlist}
\begin{explainbox}{}
\textbf{Correct Answer: B}\\[0.3em]
\textbf{Concept:} NAT Server publishes internal services to external networks.
\end{explainbox}
\end{questioncard}

\begin{questioncard}{8}{singlecol}
\qtypebadge{singlecol}{Single Choice} \hfill \textcolor{darkslate!70}{\small Topic: NAT}

Bidirectional NAT is useful when:

\begin{optlist}
\item Only source address translation is needed
\item Both source and destination must be rewritten
\item No NAT is needed
\item Only DNS traffic is involved
\end{optlist}
\begin{explainbox}{}
\textbf{Correct Answer: B}\\[0.3em]
\textbf{Concept:} It translates both source and destination addresses in a single flow.
\end{explainbox}
\end{questioncard}

\begin{questioncard}{9}{singlecol}
\qtypebadge{singlecol}{Single Choice} \hfill \textcolor{darkslate!70}{\small Topic: NAT}

Which problem can occur if NAT ALG is disabled for FTP?

\begin{optlist}
\item FTP control channel fails
\item FTP data channel may fail because dynamic ports are not opened
\item FTP becomes encrypted
\item FTP speeds up
\end{optlist}
\begin{explainbox}{}
\textbf{Correct Answer: B}\\[0.3em]
\textbf{Concept:} Without ALG, the firewall cannot open the negotiated FTP data channel.
\end{explainbox}
\end{questioncard}

\begin{questioncard}{10}{singlecol}
\qtypebadge{singlecol}{Single Choice} \hfill \textcolor{darkslate!70}{\small Topic: NAT}

Which address type is preserved by NAT No-PAT?

\begin{optlist}
\item Source port
\item Destination port
\item Both source and destination ports
\item No ports
\end{optlist}
\begin{explainbox}{}
\textbf{Correct Answer: C}\\[0.3em]
\textbf{Concept:} NAT No-PAT translates only IP addresses, preserving port numbers.
\end{explainbox}
\end{questioncard}

\begin{questioncard}{11}{singlecol}
\qtypebadge{singlecol}{Single Choice} \hfill \textcolor{darkslate!70}{\small Topic: NAT}

What is a major benefit of NAPT?

\begin{optlist}
\item One-to-one address mapping
\item Conserves public IP addresses
\item Full protocol transparency
\item No ALG needed
\end{optlist}
\begin{explainbox}{}
\textbf{Correct Answer: B}\\[0.3em]
\textbf{Concept:} NAPT conserves public IPs by sharing one address across many internal hosts.
\end{explainbox}
\end{questioncard}

\begin{questioncard}{12}{singlecol}
\qtypebadge{singlecol}{Single Choice} \hfill \textcolor{darkslate!70}{\small Topic: NAT}

A NAT Server entry is typically:

\begin{optlist}
\item Bidirectional and persistent
\item Temporary and random
\item Encrypted
\item Only for outbound traffic
\end{optlist}
\begin{explainbox}{}
\textbf{Correct Answer: A}\\[0.3em]
\textbf{Concept:} NAT Server entries are configured statically and persist for inbound access.
\end{explainbox}
\end{questioncard}

\begin{questioncard}{13}{singlecol}
\qtypebadge{singlecol}{Single Choice} \hfill \textcolor{darkslate!70}{\small Topic: NAT}

NAT ALG is often enabled together with:

\begin{optlist}
\item ASPF
\item Routing
\item DHCP
\item DNS
\end{optlist}
\begin{explainbox}{}
\textbf{Correct Answer: A}\\[0.3em]
\textbf{Concept:} ALG and ASPF work together to inspect dynamic application channels.
\end{explainbox}
\end{questioncard}

\begin{questioncard}{14}{singlecol}
\qtypebadge{singlecol}{Single Choice} \hfill \textcolor{darkslate!70}{\small Topic: NAT}

Which RFC defines private IPv4 address space?

\begin{optlist}
\item RFC 1918
\item RFC 2328
\item RFC 793
\item RFC 959
\end{optlist}
\begin{explainbox}{}
\textbf{Correct Answer: A}\\[0.3em]
\textbf{Concept:} RFC 1918 defines the private IPv4 address ranges.
\end{explainbox}
\end{questioncard}

\begin{questioncard}{15}{singlecol}
\qtypebadge{singlecol}{Single Choice} \hfill \textcolor{darkslate!70}{\small Topic: NAT}

Hairpin NAT allows internal users to access an internal server using:

\begin{optlist}
\item The server private IP only
\item The server public NAT address
\item A VPN tunnel
\item Broadcast address
\end{optlist}
\begin{explainbox}{}
\textbf{Correct Answer: B}\\[0.3em]
\textbf{Concept:} Hairpin NAT lets internal hosts reach an internal server via its public NAT address.
\end{explainbox}
\end{questioncard}

\begin{questioncard}{16}{singlecol}
\qtypebadge{singlecol}{Single Choice} \hfill \textcolor{darkslate!70}{\small Topic: NAT}

NAPT is also known as:

\begin{optlist}
\item PAT
\item Static NAT
\item NAT Server
\item Bridge NAT
\end{optlist}
\begin{explainbox}{}
\textbf{Correct Answer: A}\\[0.3em]
\textbf{Concept:} Port Address Translation (PAT) is another name for NAPT.
\end{explainbox}
\end{questioncard}

\begin{questioncard}{17}{singlecol}
\qtypebadge{singlecol}{Single Choice} \hfill \textcolor{darkslate!70}{\small Topic: NAT}

Static NAT is most often used when:

\begin{optlist}
\item A server must always use the same public IP
\item Many hosts share one IP
\item Ports must be preserved
\item Outbound only traffic exists
\end{optlist}
\begin{explainbox}{}
\textbf{Correct Answer: A}\\[0.3em]
\textbf{Concept:} Static NAT provides a consistent one-to-one public-to-private mapping.
\end{explainbox}
\end{questioncard}

\begin{questioncard}{18}{singlecol}
\qtypebadge{singlecol}{Single Choice} \hfill \textcolor{darkslate!70}{\small Topic: NAT}

Destination NAT is commonly used to:

\begin{optlist}
\item Allow internal hosts to browse the web
\item Publish internal services to external users
\item Encrypt VPN traffic
\item Assign IP addresses
\end{optlist}
\begin{explainbox}{}
\textbf{Correct Answer: B}\\[0.3em]
\textbf{Concept:} Destination NAT maps external destinations to internal servers.
\end{explainbox}
\end{questioncard}

\begin{questioncard}{19}{singlecol}
\qtypebadge{singlecol}{Single Choice} \hfill \textcolor{darkslate!70}{\small Topic: NAT}

Which feature rewrites IP addresses inside application payloads?

\begin{optlist}
\item NAT ALG
\item NAPT
\item Static NAT
\item Route lookup
\end{optlist}
\begin{explainbox}{}
\textbf{Correct Answer: A}\\[0.3em]
\textbf{Concept:} NAT ALG handles protocols that embed address information in payloads.
\end{explainbox}
\end{questioncard}

\begin{questioncard}{20}{singlecol}
\qtypebadge{singlecol}{Single Choice} \hfill \textcolor{darkslate!70}{\small Topic: NAT}

Which concept is described as: translates IP addresses between private and public?

\begin{optlist}
\item NAT ALG
\item NAT Server
\item Easy-IP
\item NAT
\end{optlist}
\begin{explainbox}{}
\textbf{Correct Answer: D}\\[0.3em]
\textbf{Concept:} NAT refers to the concept where translates IP addresses between private and public.
\end{explainbox}
\end{questioncard}

\begin{questioncard}{21}{singlecol}
\qtypebadge{singlecol}{Single Choice} \hfill \textcolor{darkslate!70}{\small Topic: NAT}

Which concept is described as: translates the source IP of outbound packets?

\begin{optlist}
\item Bidirectional NAT
\item Source NAT
\item NAT No-PAT
\item Easy-IP
\end{optlist}
\begin{explainbox}{}
\textbf{Correct Answer: B}\\[0.3em]
\textbf{Concept:} Source NAT refers to the concept where translates the source IP of outbound packets.
\end{explainbox}
\end{questioncard}

\begin{questioncard}{22}{singlecol}
\qtypebadge{singlecol}{Single Choice} \hfill \textcolor{darkslate!70}{\small Topic: NAT}

Which concept is described as: translates the destination IP of inbound packets?

\begin{optlist}
\item NAT
\item NAT No-PAT
\item Easy-IP
\item Destination NAT
\end{optlist}
\begin{explainbox}{}
\textbf{Correct Answer: D}\\[0.3em]
\textbf{Concept:} Destination NAT refers to the concept where translates the destination IP of inbound packets.
\end{explainbox}
\end{questioncard}

\begin{questioncard}{23}{singlecol}
\qtypebadge{singlecol}{Single Choice} \hfill \textcolor{darkslate!70}{\small Topic: NAT}

Which concept is described as: publishes internal services to the Internet?

\begin{optlist}
\item NAT Server
\item Bidirectional NAT
\item NAPT
\item Source NAT
\end{optlist}
\begin{explainbox}{}
\textbf{Correct Answer: A}\\[0.3em]
\textbf{Concept:} NAT Server refers to the concept where publishes internal services to the Internet.
\end{explainbox}
\end{questioncard}

\begin{questioncard}{24}{singlecol}
\qtypebadge{singlecol}{Single Choice} \hfill \textcolor{darkslate!70}{\small Topic: NAT}

Which concept is described as: translates both source and destination addresses?

\begin{optlist}
\item Bidirectional NAT
\item RFC 1918
\item Static NAT
\item Destination NAT
\end{optlist}
\begin{explainbox}{}
\textbf{Correct Answer: A}\\[0.3em]
\textbf{Concept:} Bidirectional NAT refers to the concept where translates both source and destination addresses.
\end{explainbox}
\end{questioncard}

\begin{questioncard}{25}{singlecol}
\qtypebadge{singlecol}{Single Choice} \hfill \textcolor{darkslate!70}{\small Topic: NAT}

Which concept is described as: translates many private addresses to one public IP using ports?

\begin{optlist}
\item Static NAT
\item NAT No-PAT
\item NAPT
\item NAT Server
\end{optlist}
\begin{explainbox}{}
\textbf{Correct Answer: C}\\[0.3em]
\textbf{Concept:} NAPT refers to the concept where translates many private addresses to one public IP using ports.
\end{explainbox}
\end{questioncard}

\begin{questioncard}{26}{singlecol}
\qtypebadge{singlecol}{Single Choice} \hfill \textcolor{darkslate!70}{\small Topic: NAT}

Which concept is described as: performs one-to-one address translation preserving ports?

\begin{optlist}
\item Hairpin NAT
\item NAT No-PAT
\item Source NAT
\item NAT ALG
\end{optlist}
\begin{explainbox}{}
\textbf{Correct Answer: B}\\[0.3em]
\textbf{Concept:} NAT No-PAT refers to the concept where performs one-to-one address translation preserving ports.
\end{explainbox}
\end{questioncard}

\begin{questioncard}{27}{singlecol}
\qtypebadge{singlecol}{Single Choice} \hfill \textcolor{darkslate!70}{\small Topic: NAT}

Which concept is described as: uses the outbound interface IP for translation?

\begin{optlist}
\item Hairpin NAT
\item Easy-IP
\item Static NAT
\item Destination NAT
\end{optlist}
\begin{explainbox}{}
\textbf{Correct Answer: B}\\[0.3em]
\textbf{Concept:} Easy-IP refers to the concept where uses the outbound interface IP for translation.
\end{explainbox}
\end{questioncard}

\begin{questioncard}{28}{singlecol}
\qtypebadge{singlecol}{Single Choice} \hfill \textcolor{darkslate!70}{\small Topic: NAT}

Which concept is described as: rewrites addresses inside application payloads?

\begin{optlist}
\item NAT ALG
\item Bidirectional NAT
\item Hairpin NAT
\item NAPT
\end{optlist}
\begin{explainbox}{}
\textbf{Correct Answer: A}\\[0.3em]
\textbf{Concept:} NAT ALG refers to the concept where rewrites addresses inside application payloads.
\end{explainbox}
\end{questioncard}

\begin{questioncard}{29}{singlecol}
\qtypebadge{singlecol}{Single Choice} \hfill \textcolor{darkslate!70}{\small Topic: NAT}

Which concept is described as: defines private IPv4 address ranges?

\begin{optlist}
\item Destination NAT
\item RFC 1918
\item Source NAT
\item NAT
\end{optlist}
\begin{explainbox}{}
\textbf{Correct Answer: B}\\[0.3em]
\textbf{Concept:} RFC 1918 refers to the concept where defines private IPv4 address ranges.
\end{explainbox}
\end{questioncard}

\begin{questioncard}{30}{singlecol}
\qtypebadge{singlecol}{Single Choice} \hfill \textcolor{darkslate!70}{\small Topic: NAT}

Which concept is described as: lets internal users access servers via public NAT address?

\begin{optlist}
\item Hairpin NAT
\item Easy-IP
\item NAT ALG
\item Source NAT
\end{optlist}
\begin{explainbox}{}
\textbf{Correct Answer: A}\\[0.3em]
\textbf{Concept:} Hairpin NAT refers to the concept where lets internal users access servers via public NAT address.
\end{explainbox}
\end{questioncard}

\begin{questioncard}{31}{singlecol}
\qtypebadge{singlecol}{Single Choice} \hfill \textcolor{darkslate!70}{\small Topic: NAT}

Which concept is described as: provides a persistent one-to-one mapping?

\begin{optlist}
\item Static NAT
\item NAT No-PAT
\item NAT ALG
\item RFC 1918
\end{optlist}
\begin{explainbox}{}
\textbf{Correct Answer: A}\\[0.3em]
\textbf{Concept:} Static NAT refers to the concept where provides a persistent one-to-one mapping.
\end{explainbox}
\end{questioncard}

\section{NAT Basics}
\begin{questioncard}{1}{singlecol}
\qtypebadge{singlecol}{Single Choice} \hfill \textcolor{darkslate!70}{\small Topic: NAT Basics}

What does NAT primarily translate?

\begin{optlist}
\item MAC addresses
\item IP addresses
\item Port numbers only
\item Domain names
\end{optlist}
\begin{explainbox}{}
\textbf{Correct Answer: B}\\[0.3em]
\textbf{Concept:} NAT translates IP addresses, typically between private and public addresses.
\end{explainbox}
\end{questioncard}

\section{Source NAT}
\begin{questioncard}{1}{singlecol}
\qtypebadge{singlecol}{Single Choice} \hfill \textcolor{darkslate!70}{\small Topic: Source NAT}

Which NAT type translates the source IP address of outgoing packets?

\begin{optlist}
\item Source NAT
\item Destination NAT
\item Bidirectional NAT
\item NAT Server
\end{optlist}
\begin{explainbox}{}
\textbf{Correct Answer: A}\\[0.3em]
\textbf{Concept:} Source NAT changes the source IP so internal hosts can use public addresses to access the Internet.
\end{explainbox}
\end{questioncard}

\begin{questioncard}{2}{singlecol}
\qtypebadge{singlecol}{Single Choice} \hfill \textcolor{darkslate!70}{\small Topic: Source NAT}

What is the main purpose of NAPT?

\begin{optlist}
\item Translate many private addresses to one public IP using port multiplexing
\item Hide the destination IP
\item Encrypt traffic
\item Provide VPN access
\end{optlist}
\begin{explainbox}{}
\textbf{Correct Answer: A}\\[0.3em]
\textbf{Concept:} NAPT (Network Address and Port Translation) maps multiple private addresses to a single public IP using different ports.
\end{explainbox}
\end{questioncard}

\section{Destination NAT}
\begin{questioncard}{1}{singlecol}
\qtypebadge{singlecol}{Single Choice} \hfill \textcolor{darkslate!70}{\small Topic: Destination NAT}

Which NAT type translates the destination IP address of incoming packets?

\begin{optlist}
\item Source NAT
\item Destination NAT
\item Bidirectional NAT
\item NAT Server
\end{optlist}
\begin{explainbox}{}
\textbf{Correct Answer: B}\\[0.3em]
\textbf{Concept:} Destination NAT rewrites the destination IP, often to expose internal servers to external users.
\end{explainbox}
\end{questioncard}

\section{NAT Server}
\begin{questioncard}{1}{singlecol}
\qtypebadge{singlecol}{Single Choice} \hfill \textcolor{darkslate!70}{\small Topic: NAT Server}

What is a NAT Server used for?

\begin{optlist}
\item Hiding internal server IP addresses
\item Publishing internal services to the Internet
\item Encrypting server traffic
\item Load balancing only
\end{optlist}
\begin{explainbox}{}
\textbf{Correct Answer: B}\\[0.3em]
\textbf{Concept:} NAT Server maps public IPs/ports to private server IPs/ports so external users can access internal services.
\end{explainbox}
\end{questioncard}

\section{NAT ALG}
\begin{questioncard}{1}{singlecol}
\qtypebadge{singlecol}{Single Choice} \hfill \textcolor{darkslate!70}{\small Topic: NAT ALG}

What is the purpose of NAT ALG?

\begin{optlist}
\item Encrypt NAT traffic
\item Translate application-layer addresses embedded in protocol payloads
\item Replace firewalls
\item Assign IP addresses
\end{optlist}
\begin{explainbox}{}
\textbf{Correct Answer: B}\\[0.3em]
\textbf{Concept:} NAT Application Level Gateway inspects and rewrites IP addresses inside application-layer protocols (e.g., FTP, SIP).
\end{explainbox}
\end{questioncard}

\section{Firewall Hot Standby}
\begin{questioncard}{1}{singlecol}
\qtypebadge{singlecol}{Single Choice} \hfill \textcolor{darkslate!70}{\small Topic: Firewall Hot Standby}

In active/standby mode, how many firewalls forward traffic at the same time?

\begin{optlist}
\item One
\item Two
\item All
\item None
\end{optlist}
\begin{explainbox}{}
\textbf{Correct Answer: A}\\[0.3em]
\textbf{Concept:} In active/standby mode, one firewall actively forwards while the other remains standby.
\end{explainbox}
\end{questioncard}

\begin{questioncard}{2}{singlecol}
\qtypebadge{singlecol}{Single Choice} \hfill \textcolor{darkslate!70}{\small Topic: Firewall Hot Standby}

In load-sharing mode, how many firewalls forward traffic?

\begin{optlist}
\item One
\item Two or more
\item None
\item Only the backup
\end{optlist}
\begin{explainbox}{}
\textbf{Correct Answer: B}\\[0.3em]
\textbf{Concept:} Load-sharing mode allows multiple firewalls to forward traffic simultaneously.
\end{explainbox}
\end{questioncard}

\begin{questioncard}{3}{singlecol}
\qtypebadge{singlecol}{Single Choice} \hfill \textcolor{darkslate!70}{\small Topic: Firewall Hot Standby}

What is the default VRRP advertisement interval?

\begin{optlist}
\item 1 second
\item 3 seconds
\item 10 seconds
\item 30 seconds
\end{optlist}
\begin{explainbox}{}
\textbf{Correct Answer: A}\\[0.3em]
\textbf{Concept:} VRRP master sends advertisements approximately every 1 second by default.
\end{explainbox}
\end{questioncard}

\begin{questioncard}{4}{singlecol}
\qtypebadge{singlecol}{Single Choice} \hfill \textcolor{darkslate!70}{\small Topic: Firewall Hot Standby}

Which VRRP priority value is highest configurable?

\begin{optlist}
\item 1
\item 100
\item 254
\item 255
\end{optlist}
\begin{explainbox}{}
\textbf{Correct Answer: C}\\[0.3em]
\textbf{Concept:} VRRP priority ranges from 1 to 254 configurable; 255 is reserved for the IP owner.
\end{explainbox}
\end{questioncard}

\begin{questioncard}{5}{singlecol}
\qtypebadge{singlecol}{Single Choice} \hfill \textcolor{darkslate!70}{\small Topic: Firewall Hot Standby}

What happens if a backup router stops receiving VRRP advertisements?

\begin{optlist}
\item It remains backup
\item It preempts and becomes master
\item It shuts down
\item It sends ARP requests
\end{optlist}
\begin{explainbox}{}
\textbf{Correct Answer: B}\\[0.3em]
\textbf{Concept:} Missing advertisements trigger a backup to assume the master role.
\end{explainbox}
\end{questioncard}

\begin{questioncard}{6}{singlecol}
\qtypebadge{singlecol}{Single Choice} \hfill \textcolor{darkslate!70}{\small Topic: Firewall Hot Standby}

What does VGMP group priority determine?

\begin{optlist}
\item NAT pool size
\item Which firewall becomes active
\item Number of interfaces
\item DHCP lease time
\end{optlist}
\begin{explainbox}{}
\textbf{Correct Answer: B}\\[0.3em]
\textbf{Concept:} VGMP priority determines the active/standby role among redundant firewalls.
\end{explainbox}
\end{questioncard}

\begin{questioncard}{7}{singlecol}
\qtypebadge{singlecol}{Single Choice} \hfill \textcolor{darkslate!70}{\small Topic: Firewall Hot Standby}

Which HRP state indicates the firewall is currently forwarding traffic?

\begin{optlist}
\item Standby
\item Active
\item Initialize
\item Backup
\end{optlist}
\begin{explainbox}{}
\textbf{Correct Answer: B}\\[0.3em]
\textbf{Concept:} The active HRP peer forwards traffic and synchronizes state to the standby.
\end{explainbox}
\end{questioncard}

\begin{questioncard}{8}{singlecol}
\qtypebadge{singlecol}{Single Choice} \hfill \textcolor{darkslate!70}{\small Topic: Firewall Hot Standby}

Which HRP command enables automatic configuration synchronization?

\begin{optlist}
\item hrp enable
\item hrp auto-sync config
\item hrp mirror session
\item hrp vrrp enable
\end{optlist}
\begin{explainbox}{}
\textbf{Correct Answer: B}\\[0.3em]
\textbf{Concept:} hrp auto-sync config enables automatic configuration synchronization between peers.
\end{explainbox}
\end{questioncard}

\begin{questioncard}{9}{singlecol}
\qtypebadge{singlecol}{Single Choice} \hfill \textcolor{darkslate!70}{\small Topic: Firewall Hot Standby}

Which deployment places both firewalls inline with the same traffic path?

\begin{optlist}
\item Dual-hot standby
\item Active/standby inline
\item Route backup only
\item Transparent bridge
\end{optlist}
\begin{explainbox}{}
\textbf{Correct Answer: B}\\[0.3em]
\textbf{Concept:} Inline active/standby deployment places firewalls in the same forwarding path.
\end{explainbox}
\end{questioncard}

\begin{questioncard}{10}{singlecol}
\qtypebadge{singlecol}{Single Choice} \hfill \textcolor{darkslate!70}{\small Topic: Firewall Hot Standby}

What is a common requirement for HRP heartbeat links?

\begin{optlist}
\item Use the same interface as user traffic
\item Use a dedicated high-reliability link
\item Use wireless link only
\item No link required
\end{optlist}
\begin{explainbox}{}
\textbf{Correct Answer: B}\\[0.3em]
\textbf{Concept:} Heartbeat links should be dedicated and reliable to avoid false failovers.
\end{explainbox}
\end{questioncard}

\begin{questioncard}{11}{singlecol}
\qtypebadge{singlecol}{Single Choice} \hfill \textcolor{darkslate!70}{\small Topic: Firewall Hot Standby}

Which mechanism reduces unnecessary VRRP state changes when interface flaps briefly?

\begin{optlist}
\item Preempt delay
\item Track interface
\item Authentication
\item Priority 255
\end{optlist}
\begin{explainbox}{}
\textbf{Correct Answer: A}\\[0.3em]
\textbf{Concept:} Preempt delay prevents rapid master/backup flapping.
\end{explainbox}
\end{questioncard}

\begin{questioncard}{12}{singlecol}
\qtypebadge{singlecol}{Single Choice} \hfill \textcolor{darkslate!70}{\small Topic: Firewall Hot Standby}

Which HRP feature ensures sessions survive a failover?

\begin{optlist}
\item Session backup
\item Route backup
\item NAT ALG
\item URL filtering
\end{optlist}
\begin{explainbox}{}
\textbf{Correct Answer: A}\\[0.3em]
\textbf{Concept:} HRP session backup keeps session state synchronized so connections continue after failover.
\end{explainbox}
\end{questioncard}

\begin{questioncard}{13}{singlecol}
\qtypebadge{singlecol}{Single Choice} \hfill \textcolor{darkslate!70}{\small Topic: Firewall Hot Standby}

VRRP master election is based primarily on:

\begin{optlist}
\item Priority value
\item Router hostname
\item Number of interfaces
\item Uptime only
\end{optlist}
\begin{explainbox}{}
\textbf{Correct Answer: A}\\[0.3em]
\textbf{Concept:} The router with the highest priority becomes the VRRP master.
\end{explainbox}
\end{questioncard}

\begin{questioncard}{14}{singlecol}
\qtypebadge{singlecol}{Single Choice} \hfill \textcolor{darkslate!70}{\small Topic: Firewall Hot Standby}

VGMP group management prevents:

\begin{optlist}
\item Split-brain VRRP states
\item NAT exhaustion
\item DNS failures
\item DHCP conflicts
\end{optlist}
\begin{explainbox}{}
\textbf{Correct Answer: A}\\[0.3em]
\textbf{Concept:} VGMP ensures VRRP groups fail over consistently.
\end{explainbox}
\end{questioncard}

\begin{questioncard}{15}{singlecol}
\qtypebadge{singlecol}{Single Choice} \hfill \textcolor{darkslate!70}{\small Topic: Firewall Hot Standby}

HRP is used to synchronize:

\begin{optlist}
\item Session and configuration data
\item Email messages
\item Web pages
\item Printer drivers
\end{optlist}
\begin{explainbox}{}
\textbf{Correct Answer: A}\\[0.3em]
\textbf{Concept:} HRP synchronizes sessions, NAT, server-map entries, and configurations.
\end{explainbox}
\end{questioncard}

\begin{questioncard}{16}{singlecol}
\qtypebadge{singlecol}{Single Choice} \hfill \textcolor{darkslate!70}{\small Topic: Firewall Hot Standby}

In active/standby mode, the standby firewall:

\begin{optlist}
\item Does not forward traffic until failover
\item Forwards half the traffic
\item Processes all traffic
\item Acts as a switch
\end{optlist}
\begin{explainbox}{}
\textbf{Correct Answer: A}\\[0.3em]
\textbf{Concept:} The standby firewall remains ready but inactive until failover.
\end{explainbox}
\end{questioncard}

\begin{questioncard}{17}{singlecol}
\qtypebadge{singlecol}{Single Choice} \hfill \textcolor{darkslate!70}{\small Topic: Firewall Hot Standby}

Which concept is described as: provides gateway redundancy with a virtual IP?

\begin{optlist}
\item Load-sharing mode
\item VGMP
\item VRRP
\item Preempt delay
\end{optlist}
\begin{explainbox}{}
\textbf{Correct Answer: C}\\[0.3em]
\textbf{Concept:} VRRP refers to the concept where provides gateway redundancy with a virtual IP.
\end{explainbox}
\end{questioncard}

\begin{questioncard}{18}{singlecol}
\qtypebadge{singlecol}{Single Choice} \hfill \textcolor{darkslate!70}{\small Topic: Firewall Hot Standby}

Which concept is described as: forwards traffic for the virtual IP?

\begin{optlist}
\item Heartbeat link
\item VGMP
\item HRP
\item VRRP master
\end{optlist}
\begin{explainbox}{}
\textbf{Correct Answer: D}\\[0.3em]
\textbf{Concept:} VRRP master refers to the concept where forwards traffic for the virtual IP.
\end{explainbox}
\end{questioncard}

\begin{questioncard}{19}{singlecol}
\qtypebadge{singlecol}{Single Choice} \hfill \textcolor{darkslate!70}{\small Topic: Firewall Hot Standby}

Which concept is described as: takes over if the master fails?

\begin{optlist}
\item VGMP
\item VRRP backup
\item HRP session backup
\item Load-sharing mode
\end{optlist}
\begin{explainbox}{}
\textbf{Correct Answer: B}\\[0.3em]
\textbf{Concept:} VRRP backup refers to the concept where takes over if the master fails.
\end{explainbox}
\end{questioncard}

\begin{questioncard}{20}{singlecol}
\qtypebadge{singlecol}{Single Choice} \hfill \textcolor{darkslate!70}{\small Topic: Firewall Hot Standby}

Which concept is described as: manages VRRP group states consistently?

\begin{optlist}
\item VGMP
\item HRP
\item VRRP
\item VRRP master
\end{optlist}
\begin{explainbox}{}
\textbf{Correct Answer: A}\\[0.3em]
\textbf{Concept:} VGMP refers to the concept where manages VRRP group states consistently.
\end{explainbox}
\end{questioncard}

\begin{questioncard}{21}{singlecol}
\qtypebadge{singlecol}{Single Choice} \hfill \textcolor{darkslate!70}{\small Topic: Firewall Hot Standby}

Which concept is described as: synchronizes sessions and configurations?

\begin{optlist}
\item VRRP master
\item Heartbeat link
\item VRRP
\item HRP
\end{optlist}
\begin{explainbox}{}
\textbf{Correct Answer: D}\\[0.3em]
\textbf{Concept:} HRP refers to the concept where synchronizes sessions and configurations.
\end{explainbox}
\end{questioncard}

\begin{questioncard}{22}{singlecol}
\qtypebadge{singlecol}{Single Choice} \hfill \textcolor{darkslate!70}{\small Topic: Firewall Hot Standby}

Which concept is described as: has one firewall forwarding traffic?

\begin{optlist}
\item Preempt delay
\item HRP session backup
\item Active/standby mode
\item Load-sharing mode
\end{optlist}
\begin{explainbox}{}
\textbf{Correct Answer: C}\\[0.3em]
\textbf{Concept:} Active/standby mode refers to the concept where has one firewall forwarding traffic.
\end{explainbox}
\end{questioncard}

\begin{questioncard}{23}{singlecol}
\qtypebadge{singlecol}{Single Choice} \hfill \textcolor{darkslate!70}{\small Topic: Firewall Hot Standby}

Which concept is described as: has multiple firewalls forwarding traffic?

\begin{optlist}
\item Load-sharing mode
\item VRRP backup
\item Heartbeat link
\item HRP
\end{optlist}
\begin{explainbox}{}
\textbf{Correct Answer: A}\\[0.3em]
\textbf{Concept:} Load-sharing mode refers to the concept where has multiple firewalls forwarding traffic.
\end{explainbox}
\end{questioncard}

\begin{questioncard}{24}{singlecol}
\qtypebadge{singlecol}{Single Choice} \hfill \textcolor{darkslate!70}{\small Topic: Firewall Hot Standby}

Which concept is described as: should be dedicated and reliable?

\begin{optlist}
\item Heartbeat link
\item Preempt delay
\item VRRP
\item VRRP backup
\end{optlist}
\begin{explainbox}{}
\textbf{Correct Answer: A}\\[0.3em]
\textbf{Concept:} Heartbeat link refers to the concept where should be dedicated and reliable.
\end{explainbox}
\end{questioncard}

\begin{questioncard}{25}{singlecol}
\qtypebadge{singlecol}{Single Choice} \hfill \textcolor{darkslate!70}{\small Topic: Firewall Hot Standby}

Which concept is described as: ensures connections survive failover?

\begin{optlist}
\item Load-sharing mode
\item VGMP
\item HRP session backup
\item Heartbeat link
\end{optlist}
\begin{explainbox}{}
\textbf{Correct Answer: C}\\[0.3em]
\textbf{Concept:} HRP session backup refers to the concept where ensures connections survive failover.
\end{explainbox}
\end{questioncard}

\begin{questioncard}{26}{singlecol}
\qtypebadge{singlecol}{Single Choice} \hfill \textcolor{darkslate!70}{\small Topic: Firewall Hot Standby}

Which concept is described as: reduces rapid master-backup flapping?

\begin{optlist}
\item Preempt delay
\item VRRP master
\item Heartbeat link
\item Active/standby mode
\end{optlist}
\begin{explainbox}{}
\textbf{Correct Answer: A}\\[0.3em]
\textbf{Concept:} Preempt delay refers to the concept where reduces rapid master-backup flapping.
\end{explainbox}
\end{questioncard}

\section{VRRP}
\begin{questioncard}{1}{singlecol}
\qtypebadge{singlecol}{Single Choice} \hfill \textcolor{darkslate!70}{\small Topic: VRRP}

What does VRRP stand for?

\begin{optlist}
\item Virtual Router Redundancy Protocol
\item Virtual Routing Resolution Protocol
\item Virtual Redundant Routing Path
\item Virtual Routing Relay Protocol
\end{optlist}
\begin{explainbox}{}
\textbf{Correct Answer: A}\\[0.3em]
\textbf{Concept:} VRRP provides gateway redundancy by allowing multiple routers to share a virtual IP address.
\end{explainbox}
\end{questioncard}

\begin{questioncard}{2}{singlecol}
\qtypebadge{singlecol}{Single Choice} \hfill \textcolor{darkslate!70}{\small Topic: VRRP}

In VRRP, which router forwards traffic for the virtual IP?

\begin{optlist}
\item Backup router
\item Master router
\item All routers simultaneously
\item Only the physical IP owner
\end{optlist}
\begin{explainbox}{}
\textbf{Correct Answer: B}\\[0.3em]
\textbf{Concept:} The VRRP master router owns the virtual IP and forwards traffic; backups take over if it fails.
\end{explainbox}
\end{questioncard}

\section{VGMP}
\begin{questioncard}{1}{singlecol}
\qtypebadge{singlecol}{Single Choice} \hfill \textcolor{darkslate!70}{\small Topic: VGMP}

What does VGMP manage on Huawei firewalls?

\begin{optlist}
\item VRRP groups
\item NAT pools
\item Routing tables
\item DHCP scopes
\end{optlist}
\begin{explainbox}{}
\textbf{Correct Answer: A}\\[0.3em]
\textbf{Concept:} VGMP (VRRP Group Management Protocol) manages the state of multiple VRRP groups consistently.
\end{explainbox}
\end{questioncard}

\section{HRP}
\begin{questioncard}{1}{singlecol}
\qtypebadge{singlecol}{Single Choice} \hfill \textcolor{darkslate!70}{\small Topic: HRP}

What is the primary purpose of HRP?

\begin{optlist}
\item Encrypt traffic
\item Synchronize sessions and configurations between firewalls
\item Assign IP addresses
\item Filter URLs
\end{optlist}
\begin{explainbox}{}
\textbf{Correct Answer: B}\\[0.3em]
\textbf{Concept:} HRP (Huawei Redundancy Protocol) synchronizes sessions, configurations, and state between hot-standby firewalls.
\end{explainbox}
\end{questioncard}

\section{Intrusion Prevention and Antivirus}
\begin{questioncard}{1}{singlecol}
\qtypebadge{singlecol}{Single Choice} \hfill \textcolor{darkslate!70}{\small Topic: Intrusion Prevention and Antivirus}

Which phase of an attack involves identifying vulnerable services?

\begin{optlist}
\item Exploitation
\item Reconnaissance
\item Installation
\item Exfiltration
\end{optlist}
\begin{explainbox}{}
\textbf{Correct Answer: B}\\[0.3em]
\textbf{Concept:} Reconnaissance gathers information about targets before exploitation.
\end{explainbox}
\end{questioncard}

\begin{questioncard}{2}{singlecol}
\qtypebadge{singlecol}{Single Choice} \hfill \textcolor{darkslate!70}{\small Topic: Intrusion Prevention and Antivirus}

Which phase involves gaining unauthorized access?

\begin{optlist}
\item Reconnaissance
\item Exploitation
\item Covering tracks
\item Reporting
\end{optlist}
\begin{explainbox}{}
\textbf{Correct Answer: B}\\[0.3em]
\textbf{Concept:} Exploitation uses vulnerabilities to gain access to the target.
\end{explainbox}
\end{questioncard}

\begin{questioncard}{3}{singlecol}
\qtypebadge{singlecol}{Single Choice} \hfill \textcolor{darkslate!70}{\small Topic: Intrusion Prevention and Antivirus}

Which deployment mode allows IDS to monitor traffic without affecting it?

\begin{optlist}
\item Inline
\item Tap/SPAN
\item Bridge
\item Gateway
\end{optlist}
\begin{explainbox}{}
\textbf{Correct Answer: B}\\[0.3em]
\textbf{Concept:} IDS commonly connects to a tap or switch SPAN port for passive monitoring.
\end{explainbox}
\end{questioncard}

\begin{questioncard}{4}{singlecol}
\qtypebadge{singlecol}{Single Choice} \hfill \textcolor{darkslate!70}{\small Topic: Intrusion Prevention and Antivirus}

Which signature action drops the offending packet?

\begin{optlist}
\item Alert
\item Block
\item Permit
\item Mirror
\end{optlist}
\begin{explainbox}{}
\textbf{Correct Answer: B}\\[0.3em]
\textbf{Concept:} Block actions discard malicious packets.
\end{explainbox}
\end{questioncard}

\begin{questioncard}{5}{singlecol}
\qtypebadge{singlecol}{Single Choice} \hfill \textcolor{darkslate!70}{\small Topic: Intrusion Prevention and Antivirus}

What is a false positive in IPS?

\begin{optlist}
\item Attack missed
\item Legitimate traffic blocked as malicious
\item Correct detection
\item Signature update failure
\end{optlist}
\begin{explainbox}{}
\textbf{Correct Answer: B}\\[0.3em]
\textbf{Concept:} False positives incorrectly classify legitimate traffic as malicious.
\end{explainbox}
\end{questioncard}

\begin{questioncard}{6}{singlecol}
\qtypebadge{singlecol}{Single Choice} \hfill \textcolor{darkslate!70}{\small Topic: Intrusion Prevention and Antivirus}

What is a false negative in IPS?

\begin{optlist}
\item Legitimate traffic blocked
\item Attack missed by IPS
\item Correct block
\item Logging error
\end{optlist}
\begin{explainbox}{}
\textbf{Correct Answer: B}\\[0.3em]
\textbf{Concept:} False negatives occur when an actual attack is not detected.
\end{explainbox}
\end{questioncard}

\begin{questioncard}{7}{singlecol}
\qtypebadge{singlecol}{Single Choice} \hfill \textcolor{darkslate!70}{\small Topic: Intrusion Prevention and Antivirus}

Which malware type encrypts user files and demands ransom?

\begin{optlist}
\item Virus
\item Worm
\item Trojan
\item Ransomware
\end{optlist}
\begin{explainbox}{}
\textbf{Correct Answer: D}\\[0.3em]
\textbf{Concept:} Ransomware encrypts files and demands a ransom for decryption.
\end{explainbox}
\end{questioncard}

\begin{questioncard}{8}{singlecol}
\qtypebadge{singlecol}{Single Choice} \hfill \textcolor{darkslate!70}{\small Topic: Intrusion Prevention and Antivirus}

Which approach only allows known-good applications to run?

\begin{optlist}
\item Blacklisting
\item Whitelisting
\item Heuristics
\item Sandboxing
\end{optlist}
\begin{explainbox}{}
\textbf{Correct Answer: B}\\[0.3em]
\textbf{Concept:} Whitelisting permits only approved applications.
\end{explainbox}
\end{questioncard}

\begin{questioncard}{9}{singlecol}
\qtypebadge{singlecol}{Single Choice} \hfill \textcolor{darkslate!70}{\small Topic: Intrusion Prevention and Antivirus}

Which technique detects malware by examining code behavior rather than signatures?

\begin{optlist}
\item Signature scanning
\item Heuristic analysis
\item Whitelisting
\item Patching
\end{optlist}
\begin{explainbox}{}
\textbf{Correct Answer: B}\\[0.3em]
\textbf{Concept:} Heuristics analyze code behavior to identify potentially malicious activity.
\end{explainbox}
\end{questioncard}

\begin{questioncard}{10}{singlecol}
\qtypebadge{singlecol}{Single Choice} \hfill \textcolor{darkslate!70}{\small Topic: Intrusion Prevention and Antivirus}

Which protocol is commonly targeted by brute-force attacks against remote access?

\begin{optlist}
\item HTTP
\item SSH
\item DNS
\item NTP
\end{optlist}
\begin{explainbox}{}
\textbf{Correct Answer: B}\\[0.3em]
\textbf{Concept:} SSH and similar remote-access protocols are frequent brute-force targets.
\end{explainbox}
\end{questioncard}

\begin{questioncard}{11}{singlecol}
\qtypebadge{singlecol}{Single Choice} \hfill \textcolor{darkslate!70}{\small Topic: Intrusion Prevention and Antivirus}

What does an IPS signature database contain?

\begin{optlist}
\item User passwords
\item Known attack patterns
\item Routing tables
\item DNS records
\end{optlist}
\begin{explainbox}{}
\textbf{Correct Answer: B}\\[0.3em]
\textbf{Concept:} Signature databases contain patterns used to identify known attacks.
\end{explainbox}
\end{questioncard}

\begin{questioncard}{12}{singlecol}
\qtypebadge{singlecol}{Single Choice} \hfill \textcolor{darkslate!70}{\small Topic: Intrusion Prevention and Antivirus}

What is the role of the antivirus engine on a firewall?

\begin{optlist}
\item Route packets
\item Scan files and traffic for malware
\item Assign IP addresses
\item Authenticate users
\end{optlist}
\begin{explainbox}{}
\textbf{Correct Answer: B}\\[0.3em]
\textbf{Concept:} The antivirus engine inspects files and traffic to detect and block malware.
\end{explainbox}
\end{questioncard}

\begin{questioncard}{13}{singlecol}
\qtypebadge{singlecol}{Single Choice} \hfill \textcolor{darkslate!70}{\small Topic: Intrusion Prevention and Antivirus}

An IPS differs from an IDS mainly because an IPS can:

\begin{optlist}
\item Block traffic
\item Only log events
\item Operate out-of-band only
\item Never generate alerts
\end{optlist}
\begin{explainbox}{}
\textbf{Correct Answer: A}\\[0.3em]
\textbf{Concept:} IPS can take active blocking actions, while IDS primarily detects.
\end{explainbox}
\end{questioncard}

\begin{questioncard}{14}{singlecol}
\qtypebadge{singlecol}{Single Choice} \hfill \textcolor{darkslate!70}{\small Topic: Intrusion Prevention and Antivirus}

Signature-based detection cannot easily detect:

\begin{optlist}
\item Known attacks
\item Zero-day attacks
\item Old malware
\item Common exploits
\end{optlist}
\begin{explainbox}{}
\textbf{Correct Answer: B}\\[0.3em]
\textbf{Concept:} Zero-day attacks have no known signatures.
\end{explainbox}
\end{questioncard}

\begin{questioncard}{15}{singlecol}
\qtypebadge{singlecol}{Single Choice} \hfill \textcolor{darkslate!70}{\small Topic: Intrusion Prevention and Antivirus}

Sandboxing helps detect malware by:

\begin{optlist}
\item Observing behavior in an isolated environment
\item Matching signatures
\item Comparing file sizes
\item Checking file names
\end{optlist}
\begin{explainbox}{}
\textbf{Correct Answer: A}\\[0.3em]
\textbf{Concept:} Sandboxing executes suspicious files safely to observe malicious behavior.
\end{explainbox}
\end{questioncard}

\begin{questioncard}{16}{singlecol}
\qtypebadge{singlecol}{Single Choice} \hfill \textcolor{darkslate!70}{\small Topic: Intrusion Prevention and Antivirus}

Antivirus engines require regular updates to:

\begin{optlist}
\item Recognize new malware variants
\item Improve routing
\item Increase bandwidth
\item Disable firewalls
\end{optlist}
\begin{explainbox}{}
\textbf{Correct Answer: A}\\[0.3em]
\textbf{Concept:} Signature and detection databases must be updated to detect emerging threats.
\end{explainbox}
\end{questioncard}

\begin{questioncard}{17}{singlecol}
\qtypebadge{singlecol}{Single Choice} \hfill \textcolor{darkslate!70}{\small Topic: Intrusion Prevention and Antivirus}

Which concept is described as: detects intrusions and generates alerts?

\begin{optlist}
\item Antivirus signatures
\item False negative
\item Signature detection
\item IDS
\end{optlist}
\begin{explainbox}{}
\textbf{Correct Answer: D}\\[0.3em]
\textbf{Concept:} IDS refers to the concept where detects intrusions and generates alerts.
\end{explainbox}
\end{questioncard}

\begin{questioncard}{18}{singlecol}
\qtypebadge{singlecol}{Single Choice} \hfill \textcolor{darkslate!70}{\small Topic: Intrusion Prevention and Antivirus}

Which concept is described as: detects and blocks intrusions inline?

\begin{optlist}
\item False positive
\item IPS
\item Anomaly detection
\item NIDS
\end{optlist}
\begin{explainbox}{}
\textbf{Correct Answer: B}\\[0.3em]
\textbf{Concept:} IPS refers to the concept where detects and blocks intrusions inline.
\end{explainbox}
\end{questioncard}

\begin{questioncard}{19}{singlecol}
\qtypebadge{singlecol}{Single Choice} \hfill \textcolor{darkslate!70}{\small Topic: Intrusion Prevention and Antivirus}

Which concept is described as: matches known attack patterns?

\begin{optlist}
\item Signature detection
\item Whitelisting
\item IPS
\item Exploitation
\end{optlist}
\begin{explainbox}{}
\textbf{Correct Answer: A}\\[0.3em]
\textbf{Concept:} Signature detection refers to the concept where matches known attack patterns.
\end{explainbox}
\end{questioncard}

\begin{questioncard}{20}{singlecol}
\qtypebadge{singlecol}{Single Choice} \hfill \textcolor{darkslate!70}{\small Topic: Intrusion Prevention and Antivirus}

Which concept is described as: flags deviations from normal behavior?

\begin{optlist}
\item NIDS
\item Anomaly detection
\item Exploitation
\item Whitelisting
\end{optlist}
\begin{explainbox}{}
\textbf{Correct Answer: B}\\[0.3em]
\textbf{Concept:} Anomaly detection refers to the concept where flags deviations from normal behavior.
\end{explainbox}
\end{questioncard}

\begin{questioncard}{21}{singlecol}
\qtypebadge{singlecol}{Single Choice} \hfill \textcolor{darkslate!70}{\small Topic: Intrusion Prevention and Antivirus}

Which concept is described as: incorrectly blocks legitimate traffic?

\begin{optlist}
\item Heuristic analysis
\item Anomaly detection
\item Exploitation
\item False positive
\end{optlist}
\begin{explainbox}{}
\textbf{Correct Answer: D}\\[0.3em]
\textbf{Concept:} False positive refers to the concept where incorrectly blocks legitimate traffic.
\end{explainbox}
\end{questioncard}

\begin{questioncard}{22}{singlecol}
\qtypebadge{singlecol}{Single Choice} \hfill \textcolor{darkslate!70}{\small Topic: Intrusion Prevention and Antivirus}

Which concept is described as: misses an actual attack?

\begin{optlist}
\item Sandboxing
\item Anomaly detection
\item Whitelisting
\item False negative
\end{optlist}
\begin{explainbox}{}
\textbf{Correct Answer: D}\\[0.3em]
\textbf{Concept:} False negative refers to the concept where misses an actual attack.
\end{explainbox}
\end{questioncard}

\begin{questioncard}{23}{singlecol}
\qtypebadge{singlecol}{Single Choice} \hfill \textcolor{darkslate!70}{\small Topic: Intrusion Prevention and Antivirus}

Which concept is described as: observes file behavior in an isolated environment?

\begin{optlist}
\item False negative
\item Sandboxing
\item Heuristic analysis
\item NIDS
\end{optlist}
\begin{explainbox}{}
\textbf{Correct Answer: B}\\[0.3em]
\textbf{Concept:} Sandboxing refers to the concept where observes file behavior in an isolated environment.
\end{explainbox}
\end{questioncard}

\begin{questioncard}{24}{singlecol}
\qtypebadge{singlecol}{Single Choice} \hfill \textcolor{darkslate!70}{\small Topic: Intrusion Prevention and Antivirus}

Which concept is described as: detects malware by examining behavior?

\begin{optlist}
\item Reconnaissance
\item Heuristic analysis
\item Whitelisting
\item Signature detection
\end{optlist}
\begin{explainbox}{}
\textbf{Correct Answer: B}\\[0.3em]
\textbf{Concept:} Heuristic analysis refers to the concept where detects malware by examining behavior.
\end{explainbox}
\end{questioncard}

\begin{questioncard}{25}{singlecol}
\qtypebadge{singlecol}{Single Choice} \hfill \textcolor{darkslate!70}{\small Topic: Intrusion Prevention and Antivirus}

Which concept is described as: allows only approved applications?

\begin{optlist}
\item IDS
\item HIDS
\item Whitelisting
\item Sandboxing
\end{optlist}
\begin{explainbox}{}
\textbf{Correct Answer: C}\\[0.3em]
\textbf{Concept:} Whitelisting refers to the concept where allows only approved applications.
\end{explainbox}
\end{questioncard}

\begin{questioncard}{26}{singlecol}
\qtypebadge{singlecol}{Single Choice} \hfill \textcolor{darkslate!70}{\small Topic: Intrusion Prevention and Antivirus}

Which concept is described as: must be updated regularly?

\begin{optlist}
\item IDS
\item Antivirus signatures
\item False positive
\item NIDS
\end{optlist}
\begin{explainbox}{}
\textbf{Correct Answer: B}\\[0.3em]
\textbf{Concept:} Antivirus signatures refers to the concept where must be updated regularly.
\end{explainbox}
\end{questioncard}

\begin{questioncard}{27}{singlecol}
\qtypebadge{singlecol}{Single Choice} \hfill \textcolor{darkslate!70}{\small Topic: Intrusion Prevention and Antivirus}

Which concept is described as: monitors network traffic for intrusions?

\begin{optlist}
\item NIDS
\item HIDS
\item Whitelisting
\item Anomaly detection
\end{optlist}
\begin{explainbox}{}
\textbf{Correct Answer: A}\\[0.3em]
\textbf{Concept:} NIDS refers to the concept where monitors network traffic for intrusions.
\end{explainbox}
\end{questioncard}

\begin{questioncard}{28}{singlecol}
\qtypebadge{singlecol}{Single Choice} \hfill \textcolor{darkslate!70}{\small Topic: Intrusion Prevention and Antivirus}

Which concept is described as: monitors host activity for intrusions?

\begin{optlist}
\item Antivirus signatures
\item HIDS
\item IDS
\item False negative
\end{optlist}
\begin{explainbox}{}
\textbf{Correct Answer: B}\\[0.3em]
\textbf{Concept:} HIDS refers to the concept where monitors host activity for intrusions.
\end{explainbox}
\end{questioncard}

\begin{questioncard}{29}{singlecol}
\qtypebadge{singlecol}{Single Choice} \hfill \textcolor{darkslate!70}{\small Topic: Intrusion Prevention and Antivirus}

Which concept is described as: is an early attack phase?

\begin{optlist}
\item IDS
\item Reconnaissance
\item False positive
\item Heuristic analysis
\end{optlist}
\begin{explainbox}{}
\textbf{Correct Answer: B}\\[0.3em]
\textbf{Concept:} Reconnaissance refers to the concept where is an early attack phase.
\end{explainbox}
\end{questioncard}

\begin{questioncard}{30}{singlecol}
\qtypebadge{singlecol}{Single Choice} \hfill \textcolor{darkslate!70}{\small Topic: Intrusion Prevention and Antivirus}

Which concept is described as: gains unauthorized access using vulnerabilities?

\begin{optlist}
\item Antivirus signatures
\item Exploitation
\item Sandboxing
\item HIDS
\end{optlist}
\begin{explainbox}{}
\textbf{Correct Answer: B}\\[0.3em]
\textbf{Concept:} Exploitation refers to the concept where gains unauthorized access using vulnerabilities.
\end{explainbox}
\end{questioncard}

\section{Intrusion Overview}
\begin{questioncard}{1}{singlecol}
\qtypebadge{singlecol}{Single Choice} \hfill \textcolor{darkslate!70}{\small Topic: Intrusion Overview}

What is an intrusion?

\begin{optlist}
\item Authorized access
\item Unauthorized access or attack on a network/system
\item A software patch
\item A backup operation
\end{optlist}
\begin{explainbox}{}
\textbf{Correct Answer: B}\\[0.3em]
\textbf{Concept:} An intrusion is any unauthorized attempt to access, manipulate, or harm a system or network.
\end{explainbox}
\end{questioncard}

\section{Intrusion Prevention}
\begin{questioncard}{1}{singlecol}
\qtypebadge{singlecol}{Single Choice} \hfill \textcolor{darkslate!70}{\small Topic: Intrusion Prevention}

What is the main difference between IDS and IPS?

\begin{optlist}
\item IDS only detects; IPS detects and blocks
\item IDS is faster
\item IPS cannot log
\item IDS operates inline
\end{optlist}
\begin{explainbox}{}
\textbf{Correct Answer: A}\\[0.3em]
\textbf{Concept:} IDS detects intrusions and alerts; IPS can also take action to block malicious traffic.
\end{explainbox}
\end{questioncard}

\begin{questioncard}{2}{singlecol}
\qtypebadge{singlecol}{Single Choice} \hfill \textcolor{darkslate!70}{\small Topic: Intrusion Prevention}

Which detection method compares traffic against known attack signatures?

\begin{optlist}
\item Anomaly detection
\item Signature detection
\item Behavioral analysis
\item Heuristic detection
\end{optlist}
\begin{explainbox}{}
\textbf{Correct Answer: B}\\[0.3em]
\textbf{Concept:} Signature-based detection matches traffic patterns against a database of known attack signatures.
\end{explainbox}
\end{questioncard}

\begin{questioncard}{3}{singlecol}
\qtypebadge{singlecol}{Single Choice} \hfill \textcolor{darkslate!70}{\small Topic: Intrusion Prevention}

Which detection method builds a baseline of normal behavior and flags deviations?

\begin{optlist}
\item Signature detection
\item Anomaly detection
\item Static analysis
\item Port scanning
\end{optlist}
\begin{explainbox}{}
\textbf{Correct Answer: B}\\[0.3em]
\textbf{Concept:} Anomaly detection identifies traffic that deviates from a learned baseline.
\end{explainbox}
\end{questioncard}

\begin{questioncard}{4}{singlecol}
\qtypebadge{singlecol}{Single Choice} \hfill \textcolor{darkslate!70}{\small Topic: Intrusion Prevention}

Where is an IPS typically deployed?

\begin{optlist}
\item Out-of-band only
\item Inline in the traffic path
\item On user desktops only
\item Inside a printer
\end{optlist}
\begin{explainbox}{}
\textbf{Correct Answer: B}\\[0.3em]
\textbf{Concept:} IPS is deployed inline so it can block malicious traffic in real time.
\end{explainbox}
\end{questioncard}

\section{Antivirus}
\begin{questioncard}{1}{singlecol}
\qtypebadge{singlecol}{Single Choice} \hfill \textcolor{darkslate!70}{\small Topic: Antivirus}

What does antivirus software primarily detect and remove?

\begin{optlist}
\item Routing loops
\item Malware
\item Cable faults
\item DNS errors
\end{optlist}
\begin{explainbox}{}
\textbf{Correct Answer: B}\\[0.3em]
\textbf{Concept:} Antivirus detects and removes malware such as viruses, worms, and Trojans.
\end{explainbox}
\end{questioncard}

\begin{questioncard}{2}{singlecol}
\qtypebadge{singlecol}{Single Choice} \hfill \textcolor{darkslate!70}{\small Topic: Antivirus}

Which antivirus technique looks for known byte sequences of malware?

\begin{optlist}
\item Heuristic analysis
\item Signature scanning
\item Sandboxing
\item Whitelisting
\end{optlist}
\begin{explainbox}{}
\textbf{Correct Answer: B}\\[0.3em]
\textbf{Concept:} Signature scanning matches files/traffic against known malware signatures.
\end{explainbox}
\end{questioncard}

\begin{questioncard}{3}{singlecol}
\qtypebadge{singlecol}{Single Choice} \hfill \textcolor{darkslate!70}{\small Topic: Antivirus}

Which method executes suspicious files in an isolated environment to observe behavior?

\begin{optlist}
\item Signature scanning
\item Heuristic analysis
\item Sandboxing
\item Whitelisting
\end{optlist}
\begin{explainbox}{}
\textbf{Correct Answer: C}\\[0.3em]
\textbf{Concept:} Sandboxing runs files in an isolated environment to detect malicious behavior.
\end{explainbox}
\end{questioncard}

\section{AAA and User Authentication}
\begin{questioncard}{1}{singlecol}
\qtypebadge{singlecol}{Single Choice} \hfill \textcolor{darkslate!70}{\small Topic: AAA and User Authentication}

Which AAA protocol commonly uses UDP ports 1812 and 1813?

\begin{optlist}
\item TACACS+
\item RADIUS
\item LDAP
\item Kerberos
\end{optlist}
\begin{explainbox}{}
\textbf{Correct Answer: B}\\[0.3em]
\textbf{Concept:} RADIUS commonly uses UDP ports 1812/1813 (or 1645/1646 legacy).
\end{explainbox}
\end{questioncard}

\begin{questioncard}{2}{singlecol}
\qtypebadge{singlecol}{Single Choice} \hfill \textcolor{darkslate!70}{\small Topic: AAA and User Authentication}

Which AAA protocol uses TCP port 49 by default?

\begin{optlist}
\item RADIUS
\item TACACS+
\item LDAP
\item DHCP
\end{optlist}
\begin{explainbox}{}
\textbf{Correct Answer: B}\\[0.3em]
\textbf{Concept:} TACACS+ uses TCP port 49 by default.
\end{explainbox}
\end{questioncard}

\begin{questioncard}{3}{singlecol}
\qtypebadge{singlecol}{Single Choice} \hfill \textcolor{darkslate!70}{\small Topic: AAA and User Authentication}

In AAA, which process maps users to privileges?

\begin{optlist}
\item Authentication
\item Authorization
\item Accounting
\item Admission
\end{optlist}
\begin{explainbox}{}
\textbf{Correct Answer: B}\\[0.3em]
\textbf{Concept:} Authorization maps authenticated identities to permitted actions/resources.
\end{explainbox}
\end{questioncard}

\begin{questioncard}{4}{singlecol}
\qtypebadge{singlecol}{Single Choice} \hfill \textcolor{darkslate!70}{\small Topic: AAA and User Authentication}

Which authentication type identifies users based on source IP?

\begin{optlist}
\item User authentication
\item Access user authentication
\item Device authentication
\item IP/MAC binding
\end{optlist}
\begin{explainbox}{}
\textbf{Correct Answer: B}\\[0.3em]
\textbf{Concept:} Access user authentication can identify users by IP/MAC for transparent access control.
\end{explainbox}
\end{questioncard}

\begin{questioncard}{5}{singlecol}
\qtypebadge{singlecol}{Single Choice} \hfill \textcolor{darkslate!70}{\small Topic: AAA and User Authentication}

Which user type is created directly on the firewall?

\begin{optlist}
\item RADIUS user
\item LDAP user
\item Local user
\item Domain user
\end{optlist}
\begin{explainbox}{}
\textbf{Correct Answer: C}\\[0.3em]
\textbf{Concept:} Local users are defined directly in the firewall local user database.
\end{explainbox}
\end{questioncard}

\begin{questioncard}{6}{singlecol}
\qtypebadge{singlecol}{Single Choice} \hfill \textcolor{darkslate!70}{\small Topic: AAA and User Authentication}

Which component is queried when using LDAP authentication?

\begin{optlist}
\item DNS server
\item LDAP directory server
\item NTP server
\item SNMP server
\end{optlist}
\begin{explainbox}{}
\textbf{Correct Answer: B}\\[0.3em]
\textbf{Concept:} The firewall queries an LDAP directory server to validate credentials.
\end{explainbox}
\end{questioncard}

\begin{questioncard}{7}{singlecol}
\qtypebadge{singlecol}{Single Choice} \hfill \textcolor{darkslate!70}{\small Topic: AAA and User Authentication}

What is the purpose of a user group?

\begin{optlist}
\item Increase Internet speed
\item Apply policies to multiple users collectively
\item Replace firewalls
\item Encrypt traffic
\end{optlist}
\begin{explainbox}{}
\textbf{Correct Answer: B}\\[0.3em]
\textbf{Concept:} User groups simplify policy management by grouping users with similar access needs.
\end{explainbox}
\end{questioncard}

\begin{questioncard}{8}{singlecol}
\qtypebadge{singlecol}{Single Choice} \hfill \textcolor{darkslate!70}{\small Topic: AAA and User Authentication}

Which authentication method is often used for wired/wireless network access control?

\begin{optlist}
\item Portal
\item 802.1X
\item Local
\item RADIUS accounting
\end{optlist}
\begin{explainbox}{}
\textbf{Correct Answer: B}\\[0.3em]
\textbf{Concept:} 802.1X is commonly used for port-based network access control on wired and wireless networks.
\end{explainbox}
\end{questioncard}

\begin{questioncard}{9}{singlecol}
\qtypebadge{singlecol}{Single Choice} \hfill \textcolor{darkslate!70}{\small Topic: AAA and User Authentication}

What does single sign-on (SSO) provide?

\begin{optlist}
\item Multiple passwords
\item One authentication for multiple services
\item No authentication
\item Local only access
\end{optlist}
\begin{explainbox}{}
\textbf{Correct Answer: B}\\[0.3em]
\textbf{Concept:} SSO allows users to authenticate once and access multiple services.
\end{explainbox}
\end{questioncard}

\begin{questioncard}{10}{singlecol}
\qtypebadge{singlecol}{Single Choice} \hfill \textcolor{darkslate!70}{\small Topic: AAA and User Authentication}

Which user authentication policy condition can match traffic?

\begin{optlist}
\item Source zone
\item Time range
\item Source/destination address
\item All of the above
\end{optlist}
\begin{explainbox}{}
\textbf{Correct Answer: D}\\[0.3em]
\textbf{Concept:} Authentication policies can match zones, addresses, time ranges, and other conditions.
\end{explainbox}
\end{questioncard}

\begin{questioncard}{11}{singlecol}
\qtypebadge{singlecol}{Single Choice} \hfill \textcolor{darkslate!70}{\small Topic: AAA and User Authentication}

The first A in AAA stands for:

\begin{optlist}
\item Authentication
\item Authorization
\item Accounting
\item Auditing
\end{optlist}
\begin{explainbox}{}
\textbf{Correct Answer: A}\\[0.3em]
\textbf{Concept:} AAA stands for Authentication, Authorization, and Accounting.
\end{explainbox}
\end{questioncard}

\begin{questioncard}{12}{singlecol}
\qtypebadge{singlecol}{Single Choice} \hfill \textcolor{darkslate!70}{\small Topic: AAA and User Authentication}

RADIUS commonly runs over:

\begin{optlist}
\item UDP
\item TCP
\item SCTP
\item ICMP
\end{optlist}
\begin{explainbox}{}
\textbf{Correct Answer: A}\\[0.3em]
\textbf{Concept:} RADIUS typically uses UDP ports 1812/1813.
\end{explainbox}
\end{questioncard}

\begin{questioncard}{13}{singlecol}
\qtypebadge{singlecol}{Single Choice} \hfill \textcolor{darkslate!70}{\small Topic: AAA and User Authentication}

TACACS+ commonly runs over:

\begin{optlist}
\item UDP
\item TCP
\item HTTP
\item FTP
\end{optlist}
\begin{explainbox}{}
\textbf{Correct Answer: B}\\[0.3em]
\textbf{Concept:} TACACS+ uses TCP port 49.
\end{explainbox}
\end{questioncard}

\begin{questioncard}{14}{singlecol}
\qtypebadge{singlecol}{Single Choice} \hfill \textcolor{darkslate!70}{\small Topic: AAA and User Authentication}

Portal authentication is typically used for:

\begin{optlist}
\item Web-based guest access
\item Cable modem registration
\item Switch MAC learning
\item DNS resolution
\end{optlist}
\begin{explainbox}{}
\textbf{Correct Answer: A}\\[0.3em]
\textbf{Concept:} Portal authentication presents a web login page.
\end{explainbox}
\end{questioncard}

\begin{questioncard}{15}{singlecol}
\qtypebadge{singlecol}{Single Choice} \hfill \textcolor{darkslate!70}{\small Topic: AAA and User Authentication}

802.1X authentication occurs at:

\begin{optlist}
\item Layer 2 port level
\item Application layer only
\item Physical layer only
\item Transport layer only
\end{optlist}
\begin{explainbox}{}
\textbf{Correct Answer: A}\\[0.3em]
\textbf{Concept:} 802.1X provides port-based network access control at Layer 2.
\end{explainbox}
\end{questioncard}

\begin{questioncard}{16}{singlecol}
\qtypebadge{singlecol}{Single Choice} \hfill \textcolor{darkslate!70}{\small Topic: AAA and User Authentication}

Which concept is described as: verifies user identity?

\begin{optlist}
\item Accounting
\item Authorization
\item RADIUS
\item Authentication
\end{optlist}
\begin{explainbox}{}
\textbf{Correct Answer: D}\\[0.3em]
\textbf{Concept:} Authentication refers to the concept where verifies user identity.
\end{explainbox}
\end{questioncard}

\begin{questioncard}{17}{singlecol}
\qtypebadge{singlecol}{Single Choice} \hfill \textcolor{darkslate!70}{\small Topic: AAA and User Authentication}

Which concept is described as: determines what a user can access?

\begin{optlist}
\item Accounting
\item Portal authentication
\item Local authentication
\item Authorization
\end{optlist}
\begin{explainbox}{}
\textbf{Correct Answer: D}\\[0.3em]
\textbf{Concept:} Authorization refers to the concept where determines what a user can access.
\end{explainbox}
\end{questioncard}

\begin{questioncard}{18}{singlecol}
\qtypebadge{singlecol}{Single Choice} \hfill \textcolor{darkslate!70}{\small Topic: AAA and User Authentication}

Which concept is described as: records user activities and resource usage?

\begin{optlist}
\item Accounting
\item Authorization
\item User group
\item Portal authentication
\end{optlist}
\begin{explainbox}{}
\textbf{Correct Answer: A}\\[0.3em]
\textbf{Concept:} Accounting refers to the concept where records user activities and resource usage.
\end{explainbox}
\end{questioncard}

\begin{questioncard}{19}{singlecol}
\qtypebadge{singlecol}{Single Choice} \hfill \textcolor{darkslate!70}{\small Topic: AAA and User Authentication}

Which concept is described as: uses UDP ports 1812 and 1813?

\begin{optlist}
\item Portal authentication
\item RADIUS
\item SSO
\item Local authentication
\end{optlist}
\begin{explainbox}{}
\textbf{Correct Answer: B}\\[0.3em]
\textbf{Concept:} RADIUS refers to the concept where uses UDP ports 1812 and 1813.
\end{explainbox}
\end{questioncard}

\begin{questioncard}{20}{singlecol}
\qtypebadge{singlecol}{Single Choice} \hfill \textcolor{darkslate!70}{\small Topic: AAA and User Authentication}

Which concept is described as: uses TCP port 49?

\begin{optlist}
\item RADIUS
\item TACACS+
\item 802.1X
\item Authentication
\end{optlist}
\begin{explainbox}{}
\textbf{Correct Answer: B}\\[0.3em]
\textbf{Concept:} TACACS+ refers to the concept where uses TCP port 49.
\end{explainbox}
\end{questioncard}

\begin{questioncard}{21}{singlecol}
\qtypebadge{singlecol}{Single Choice} \hfill \textcolor{darkslate!70}{\small Topic: AAA and User Authentication}

Which concept is described as: uses a web login page?

\begin{optlist}
\item Authorization
\item User group
\item Local authentication
\item Portal authentication
\end{optlist}
\begin{explainbox}{}
\textbf{Correct Answer: D}\\[0.3em]
\textbf{Concept:} Portal authentication refers to the concept where uses a web login page.
\end{explainbox}
\end{questioncard}

\begin{questioncard}{22}{singlecol}
\qtypebadge{singlecol}{Single Choice} \hfill \textcolor{darkslate!70}{\small Topic: AAA and User Authentication}

Which concept is described as: provides port-based network access control?

\begin{optlist}
\item User group
\item Accounting
\item SSO
\item 802.1X
\end{optlist}
\begin{explainbox}{}
\textbf{Correct Answer: D}\\[0.3em]
\textbf{Concept:} 802.1X refers to the concept where provides port-based network access control.
\end{explainbox}
\end{questioncard}

\begin{questioncard}{23}{singlecol}
\qtypebadge{singlecol}{Single Choice} \hfill \textcolor{darkslate!70}{\small Topic: AAA and User Authentication}

Which concept is described as: uses the device user database?

\begin{optlist}
\item Authentication
\item LDAP
\item Accounting
\item Local authentication
\end{optlist}
\begin{explainbox}{}
\textbf{Correct Answer: D}\\[0.3em]
\textbf{Concept:} Local authentication refers to the concept where uses the device user database.
\end{explainbox}
\end{questioncard}

\begin{questioncard}{24}{singlecol}
\qtypebadge{singlecol}{Single Choice} \hfill \textcolor{darkslate!70}{\small Topic: AAA and User Authentication}

Which concept is described as: queries a directory server for authentication?

\begin{optlist}
\item 802.1X
\item RADIUS
\item LDAP
\item Authentication
\end{optlist}
\begin{explainbox}{}
\textbf{Correct Answer: C}\\[0.3em]
\textbf{Concept:} LDAP refers to the concept where queries a directory server for authentication.
\end{explainbox}
\end{questioncard}

\begin{questioncard}{25}{singlecol}
\qtypebadge{singlecol}{Single Choice} \hfill \textcolor{darkslate!70}{\small Topic: AAA and User Authentication}

Which concept is described as: applies policies to multiple users collectively?

\begin{optlist}
\item Portal authentication
\item User group
\item 802.1X
\item TACACS+
\end{optlist}
\begin{explainbox}{}
\textbf{Correct Answer: B}\\[0.3em]
\textbf{Concept:} User group refers to the concept where applies policies to multiple users collectively.
\end{explainbox}
\end{questioncard}

\begin{questioncard}{26}{singlecol}
\qtypebadge{singlecol}{Single Choice} \hfill \textcolor{darkslate!70}{\small Topic: AAA and User Authentication}

Which concept is described as: allows one authentication for multiple services?

\begin{optlist}
\item 802.1X
\item LDAP
\item SSO
\item RADIUS
\end{optlist}
\begin{explainbox}{}
\textbf{Correct Answer: C}\\[0.3em]
\textbf{Concept:} SSO refers to the concept where allows one authentication for multiple services.
\end{explainbox}
\end{questioncard}

\section{AAA Basics}
\begin{questioncard}{1}{singlecol}
\qtypebadge{singlecol}{Single Choice} \hfill \textcolor{darkslate!70}{\small Topic: AAA Basics}

What does AAA stand for?

\begin{optlist}
\item Authentication, Authorization, Accounting
\item Access, Audit, Alert
\item Attack, Analysis, Action
\item Application, Authorization, Audit
\end{optlist}
\begin{explainbox}{}
\textbf{Correct Answer: A}\\[0.3em]
\textbf{Concept:} AAA provides authentication (who are you), authorization (what can you do), and accounting (what did you do).
\end{explainbox}
\end{questioncard}

\begin{questioncard}{2}{singlecol}
\qtypebadge{singlecol}{Single Choice} \hfill \textcolor{darkslate!70}{\small Topic: AAA Basics}

Which AAA component verifies a user identity?

\begin{optlist}
\item Authentication
\item Authorization
\item Accounting
\item Auditing
\end{optlist}
\begin{explainbox}{}
\textbf{Correct Answer: A}\\[0.3em]
\textbf{Concept:} Authentication verifies identity, for example, by username/password or certificate.
\end{explainbox}
\end{questioncard}

\begin{questioncard}{3}{singlecol}
\qtypebadge{singlecol}{Single Choice} \hfill \textcolor{darkslate!70}{\small Topic: AAA Basics}

Which AAA component determines what resources a user can access?

\begin{optlist}
\item Authentication
\item Authorization
\item Accounting
\item Auditing
\end{optlist}
\begin{explainbox}{}
\textbf{Correct Answer: B}\\[0.3em]
\textbf{Concept:} Authorization defines permissions and access rights after authentication.
\end{explainbox}
\end{questioncard}

\begin{questioncard}{4}{singlecol}
\qtypebadge{singlecol}{Single Choice} \hfill \textcolor{darkslate!70}{\small Topic: AAA Basics}

Which AAA component records user activities for billing or auditing?

\begin{optlist}
\item Authentication
\item Authorization
\item Accounting
\item Auditing
\end{optlist}
\begin{explainbox}{}
\textbf{Correct Answer: C}\\[0.3em]
\textbf{Concept:} Accounting logs user activities such as login time, data usage, and commands executed.
\end{explainbox}
\end{questioncard}

\section{Firewall User Authentication}
\begin{questioncard}{1}{singlecol}
\qtypebadge{singlecol}{Single Choice} \hfill \textcolor{darkslate!70}{\small Topic: Firewall User Authentication}

Which authentication method redirects a user to a web portal for login?

\begin{optlist}
\item Local authentication
\item RADIUS authentication
\item Portal authentication
\item LDAP authentication
\end{optlist}
\begin{explainbox}{}
\textbf{Correct Answer: C}\\[0.3em]
\textbf{Concept:} Portal authentication presents a web page to users for credential entry.
\end{explainbox}
\end{questioncard}

\begin{questioncard}{2}{singlecol}
\qtypebadge{singlecol}{Single Choice} \hfill \textcolor{darkslate!70}{\small Topic: Firewall User Authentication}

Which authentication method uses a client-side agent to authenticate users before network access?

\begin{optlist}
\item Portal authentication
\item 802.1X authentication
\item Local authentication
\item Certificate authentication
\end{optlist}
\begin{explainbox}{}
\textbf{Correct Answer: B}\\[0.3em]
\textbf{Concept:} 802.1X uses a supplicant client to authenticate devices before granting network access.
\end{explainbox}
\end{questioncard}

\begin{questioncard}{3}{singlecol}
\qtypebadge{singlecol}{Single Choice} \hfill \textcolor{darkslate!70}{\small Topic: Firewall User Authentication}

Which protocol is commonly used for centralized authentication and accounting?

\begin{optlist}
\item RADIUS
\item SMTP
\item SNMP
\item NTP
\end{optlist}
\begin{explainbox}{}
\textbf{Correct Answer: A}\\[0.3em]
\textbf{Concept:} RADIUS is a widely used AAA protocol for centralized authentication, authorization, and accounting.
\end{explainbox}
\end{questioncard}

\begin{questioncard}{4}{singlecol}
\qtypebadge{singlecol}{Single Choice} \hfill \textcolor{darkslate!70}{\small Topic: Firewall User Authentication}

Which protocol provides authentication with stronger security and uses TCP?

\begin{optlist}
\item RADIUS
\item TACACS+
\item LDAP
\item DHCP
\end{optlist}
\begin{explainbox}{}
\textbf{Correct Answer: B}\\[0.3em]
\textbf{Concept:} TACACS+ uses TCP and separates authentication, authorization, and accounting functions.
\end{explainbox}
\end{questioncard}

\section{Cryptography and PKI}
\begin{questioncard}{1}{singlecol}
\qtypebadge{singlecol}{Single Choice} \hfill \textcolor{darkslate!70}{\small Topic: Cryptography and PKI}

Which symmetric algorithm uses variable key lengths of 128, 192, or 256 bits?

\begin{optlist}
\item DES
\item AES
\item RSA
\item MD5
\end{optlist}
\begin{explainbox}{}
\textbf{Correct Answer: B}\\[0.3em]
\textbf{Concept:} AES supports 128, 192, and 256-bit keys.
\end{explainbox}
\end{questioncard}

\begin{questioncard}{2}{singlecol}
\qtypebadge{singlecol}{Single Choice} \hfill \textcolor{darkslate!70}{\small Topic: Cryptography and PKI}

How many bits is a DES key?

\begin{optlist}
\item 56
\item 64
\item 128
\item 256
\end{optlist}
\begin{explainbox}{}
\textbf{Correct Answer: A}\\[0.3em]
\textbf{Concept:} DES uses a 56-bit effective key length.
\end{explainbox}
\end{questioncard}

\begin{questioncard}{3}{singlecol}
\qtypebadge{singlecol}{Single Choice} \hfill \textcolor{darkslate!70}{\small Topic: Cryptography and PKI}

Which asymmetric algorithm is commonly used for key exchange in TLS?

\begin{optlist}
\item AES
\item RSA
\item SHA-256
\item HMAC
\end{optlist}
\begin{explainbox}{}
\textbf{Correct Answer: B}\\[0.3em]
\textbf{Concept:} RSA or Diffie-Hellman are commonly used for TLS key exchange.
\end{explainbox}
\end{questioncard}

\begin{questioncard}{4}{singlecol}
\qtypebadge{singlecol}{Single Choice} \hfill \textcolor{darkslate!70}{\small Topic: Cryptography and PKI}

What is the main disadvantage of symmetric encryption?

\begin{optlist}
\item Slow speed
\item Key distribution challenge
\item Large key sizes
\item No confidentiality
\end{optlist}
\begin{explainbox}{}
\textbf{Correct Answer: B}\\[0.3em]
\textbf{Concept:} Distributing the shared secret key securely is the main challenge.
\end{explainbox}
\end{questioncard}

\begin{questioncard}{5}{singlecol}
\qtypebadge{singlecol}{Single Choice} \hfill \textcolor{darkslate!70}{\small Topic: Cryptography and PKI}

Which property ensures a small input change produces a very different hash?

\begin{optlist}
\item Determinism
\item Avalanche effect
\item Fixed output
\item Reversibility
\end{optlist}
\begin{explainbox}{}
\textbf{Correct Answer: B}\\[0.3em]
\textbf{Concept:} The avalanche effect means tiny input changes cause large output changes.
\end{explainbox}
\end{questioncard}

\begin{questioncard}{6}{singlecol}
\qtypebadge{singlecol}{Single Choice} \hfill \textcolor{darkslate!70}{\small Topic: Cryptography and PKI}

Which hash algorithm produces a 128-bit digest?

\begin{optlist}
\item MD5
\item SHA-1
\item SHA-256
\item SHA-512
\end{optlist}
\begin{explainbox}{}
\textbf{Correct Answer: A}\\[0.3em]
\textbf{Concept:} MD5 produces a 128-bit digest.
\end{explainbox}
\end{questioncard}

\begin{questioncard}{7}{singlecol}
\qtypebadge{singlecol}{Single Choice} \hfill \textcolor{darkslate!70}{\small Topic: Cryptography and PKI}

What is HMAC?

\begin{optlist}
\item A hash-based message authentication code
\item A symmetric cipher
\item A key exchange protocol
\item A firewall feature
\end{optlist}
\begin{explainbox}{}
\textbf{Correct Answer: A}\\[0.3em]
\textbf{Concept:} HMAC combines a hash function with a secret key to provide message authentication.
\end{explainbox}
\end{questioncard}

\begin{questioncard}{8}{singlecol}
\qtypebadge{singlecol}{Single Choice} \hfill \textcolor{darkslate!70}{\small Topic: Cryptography and PKI}

Which key is used to create a digital signature?

\begin{optlist}
\item Public key
\item Private key
\item Shared secret
\item Session key
\end{optlist}
\begin{explainbox}{}
\textbf{Correct Answer: B}\\[0.3em]
\textbf{Concept:} The signer uses their private key to create a digital signature.
\end{explainbox}
\end{questioncard}

\begin{questioncard}{9}{singlecol}
\qtypebadge{singlecol}{Single Choice} \hfill \textcolor{darkslate!70}{\small Topic: Cryptography and PKI}

Which key is used to verify a digital signature?

\begin{optlist}
\item Public key
\item Private key
\item Shared secret
\item Master key
\end{optlist}
\begin{explainbox}{}
\textbf{Correct Answer: A}\\[0.3em]
\textbf{Concept:} The verifier uses the signer public key to verify the signature.
\end{explainbox}
\end{questioncard}

\begin{questioncard}{10}{singlecol}
\qtypebadge{singlecol}{Single Choice} \hfill \textcolor{darkslate!70}{\small Topic: Cryptography and PKI}

Which standard defines the format of digital certificates?

\begin{optlist}
\item X.509
\item X.500
\item LDAP
\item DNS
\end{optlist}
\begin{explainbox}{}
\textbf{Correct Answer: A}\\[0.3em]
\textbf{Concept:} X.509 is the standard format for public key certificates.
\end{explainbox}
\end{questioncard}

\begin{questioncard}{11}{singlecol}
\qtypebadge{singlecol}{Single Choice} \hfill \textcolor{darkslate!70}{\small Topic: Cryptography and PKI}

What does OCSP provide?

\begin{optlist}
\item Online certificate status checking
\item Certificate issuance
\item Key generation
\item Email encryption
\end{optlist}
\begin{explainbox}{}
\textbf{Correct Answer: A}\\[0.3em]
\textbf{Concept:} OCSP checks certificate revocation status in real time.
\end{explainbox}
\end{questioncard}

\begin{questioncard}{12}{singlecol}
\qtypebadge{singlecol}{Single Choice} \hfill \textcolor{darkslate!70}{\small Topic: Cryptography and PKI}

In a PKI hierarchy, which certificate is self-signed?

\begin{optlist}
\item End-entity certificate
\item Root CA certificate
\item Intermediate CA certificate
\item Server certificate
\end{optlist}
\begin{explainbox}{}
\textbf{Correct Answer: B}\\[0.3em]
\textbf{Concept:} Root CA certificates are self-signed because they are the trust anchor.
\end{explainbox}
\end{questioncard}

\begin{questioncard}{13}{singlecol}
\qtypebadge{singlecol}{Single Choice} \hfill \textcolor{darkslate!70}{\small Topic: Cryptography and PKI}

Which entity stores issued certificates for retrieval?

\begin{optlist}
\item CA
\item RA
\item Repository
\item CRL
\end{optlist}
\begin{explainbox}{}
\textbf{Correct Answer: C}\\[0.3em]
\textbf{Concept:} A repository distributes certificates and CRLs to relying parties.
\end{explainbox}
\end{questioncard}

\begin{questioncard}{14}{singlecol}
\qtypebadge{singlecol}{Single Choice} \hfill \textcolor{darkslate!70}{\small Topic: Cryptography and PKI}

Which cipher mode provides both confidentiality and authentication?

\begin{optlist}
\item ECB
\item CBC
\item GCM
\item CTR
\end{optlist}
\begin{explainbox}{}
\textbf{Correct Answer: C}\\[0.3em]
\textbf{Concept:} Galois/Counter Mode (GCM) provides authenticated encryption.
\end{explainbox}
\end{questioncard}

\begin{questioncard}{15}{singlecol}
\qtypebadge{singlecol}{Single Choice} \hfill \textcolor{darkslate!70}{\small Topic: Cryptography and PKI}

Which mode should be avoided because identical plaintext blocks produce identical ciphertext?

\begin{optlist}
\item ECB
\item CBC
\item GCM
\item CFB
\end{optlist}
\begin{explainbox}{}
\textbf{Correct Answer: A}\\[0.3em]
\textbf{Concept:} ECB reveals patterns in plaintext and is generally insecure.
\end{explainbox}
\end{questioncard}

\begin{questioncard}{16}{singlecol}
\qtypebadge{singlecol}{Single Choice} \hfill \textcolor{darkslate!70}{\small Topic: Cryptography and PKI}

AES is an example of:

\begin{optlist}
\item Symmetric encryption
\item Asymmetric encryption
\item Hashing
\item Steganography
\end{optlist}
\begin{explainbox}{}
\textbf{Correct Answer: A}\\[0.3em]
\textbf{Concept:} AES uses a single shared key for encryption and decryption.
\end{explainbox}
\end{questioncard}

\begin{questioncard}{17}{singlecol}
\qtypebadge{singlecol}{Single Choice} \hfill \textcolor{darkslate!70}{\small Topic: Cryptography and PKI}

RSA is an example of:

\begin{optlist}
\item Asymmetric encryption
\item Symmetric encryption
\item Hash algorithm
\item Compression algorithm
\end{optlist}
\begin{explainbox}{}
\textbf{Correct Answer: A}\\[0.3em]
\textbf{Concept:} RSA uses public and private key pairs.
\end{explainbox}
\end{questioncard}

\begin{questioncard}{18}{singlecol}
\qtypebadge{singlecol}{Single Choice} \hfill \textcolor{darkslate!70}{\small Topic: Cryptography and PKI}

A hash function produces:

\begin{optlist}
\item A fixed-length digest
\item The original input
\item An encrypted file
\item A random key
\end{optlist}
\begin{explainbox}{}
\textbf{Correct Answer: A}\\[0.3em]
\textbf{Concept:} Hash functions map variable input to a fixed-size digest.
\end{explainbox}
\end{questioncard}

\begin{questioncard}{19}{singlecol}
\qtypebadge{singlecol}{Single Choice} \hfill \textcolor{darkslate!70}{\small Topic: Cryptography and PKI}

MD5 is considered insecure because:

\begin{optlist}
\item Collision vulnerabilities exist
\item It is too slow
\item It produces too long output
\item It requires a public key
\end{optlist}
\begin{explainbox}{}
\textbf{Correct Answer: A}\\[0.3em]
\textbf{Concept:} MD5 collision attacks make it unsuitable for security purposes.
\end{explainbox}
\end{questioncard}

\begin{questioncard}{20}{singlecol}
\qtypebadge{singlecol}{Single Choice} \hfill \textcolor{darkslate!70}{\small Topic: Cryptography and PKI}

A digital certificate binds a public key to:

\begin{optlist}
\item An identity
\item A password
\item A MAC address
\item A VLAN
\end{optlist}
\begin{explainbox}{}
\textbf{Correct Answer: A}\\[0.3em]
\textbf{Concept:} Certificates associate a public key with an entity identity.
\end{explainbox}
\end{questioncard}

\begin{questioncard}{21}{singlecol}
\qtypebadge{singlecol}{Single Choice} \hfill \textcolor{darkslate!70}{\small Topic: Cryptography and PKI}

Certificate revocation status can be checked using:

\begin{optlist}
\item CRL or OCSP
\item DNS only
\item DHCP only
\item SNMP only
\end{optlist}
\begin{explainbox}{}
\textbf{Correct Answer: A}\\[0.3em]
\textbf{Concept:} CRLs and OCSP provide revocation status information.
\end{explainbox}
\end{questioncard}

\begin{questioncard}{22}{singlecol}
\qtypebadge{singlecol}{Single Choice} \hfill \textcolor{darkslate!70}{\small Topic: Cryptography and PKI}

Which concept is described as: uses the same key for encryption and decryption?

\begin{optlist}
\item Root CA certificate
\item Public key
\item Private key
\item Symmetric encryption
\end{optlist}
\begin{explainbox}{}
\textbf{Correct Answer: D}\\[0.3em]
\textbf{Concept:} Symmetric encryption refers to the concept where uses the same key for encryption and decryption.
\end{explainbox}
\end{questioncard}

\begin{questioncard}{23}{singlecol}
\qtypebadge{singlecol}{Single Choice} \hfill \textcolor{darkslate!70}{\small Topic: Cryptography and PKI}

Which concept is described as: uses public and private key pairs?

\begin{optlist}
\item Asymmetric encryption
\item X.509
\item Root CA certificate
\item Hash function
\end{optlist}
\begin{explainbox}{}
\textbf{Correct Answer: A}\\[0.3em]
\textbf{Concept:} Asymmetric encryption refers to the concept where uses public and private key pairs.
\end{explainbox}
\end{questioncard}

\begin{questioncard}{24}{singlecol}
\qtypebadge{singlecol}{Single Choice} \hfill \textcolor{darkslate!70}{\small Topic: Cryptography and PKI}

Which concept is described as: is a symmetric encryption standard?

\begin{optlist}
\item GCM
\item SHA-256
\item AES
\item RA
\end{optlist}
\begin{explainbox}{}
\textbf{Correct Answer: C}\\[0.3em]
\textbf{Concept:} AES refers to the concept where is a symmetric encryption standard.
\end{explainbox}
\end{questioncard}

\begin{questioncard}{25}{singlecol}
\qtypebadge{singlecol}{Single Choice} \hfill \textcolor{darkslate!70}{\small Topic: Cryptography and PKI}

Which concept is described as: is an asymmetric encryption algorithm?

\begin{optlist}
\item PKI
\item Asymmetric encryption
\item RSA
\item HMAC
\end{optlist}
\begin{explainbox}{}
\textbf{Correct Answer: C}\\[0.3em]
\textbf{Concept:} RSA refers to the concept where is an asymmetric encryption algorithm.
\end{explainbox}
\end{questioncard}

\begin{questioncard}{26}{singlecol}
\qtypebadge{singlecol}{Single Choice} \hfill \textcolor{darkslate!70}{\small Topic: Cryptography and PKI}

Which concept is described as: uses a 56-bit effective key?

\begin{optlist}
\item CA
\item OCSP
\item DES
\item SHA-256
\end{optlist}
\begin{explainbox}{}
\textbf{Correct Answer: C}\\[0.3em]
\textbf{Concept:} DES refers to the concept where uses a 56-bit effective key.
\end{explainbox}
\end{questioncard}

\begin{questioncard}{27}{singlecol}
\qtypebadge{singlecol}{Single Choice} \hfill \textcolor{darkslate!70}{\small Topic: Cryptography and PKI}

Which concept is described as: produces a fixed-length digest?

\begin{optlist}
\item Root CA certificate
\item Hash function
\item CA
\item RA
\end{optlist}
\begin{explainbox}{}
\textbf{Correct Answer: B}\\[0.3em]
\textbf{Concept:} Hash function refers to the concept where produces a fixed-length digest.
\end{explainbox}
\end{questioncard}

\begin{questioncard}{28}{singlecol}
\qtypebadge{singlecol}{Single Choice} \hfill \textcolor{darkslate!70}{\small Topic: Cryptography and PKI}

Which concept is described as: is considered broken due to collision attacks?

\begin{optlist}
\item Hash function
\item MD5
\item AES
\item CRL
\end{optlist}
\begin{explainbox}{}
\textbf{Correct Answer: B}\\[0.3em]
\textbf{Concept:} MD5 refers to the concept where is considered broken due to collision attacks.
\end{explainbox}
\end{questioncard}

\begin{questioncard}{29}{singlecol}
\qtypebadge{singlecol}{Single Choice} \hfill \textcolor{darkslate!70}{\small Topic: Cryptography and PKI}

Which concept is described as: produces a 256-bit digest?

\begin{optlist}
\item Digital signature
\item RA
\item Hash function
\item SHA-256
\end{optlist}
\begin{explainbox}{}
\textbf{Correct Answer: D}\\[0.3em]
\textbf{Concept:} SHA-256 refers to the concept where produces a 256-bit digest.
\end{explainbox}
\end{questioncard}

\begin{questioncard}{30}{singlecol}
\qtypebadge{singlecol}{Single Choice} \hfill \textcolor{darkslate!70}{\small Topic: Cryptography and PKI}

Which concept is described as: provides message authentication using a hash and key?

\begin{optlist}
\item PKI
\item CRL
\item HMAC
\item Hash function
\end{optlist}
\begin{explainbox}{}
\textbf{Correct Answer: C}\\[0.3em]
\textbf{Concept:} HMAC refers to the concept where provides message authentication using a hash and key.
\end{explainbox}
\end{questioncard}

\begin{questioncard}{31}{singlecol}
\qtypebadge{singlecol}{Single Choice} \hfill \textcolor{darkslate!70}{\small Topic: Cryptography and PKI}

Which concept is described as: verifies authenticity and integrity?

\begin{optlist}
\item Public key
\item OCSP
\item Digital signature
\item Asymmetric encryption
\end{optlist}
\begin{explainbox}{}
\textbf{Correct Answer: C}\\[0.3em]
\textbf{Concept:} Digital signature refers to the concept where verifies authenticity and integrity.
\end{explainbox}
\end{questioncard}

\begin{questioncard}{32}{singlecol}
\qtypebadge{singlecol}{Single Choice} \hfill \textcolor{darkslate!70}{\small Topic: Cryptography and PKI}

Which concept is described as: is used to verify a digital signature?

\begin{optlist}
\item Root CA certificate
\item Public key
\item HMAC
\item Hash function
\end{optlist}
\begin{explainbox}{}
\textbf{Correct Answer: B}\\[0.3em]
\textbf{Concept:} Public key refers to the concept where is used to verify a digital signature.
\end{explainbox}
\end{questioncard}

\begin{questioncard}{33}{singlecol}
\qtypebadge{singlecol}{Single Choice} \hfill \textcolor{darkslate!70}{\small Topic: Cryptography and PKI}

Which concept is described as: is used to create a digital signature?

\begin{optlist}
\item Asymmetric encryption
\item CA
\item Private key
\item OCSP
\end{optlist}
\begin{explainbox}{}
\textbf{Correct Answer: C}\\[0.3em]
\textbf{Concept:} Private key refers to the concept where is used to create a digital signature.
\end{explainbox}
\end{questioncard}

\begin{questioncard}{34}{singlecol}
\qtypebadge{singlecol}{Single Choice} \hfill \textcolor{darkslate!70}{\small Topic: Cryptography and PKI}

Which concept is described as: manages public keys and certificates?

\begin{optlist}
\item PKI
\item HMAC
\item Public key
\item OCSP
\end{optlist}
\begin{explainbox}{}
\textbf{Correct Answer: A}\\[0.3em]
\textbf{Concept:} PKI refers to the concept where manages public keys and certificates.
\end{explainbox}
\end{questioncard}

\begin{questioncard}{35}{singlecol}
\qtypebadge{singlecol}{Single Choice} \hfill \textcolor{darkslate!70}{\small Topic: Cryptography and PKI}

Which concept is described as: issues and signs digital certificates?

\begin{optlist}
\item CA
\item Symmetric encryption
\item CRL
\item SHA-256
\end{optlist}
\begin{explainbox}{}
\textbf{Correct Answer: A}\\[0.3em]
\textbf{Concept:} CA refers to the concept where issues and signs digital certificates.
\end{explainbox}
\end{questioncard}

\begin{questioncard}{36}{singlecol}
\qtypebadge{singlecol}{Single Choice} \hfill \textcolor{darkslate!70}{\small Topic: Cryptography and PKI}

Which concept is described as: validates identity before certificate issuance?

\begin{optlist}
\item Public key
\item RA
\item Private key
\item PKI
\end{optlist}
\begin{explainbox}{}
\textbf{Correct Answer: B}\\[0.3em]
\textbf{Concept:} RA refers to the concept where validates identity before certificate issuance.
\end{explainbox}
\end{questioncard}

\begin{questioncard}{37}{singlecol}
\qtypebadge{singlecol}{Single Choice} \hfill \textcolor{darkslate!70}{\small Topic: Cryptography and PKI}

Which concept is described as: lists revoked certificates?

\begin{optlist}
\item HMAC
\item CRL
\item Digital signature
\item ECB
\end{optlist}
\begin{explainbox}{}
\textbf{Correct Answer: B}\\[0.3em]
\textbf{Concept:} CRL refers to the concept where lists revoked certificates.
\end{explainbox}
\end{questioncard}

\begin{questioncard}{38}{singlecol}
\qtypebadge{singlecol}{Single Choice} \hfill \textcolor{darkslate!70}{\small Topic: Cryptography and PKI}

Which concept is described as: provides online certificate status checking?

\begin{optlist}
\item X.509
\item OCSP
\item Hash function
\item HMAC
\end{optlist}
\begin{explainbox}{}
\textbf{Correct Answer: B}\\[0.3em]
\textbf{Concept:} OCSP refers to the concept where provides online certificate status checking.
\end{explainbox}
\end{questioncard}

\begin{questioncard}{39}{singlecol}
\qtypebadge{singlecol}{Single Choice} \hfill \textcolor{darkslate!70}{\small Topic: Cryptography and PKI}

Which concept is described as: is the standard certificate format?

\begin{optlist}
\item AES
\item X.509
\item Symmetric encryption
\item Digital signature
\end{optlist}
\begin{explainbox}{}
\textbf{Correct Answer: B}\\[0.3em]
\textbf{Concept:} X.509 refers to the concept where is the standard certificate format.
\end{explainbox}
\end{questioncard}

\begin{questioncard}{40}{singlecol}
\qtypebadge{singlecol}{Single Choice} \hfill \textcolor{darkslate!70}{\small Topic: Cryptography and PKI}

Which concept is described as: is self-signed and acts as a trust anchor?

\begin{optlist}
\item Root CA certificate
\item Digital signature
\item ECB
\item CA
\end{optlist}
\begin{explainbox}{}
\textbf{Correct Answer: A}\\[0.3em]
\textbf{Concept:} Root CA certificate refers to the concept where is self-signed and acts as a trust anchor.
\end{explainbox}
\end{questioncard}

\begin{questioncard}{41}{singlecol}
\qtypebadge{singlecol}{Single Choice} \hfill \textcolor{darkslate!70}{\small Topic: Cryptography and PKI}

Which concept is described as: provides authenticated encryption?

\begin{optlist}
\item GCM
\item Root CA certificate
\item X.509
\item OCSP
\end{optlist}
\begin{explainbox}{}
\textbf{Correct Answer: A}\\[0.3em]
\textbf{Concept:} GCM refers to the concept where provides authenticated encryption.
\end{explainbox}
\end{questioncard}

\begin{questioncard}{42}{singlecol}
\qtypebadge{singlecol}{Single Choice} \hfill \textcolor{darkslate!70}{\small Topic: Cryptography and PKI}

Which concept is described as: should be avoided because it reveals plaintext patterns?

\begin{optlist}
\item RA
\item Root CA certificate
\item DES
\item ECB
\end{optlist}
\begin{explainbox}{}
\textbf{Correct Answer: D}\\[0.3em]
\textbf{Concept:} ECB refers to the concept where should be avoided because it reveals plaintext patterns.
\end{explainbox}
\end{questioncard}

\section{Cryptography Basics}
\begin{questioncard}{1}{singlecol}
\qtypebadge{singlecol}{Single Choice} \hfill \textcolor{darkslate!70}{\small Topic: Cryptography Basics}

Which encryption type uses the same key for encryption and decryption?

\begin{optlist}
\item Symmetric encryption
\item Asymmetric encryption
\item Hashing
\item Steganography
\end{optlist}
\begin{explainbox}{}
\textbf{Correct Answer: A}\\[0.3em]
\textbf{Concept:} Symmetric encryption uses a single shared key for both encryption and decryption.
\end{explainbox}
\end{questioncard}

\begin{questioncard}{2}{singlecol}
\qtypebadge{singlecol}{Single Choice} \hfill \textcolor{darkslate!70}{\small Topic: Cryptography Basics}

Which encryption type uses a key pair: public key and private key?

\begin{optlist}
\item Symmetric encryption
\item Asymmetric encryption
\item Hashing
\item Encoding
\end{optlist}
\begin{explainbox}{}
\textbf{Correct Answer: B}\\[0.3em]
\textbf{Concept:} Asymmetric encryption uses a public key to encrypt and a private key to decrypt (or vice versa).
\end{explainbox}
\end{questioncard}

\begin{questioncard}{3}{singlecol}
\qtypebadge{singlecol}{Single Choice} \hfill \textcolor{darkslate!70}{\small Topic: Cryptography Basics}

Which algorithm is a symmetric encryption standard?

\begin{optlist}
\item RSA
\item AES
\item ECC
\item DSA
\end{optlist}
\begin{explainbox}{}
\textbf{Correct Answer: B}\\[0.3em]
\textbf{Concept:} AES (Advanced Encryption Standard) is a widely used symmetric encryption algorithm.
\end{explainbox}
\end{questioncard}

\begin{questioncard}{4}{singlecol}
\qtypebadge{singlecol}{Single Choice} \hfill \textcolor{darkslate!70}{\small Topic: Cryptography Basics}

Which algorithm is an asymmetric encryption algorithm?

\begin{optlist}
\item AES
\item DES
\item RSA
\item 3DES
\end{optlist}
\begin{explainbox}{}
\textbf{Correct Answer: C}\\[0.3em]
\textbf{Concept:} RSA is an asymmetric algorithm used for encryption, digital signatures, and key exchange.
\end{explainbox}
\end{questioncard}

\section{Hash Algorithms}
\begin{questioncard}{1}{singlecol}
\qtypebadge{singlecol}{Single Choice} \hfill \textcolor{darkslate!70}{\small Topic: Hash Algorithms}

What is a key property of a cryptographic hash function?

\begin{optlist}
\item Reversible output
\item Fixed-length output regardless of input size
\item Requires a public key
\item Uses the same key for hashing and verification
\end{optlist}
\begin{explainbox}{}
\textbf{Correct Answer: B}\\[0.3em]
\textbf{Concept:} Hash functions produce a fixed-length digest for any input size.
\end{explainbox}
\end{questioncard}

\begin{questioncard}{2}{singlecol}
\qtypebadge{singlecol}{Single Choice} \hfill \textcolor{darkslate!70}{\small Topic: Hash Algorithms}

Which hash algorithm produces a 256-bit digest?

\begin{optlist}
\item MD5
\item SHA-1
\item SHA-256
\item SHA-512
\end{optlist}
\begin{explainbox}{}
\textbf{Correct Answer: C}\\[0.3em]
\textbf{Concept:} SHA-256 produces a 256-bit hash digest.
\end{explainbox}
\end{questioncard}

\begin{questioncard}{3}{singlecol}
\qtypebadge{singlecol}{Single Choice} \hfill \textcolor{darkslate!70}{\small Topic: Hash Algorithms}

Which hash algorithm is considered broken and unsuitable for security use?

\begin{optlist}
\item SHA-256
\item MD5
\item SHA-3
\item BLAKE2
\end{optlist}
\begin{explainbox}{}
\textbf{Correct Answer: B}\\[0.3em]
\textbf{Concept:} MD5 is vulnerable to collision attacks and should not be used for security-sensitive applications.
\end{explainbox}
\end{questioncard}

\section{Digital Signatures}
\begin{questioncard}{1}{singlecol}
\qtypebadge{singlecol}{Single Choice} \hfill \textcolor{darkslate!70}{\small Topic: Digital Signatures}

What is the purpose of a digital signature?

\begin{optlist}
\item Encrypt all data
\item Verify authenticity and integrity
\item Compress files
\item Hide sender identity
\end{optlist}
\begin{explainbox}{}
\textbf{Correct Answer: B}\\[0.3em]
\textbf{Concept:} Digital signatures verify the sender identity and ensure data has not been altered.
\end{explainbox}
\end{questioncard}

\section{PKI}
\begin{questioncard}{1}{singlecol}
\qtypebadge{singlecol}{Single Choice} \hfill \textcolor{darkslate!70}{\small Topic: PKI}

What does PKI stand for?

\begin{optlist}
\item Public Key Infrastructure
\item Private Key Interface
\item Packet Knowledge Index
\item Policy Key Identifier
\end{optlist}
\begin{explainbox}{}
\textbf{Correct Answer: A}\\[0.3em]
\textbf{Concept:} PKI provides a framework for managing public keys, certificates, and trust relationships.
\end{explainbox}
\end{questioncard}

\begin{questioncard}{2}{singlecol}
\qtypebadge{singlecol}{Single Choice} \hfill \textcolor{darkslate!70}{\small Topic: PKI}

Which entity issues digital certificates in a PKI?

\begin{optlist}
\item Registration Authority
\item Certificate Authority
\item End user
\item Repository
\end{optlist}
\begin{explainbox}{}
\textbf{Correct Answer: B}\\[0.3em]
\textbf{Concept:} A Certificate Authority (CA) issues and signs digital certificates.
\end{explainbox}
\end{questioncard}

\begin{questioncard}{3}{singlecol}
\qtypebadge{singlecol}{Single Choice} \hfill \textcolor{darkslate!70}{\small Topic: PKI}

Which PKI component validates an applicant identity before certificate issuance?

\begin{optlist}
\item Certificate Authority
\item Registration Authority
\item Repository
\item CRL
\end{optlist}
\begin{explainbox}{}
\textbf{Correct Answer: B}\\[0.3em]
\textbf{Concept:} A Registration Authority (RA) verifies identities and forwards requests to the CA.
\end{explainbox}
\end{questioncard}

\begin{questioncard}{4}{singlecol}
\qtypebadge{singlecol}{Single Choice} \hfill \textcolor{darkslate!70}{\small Topic: PKI}

What is a CRL used for?

\begin{optlist}
\item Storing public keys
\item Listing revoked certificates
\item Encrypting emails
\item Generating hash values
\end{optlist}
\begin{explainbox}{}
\textbf{Correct Answer: B}\\[0.3em]
\textbf{Concept:} A Certificate Revocation List (CRL) contains certificates that have been revoked before expiration.
\end{explainbox}
\end{questioncard}

\section{VPN}
\begin{questioncard}{1}{singlecol}
\qtypebadge{singlecol}{Single Choice} \hfill \textcolor{darkslate!70}{\small Topic: VPN}

Which VPN deployment connects two fixed sites?

\begin{optlist}
\item Remote access VPN
\item Site-to-site VPN
\item Clientless VPN
\item Host-to-host VPN
\end{optlist}
\begin{explainbox}{}
\textbf{Correct Answer: B}\\[0.3em]
\textbf{Concept:} Site-to-site VPNs connect entire networks at different locations.
\end{explainbox}
\end{questioncard}

\begin{questioncard}{2}{singlecol}
\qtypebadge{singlecol}{Single Choice} \hfill \textcolor{darkslate!70}{\small Topic: VPN}

What is a key advantage of GRE?

\begin{optlist}
\item Built-in encryption
\item Can encapsulate multicast and non-IP protocols
\item Uses TCP
\item Automatic key exchange
\end{optlist}
\begin{explainbox}{}
\textbf{Correct Answer: B}\\[0.3em]
\textbf{Concept:} GRE can tunnel multicast and non-IP traffic, making it useful for routing protocols.
\end{explainbox}
\end{questioncard}

\begin{questioncard}{3}{singlecol}
\qtypebadge{singlecol}{Single Choice} \hfill \textcolor{darkslate!70}{\small Topic: VPN}

Which protocol number is assigned to GRE?

\begin{optlist}
\item 47
\item 50
\item 51
\item 89
\end{optlist}
\begin{explainbox}{}
\textbf{Correct Answer: A}\\[0.3em]
\textbf{Concept:} GRE is IP protocol 47.
\end{explainbox}
\end{questioncard}

\begin{questioncard}{4}{singlecol}
\qtypebadge{singlecol}{Single Choice} \hfill \textcolor{darkslate!70}{\small Topic: VPN}

What does IKE Phase 1 establish?

\begin{optlist}
\item IPsec SA
\item IKE SA
\item Routing table
\item DHCP lease
\end{optlist}
\begin{explainbox}{}
\textbf{Correct Answer: B}\\[0.3em]
\textbf{Concept:} IKE Phase 1 establishes a secure IKE SA for protected negotiation.
\end{explainbox}
\end{questioncard}

\begin{questioncard}{5}{singlecol}
\qtypebadge{singlecol}{Single Choice} \hfill \textcolor{darkslate!70}{\small Topic: VPN}

What does IKE Phase 2 establish?

\begin{optlist}
\item IKE SA
\item IPsec SA
\item SSL session
\item ARP entry
\end{optlist}
\begin{explainbox}{}
\textbf{Correct Answer: B}\\[0.3em]
\textbf{Concept:} IKE Phase 2 negotiates the IPsec SA that protects actual data traffic.
\end{explainbox}
\end{questioncard}

\begin{questioncard}{6}{singlecol}
\qtypebadge{singlecol}{Single Choice} \hfill \textcolor{darkslate!70}{\small Topic: VPN}

Which authentication method does IKE support? (single best answer)

\begin{optlist}
\item Pre-shared key
\item Digital signatures
\item Both pre-shared key and digital signatures
\item None
\end{optlist}
\begin{explainbox}{}
\textbf{Correct Answer: C}\\[0.3em]
\textbf{Concept:} IKE supports both pre-shared key and digital signature authentication.
\end{explainbox}
\end{questioncard}

\begin{questioncard}{7}{singlecol}
\qtypebadge{singlecol}{Single Choice} \hfill \textcolor{darkslate!70}{\small Topic: VPN}

Which Diffie-Hellman group provides stronger key exchange?

\begin{optlist}
\item Group 1
\item Group 2
\item Group 5
\item Group 14 or higher
\end{optlist}
\begin{explainbox}{}
\textbf{Correct Answer: D}\\[0.3em]
\textbf{Concept:} Higher Diffie-Hellman groups such as 14 or above provide stronger security.
\end{explainbox}
\end{questioncard}

\begin{questioncard}{8}{singlecol}
\qtypebadge{singlecol}{Single Choice} \hfill \textcolor{darkslate!70}{\small Topic: VPN}

Which IKE version provides faster SA setup and improved reliability?

\begin{optlist}
\item IKEv1
\item IKEv2
\item IKEv3
\item ISAKMP
\end{optlist}
\begin{explainbox}{}
\textbf{Correct Answer: B}\\[0.3em]
\textbf{Concept:} IKEv2 simplifies negotiation and improves resilience compared to IKEv1.
\end{explainbox}
\end{questioncard}

\begin{questioncard}{9}{singlecol}
\qtypebadge{singlecol}{Single Choice} \hfill \textcolor{darkslate!70}{\small Topic: VPN}

Which UDP ports does L2TP use?

\begin{optlist}
\item 500 and 4500
\item 1701
\item 443
\item 50 and 51
\end{optlist}
\begin{explainbox}{}
\textbf{Correct Answer: B}\\[0.3em]
\textbf{Concept:} L2TP uses UDP port 1701.
\end{explainbox}
\end{questioncard}

\begin{questioncard}{10}{singlecol}
\qtypebadge{singlecol}{Single Choice} \hfill \textcolor{darkslate!70}{\small Topic: VPN}

Why is L2TP often combined with IPsec?

\begin{optlist}
\item To provide encryption
\item To increase speed
\item To reduce overhead
\item To avoid tunnels
\end{optlist}
\begin{explainbox}{}
\textbf{Correct Answer: A}\\[0.3em]
\textbf{Concept:} L2TP lacks encryption, so IPsec is added for confidentiality and integrity.
\end{explainbox}
\end{questioncard}

\begin{questioncard}{11}{singlecol}
\qtypebadge{singlecol}{Single Choice} \hfill \textcolor{darkslate!70}{\small Topic: VPN}

Which SSL VPN mode grants the remote host a virtual IP on the internal network?

\begin{optlist}
\item Web proxy
\item Network extension
\item File sharing
\item Port forwarding
\end{optlist}
\begin{explainbox}{}
\textbf{Correct Answer: B}\\[0.3em]
\textbf{Concept:} Network extension mode assigns a virtual IP and routes traffic into the corporate network.
\end{explainbox}
\end{questioncard}

\begin{questioncard}{12}{singlecol}
\qtypebadge{singlecol}{Single Choice} \hfill \textcolor{darkslate!70}{\small Topic: VPN}

Which SSL VPN mode is best for accessing internal web apps without installing a client?

\begin{optlist}
\item Web proxy
\item Network extension
\item IPsec
\item L2TP
\end{optlist}
\begin{explainbox}{}
\textbf{Correct Answer: A}\\[0.3em]
\textbf{Concept:} Web proxy mode requires only a browser.
\end{explainbox}
\end{questioncard}

\begin{questioncard}{13}{singlecol}
\qtypebadge{singlecol}{Single Choice} \hfill \textcolor{darkslate!70}{\small Topic: VPN}

Which encapsulation mode is typically used for site-to-site VPNs?

\begin{optlist}
\item Transport mode
\item Tunnel mode
\item Bridge mode
\item Trunk mode
\end{optlist}
\begin{explainbox}{}
\textbf{Correct Answer: B}\\[0.3em]
\textbf{Concept:} Site-to-site VPNs typically use tunnel mode to protect the entire original packet.
\end{explainbox}
\end{questioncard}

\begin{questioncard}{14}{singlecol}
\qtypebadge{singlecol}{Single Choice} \hfill \textcolor{darkslate!70}{\small Topic: VPN}

Which tunneling protocol is considered a Layer 2 VPN technology?

\begin{optlist}
\item GRE
\item IPsec
\item L2TP
\item SSL VPN
\end{optlist}
\begin{explainbox}{}
\textbf{Correct Answer: C}\\[0.3em]
\textbf{Concept:} L2TP operates at Layer 2, while GRE, IPsec, and SSL VPN operate at higher layers.
\end{explainbox}
\end{questioncard}

\begin{questioncard}{15}{singlecol}
\qtypebadge{singlecol}{Single Choice} \hfill \textcolor{darkslate!70}{\small Topic: VPN}

Which component defines the security parameters for an IPsec connection?

\begin{optlist}
\item SPD
\item SA
\item NAT
\item ARP
\end{optlist}
\begin{explainbox}{}
\textbf{Correct Answer: B}\\[0.3em]
\textbf{Concept:} A Security Association (SA) defines algorithms, keys, and lifetimes for IPsec.
\end{explainbox}
\end{questioncard}

\begin{questioncard}{16}{singlecol}
\qtypebadge{singlecol}{Single Choice} \hfill \textcolor{darkslate!70}{\small Topic: VPN}

The main goal of a VPN is to provide:

\begin{optlist}
\item Secure communication over public networks
\item Faster Internet speed
\item Free public Wi-Fi
\item Unencrypted remote access
\end{optlist}
\begin{explainbox}{}
\textbf{Correct Answer: A}\\[0.3em]
\textbf{Concept:} VPNs create secure encrypted tunnels over untrusted networks.
\end{explainbox}
\end{questioncard}

\begin{questioncard}{17}{singlecol}
\qtypebadge{singlecol}{Single Choice} \hfill \textcolor{darkslate!70}{\small Topic: VPN}

GRE protocol number is:

\begin{optlist}
\item 47
\item 50
\item 51
\item 89
\end{optlist}
\begin{explainbox}{}
\textbf{Correct Answer: A}\\[0.3em]
\textbf{Concept:} GRE is encapsulated using IP protocol 47.
\end{explainbox}
\end{questioncard}

\begin{questioncard}{18}{singlecol}
\qtypebadge{singlecol}{Single Choice} \hfill \textcolor{darkslate!70}{\small Topic: VPN}

IPsec ESP provides:

\begin{optlist}
\item Encryption and authentication
\item Only routing
\item Only NAT
\item Only compression
\end{optlist}
\begin{explainbox}{}
\textbf{Correct Answer: A}\\[0.3em]
\textbf{Concept:} ESP provides confidentiality, integrity, and authentication.
\end{explainbox}
\end{questioncard}

\begin{questioncard}{19}{singlecol}
\qtypebadge{singlecol}{Single Choice} \hfill \textcolor{darkslate!70}{\small Topic: VPN}

IPsec AH provides:

\begin{optlist}
\item Authentication and integrity without encryption
\item Encryption only
\item Compression only
\item Key exchange
\end{optlist}
\begin{explainbox}{}
\textbf{Correct Answer: A}\\[0.3em]
\textbf{Concept:} AH authenticates and protects integrity but does not encrypt.
\end{explainbox}
\end{questioncard}

\begin{questioncard}{20}{singlecol}
\qtypebadge{singlecol}{Single Choice} \hfill \textcolor{darkslate!70}{\small Topic: VPN}

IKE negotiates:

\begin{optlist}
\item Security associations and keys for IPsec
\item DHCP leases
\item DNS records
\item MAC addresses
\end{optlist}
\begin{explainbox}{}
\textbf{Correct Answer: A}\\[0.3em]
\textbf{Concept:} IKE handles key exchange and SA negotiation.
\end{explainbox}
\end{questioncard}

\begin{questioncard}{21}{singlecol}
\qtypebadge{singlecol}{Single Choice} \hfill \textcolor{darkslate!70}{\small Topic: VPN}

SSL VPN commonly uses port:

\begin{optlist}
\item 80
\item 443
\item 22
\item 1701
\end{optlist}
\begin{explainbox}{}
\textbf{Correct Answer: B}\\[0.3em]
\textbf{Concept:} SSL VPN uses HTTPS/TCP 443.
\end{explainbox}
\end{questioncard}

\begin{questioncard}{22}{singlecol}
\qtypebadge{singlecol}{Single Choice} \hfill \textcolor{darkslate!70}{\small Topic: VPN}

L2TP is commonly paired with which protocol for encryption?

\begin{optlist}
\item IPsec
\item GRE
\item OSPF
\item BGP
\end{optlist}
\begin{explainbox}{}
\textbf{Correct Answer: A}\\[0.3em]
\textbf{Concept:} L2TP provides tunneling and is paired with IPsec for encryption.
\end{explainbox}
\end{questioncard}

\begin{questioncard}{23}{singlecol}
\qtypebadge{singlecol}{Single Choice} \hfill \textcolor{darkslate!70}{\small Topic: VPN}

Which concept is described as: provides secure communication over public networks?

\begin{optlist}
\item Remote-access VPN
\item AH
\item VPN
\item SSL VPN
\end{optlist}
\begin{explainbox}{}
\textbf{Correct Answer: C}\\[0.3em]
\textbf{Concept:} VPN refers to the concept where provides secure communication over public networks.
\end{explainbox}
\end{questioncard}

\begin{questioncard}{24}{singlecol}
\qtypebadge{singlecol}{Single Choice} \hfill \textcolor{darkslate!70}{\small Topic: VPN}

Which concept is described as: connects entire networks at different locations?

\begin{optlist}
\item Site-to-site VPN
\item IPsec
\item Tunnel mode
\item GRE
\end{optlist}
\begin{explainbox}{}
\textbf{Correct Answer: A}\\[0.3em]
\textbf{Concept:} Site-to-site VPN refers to the concept where connects entire networks at different locations.
\end{explainbox}
\end{questioncard}

\begin{questioncard}{25}{singlecol}
\qtypebadge{singlecol}{Single Choice} \hfill \textcolor{darkslate!70}{\small Topic: VPN}

Which concept is described as: connects individual users to a corporate network?

\begin{optlist}
\item Remote-access VPN
\item IKEv2
\item Tunnel mode
\item Network extension mode
\end{optlist}
\begin{explainbox}{}
\textbf{Correct Answer: A}\\[0.3em]
\textbf{Concept:} Remote-access VPN refers to the concept where connects individual users to a corporate network.
\end{explainbox}
\end{questioncard}

\begin{questioncard}{26}{singlecol}
\qtypebadge{singlecol}{Single Choice} \hfill \textcolor{darkslate!70}{\small Topic: VPN}

Which concept is described as: is IP protocol 47 and lacks built-in encryption?

\begin{optlist}
\item ESP
\item GRE
\item IKE
\item IPsec
\end{optlist}
\begin{explainbox}{}
\textbf{Correct Answer: B}\\[0.3em]
\textbf{Concept:} GRE refers to the concept where is IP protocol 47 and lacks built-in encryption.
\end{explainbox}
\end{questioncard}

\begin{questioncard}{27}{singlecol}
\qtypebadge{singlecol}{Single Choice} \hfill \textcolor{darkslate!70}{\small Topic: VPN}

Which concept is described as: can tunnel multicast and non-IP protocols?

\begin{optlist}
\item Tunnel mode
\item Transport mode
\item GRE
\item Remote-access VPN
\end{optlist}
\begin{explainbox}{}
\textbf{Correct Answer: C}\\[0.3em]
\textbf{Concept:} GRE refers to the concept where can tunnel multicast and non-IP protocols.
\end{explainbox}
\end{questioncard}

\begin{questioncard}{28}{singlecol}
\qtypebadge{singlecol}{Single Choice} \hfill \textcolor{darkslate!70}{\small Topic: VPN}

Which concept is described as: provides confidentiality, integrity, and authentication?

\begin{optlist}
\item Network extension mode
\item Remote-access VPN
\item IPsec
\item IKEv2
\end{optlist}
\begin{explainbox}{}
\textbf{Correct Answer: C}\\[0.3em]
\textbf{Concept:} IPsec refers to the concept where provides confidentiality, integrity, and authentication.
\end{explainbox}
\end{questioncard}

\begin{questioncard}{29}{singlecol}
\qtypebadge{singlecol}{Single Choice} \hfill \textcolor{darkslate!70}{\small Topic: VPN}

Which concept is described as: provides encryption and authentication?

\begin{optlist}
\item ESP
\item IKEv2
\item SSL VPN
\item GRE
\end{optlist}
\begin{explainbox}{}
\textbf{Correct Answer: A}\\[0.3em]
\textbf{Concept:} ESP refers to the concept where provides encryption and authentication.
\end{explainbox}
\end{questioncard}

\begin{questioncard}{30}{singlecol}
\qtypebadge{singlecol}{Single Choice} \hfill \textcolor{darkslate!70}{\small Topic: VPN}

Which concept is described as: provides authentication without encryption?

\begin{optlist}
\item ESP
\item GRE
\item AH
\item IPsec
\end{optlist}
\begin{explainbox}{}
\textbf{Correct Answer: C}\\[0.3em]
\textbf{Concept:} AH refers to the concept where provides authentication without encryption.
\end{explainbox}
\end{questioncard}

\begin{questioncard}{31}{singlecol}
\qtypebadge{singlecol}{Single Choice} \hfill \textcolor{darkslate!70}{\small Topic: VPN}

Which concept is described as: negotiates IPsec keys and security associations?

\begin{optlist}
\item GRE
\item Tunnel mode
\item IKEv2
\item IKE
\end{optlist}
\begin{explainbox}{}
\textbf{Correct Answer: D}\\[0.3em]
\textbf{Concept:} IKE refers to the concept where negotiates IPsec keys and security associations.
\end{explainbox}
\end{questioncard}

\begin{questioncard}{32}{singlecol}
\qtypebadge{singlecol}{Single Choice} \hfill \textcolor{darkslate!70}{\small Topic: VPN}

Which concept is described as: encrypts the entire original IP packet?

\begin{optlist}
\item IKEv2
\item L2TP
\item IKE
\item Tunnel mode
\end{optlist}
\begin{explainbox}{}
\textbf{Correct Answer: D}\\[0.3em]
\textbf{Concept:} Tunnel mode refers to the concept where encrypts the entire original IP packet.
\end{explainbox}
\end{questioncard}

\begin{questioncard}{33}{singlecol}
\qtypebadge{singlecol}{Single Choice} \hfill \textcolor{darkslate!70}{\small Topic: VPN}

Which concept is described as: encrypts only the IP payload?

\begin{optlist}
\item Remote-access VPN
\item Web proxy mode
\item Transport mode
\item GRE
\end{optlist}
\begin{explainbox}{}
\textbf{Correct Answer: C}\\[0.3em]
\textbf{Concept:} Transport mode refers to the concept where encrypts only the IP payload.
\end{explainbox}
\end{questioncard}

\begin{questioncard}{34}{singlecol}
\qtypebadge{singlecol}{Single Choice} \hfill \textcolor{darkslate!70}{\small Topic: VPN}

Which concept is described as: uses UDP port 1701 and lacks encryption?

\begin{optlist}
\item IKEv2
\item VPN
\item L2TP over IPsec
\item L2TP
\end{optlist}
\begin{explainbox}{}
\textbf{Correct Answer: D}\\[0.3em]
\textbf{Concept:} L2TP refers to the concept where uses UDP port 1701 and lacks encryption.
\end{explainbox}
\end{questioncard}

\begin{questioncard}{35}{singlecol}
\qtypebadge{singlecol}{Single Choice} \hfill \textcolor{darkslate!70}{\small Topic: VPN}

Which concept is described as: combines L2TP tunneling with IPsec encryption?

\begin{optlist}
\item Network extension mode
\item L2TP
\item L2TP over IPsec
\item Site-to-site VPN
\end{optlist}
\begin{explainbox}{}
\textbf{Correct Answer: C}\\[0.3em]
\textbf{Concept:} L2TP over IPsec refers to the concept where combines L2TP tunneling with IPsec encryption.
\end{explainbox}
\end{questioncard}

\begin{questioncard}{36}{singlecol}
\qtypebadge{singlecol}{Single Choice} \hfill \textcolor{darkslate!70}{\small Topic: VPN}

Which concept is described as: commonly uses TCP port 443?

\begin{optlist}
\item SSL VPN
\item IKE
\item IPsec
\item Web proxy mode
\end{optlist}
\begin{explainbox}{}
\textbf{Correct Answer: A}\\[0.3em]
\textbf{Concept:} SSL VPN refers to the concept where commonly uses TCP port 443.
\end{explainbox}
\end{questioncard}

\begin{questioncard}{37}{singlecol}
\qtypebadge{singlecol}{Single Choice} \hfill \textcolor{darkslate!70}{\small Topic: VPN}

Which concept is described as: provides browser-based access to web apps?

\begin{optlist}
\item Network extension mode
\item Web proxy mode
\item IKEv2
\item GRE
\end{optlist}
\begin{explainbox}{}
\textbf{Correct Answer: B}\\[0.3em]
\textbf{Concept:} Web proxy mode refers to the concept where provides browser-based access to web apps.
\end{explainbox}
\end{questioncard}

\begin{questioncard}{38}{singlecol}
\qtypebadge{singlecol}{Single Choice} \hfill \textcolor{darkslate!70}{\small Topic: VPN}

Which concept is described as: assigns a virtual IP to the remote host?

\begin{optlist}
\item Network extension mode
\item Site-to-site VPN
\item Web proxy mode
\item VPN
\end{optlist}
\begin{explainbox}{}
\textbf{Correct Answer: A}\\[0.3em]
\textbf{Concept:} Network extension mode refers to the concept where assigns a virtual IP to the remote host.
\end{explainbox}
\end{questioncard}

\begin{questioncard}{39}{singlecol}
\qtypebadge{singlecol}{Single Choice} \hfill \textcolor{darkslate!70}{\small Topic: VPN}

Which concept is described as: is faster and more reliable than IKEv1?

\begin{optlist}
\item IKEv2
\item Web proxy mode
\item IPsec
\item Site-to-site VPN
\end{optlist}
\begin{explainbox}{}
\textbf{Correct Answer: A}\\[0.3em]
\textbf{Concept:} IKEv2 refers to the concept where is faster and more reliable than IKEv1.
\end{explainbox}
\end{questioncard}

\section{VPN Overview}
\begin{questioncard}{1}{singlecol}
\qtypebadge{singlecol}{Single Choice} \hfill \textcolor{darkslate!70}{\small Topic: VPN Overview}

What does VPN stand for?

\begin{optlist}
\item Virtual Private Network
\item Virtual Public Network
\item Verified Private Node
\item Virtual Protocol Network
\end{optlist}
\begin{explainbox}{}
\textbf{Correct Answer: A}\\[0.3em]
\textbf{Concept:} A VPN extends a private network across a public network using encryption and tunneling.
\end{explainbox}
\end{questioncard}

\begin{questioncard}{2}{singlecol}
\qtypebadge{singlecol}{Single Choice} \hfill \textcolor{darkslate!70}{\small Topic: VPN Overview}

Which benefit does a VPN provide over the public Internet?

\begin{optlist}
\item Higher physical speed
\item Confidentiality and secure remote access
\item Free bandwidth
\item No encryption needed
\end{optlist}
\begin{explainbox}{}
\textbf{Correct Answer: B}\\[0.3em]
\textbf{Concept:} VPNs provide confidentiality, integrity, and secure access over untrusted networks.
\end{explainbox}
\end{questioncard}

\begin{questioncard}{3}{singlecol}
\qtypebadge{singlecol}{Single Choice} \hfill \textcolor{darkslate!70}{\small Topic: VPN Overview}

Which VPN type connects individual remote users to a corporate network?

\begin{optlist}
\item Site-to-site VPN
\item Remote-access VPN
\item Intranet VPN
\item Extranet VPN
\end{optlist}
\begin{explainbox}{}
\textbf{Correct Answer: B}\\[0.3em]
\textbf{Concept:} Remote-access VPNs allow individual users to securely connect to a corporate network.
\end{explainbox}
\end{questioncard}

\section{GRE VPN}
\begin{questioncard}{1}{singlecol}
\qtypebadge{singlecol}{Single Choice} \hfill \textcolor{darkslate!70}{\small Topic: GRE VPN}

Which protocol does GRE use as its transport protocol?

\begin{optlist}
\item TCP
\item UDP
\item IP protocol 47
\item IP protocol 50
\end{optlist}
\begin{explainbox}{}
\textbf{Correct Answer: C}\\[0.3em]
\textbf{Concept:} GRE encapsulates packets using IP protocol number 47.
\end{explainbox}
\end{questioncard}

\begin{questioncard}{2}{singlecol}
\qtypebadge{singlecol}{Single Choice} \hfill \textcolor{darkslate!70}{\small Topic: GRE VPN}

Which statement about GRE is correct?

\begin{optlist}
\item GRE provides encryption by default
\item GRE is a simple tunneling protocol without built-in encryption
\item GRE cannot tunnel multicast
\item GRE uses UDP port 500
\end{optlist}
\begin{explainbox}{}
\textbf{Correct Answer: B}\\[0.3em]
\textbf{Concept:} GRE provides encapsulation and tunneling but does not provide encryption or authentication by itself.
\end{explainbox}
\end{questioncard}

\section{IPsec VPN}
\begin{questioncard}{1}{singlecol}
\qtypebadge{singlecol}{Single Choice} \hfill \textcolor{darkslate!70}{\small Topic: IPsec VPN}

Which IPsec protocol provides encryption and confidentiality?

\begin{optlist}
\item AH
\item ESP
\item IKE
\item GRE
\end{optlist}
\begin{explainbox}{}
\textbf{Correct Answer: B}\\[0.3em]
\textbf{Concept:} ESP (Encapsulating Security Payload) provides encryption, authentication, and integrity.
\end{explainbox}
\end{questioncard}

\begin{questioncard}{2}{singlecol}
\qtypebadge{singlecol}{Single Choice} \hfill \textcolor{darkslate!70}{\small Topic: IPsec VPN}

Which IPsec protocol provides authentication and integrity but not encryption?

\begin{optlist}
\item AH
\item ESP
\item IKE
\item L2TP
\end{optlist}
\begin{explainbox}{}
\textbf{Correct Answer: A}\\[0.3em]
\textbf{Concept:} AH (Authentication Header) provides source authentication and integrity but does not encrypt payloads.
\end{explainbox}
\end{questioncard}

\begin{questioncard}{3}{singlecol}
\qtypebadge{singlecol}{Single Choice} \hfill \textcolor{darkslate!70}{\small Topic: IPsec VPN}

Which protocol is used to establish and manage IPsec security associations?

\begin{optlist}
\item AH
\item ESP
\item IKE
\item GRE
\end{optlist}
\begin{explainbox}{}
\textbf{Correct Answer: C}\\[0.3em]
\textbf{Concept:} IKE (Internet Key Exchange) negotiates keys and Security Associations for IPsec.
\end{explainbox}
\end{questioncard}

\begin{questioncard}{4}{singlecol}
\qtypebadge{singlecol}{Single Choice} \hfill \textcolor{darkslate!70}{\small Topic: IPsec VPN}

Which IPsec mode encrypts the entire original IP packet and adds a new IP header?

\begin{optlist}
\item Transport mode
\item Tunnel mode
\item Gre mode
\item Bridge mode
\end{optlist}
\begin{explainbox}{}
\textbf{Correct Answer: B}\\[0.3em]
\textbf{Concept:} Tunnel mode encapsulates and encrypts the entire original packet, adding a new IP header.
\end{explainbox}
\end{questioncard}

\begin{questioncard}{5}{singlecol}
\qtypebadge{singlecol}{Single Choice} \hfill \textcolor{darkslate!70}{\small Topic: IPsec VPN}

Which IPsec mode encrypts only the payload and is used for host-to-host communications?

\begin{optlist}
\item Transport mode
\item Tunnel mode
\item Bridge mode
\item Trunk mode
\end{optlist}
\begin{explainbox}{}
\textbf{Correct Answer: A}\\[0.3em]
\textbf{Concept:} Transport mode encrypts only the IP payload, leaving the original IP header intact.
\end{explainbox}
\end{questioncard}

\section{L2TP VPN}
\begin{questioncard}{1}{singlecol}
\qtypebadge{singlecol}{Single Choice} \hfill \textcolor{darkslate!70}{\small Topic: L2TP VPN}

Which layer-2 tunneling protocol is often combined with IPsec for encryption?

\begin{optlist}
\item PPTP
\item L2TP
\item GRE
\item MPLS
\end{optlist}
\begin{explainbox}{}
\textbf{Correct Answer: B}\\[0.3em]
\textbf{Concept:} L2TP provides tunneling and is commonly paired with IPsec to provide encryption.
\end{explainbox}
\end{questioncard}

\section{SSL VPN}
\begin{questioncard}{1}{singlecol}
\qtypebadge{singlecol}{Single Choice} \hfill \textcolor{darkslate!70}{\small Topic: SSL VPN}

Which VPN type is commonly accessed through a web browser without a dedicated client?

\begin{optlist}
\item IPsec VPN
\item SSL VPN
\item GRE VPN
\item L2TP VPN
\end{optlist}
\begin{explainbox}{}
\textbf{Correct Answer: B}\\[0.3em]
\textbf{Concept:} SSL VPN can be accessed via a web browser, making it convenient for remote users.
\end{explainbox}
\end{questioncard}

\begin{questioncard}{2}{singlecol}
\qtypebadge{singlecol}{Single Choice} \hfill \textcolor{darkslate!70}{\small Topic: SSL VPN}

Which SSL VPN mode provides access to specific web applications through a browser?

\begin{optlist}
\item Web proxy mode
\item Network extension mode
\item File sharing mode
\item Full tunnel mode
\end{optlist}
\begin{explainbox}{}
\textbf{Correct Answer: A}\\[0.3em]
\textbf{Concept:} Web proxy mode allows browser-based access to internal web applications.
\end{explainbox}
\end{questioncard}

\part{Multiple-Answer Questions}
\section{Network Security Concepts}
\begin{questioncard}{1}{multicol}
\qtypebadge{multicol}{Multiple Choice} \hfill \textcolor{darkslate!70}{\small Topic: Network Security Concepts}

Which of the following are characteristics of network security? (Choose all that apply)

\begin{optlist}
\item Confidentiality
\item Integrity
\item Availability
\item Controllability
\item Non-repudiation
\end{optlist}
\begin{explainbox}{}
\textbf{Correct Answer: A,B,C,D,E}\\[0.3em]
\textbf{Concept:} Network security commonly encompasses confidentiality, integrity, availability, controllability, and non-repudiation.
\end{explainbox}
\end{questioncard}

\begin{questioncard}{2}{multicol}
\qtypebadge{multicol}{Multiple Choice} \hfill \textcolor{darkslate!70}{\small Topic: Network Security Concepts}

Which are stages in the development history of network security? (Choose all that apply)

\begin{optlist}
\item Communication security period
\item Information security period
\item Information assurance period
\item Cyberspace security period
\end{optlist}
\begin{explainbox}{}
\textbf{Correct Answer: A,B,C,D}\\[0.3em]
\textbf{Concept:} Network security has evolved through communication security, information security, information assurance, and cyberspace security periods.
\end{explainbox}
\end{questioncard}

\begin{questioncard}{3}{multicol}
\qtypebadge{multicol}{Multiple Choice} \hfill \textcolor{darkslate!70}{\small Topic: Network Security Concepts}

Which are core information security properties? (Choose all that apply)

\begin{optlist}
\item Confidentiality
\item Integrity
\item Availability
\item Controllability
\item Non-repudiation
\end{optlist}
\begin{explainbox}{}
\textbf{Correct Answer: A,B,C,D,E}\\[0.3em]
\textbf{Concept:} Network security includes confidentiality, integrity, availability, controllability, and non-repudiation.
\end{explainbox}
\end{questioncard}

\section{Network Security Standards}
\begin{questioncard}{1}{multicol}
\qtypebadge{multicol}{Multiple Choice} \hfill \textcolor{darkslate!70}{\small Topic: Network Security Standards}

Which are typical Classified Protection 2.0 lifecycle phases? (Choose all that apply)

\begin{optlist}
\item Classification
\item Filing
\item Construction and rectification
\item Level assessment
\item Supervision and inspection
\end{optlist}
\begin{explainbox}{}
\textbf{Correct Answer: A,B,C,D,E}\\[0.3em]
\textbf{Concept:} The full lifecycle includes classification, filing, construction/rectification, level assessment, and supervision/inspection.
\end{explainbox}
\end{questioncard}

\begin{questioncard}{2}{multicol}
\qtypebadge{multicol}{Multiple Choice} \hfill \textcolor{darkslate!70}{\small Topic: Network Security Standards}

Which are benefits of implementing ISO 27001? (Choose all that apply)

\begin{optlist}
\item Systematic risk management
\item Improved stakeholder confidence
\item Legal/regulatory compliance
\item Guaranteed zero attacks
\end{optlist}
\begin{explainbox}{}
\textbf{Correct Answer: A,B,C}\\[0.3em]
\textbf{Concept:} ISO 27001 helps manage risk, build confidence, and comply with regulations, but cannot guarantee zero attacks.
\end{explainbox}
\end{questioncard}

\begin{questioncard}{3}{multicol}
\qtypebadge{multicol}{Multiple Choice} \hfill \textcolor{darkslate!70}{\small Topic: Network Security Standards}

Which are typical Classified Protection 2.0 security extension requirements? (Choose all that apply)

\begin{optlist}
\item Cloud computing
\item Mobile Internet
\item Internet of Things
\item Industrial control systems
\end{optlist}
\begin{explainbox}{}
\textbf{Correct Answer: A,B,C,D}\\[0.3em]
\textbf{Concept:} Classified Protection 2.0 extends coverage to cloud, mobile Internet, IoT, and industrial control scenarios.
\end{explainbox}
\end{questioncard}

\begin{questioncard}{4}{multicol}
\qtypebadge{multicol}{Multiple Choice} \hfill \textcolor{darkslate!70}{\small Topic: Network Security Standards}

Which belong to the Classified Protection lifecycle? (Choose all that apply)

\begin{optlist}
\item Classification
\item Filing
\item Construction and rectification
\item Level assessment
\item Supervision and inspection
\end{optlist}
\begin{explainbox}{}
\textbf{Correct Answer: A,B,C,D,E}\\[0.3em]
\textbf{Concept:} The full lifecycle includes classification, filing, construction, level assessment, and supervision.
\end{explainbox}
\end{questioncard}

\begin{questioncard}{5}{multicol}
\qtypebadge{multicol}{Multiple Choice} \hfill \textcolor{darkslate!70}{\small Topic: Network Security Standards}

Which are ISO 27000 family standards? (Choose all that apply)

\begin{optlist}
\item ISO 27001
\item ISO 27002
\item ISO 27005
\item ISO 9001
\end{optlist}
\begin{explainbox}{}
\textbf{Correct Answer: A,B,C}\\[0.3em]
\textbf{Concept:} ISO 27001, 27002, and 27005 are information security standards; ISO 9001 is quality management.
\end{explainbox}
\end{questioncard}

\begin{questioncard}{6}{multicol}
\qtypebadge{multicol}{Multiple Choice} \hfill \textcolor{darkslate!70}{\small Topic: Network Security Standards}

Which of the following relate to ISO? (Choose all that apply)

\begin{optlist}
\item ISO 27001
\item ISO 27002
\item ISO 27001:2022
\end{optlist}
\begin{explainbox}{}
\textbf{Correct Answer: A,B,C}\\[0.3em]
\textbf{Concept:} These concepts all relate to ISO in Network Security Standards.
\end{explainbox}
\end{questioncard}

\begin{questioncard}{7}{multicol}
\qtypebadge{multicol}{Multiple Choice} \hfill \textcolor{darkslate!70}{\small Topic: Network Security Standards}

Which of the following relate to Classified? (Choose all that apply)

\begin{optlist}
\item Classified Protection 2.0
\item Classified Protection levels
\item Classified Protection lifecycle
\end{optlist}
\begin{explainbox}{}
\textbf{Correct Answer: A,B,C}\\[0.3em]
\textbf{Concept:} These concepts all relate to Classified in Network Security Standards.
\end{explainbox}
\end{questioncard}

\section{Network Fundamentals}
\begin{questioncard}{1}{multicol}
\qtypebadge{multicol}{Multiple Choice} \hfill \textcolor{darkslate!70}{\small Topic: Network Fundamentals}

Which layers are included in the TCP/IP equivalent model? (Choose all that apply)

\begin{optlist}
\item Application
\item Transport
\item Network
\item Data link
\item Physical
\end{optlist}
\begin{explainbox}{}
\textbf{Correct Answer: A,B,C,D,E}\\[0.3em]
\textbf{Concept:} The TCP/IP equivalent model includes application, transport, network, data link, and physical layers.
\end{explainbox}
\end{questioncard}

\begin{questioncard}{2}{multicol}
\qtypebadge{multicol}{Multiple Choice} \hfill \textcolor{darkslate!70}{\small Topic: Network Fundamentals}

Which are connectionless protocols? (Choose all that apply)

\begin{optlist}
\item UDP
\item IP
\item ICMP
\item TCP
\end{optlist}
\begin{explainbox}{}
\textbf{Correct Answer: A,B,C}\\[0.3em]
\textbf{Concept:} UDP, IP, and ICMP are connectionless; TCP is connection-oriented.
\end{explainbox}
\end{questioncard}

\begin{questioncard}{3}{multicol}
\qtypebadge{multicol}{Multiple Choice} \hfill \textcolor{darkslate!70}{\small Topic: Network Fundamentals}

Which fields are in a TCP header? (Choose all that apply)

\begin{optlist}
\item Source port
\item Destination port
\item Sequence number
\item Acknowledgment number
\end{optlist}
\begin{explainbox}{}
\textbf{Correct Answer: A,B,C,D}\\[0.3em]
\textbf{Concept:} TCP headers include source/destination ports, sequence and acknowledgment numbers, flags, window size, etc.
\end{explainbox}
\end{questioncard}

\begin{questioncard}{4}{multicol}
\qtypebadge{multicol}{Multiple Choice} \hfill \textcolor{darkslate!70}{\small Topic: Network Fundamentals}

Which devices operate primarily at Layer 1? (Choose all that apply)

\begin{optlist}
\item Hub
\item Repeater
\item Switch
\item Router
\end{optlist}
\begin{explainbox}{}
\textbf{Correct Answer: A,B}\\[0.3em]
\textbf{Concept:} Hubs and repeaters operate at the physical layer.
\end{explainbox}
\end{questioncard}

\begin{questioncard}{5}{multicol}
\qtypebadge{multicol}{Multiple Choice} \hfill \textcolor{darkslate!70}{\small Topic: Network Fundamentals}

Which protocols use TCP? (Choose all that apply)

\begin{optlist}
\item HTTP
\item FTP
\item SMTP
\item DNS (sometimes)
\end{optlist}
\begin{explainbox}{}
\textbf{Correct Answer: A,B,C,D}\\[0.3em]
\textbf{Concept:} HTTP, FTP, and SMTP use TCP; DNS primarily uses UDP but can use TCP for zone transfers and large responses.
\end{explainbox}
\end{questioncard}

\begin{questioncard}{6}{multicol}
\qtypebadge{multicol}{Multiple Choice} \hfill \textcolor{darkslate!70}{\small Topic: Network Fundamentals}

Which protocols use UDP? (Choose all that apply)

\begin{optlist}
\item DNS
\item TFTP
\item SNMP
\item DHCP
\end{optlist}
\begin{explainbox}{}
\textbf{Correct Answer: A,B,C,D}\\[0.3em]
\textbf{Concept:} DNS, TFTP, SNMP, and DHCP commonly use UDP.
\end{explainbox}
\end{questioncard}

\begin{questioncard}{7}{multicol}
\qtypebadge{multicol}{Multiple Choice} \hfill \textcolor{darkslate!70}{\small Topic: Network Fundamentals}

Which are valid IPv4 private address ranges? (Choose all that apply)

\begin{optlist}
\item 10.0.0.0/8
\item 172.16.0.0/12
\item 192.168.0.0/16
\item 224.0.0.0/4
\end{optlist}
\begin{explainbox}{}
\textbf{Correct Answer: A,B,C}\\[0.3em]
\textbf{Concept:} RFC 1918 private ranges are 10.0.0.0/8, 172.16.0.0/12, and 192.168.0.0/16.
\end{explainbox}
\end{questioncard}

\begin{questioncard}{8}{multicol}
\qtypebadge{multicol}{Multiple Choice} \hfill \textcolor{darkslate!70}{\small Topic: Network Fundamentals}

Which are OSI layers? (Choose all that apply)

\begin{optlist}
\item Physical
\item Data link
\item Network
\item Transport
\item Application
\end{optlist}
\begin{explainbox}{}
\textbf{Correct Answer: A,B,C,D,E}\\[0.3em]
\textbf{Concept:} The OSI model includes these layers plus session and presentation.
\end{explainbox}
\end{questioncard}

\begin{questioncard}{9}{multicol}
\qtypebadge{multicol}{Multiple Choice} \hfill \textcolor{darkslate!70}{\small Topic: Network Fundamentals}

Which are TCP/IP equivalent model layers? (Choose all that apply)

\begin{optlist}
\item Application
\item Transport
\item Network
\item Data link
\item Physical
\end{optlist}
\begin{explainbox}{}
\textbf{Correct Answer: A,B,C,D,E}\\[0.3em]
\textbf{Concept:} The TCP/IP equivalent model includes all five layers.
\end{explainbox}
\end{questioncard}

\begin{questioncard}{10}{multicol}
\qtypebadge{multicol}{Multiple Choice} \hfill \textcolor{darkslate!70}{\small Topic: Network Fundamentals}

Which protocols use TCP? (Choose all that apply)

\begin{optlist}
\item HTTP
\item FTP
\item SMTP
\item SSH
\end{optlist}
\begin{explainbox}{}
\textbf{Correct Answer: A,B,C,D}\\[0.3em]
\textbf{Concept:} HTTP, FTP, SMTP, and SSH all use TCP.
\end{explainbox}
\end{questioncard}

\begin{questioncard}{11}{multicol}
\qtypebadge{multicol}{Multiple Choice} \hfill \textcolor{darkslate!70}{\small Topic: Network Fundamentals}

Which devices operate at Layer 2? (Choose all that apply)

\begin{optlist}
\item Switch
\item Bridge
\item Router
\item Hub
\end{optlist}
\begin{explainbox}{}
\textbf{Correct Answer: A,B}\\[0.3em]
\textbf{Concept:} Switches and bridges operate at the data link layer.
\end{explainbox}
\end{questioncard}

\begin{questioncard}{12}{multicol}
\qtypebadge{multicol}{Multiple Choice} \hfill \textcolor{darkslate!70}{\small Topic: Network Fundamentals}

Which of the following relate to TCP/IP? (Choose all that apply)

\begin{optlist}
\item TCP/IP standard model
\item TCP/IP equivalent model
\end{optlist}
\begin{explainbox}{}
\textbf{Correct Answer: A,B}\\[0.3em]
\textbf{Concept:} These concepts all relate to TCP/IP in Network Fundamentals.
\end{explainbox}
\end{questioncard}

\section{OSI and TCP/IP Models}
\begin{questioncard}{1}{multicol}
\qtypebadge{multicol}{Multiple Choice} \hfill \textcolor{darkslate!70}{\small Topic: OSI and TCP/IP Models}

Which layers are part of the OSI model? (Choose all that apply)

\begin{optlist}
\item Physical
\item Data link
\item Network
\item Transport
\item Application
\end{optlist}
\begin{explainbox}{}
\textbf{Correct Answer: A,B,C,D,E}\\[0.3em]
\textbf{Concept:} The OSI model includes physical, data link, network, transport, session, presentation, and application layers.
\end{explainbox}
\end{questioncard}

\section{Network Protocols}
\begin{questioncard}{1}{multicol}
\qtypebadge{multicol}{Multiple Choice} \hfill \textcolor{darkslate!70}{\small Topic: Network Protocols}

Which protocols operate at the application layer? (Choose all that apply)

\begin{optlist}
\item HTTP
\item FTP
\item DNS
\item SMTP
\end{optlist}
\begin{explainbox}{}
\textbf{Correct Answer: A,B,C,D}\\[0.3em]
\textbf{Concept:} HTTP, FTP, DNS, and SMTP are all application-layer protocols.
\end{explainbox}
\end{questioncard}

\section{Network Devices}
\begin{questioncard}{1}{multicol}
\qtypebadge{multicol}{Multiple Choice} \hfill \textcolor{darkslate!70}{\small Topic: Network Devices}

Which devices can operate at Layer 2? (Choose all that apply)

\begin{optlist}
\item Switch
\item Bridge
\item Router
\item Hub
\end{optlist}
\begin{explainbox}{}
\textbf{Correct Answer: A,B}\\[0.3em]
\textbf{Concept:} Switches and bridges operate at the data link layer; routers are Layer 3 and hubs are Layer 1.
\end{explainbox}
\end{questioncard}

\section{Enterprise Security Threats}
\begin{questioncard}{1}{multicol}
\qtypebadge{multicol}{Multiple Choice} \hfill \textcolor{darkslate!70}{\small Topic: Enterprise Security Threats}

Which are examples of malware? (Choose all that apply)

\begin{optlist}
\item Virus
\item Worm
\item Trojan
\item Ransomware
\item Spyware
\end{optlist}
\begin{explainbox}{}
\textbf{Correct Answer: A,B,C,D,E}\\[0.3em]
\textbf{Concept:} All listed items are types of malware.
\end{explainbox}
\end{questioncard}

\begin{questioncard}{2}{multicol}
\qtypebadge{multicol}{Multiple Choice} \hfill \textcolor{darkslate!70}{\small Topic: Enterprise Security Threats}

Which are social engineering attacks? (Choose all that apply)

\begin{optlist}
\item Phishing
\item Pretexting
\item Baiting
\item Tailgating
\end{optlist}
\begin{explainbox}{}
\textbf{Correct Answer: A,B,C,D}\\[0.3em]
\textbf{Concept:} Social engineering manipulates people through phishing, pretexting, baiting, tailgating, etc.
\end{explainbox}
\end{questioncard}

\begin{questioncard}{3}{multicol}
\qtypebadge{multicol}{Multiple Choice} \hfill \textcolor{darkslate!70}{\small Topic: Enterprise Security Threats}

Which protect communication confidentiality? (Choose all that apply)

\begin{optlist}
\item Encryption
\item VPN
\item TLS/SSL
\item MAC filtering
\end{optlist}
\begin{explainbox}{}
\textbf{Correct Answer: A,B,C}\\[0.3em]
\textbf{Concept:} Encryption, VPNs, and TLS/SSL protect confidentiality; MAC filtering controls device access.
\end{explainbox}
\end{questioncard}

\begin{questioncard}{4}{multicol}
\qtypebadge{multicol}{Multiple Choice} \hfill \textcolor{darkslate!70}{\small Topic: Enterprise Security Threats}

Which are typical security zones? (Choose all that apply)

\begin{optlist}
\item Trust
\item Untrust
\item DMZ
\item Local
\end{optlist}
\begin{explainbox}{}
\textbf{Correct Answer: A,B,C,D}\\[0.3em]
\textbf{Concept:} Huawei firewalls define Trust, Untrust, DMZ, and Local security zones by default.
\end{explainbox}
\end{questioncard}

\begin{questioncard}{5}{multicol}
\qtypebadge{multicol}{Multiple Choice} \hfill \textcolor{darkslate!70}{\small Topic: Enterprise Security Threats}

Which are endpoint security controls? (Choose all that apply)

\begin{optlist}
\item Antivirus
\item Host firewall
\item Application whitelisting
\item Full-disk encryption
\end{optlist}
\begin{explainbox}{}
\textbf{Correct Answer: A,B,C,D}\\[0.3em]
\textbf{Concept:} Endpoint security includes antivirus, host firewalls, whitelisting, and encryption.
\end{explainbox}
\end{questioncard}

\begin{questioncard}{6}{multicol}
\qtypebadge{multicol}{Multiple Choice} \hfill \textcolor{darkslate!70}{\small Topic: Enterprise Security Threats}

Which are malware types? (Choose all that apply)

\begin{optlist}
\item Virus
\item Worm
\item Trojan
\item Ransomware
\end{optlist}
\begin{explainbox}{}
\textbf{Correct Answer: A,B,C,D}\\[0.3em]
\textbf{Concept:} All listed are categories of malware.
\end{explainbox}
\end{questioncard}

\begin{questioncard}{7}{multicol}
\qtypebadge{multicol}{Multiple Choice} \hfill \textcolor{darkslate!70}{\small Topic: Enterprise Security Threats}

Which are denial-of-service attacks? (Choose all that apply)

\begin{optlist}
\item SYN flood
\item DDoS
\item Ping of death
\item Smurf attack
\end{optlist}
\begin{explainbox}{}
\textbf{Correct Answer: A,B,C,D}\\[0.3em]
\textbf{Concept:} These attacks aim to disrupt service availability.
\end{explainbox}
\end{questioncard}

\section{Zone Border Security}
\begin{questioncard}{1}{multicol}
\qtypebadge{multicol}{Multiple Choice} \hfill \textcolor{darkslate!70}{\small Topic: Zone Border Security}

Which controls are typically applied at zone borders? (Choose all that apply)

\begin{optlist}
\item Firewall policies
\item NAT
\item IPS/IDS
\item VPN termination
\end{optlist}
\begin{explainbox}{}
\textbf{Correct Answer: A,B,C,D}\\[0.3em]
\textbf{Concept:} Zone borders commonly use firewalls, NAT, IPS/IDS, and VPN gateways for protection.
\end{explainbox}
\end{questioncard}

\section{Firewall Basics}
\begin{questioncard}{1}{multicol}
\qtypebadge{multicol}{Multiple Choice} \hfill \textcolor{darkslate!70}{\small Topic: Firewall Basics}

Which are Huawei firewall default security zones? (Choose all that apply)

\begin{optlist}
\item Trust
\item Untrust
\item DMZ
\item Local
\end{optlist}
\begin{explainbox}{}
\textbf{Correct Answer: A,B,C,D}\\[0.3em]
\textbf{Concept:} Huawei firewalls provide Trust, Untrust, DMZ, and Local zones by default.
\end{explainbox}
\end{questioncard}

\begin{questioncard}{2}{multicol}
\qtypebadge{multicol}{Multiple Choice} \hfill \textcolor{darkslate!70}{\small Topic: Firewall Basics}

Which actions can a security policy perform? (Choose all that apply)

\begin{optlist}
\item Permit
\item Deny
\item Log
\item Redirect
\end{optlist}
\begin{explainbox}{}
\textbf{Correct Answer: A,B,C}\\[0.3em]
\textbf{Concept:} Policies can permit, deny, and log traffic; redirect may be available in specific features.
\end{explainbox}
\end{questioncard}

\begin{questioncard}{3}{multicol}
\qtypebadge{multicol}{Multiple Choice} \hfill \textcolor{darkslate!70}{\small Topic: Firewall Basics}

Which are common firewall deployment modes? (Choose all that apply)

\begin{optlist}
\item Routing mode
\item Transparent mode
\item Hybrid mode
\item Bridge mode
\end{optlist}
\begin{explainbox}{}
\textbf{Correct Answer: A,B,C}\\[0.3em]
\textbf{Concept:} Firewalls commonly deploy in routing, transparent (bridge), and hybrid modes.
\end{explainbox}
\end{questioncard}

\begin{questioncard}{4}{multicol}
\qtypebadge{multicol}{Multiple Choice} \hfill \textcolor{darkslate!70}{\small Topic: Firewall Basics}

Which conditions can be matched in a security policy? (Choose all that apply)

\begin{optlist}
\item Source zone
\item Destination zone
\item Source IP
\item Destination port
\item User
\end{optlist}
\begin{explainbox}{}
\textbf{Correct Answer: A,B,C,D,E}\\[0.3em]
\textbf{Concept:} Policies can match zones, IPs, ports, users, applications, services, and time.
\end{explainbox}
\end{questioncard}

\begin{questioncard}{5}{multicol}
\qtypebadge{multicol}{Multiple Choice} \hfill \textcolor{darkslate!70}{\small Topic: Firewall Basics}

Which are default Huawei security zones? (Choose all that apply)

\begin{optlist}
\item Trust
\item Untrust
\item DMZ
\item Local
\end{optlist}
\begin{explainbox}{}
\textbf{Correct Answer: A,B,C,D}\\[0.3em]
\textbf{Concept:} Huawei firewalls provide Trust, Untrust, DMZ, and Local zones by default.
\end{explainbox}
\end{questioncard}

\begin{questioncard}{6}{multicol}
\qtypebadge{multicol}{Multiple Choice} \hfill \textcolor{darkslate!70}{\small Topic: Firewall Basics}

Which can be matched in security policies? (Choose all that apply)

\begin{optlist}
\item Source zone
\item Destination zone
\item Source IP
\item Destination port
\item User
\end{optlist}
\begin{explainbox}{}
\textbf{Correct Answer: A,B,C,D,E}\\[0.3em]
\textbf{Concept:} Policies can match zones, addresses, ports, users, services, applications, and time.
\end{explainbox}
\end{questioncard}

\begin{questioncard}{7}{multicol}
\qtypebadge{multicol}{Multiple Choice} \hfill \textcolor{darkslate!70}{\small Topic: Firewall Basics}

Which of the following relate to Security? (Choose all that apply)

\begin{optlist}
\item Security zones
\item Security policy logging
\end{optlist}
\begin{explainbox}{}
\textbf{Correct Answer: A,B}\\[0.3em]
\textbf{Concept:} These concepts all relate to Security in Firewall Basics.
\end{explainbox}
\end{questioncard}

\section{Security Zones}
\begin{questioncard}{1}{multicol}
\qtypebadge{multicol}{Multiple Choice} \hfill \textcolor{darkslate!70}{\small Topic: Security Zones}

Which are default security zones on Huawei firewalls? (Choose all that apply)

\begin{optlist}
\item Trust
\item Untrust
\item DMZ
\item Local
\end{optlist}
\begin{explainbox}{}
\textbf{Correct Answer: A,B,C,D}\\[0.3em]
\textbf{Concept:} Huawei firewalls provide four default zones: Trust, Untrust, DMZ, and Local.
\end{explainbox}
\end{questioncard}

\section{Security Policies}
\begin{questioncard}{1}{multicol}
\qtypebadge{multicol}{Multiple Choice} \hfill \textcolor{darkslate!70}{\small Topic: Security Policies}

Which elements can be referenced in a security policy? (Choose all that apply)

\begin{optlist}
\item Security zones
\item IP addresses/address groups
\item Services
\item Users and user groups
\item Time ranges
\end{optlist}
\begin{explainbox}{}
\textbf{Correct Answer: A,B,C,D,E}\\[0.3em]
\textbf{Concept:} Security policies can reference zones, addresses, services, users, applications, and time ranges.
\end{explainbox}
\end{questioncard}

\section{Stateful Inspection}
\begin{questioncard}{1}{multicol}
\qtypebadge{multicol}{Multiple Choice} \hfill \textcolor{darkslate!70}{\small Topic: Stateful Inspection}

Which are advantages of stateful inspection over simple packet filtering? (Choose all that apply)

\begin{optlist}
\item Tracks connection state
\item Allows return traffic automatically
\item Better performance for established sessions
\item Can inspect application-layer payloads
\end{optlist}
\begin{explainbox}{}
\textbf{Correct Answer: A,B,C}\\[0.3em]
\textbf{Concept:} Stateful inspection tracks connections, permits return traffic, and improves performance; deep payload inspection is more advanced.
\end{explainbox}
\end{questioncard}

\section{NAT}
\begin{questioncard}{1}{multicol}
\qtypebadge{multicol}{Multiple Choice} \hfill \textcolor{darkslate!70}{\small Topic: NAT}

Which are private IPv4 address ranges? (Choose all that apply)

\begin{optlist}
\item 10.0.0.0/8
\item 172.16.0.0/12
\item 192.168.0.0/16
\item 127.0.0.0/8
\end{optlist}
\begin{explainbox}{}
\textbf{Correct Answer: A,B,C}\\[0.3em]
\textbf{Concept:} RFC 1918 private ranges are 10/8, 172.16/12, and 192.168/16.
\end{explainbox}
\end{questioncard}

\begin{questioncard}{2}{multicol}
\qtypebadge{multicol}{Multiple Choice} \hfill \textcolor{darkslate!70}{\small Topic: NAT}

Which are valid uses of NAT Server? (Choose all that apply)

\begin{optlist}
\item Publishing a web server
\item Publishing an email server
\item Remote access to an internal host
\item Hiding internal topology
\end{optlist}
\begin{explainbox}{}
\textbf{Correct Answer: A,B,C,D}\\[0.3em]
\textbf{Concept:} NAT Server can publish services, enable remote access, and obscure internal addressing.
\end{explainbox}
\end{questioncard}

\begin{questioncard}{3}{multicol}
\qtypebadge{multicol}{Multiple Choice} \hfill \textcolor{darkslate!70}{\small Topic: NAT}

Which are NAT types? (Choose all that apply)

\begin{optlist}
\item Source NAT
\item Destination NAT
\item Bidirectional NAT
\item NAT Server
\end{optlist}
\begin{explainbox}{}
\textbf{Correct Answer: A,B,C,D}\\[0.3em]
\textbf{Concept:} Huawei firewalls support all four NAT types.
\end{explainbox}
\end{questioncard}

\begin{questioncard}{4}{multicol}
\qtypebadge{multicol}{Multiple Choice} \hfill \textcolor{darkslate!70}{\small Topic: NAT}

Which source NAT modes exist? (Choose all that apply)

\begin{optlist}
\item NAT No-PAT
\item NAPT
\item Easy-IP
\item NAT Server
\end{optlist}
\begin{explainbox}{}
\textbf{Correct Answer: A,B,C}\\[0.3em]
\textbf{Concept:} NAT Server is not a source NAT mode.
\end{explainbox}
\end{questioncard}

\begin{questioncard}{5}{multicol}
\qtypebadge{multicol}{Multiple Choice} \hfill \textcolor{darkslate!70}{\small Topic: NAT}

Which of the following relate to NAT? (Choose all that apply)

\begin{optlist}
\item NAT
\item NAT Server
\item NAT No-PAT
\item NAT ALG
\end{optlist}
\begin{explainbox}{}
\textbf{Correct Answer: A,B,C,D}\\[0.3em]
\textbf{Concept:} These concepts all relate to NAT in NAT.
\end{explainbox}
\end{questioncard}

\section{NAT Basics}
\begin{questioncard}{1}{multicol}
\qtypebadge{multicol}{Multiple Choice} \hfill \textcolor{darkslate!70}{\small Topic: NAT Basics}

Which are types of NAT supported by Huawei firewalls? (Choose all that apply)

\begin{optlist}
\item Source NAT
\item Destination NAT
\item Bidirectional NAT
\item NAT Server
\end{optlist}
\begin{explainbox}{}
\textbf{Correct Answer: A,B,C,D}\\[0.3em]
\textbf{Concept:} Huawei firewalls support source NAT, destination NAT, bidirectional NAT, and NAT server.
\end{explainbox}
\end{questioncard}

\section{Source NAT}
\begin{questioncard}{1}{multicol}
\qtypebadge{multicol}{Multiple Choice} \hfill \textcolor{darkslate!70}{\small Topic: Source NAT}

Which are common source NAT modes? (Choose all that apply)

\begin{optlist}
\item NAT No-PAT
\item NAPT
\item Easy-IP
\item NAT Server
\end{optlist}
\begin{explainbox}{}
\textbf{Correct Answer: A,B,C}\\[0.3em]
\textbf{Concept:} Source NAT includes NAT No-PAT (one-to-one), NAPT (many-to-one), and Easy-IP (interface-based).
\end{explainbox}
\end{questioncard}

\section{NAT ALG}
\begin{questioncard}{1}{multicol}
\qtypebadge{multicol}{Multiple Choice} \hfill \textcolor{darkslate!70}{\small Topic: NAT ALG}

Which protocols may require NAT ALG? (Choose all that apply)

\begin{optlist}
\item FTP
\item SIP
\item H.323
\item DNS
\end{optlist}
\begin{explainbox}{}
\textbf{Correct Answer: A,B,C}\\[0.3em]
\textbf{Concept:} Protocols embedding address/port information in payloads (FTP, SIP, H.323) need NAT ALG.
\end{explainbox}
\end{questioncard}

\section{Firewall Hot Standby}
\begin{questioncard}{1}{multicol}
\qtypebadge{multicol}{Multiple Choice} \hfill \textcolor{darkslate!70}{\small Topic: Firewall Hot Standby}

Which are VRRP components? (Choose all that apply)

\begin{optlist}
\item Virtual IP
\item Virtual MAC
\item Master router
\item Backup router
\end{optlist}
\begin{explainbox}{}
\textbf{Correct Answer: A,B,C,D}\\[0.3em]
\textbf{Concept:} VRRP uses a virtual IP, virtual MAC, master, and backup routers.
\end{explainbox}
\end{questioncard}

\begin{questioncard}{2}{multicol}
\qtypebadge{multicol}{Multiple Choice} \hfill \textcolor{darkslate!70}{\small Topic: Firewall Hot Standby}

Which are HRP synchronization modes? (Choose all that apply)

\begin{optlist}
\item Real-time
\item Fast
\item Batch
\item Auto
\end{optlist}
\begin{explainbox}{}
\textbf{Correct Answer: A,B,C}\\[0.3em]
\textbf{Concept:} HRP supports real-time, fast, and batch synchronization modes.
\end{explainbox}
\end{questioncard}

\begin{questioncard}{3}{multicol}
\qtypebadge{multicol}{Multiple Choice} \hfill \textcolor{darkslate!70}{\small Topic: Firewall Hot Standby}

Which are valid hot-standby modes? (Choose all that apply)

\begin{optlist}
\item Active/standby
\item Load sharing
\item Cluster
\item Standalone
\end{optlist}
\begin{explainbox}{}
\textbf{Correct Answer: A,B}\\[0.3em]
\textbf{Concept:} Huawei firewall hot standby supports active/standby and load-sharing modes.
\end{explainbox}
\end{questioncard}

\begin{questioncard}{4}{multicol}
\qtypebadge{multicol}{Multiple Choice} \hfill \textcolor{darkslate!70}{\small Topic: Firewall Hot Standby}

Which items should be checked when HRP is not synchronizing? (Choose all that apply)

\begin{optlist}
\item Heartbeat link status
\item HRP enable status
\item VGMP priority
\item Configuration consistency
\end{optlist}
\begin{explainbox}{}
\textbf{Correct Answer: A,B,C,D}\\[0.3em]
\textbf{Concept:} Check heartbeat links, HRP status, priorities, and configuration consistency.
\end{explainbox}
\end{questioncard}

\begin{questioncard}{5}{multicol}
\qtypebadge{multicol}{Multiple Choice} \hfill \textcolor{darkslate!70}{\small Topic: Firewall Hot Standby}

Which are HRP synchronization contents? (Choose all that apply)

\begin{optlist}
\item Session table
\item NAT entries
\item Server-map entries
\item Configuration
\end{optlist}
\begin{explainbox}{}
\textbf{Correct Answer: A,B,C,D}\\[0.3em]
\textbf{Concept:} HRP synchronizes all these items between peers.
\end{explainbox}
\end{questioncard}

\begin{questioncard}{6}{multicol}
\qtypebadge{multicol}{Multiple Choice} \hfill \textcolor{darkslate!70}{\small Topic: Firewall Hot Standby}

Which are VRRP states? (Choose all that apply)

\begin{optlist}
\item Master
\item Backup
\item Initialize
\item Failover
\end{optlist}
\begin{explainbox}{}
\textbf{Correct Answer: A,B,C}\\[0.3em]
\textbf{Concept:} VRRP states include master, backup, and initialize.
\end{explainbox}
\end{questioncard}

\begin{questioncard}{7}{multicol}
\qtypebadge{multicol}{Multiple Choice} \hfill \textcolor{darkslate!70}{\small Topic: Firewall Hot Standby}

Which of the following relate to VRRP? (Choose all that apply)

\begin{optlist}
\item VRRP
\item VRRP master
\item VRRP backup
\end{optlist}
\begin{explainbox}{}
\textbf{Correct Answer: A,B,C}\\[0.3em]
\textbf{Concept:} These concepts all relate to VRRP in Firewall Hot Standby.
\end{explainbox}
\end{questioncard}

\begin{questioncard}{8}{multicol}
\qtypebadge{multicol}{Multiple Choice} \hfill \textcolor{darkslate!70}{\small Topic: Firewall Hot Standby}

Which of the following relate to HRP? (Choose all that apply)

\begin{optlist}
\item HRP
\item HRP session backup
\end{optlist}
\begin{explainbox}{}
\textbf{Correct Answer: A,B}\\[0.3em]
\textbf{Concept:} These concepts all relate to HRP in Firewall Hot Standby.
\end{explainbox}
\end{questioncard}

\section{VRRP}
\begin{questioncard}{1}{multicol}
\qtypebadge{multicol}{Multiple Choice} \hfill \textcolor{darkslate!70}{\small Topic: VRRP}

Which are VRRP roles? (Choose all that apply)

\begin{optlist}
\item Master
\item Backup
\item Initialize
\item Failover
\end{optlist}
\begin{explainbox}{}
\textbf{Correct Answer: A,B,C}\\[0.3em]
\textbf{Concept:} VRRP devices can be in master, backup, or initialize states.
\end{explainbox}
\end{questioncard}

\section{HRP}
\begin{questioncard}{1}{multicol}
\qtypebadge{multicol}{Multiple Choice} \hfill \textcolor{darkslate!70}{\small Topic: HRP}

Which data does HRP synchronize? (Choose all that apply)

\begin{optlist}
\item Session table
\item NAT entries
\item Server-map entries
\item Configurations
\end{optlist}
\begin{explainbox}{}
\textbf{Correct Answer: A,B,C,D}\\[0.3em]
\textbf{Concept:} HRP synchronizes sessions, NAT entries, server-map entries, and configurations.
\end{explainbox}
\end{questioncard}

\section{Intrusion Prevention and Antivirus}
\begin{questioncard}{1}{multicol}
\qtypebadge{multicol}{Multiple Choice} \hfill \textcolor{darkslate!70}{\small Topic: Intrusion Prevention and Antivirus}

Which are common attack phases? (Choose all that apply)

\begin{optlist}
\item Reconnaissance
\item Exploitation
\item Lateral movement
\item Exfiltration
\end{optlist}
\begin{explainbox}{}
\textbf{Correct Answer: A,B,C,D}\\[0.3em]
\textbf{Concept:} Cyber kill chain phases include reconnaissance, exploitation, lateral movement, and exfiltration.
\end{explainbox}
\end{questioncard}

\begin{questioncard}{2}{multicol}
\qtypebadge{multicol}{Multiple Choice} \hfill \textcolor{darkslate!70}{\small Topic: Intrusion Prevention and Antivirus}

Which are types of intrusion detection/prevention systems? (Choose all that apply)

\begin{optlist}
\item NIDS
\item NIPS
\item HIDS
\item HIPS
\end{optlist}
\begin{explainbox}{}
\textbf{Correct Answer: A,B,C,D}\\[0.3em]
\textbf{Concept:} Network and host-based IDS/IPS exist for different monitoring scopes.
\end{explainbox}
\end{questioncard}

\begin{questioncard}{3}{multicol}
\qtypebadge{multicol}{Multiple Choice} \hfill \textcolor{darkslate!70}{\small Topic: Intrusion Prevention and Antivirus}

Which factors should be considered when tuning IPS signatures? (Choose all that apply)

\begin{optlist}
\item False positive rate
\item False negative rate
\item Performance impact
\item Business context
\end{optlist}
\begin{explainbox}{}
\textbf{Correct Answer: A,B,C,D}\\[0.3em]
\textbf{Concept:} Effective tuning balances detection accuracy, performance, and business needs.
\end{explainbox}
\end{questioncard}

\begin{questioncard}{4}{multicol}
\qtypebadge{multicol}{Multiple Choice} \hfill \textcolor{darkslate!70}{\small Topic: Intrusion Prevention and Antivirus}

Which are methods to improve malware detection? (Choose all that apply)

\begin{optlist}
\item Keep signatures updated
\item Use heuristics
\item Deploy sandboxing
\item User education
\end{optlist}
\begin{explainbox}{}
\textbf{Correct Answer: A,B,C,D}\\[0.3em]
\textbf{Concept:} Updated signatures, heuristics, sandboxing, and user awareness improve detection.
\end{explainbox}
\end{questioncard}

\begin{questioncard}{5}{multicol}
\qtypebadge{multicol}{Multiple Choice} \hfill \textcolor{darkslate!70}{\small Topic: Intrusion Prevention and Antivirus}

Which are IPS deployment modes? (Choose all that apply)

\begin{optlist}
\item Inline
\item Tap/SPAN
\item Out-of-band
\item Bridge
\end{optlist}
\begin{explainbox}{}
\textbf{Correct Answer: A,B,C}\\[0.3em]
\textbf{Concept:} IPS can be inline, tapped, or out-of-band.
\end{explainbox}
\end{questioncard}

\begin{questioncard}{6}{multicol}
\qtypebadge{multicol}{Multiple Choice} \hfill \textcolor{darkslate!70}{\small Topic: Intrusion Prevention and Antivirus}

Which are antivirus detection methods? (Choose all that apply)

\begin{optlist}
\item Signature-based
\item Heuristic
\item Sandboxing
\item Behavioral analysis
\end{optlist}
\begin{explainbox}{}
\textbf{Correct Answer: A,B,C,D}\\[0.3em]
\textbf{Concept:} Modern antivirus uses all these techniques.
\end{explainbox}
\end{questioncard}

\begin{questioncard}{7}{multicol}
\qtypebadge{multicol}{Multiple Choice} \hfill \textcolor{darkslate!70}{\small Topic: Intrusion Prevention and Antivirus}

Which of the following relate to False? (Choose all that apply)

\begin{optlist}
\item False positive
\item False negative
\end{optlist}
\begin{explainbox}{}
\textbf{Correct Answer: A,B}\\[0.3em]
\textbf{Concept:} These concepts all relate to False in Intrusion Prevention and Antivirus.
\end{explainbox}
\end{questioncard}

\section{Intrusion Prevention}
\begin{questioncard}{1}{multicol}
\qtypebadge{multicol}{Multiple Choice} \hfill \textcolor{darkslate!70}{\small Topic: Intrusion Prevention}

Which are common IPS actions? (Choose all that apply)

\begin{optlist}
\item Block
\item Permit
\item Reset connection
\item Alert/log
\end{optlist}
\begin{explainbox}{}
\textbf{Correct Answer: A,C,D}\\[0.3em]
\textbf{Concept:} IPS can block malicious traffic, reset connections, and alert/log events.
\end{explainbox}
\end{questioncard}

\section{Antivirus}
\begin{questioncard}{1}{multicol}
\qtypebadge{multicol}{Multiple Choice} \hfill \textcolor{darkslate!70}{\small Topic: Antivirus}

Which are antivirus detection techniques? (Choose all that apply)

\begin{optlist}
\item Signature-based
\item Heuristic
\item Sandboxing
\item Behavioral analysis
\end{optlist}
\begin{explainbox}{}
\textbf{Correct Answer: A,B,C,D}\\[0.3em]
\textbf{Concept:} Modern antivirus uses signatures, heuristics, sandboxing, and behavioral analysis.
\end{explainbox}
\end{questioncard}

\section{AAA and User Authentication}
\begin{questioncard}{1}{multicol}
\qtypebadge{multicol}{Multiple Choice} \hfill \textcolor{darkslate!70}{\small Topic: AAA and User Authentication}

Which are AAA functions? (Choose all that apply)

\begin{optlist}
\item Authentication
\item Authorization
\item Accounting
\item Alerting
\end{optlist}
\begin{explainbox}{}
\textbf{Correct Answer: A,B,C}\\[0.3em]
\textbf{Concept:} AAA provides authentication, authorization, and accounting.
\end{explainbox}
\end{questioncard}

\begin{questioncard}{2}{multicol}
\qtypebadge{multicol}{Multiple Choice} \hfill \textcolor{darkslate!70}{\small Topic: AAA and User Authentication}

Which can be authentication sources on Huawei firewalls? (Choose all that apply)

\begin{optlist}
\item Local database
\item RADIUS server
\item HWTACACS server
\item LDAP server
\end{optlist}
\begin{explainbox}{}
\textbf{Correct Answer: A,B,C,D}\\[0.3em]
\textbf{Concept:} Firewalls can authenticate against local, RADIUS, HWTACACS, and LDAP sources.
\end{explainbox}
\end{questioncard}

\begin{questioncard}{3}{multicol}
\qtypebadge{multicol}{Multiple Choice} \hfill \textcolor{darkslate!70}{\small Topic: AAA and User Authentication}

Which are differences between RADIUS and TACACS+? (Choose all that apply)

\begin{optlist}
\item RADIUS uses UDP; TACACS+ uses TCP
\item RADIUS combines auth/authz; TACACS+ separates them
\item RADIUS encrypts only password; TACACS+ encrypts entire payload
\item RADIUS is Cisco proprietary
\end{optlist}
\begin{explainbox}{}
\textbf{Correct Answer: A,B,C}\\[0.3em]
\textbf{Concept:} RADIUS uses UDP and combines auth/authz; TACACS+ uses TCP and separates AAA; only the password is encrypted in RADIUS.
\end{explainbox}
\end{questioncard}

\begin{questioncard}{4}{multicol}
\qtypebadge{multicol}{Multiple Choice} \hfill \textcolor{darkslate!70}{\small Topic: AAA and User Authentication}

Which are AAA components? (Choose all that apply)

\begin{optlist}
\item Authentication
\item Authorization
\item Accounting
\item Auditing
\end{optlist}
\begin{explainbox}{}
\textbf{Correct Answer: A,B,C}\\[0.3em]
\textbf{Concept:} AAA consists of authentication, authorization, and accounting.
\end{explainbox}
\end{questioncard}

\begin{questioncard}{5}{multicol}
\qtypebadge{multicol}{Multiple Choice} \hfill \textcolor{darkslate!70}{\small Topic: AAA and User Authentication}

Which can authenticate users on a Huawei firewall? (Choose all that apply)

\begin{optlist}
\item Local database
\item RADIUS
\item HWTACACS
\item LDAP
\end{optlist}
\begin{explainbox}{}
\textbf{Correct Answer: A,B,C,D}\\[0.3em]
\textbf{Concept:} Firewalls support all these authentication sources.
\end{explainbox}
\end{questioncard}

\section{AAA Basics}
\begin{questioncard}{1}{multicol}
\qtypebadge{multicol}{Multiple Choice} \hfill \textcolor{darkslate!70}{\small Topic: AAA Basics}

Which are components of AAA? (Choose all that apply)

\begin{optlist}
\item Authentication
\item Authorization
\item Accounting
\item Auditing
\end{optlist}
\begin{explainbox}{}
\textbf{Correct Answer: A,B,C}\\[0.3em]
\textbf{Concept:} AAA consists of authentication, authorization, and accounting.
\end{explainbox}
\end{questioncard}

\section{Firewall User Authentication}
\begin{questioncard}{1}{multicol}
\qtypebadge{multicol}{Multiple Choice} \hfill \textcolor{darkslate!70}{\small Topic: Firewall User Authentication}

Which are common firewall authentication methods? (Choose all that apply)

\begin{optlist}
\item Local authentication
\item RADIUS
\item LDAP
\item HWTACACS
\end{optlist}
\begin{explainbox}{}
\textbf{Correct Answer: A,B,C,D}\\[0.3em]
\textbf{Concept:} Firewalls support local, RADIUS, LDAP, and HWTACACS authentication.
\end{explainbox}
\end{questioncard}

\section{Cryptography and PKI}
\begin{questioncard}{1}{multicol}
\qtypebadge{multicol}{Multiple Choice} \hfill \textcolor{darkslate!70}{\small Topic: Cryptography and PKI}

Which are asymmetric algorithms? (Choose all that apply)

\begin{optlist}
\item RSA
\item ECC
\item DSA
\item AES
\end{optlist}
\begin{explainbox}{}
\textbf{Correct Answer: A,B,C}\\[0.3em]
\textbf{Concept:} RSA, ECC, and DSA are asymmetric; AES is symmetric.
\end{explainbox}
\end{questioncard}

\begin{questioncard}{2}{multicol}
\qtypebadge{multicol}{Multiple Choice} \hfill \textcolor{darkslate!70}{\small Topic: Cryptography and PKI}

Which are block cipher modes? (Choose all that apply)

\begin{optlist}
\item ECB
\item CBC
\item CTR
\item GCM
\end{optlist}
\begin{explainbox}{}
\textbf{Correct Answer: A,B,C,D}\\[0.3em]
\textbf{Concept:} ECB, CBC, CTR, and GCM are all AES block cipher modes.
\end{explainbox}
\end{questioncard}

\begin{questioncard}{3}{multicol}
\qtypebadge{multicol}{Multiple Choice} \hfill \textcolor{darkslate!70}{\small Topic: Cryptography and PKI}

Which are secure hash algorithms? (Choose all that apply)

\begin{optlist}
\item SHA-256
\item SHA-3
\item BLAKE2
\item MD5
\end{optlist}
\begin{explainbox}{}
\textbf{Correct Answer: A,B,C}\\[0.3em]
\textbf{Concept:} SHA-256, SHA-3, and BLAKE2 are considered secure; MD5 is broken.
\end{explainbox}
\end{questioncard}

\begin{questioncard}{4}{multicol}
\qtypebadge{multicol}{Multiple Choice} \hfill \textcolor{darkslate!70}{\small Topic: Cryptography and PKI}

Which are PKI trust models? (Choose all that apply)

\begin{optlist}
\item Single CA
\item Hierarchical CA
\item Cross-certification
\item Bridge CA
\end{optlist}
\begin{explainbox}{}
\textbf{Correct Answer: A,B,C,D}\\[0.3em]
\textbf{Concept:} PKI can use single, hierarchical, cross-certified, or bridge CA models.
\end{explainbox}
\end{questioncard}

\begin{questioncard}{5}{multicol}
\qtypebadge{multicol}{Multiple Choice} \hfill \textcolor{darkslate!70}{\small Topic: Cryptography and PKI}

Which fields are typically found in an X.509 certificate? (Choose all that apply)

\begin{optlist}
\item Subject
\item Issuer
\item Public key
\item Validity period
\end{optlist}
\begin{explainbox}{}
\textbf{Correct Answer: A,B,C,D}\\[0.3em]
\textbf{Concept:} X.509 certificates contain subject, issuer, public key, validity, serial number, and signature.
\end{explainbox}
\end{questioncard}

\begin{questioncard}{6}{multicol}
\qtypebadge{multicol}{Multiple Choice} \hfill \textcolor{darkslate!70}{\small Topic: Cryptography and PKI}

Which are symmetric algorithms? (Choose all that apply)

\begin{optlist}
\item AES
\item DES
\item 3DES
\item RSA
\end{optlist}
\begin{explainbox}{}
\textbf{Correct Answer: A,B,C}\\[0.3em]
\textbf{Concept:} AES, DES, and 3DES are symmetric; RSA is asymmetric.
\end{explainbox}
\end{questioncard}

\begin{questioncard}{7}{multicol}
\qtypebadge{multicol}{Multiple Choice} \hfill \textcolor{darkslate!70}{\small Topic: Cryptography and PKI}

Which are hash algorithms? (Choose all that apply)

\begin{optlist}
\item MD5
\item SHA-256
\item SHA-1
\item SHA-3
\end{optlist}
\begin{explainbox}{}
\textbf{Correct Answer: A,B,C,D}\\[0.3em]
\textbf{Concept:} These are all hash algorithms, though MD5 and SHA-1 are deprecated for security use.
\end{explainbox}
\end{questioncard}

\begin{questioncard}{8}{multicol}
\qtypebadge{multicol}{Multiple Choice} \hfill \textcolor{darkslate!70}{\small Topic: Cryptography and PKI}

Which are PKI components? (Choose all that apply)

\begin{optlist}
\item CA
\item RA
\item Repository
\item CRL/OCSP
\end{optlist}
\begin{explainbox}{}
\textbf{Correct Answer: A,B,C,D}\\[0.3em]
\textbf{Concept:} PKI includes certificate authorities, registration authorities, repositories, and revocation systems.
\end{explainbox}
\end{questioncard}

\section{Cryptography Basics}
\begin{questioncard}{1}{multicol}
\qtypebadge{multicol}{Multiple Choice} \hfill \textcolor{darkslate!70}{\small Topic: Cryptography Basics}

Which are symmetric encryption algorithms? (Choose all that apply)

\begin{optlist}
\item AES
\item DES
\item 3DES
\item RSA
\end{optlist}
\begin{explainbox}{}
\textbf{Correct Answer: A,B,C}\\[0.3em]
\textbf{Concept:} AES, DES, and 3DES are symmetric; RSA is asymmetric.
\end{explainbox}
\end{questioncard}

\section{PKI}
\begin{questioncard}{1}{multicol}
\qtypebadge{multicol}{Multiple Choice} \hfill \textcolor{darkslate!70}{\small Topic: PKI}

Which are core PKI components? (Choose all that apply)

\begin{optlist}
\item Certificate Authority
\item Registration Authority
\item Certificate repository
\item CRL/OCSP
\end{optlist}
\begin{explainbox}{}
\textbf{Correct Answer: A,B,C,D}\\[0.3em]
\textbf{Concept:} PKI includes CAs, RAs, repositories, and revocation mechanisms such as CRL and OCSP.
\end{explainbox}
\end{questioncard}

\section{VPN}
\begin{questioncard}{1}{multicol}
\qtypebadge{multicol}{Multiple Choice} \hfill \textcolor{darkslate!70}{\small Topic: VPN}

Which are benefits of using a VPN? (Choose all that apply)

\begin{optlist}
\item Confidentiality
\item Data integrity
\item Authentication
\item Secure remote access
\end{optlist}
\begin{explainbox}{}
\textbf{Correct Answer: A,B,C,D}\\[0.3em]
\textbf{Concept:} VPNs provide confidentiality, integrity, authentication, and secure remote access.
\end{explainbox}
\end{questioncard}

\begin{questioncard}{2}{multicol}
\qtypebadge{multicol}{Multiple Choice} \hfill \textcolor{darkslate!70}{\small Topic: VPN}

Which are IPsec security services? (Choose all that apply)

\begin{optlist}
\item Confidentiality
\item Integrity
\item Authentication
\item Anti-replay
\end{optlist}
\begin{explainbox}{}
\textbf{Correct Answer: A,B,C,D}\\[0.3em]
\textbf{Concept:} IPsec provides confidentiality, integrity, authentication, and anti-replay protection.
\end{explainbox}
\end{questioncard}

\begin{questioncard}{3}{multicol}
\qtypebadge{multicol}{Multiple Choice} \hfill \textcolor{darkslate!70}{\small Topic: VPN}

Which are phases of IKEv1? (Choose all that apply)

\begin{optlist}
\item Phase 1
\item Phase 2
\item Phase 3
\item Phase 1.5
\end{optlist}
\begin{explainbox}{}
\textbf{Correct Answer: A,B}\\[0.3em]
\textbf{Concept:} IKEv1 has two phases: Phase 1 (IKE SA) and Phase 2 (IPsec SA).
\end{explainbox}
\end{questioncard}

\begin{questioncard}{4}{multicol}
\qtypebadge{multicol}{Multiple Choice} \hfill \textcolor{darkslate!70}{\small Topic: VPN}

Which protocols operate at the IP layer? (Choose all that apply)

\begin{optlist}
\item AH
\item ESP
\item GRE
\item L2TP
\end{optlist}
\begin{explainbox}{}
\textbf{Correct Answer: A,B}\\[0.3em]
\textbf{Concept:} AH and ESP are IP-layer protocols; GRE is also IP-layer but L2TP is not.
\end{explainbox}
\end{questioncard}

\begin{questioncard}{5}{multicol}
\qtypebadge{multicol}{Multiple Choice} \hfill \textcolor{darkslate!70}{\small Topic: VPN}

Which are VPN types? (Choose all that apply)

\begin{optlist}
\item Site-to-site
\item Remote-access
\item Intranet
\item Extranet
\end{optlist}
\begin{explainbox}{}
\textbf{Correct Answer: A,B,C,D}\\[0.3em]
\textbf{Concept:} VPNs can be categorized by these deployment models.
\end{explainbox}
\end{questioncard}

\begin{questioncard}{6}{multicol}
\qtypebadge{multicol}{Multiple Choice} \hfill \textcolor{darkslate!70}{\small Topic: VPN}

Which are IPsec protocols? (Choose all that apply)

\begin{optlist}
\item AH
\item ESP
\item IKE
\item GRE
\end{optlist}
\begin{explainbox}{}
\textbf{Correct Answer: A,B,C}\\[0.3em]
\textbf{Concept:} AH, ESP, and IKE are IPsec protocols; GRE is separate.
\end{explainbox}
\end{questioncard}

\begin{questioncard}{7}{multicol}
\qtypebadge{multicol}{Multiple Choice} \hfill \textcolor{darkslate!70}{\small Topic: VPN}

Which are SSL VPN access modes? (Choose all that apply)

\begin{optlist}
\item Web proxy
\item File sharing
\item Network extension
\item Port forwarding
\end{optlist}
\begin{explainbox}{}
\textbf{Correct Answer: A,B,C,D}\\[0.3em]
\textbf{Concept:} SSL VPN commonly supports all these modes.
\end{explainbox}
\end{questioncard}

\begin{questioncard}{8}{multicol}
\qtypebadge{multicol}{Multiple Choice} \hfill \textcolor{darkslate!70}{\small Topic: VPN}

Which of the following relate to GRE? (Choose all that apply)

\begin{optlist}
\item GRE
\item GRE
\end{optlist}
\begin{explainbox}{}
\textbf{Correct Answer: A,B}\\[0.3em]
\textbf{Concept:} These concepts all relate to GRE in VPN.
\end{explainbox}
\end{questioncard}

\begin{questioncard}{9}{multicol}
\qtypebadge{multicol}{Multiple Choice} \hfill \textcolor{darkslate!70}{\small Topic: VPN}

Which of the following relate to L2TP? (Choose all that apply)

\begin{optlist}
\item L2TP
\item L2TP over IPsec
\end{optlist}
\begin{explainbox}{}
\textbf{Correct Answer: A,B}\\[0.3em]
\textbf{Concept:} These concepts all relate to L2TP in VPN.
\end{explainbox}
\end{questioncard}

\section{VPN Overview}
\begin{questioncard}{1}{multicol}
\qtypebadge{multicol}{Multiple Choice} \hfill \textcolor{darkslate!70}{\small Topic: VPN Overview}

Which are common VPN types? (Choose all that apply)

\begin{optlist}
\item Site-to-site VPN
\item Remote-access VPN
\item Intranet VPN
\item Extranet VPN
\end{optlist}
\begin{explainbox}{}
\textbf{Correct Answer: A,B,C,D}\\[0.3em]
\textbf{Concept:} VPNs can be site-to-site, remote-access, intranet, or extranet based on deployment.
\end{explainbox}
\end{questioncard}

\section{IPsec VPN}
\begin{questioncard}{1}{multicol}
\qtypebadge{multicol}{Multiple Choice} \hfill \textcolor{darkslate!70}{\small Topic: IPsec VPN}

Which are IPsec protocols? (Choose all that apply)

\begin{optlist}
\item AH
\item ESP
\item IKE
\item GRE
\end{optlist}
\begin{explainbox}{}
\textbf{Correct Answer: A,B,C}\\[0.3em]
\textbf{Concept:} AH, ESP, and IKE are part of IPsec; GRE is a separate tunneling protocol.
\end{explainbox}
\end{questioncard}

\section{SSL VPN}
\begin{questioncard}{1}{multicol}
\qtypebadge{multicol}{Multiple Choice} \hfill \textcolor{darkslate!70}{\small Topic: SSL VPN}

Which are common SSL VPN access modes? (Choose all that apply)

\begin{optlist}
\item Web proxy
\item File sharing
\item Network extension
\item Port forwarding
\end{optlist}
\begin{explainbox}{}
\textbf{Correct Answer: A,B,C,D}\\[0.3em]
\textbf{Concept:} SSL VPN can offer web proxy, file sharing, network extension, and port forwarding.
\end{explainbox}
\end{questioncard}

\part{True / False Questions}
\section{Network Security Concepts}
\begin{questioncard}{1}{tfcol}
\qtypebadge{tfcol}{True / False} \hfill \textcolor{darkslate!70}{\small Topic: Network Security Concepts}

Network security in the narrow sense only refers to national cybersecurity laws.

\begin{optlist}
\item True
\item False
\end{optlist}
\begin{explainbox}{}
\textbf{Correct Answer: False}\\[0.3em]
\textbf{Concept:} Narrow-sense network security covers vendor solutions and enterprise network protection, not only national laws.
\end{explainbox}
\end{questioncard}

\begin{questioncard}{2}{tfcol}
\qtypebadge{tfcol}{True / False} \hfill \textcolor{darkslate!70}{\small Topic: Network Security Concepts}

The CIA triad includes confidentiality, integrity, and availability.

\begin{optlist}
\item True
\item False
\end{optlist}
\begin{explainbox}{}
\textbf{Correct Answer: True}\\[0.3em]
\textbf{Concept:} Confidentiality, integrity, and availability are the three foundational goals of information security.
\end{explainbox}
\end{questioncard}

\begin{questioncard}{3}{tfcol}
\qtypebadge{tfcol}{True / False} \hfill \textcolor{darkslate!70}{\small Topic: Network Security Concepts}

Building network security helps enterprises reduce losses from attacks.

\begin{optlist}
\item True
\item False
\end{optlist}
\begin{explainbox}{}
\textbf{Correct Answer: True}\\[0.3em]
\textbf{Concept:} Security controls reduce the impact of breaches, downtime, and data loss, thereby lowering financial losses.
\end{explainbox}
\end{questioncard}

\begin{questioncard}{4}{tfcol}
\qtypebadge{tfcol}{True / False} \hfill \textcolor{darkslate!70}{\small Topic: Network Security Concepts}

Non-repudiation ensures that senders or receivers cannot deny their actions.

\begin{optlist}
\item True
\item False
\end{optlist}
\begin{explainbox}{}
\textbf{Correct Answer: True}\\[0.3em]
\textbf{Concept:} Non-repudiation provides evidence to prevent denial of participation.
\end{explainbox}
\end{questioncard}

\begin{questioncard}{5}{tfcol}
\qtypebadge{tfcol}{True / False} \hfill \textcolor{darkslate!70}{\small Topic: Network Security Concepts}

Availability means anyone can access the system at any time.

\begin{optlist}
\item True
\item False
\end{optlist}
\begin{explainbox}{}
\textbf{Correct Answer: False}\\[0.3em]
\textbf{Concept:} Availability means authorized users can access information when required, not unlimited access.
\end{explainbox}
\end{questioncard}

\begin{questioncard}{6}{tfcol}
\qtypebadge{tfcol}{True / False} \hfill \textcolor{darkslate!70}{\small Topic: Network Security Concepts}

Zero trust means everything inside the network is trusted by default.

\begin{optlist}
\item True
\item False
\end{optlist}
\begin{explainbox}{}
\textbf{Correct Answer: False}\\[0.3em]
\textbf{Concept:} Zero trust never assumes trust based on network location and verifies every request.
\end{explainbox}
\end{questioncard}

\begin{questioncard}{7}{tfcol}
\qtypebadge{tfcol}{True / False} \hfill \textcolor{darkslate!70}{\small Topic: Network Security Concepts}

Integrity ensures that data is not tampered with during transmission.

\begin{optlist}
\item True
\item False
\end{optlist}
\begin{explainbox}{}
\textbf{Correct Answer: True}\\[0.3em]
\textbf{Concept:} Integrity protects against unauthorized modification or destruction of data.
\end{explainbox}
\end{questioncard}

\begin{questioncard}{8}{tfcol}
\qtypebadge{tfcol}{True / False} \hfill \textcolor{darkslate!70}{\small Topic: Network Security Concepts}

Network security broadly refers to cyberspace security at a national level.

\begin{optlist}
\item True
\item False
\end{optlist}
\begin{explainbox}{}
\textbf{Correct Answer: True}\\[0.3em]
\textbf{Concept:} Broad network security covers national cyberspace protection.
\end{explainbox}
\end{questioncard}

\begin{questioncard}{9}{tfcol}
\qtypebadge{tfcol}{True / False} \hfill \textcolor{darkslate!70}{\small Topic: Network Security Concepts}

The CIA triad excludes availability.

\begin{optlist}
\item True
\item False
\end{optlist}
\begin{explainbox}{}
\textbf{Correct Answer: False}\\[0.3em]
\textbf{Concept:} The CIA triad includes confidentiality, integrity, and availability.
\end{explainbox}
\end{questioncard}

\begin{questioncard}{10}{tfcol}
\qtypebadge{tfcol}{True / False} \hfill \textcolor{darkslate!70}{\small Topic: Network Security Concepts}

Risk is the combination of threat likelihood and impact.

\begin{optlist}
\item True
\item False
\end{optlist}
\begin{explainbox}{}
\textbf{Correct Answer: True}\\[0.3em]
\textbf{Concept:} Risk assesses probability and consequence of security incidents.
\end{explainbox}
\end{questioncard}

\begin{questioncard}{11}{tfcol}
\qtypebadge{tfcol}{True / False} \hfill \textcolor{darkslate!70}{\small Topic: Network Security Concepts}

Broad network security means that is synonymous with national cyberspace security.

\begin{optlist}
\item True
\item False
\end{optlist}
\begin{explainbox}{}
\textbf{Correct Answer: True}\\[0.3em]
\textbf{Concept:} This statement correctly describes Broad network security.
\end{explainbox}
\end{questioncard}

\begin{questioncard}{12}{tfcol}
\qtypebadge{tfcol}{True / False} \hfill \textcolor{darkslate!70}{\small Topic: Network Security Concepts}

Broad network security is unrelated to is synonymous with national cyberspace security.

\begin{optlist}
\item True
\item False
\end{optlist}
\begin{explainbox}{}
\textbf{Correct Answer: False}\\[0.3em]
\textbf{Concept:} Broad network security is directly related to is synonymous with national cyberspace security.
\end{explainbox}
\end{questioncard}

\begin{questioncard}{13}{tfcol}
\qtypebadge{tfcol}{True / False} \hfill \textcolor{darkslate!70}{\small Topic: Network Security Concepts}

Narrow network security means that focuses on enterprise network protection solutions.

\begin{optlist}
\item True
\item False
\end{optlist}
\begin{explainbox}{}
\textbf{Correct Answer: True}\\[0.3em]
\textbf{Concept:} This statement correctly describes Narrow network security.
\end{explainbox}
\end{questioncard}

\begin{questioncard}{14}{tfcol}
\qtypebadge{tfcol}{True / False} \hfill \textcolor{darkslate!70}{\small Topic: Network Security Concepts}

Narrow network security is unrelated to focuses on enterprise network protection solutions.

\begin{optlist}
\item True
\item False
\end{optlist}
\begin{explainbox}{}
\textbf{Correct Answer: False}\\[0.3em]
\textbf{Concept:} Narrow network security is directly related to focuses on enterprise network protection solutions.
\end{explainbox}
\end{questioncard}

\begin{questioncard}{15}{tfcol}
\qtypebadge{tfcol}{True / False} \hfill \textcolor{darkslate!70}{\small Topic: Network Security Concepts}

CIA triad means that consists of confidentiality, integrity, and availability.

\begin{optlist}
\item True
\item False
\end{optlist}
\begin{explainbox}{}
\textbf{Correct Answer: True}\\[0.3em]
\textbf{Concept:} This statement correctly describes CIA triad.
\end{explainbox}
\end{questioncard}

\begin{questioncard}{16}{tfcol}
\qtypebadge{tfcol}{True / False} \hfill \textcolor{darkslate!70}{\small Topic: Network Security Concepts}

CIA triad is unrelated to consists of confidentiality, integrity, and availability.

\begin{optlist}
\item True
\item False
\end{optlist}
\begin{explainbox}{}
\textbf{Correct Answer: False}\\[0.3em]
\textbf{Concept:} CIA triad is directly related to consists of confidentiality, integrity, and availability.
\end{explainbox}
\end{questioncard}

\begin{questioncard}{17}{tfcol}
\qtypebadge{tfcol}{True / False} \hfill \textcolor{darkslate!70}{\small Topic: Network Security Concepts}

Information assurance period means that added controllability and non-repudiation.

\begin{optlist}
\item True
\item False
\end{optlist}
\begin{explainbox}{}
\textbf{Correct Answer: True}\\[0.3em]
\textbf{Concept:} This statement correctly describes Information assurance period.
\end{explainbox}
\end{questioncard}

\begin{questioncard}{18}{tfcol}
\qtypebadge{tfcol}{True / False} \hfill \textcolor{darkslate!70}{\small Topic: Network Security Concepts}

Information assurance period is unrelated to added controllability and non-repudiation.

\begin{optlist}
\item True
\item False
\end{optlist}
\begin{explainbox}{}
\textbf{Correct Answer: False}\\[0.3em]
\textbf{Concept:} Information assurance period is directly related to added controllability and non-repudiation.
\end{explainbox}
\end{questioncard}

\begin{questioncard}{19}{tfcol}
\qtypebadge{tfcol}{True / False} \hfill \textcolor{darkslate!70}{\small Topic: Network Security Concepts}

Cyberspace security means that protects facilities, data, users, and operations.

\begin{optlist}
\item True
\item False
\end{optlist}
\begin{explainbox}{}
\textbf{Correct Answer: True}\\[0.3em]
\textbf{Concept:} This statement correctly describes Cyberspace security.
\end{explainbox}
\end{questioncard}

\begin{questioncard}{20}{tfcol}
\qtypebadge{tfcol}{True / False} \hfill \textcolor{darkslate!70}{\small Topic: Network Security Concepts}

Cyberspace security is unrelated to protects facilities, data, users, and operations.

\begin{optlist}
\item True
\item False
\end{optlist}
\begin{explainbox}{}
\textbf{Correct Answer: False}\\[0.3em]
\textbf{Concept:} Cyberspace security is directly related to protects facilities, data, users, and operations.
\end{explainbox}
\end{questioncard}

\begin{questioncard}{21}{tfcol}
\qtypebadge{tfcol}{True / False} \hfill \textcolor{darkslate!70}{\small Topic: Network Security Concepts}

Confidentiality means that ensures information is accessible only to authorized users.

\begin{optlist}
\item True
\item False
\end{optlist}
\begin{explainbox}{}
\textbf{Correct Answer: True}\\[0.3em]
\textbf{Concept:} This statement correctly describes Confidentiality.
\end{explainbox}
\end{questioncard}

\begin{questioncard}{22}{tfcol}
\qtypebadge{tfcol}{True / False} \hfill \textcolor{darkslate!70}{\small Topic: Network Security Concepts}

Confidentiality is unrelated to ensures information is accessible only to authorized users.

\begin{optlist}
\item True
\item False
\end{optlist}
\begin{explainbox}{}
\textbf{Correct Answer: False}\\[0.3em]
\textbf{Concept:} Confidentiality is directly related to ensures information is accessible only to authorized users.
\end{explainbox}
\end{questioncard}

\begin{questioncard}{23}{tfcol}
\qtypebadge{tfcol}{True / False} \hfill \textcolor{darkslate!70}{\small Topic: Network Security Concepts}

Integrity means that ensures data is not tampered with during transmission or storage.

\begin{optlist}
\item True
\item False
\end{optlist}
\begin{explainbox}{}
\textbf{Correct Answer: True}\\[0.3em]
\textbf{Concept:} This statement correctly describes Integrity.
\end{explainbox}
\end{questioncard}

\begin{questioncard}{24}{tfcol}
\qtypebadge{tfcol}{True / False} \hfill \textcolor{darkslate!70}{\small Topic: Network Security Concepts}

Integrity is unrelated to ensures data is not tampered with during transmission or storage.

\begin{optlist}
\item True
\item False
\end{optlist}
\begin{explainbox}{}
\textbf{Correct Answer: False}\\[0.3em]
\textbf{Concept:} Integrity is directly related to ensures data is not tampered with during transmission or storage.
\end{explainbox}
\end{questioncard}

\begin{questioncard}{25}{tfcol}
\qtypebadge{tfcol}{True / False} \hfill \textcolor{darkslate!70}{\small Topic: Network Security Concepts}

Availability means that ensures authorized users can access information when needed.

\begin{optlist}
\item True
\item False
\end{optlist}
\begin{explainbox}{}
\textbf{Correct Answer: True}\\[0.3em]
\textbf{Concept:} This statement correctly describes Availability.
\end{explainbox}
\end{questioncard}

\begin{questioncard}{26}{tfcol}
\qtypebadge{tfcol}{True / False} \hfill \textcolor{darkslate!70}{\small Topic: Network Security Concepts}

Availability is unrelated to ensures authorized users can access information when needed.

\begin{optlist}
\item True
\item False
\end{optlist}
\begin{explainbox}{}
\textbf{Correct Answer: False}\\[0.3em]
\textbf{Concept:} Availability is directly related to ensures authorized users can access information when needed.
\end{explainbox}
\end{questioncard}

\begin{questioncard}{27}{tfcol}
\qtypebadge{tfcol}{True / False} \hfill \textcolor{darkslate!70}{\small Topic: Network Security Concepts}

Controllability means that enables monitoring and control of information and systems.

\begin{optlist}
\item True
\item False
\end{optlist}
\begin{explainbox}{}
\textbf{Correct Answer: True}\\[0.3em]
\textbf{Concept:} This statement correctly describes Controllability.
\end{explainbox}
\end{questioncard}

\begin{questioncard}{28}{tfcol}
\qtypebadge{tfcol}{True / False} \hfill \textcolor{darkslate!70}{\small Topic: Network Security Concepts}

Controllability is unrelated to enables monitoring and control of information and systems.

\begin{optlist}
\item True
\item False
\end{optlist}
\begin{explainbox}{}
\textbf{Correct Answer: False}\\[0.3em]
\textbf{Concept:} Controllability is directly related to enables monitoring and control of information and systems.
\end{explainbox}
\end{questioncard}

\begin{questioncard}{29}{tfcol}
\qtypebadge{tfcol}{True / False} \hfill \textcolor{darkslate!70}{\small Topic: Network Security Concepts}

Non-repudiation means that prevents senders or receivers from denying actions.

\begin{optlist}
\item True
\item False
\end{optlist}
\begin{explainbox}{}
\textbf{Correct Answer: True}\\[0.3em]
\textbf{Concept:} This statement correctly describes Non-repudiation.
\end{explainbox}
\end{questioncard}

\begin{questioncard}{30}{tfcol}
\qtypebadge{tfcol}{True / False} \hfill \textcolor{darkslate!70}{\small Topic: Network Security Concepts}

Non-repudiation is unrelated to prevents senders or receivers from denying actions.

\begin{optlist}
\item True
\item False
\end{optlist}
\begin{explainbox}{}
\textbf{Correct Answer: False}\\[0.3em]
\textbf{Concept:} Non-repudiation is directly related to prevents senders or receivers from denying actions.
\end{explainbox}
\end{questioncard}

\begin{questioncard}{31}{tfcol}
\qtypebadge{tfcol}{True / False} \hfill \textcolor{darkslate!70}{\small Topic: Network Security Concepts}

Risk means that is the combination of threat likelihood and impact.

\begin{optlist}
\item True
\item False
\end{optlist}
\begin{explainbox}{}
\textbf{Correct Answer: True}\\[0.3em]
\textbf{Concept:} This statement correctly describes Risk.
\end{explainbox}
\end{questioncard}

\begin{questioncard}{32}{tfcol}
\qtypebadge{tfcol}{True / False} \hfill \textcolor{darkslate!70}{\small Topic: Network Security Concepts}

Risk is unrelated to is the combination of threat likelihood and impact.

\begin{optlist}
\item True
\item False
\end{optlist}
\begin{explainbox}{}
\textbf{Correct Answer: False}\\[0.3em]
\textbf{Concept:} Risk is directly related to is the combination of threat likelihood and impact.
\end{explainbox}
\end{questioncard}

\begin{questioncard}{33}{tfcol}
\qtypebadge{tfcol}{True / False} \hfill \textcolor{darkslate!70}{\small Topic: Network Security Concepts}

Vulnerability means that is a weakness that can be exploited by threats.

\begin{optlist}
\item True
\item False
\end{optlist}
\begin{explainbox}{}
\textbf{Correct Answer: True}\\[0.3em]
\textbf{Concept:} This statement correctly describes Vulnerability.
\end{explainbox}
\end{questioncard}

\begin{questioncard}{34}{tfcol}
\qtypebadge{tfcol}{True / False} \hfill \textcolor{darkslate!70}{\small Topic: Network Security Concepts}

Vulnerability is unrelated to is a weakness that can be exploited by threats.

\begin{optlist}
\item True
\item False
\end{optlist}
\begin{explainbox}{}
\textbf{Correct Answer: False}\\[0.3em]
\textbf{Concept:} Vulnerability is directly related to is a weakness that can be exploited by threats.
\end{explainbox}
\end{questioncard}

\begin{questioncard}{35}{tfcol}
\qtypebadge{tfcol}{True / False} \hfill \textcolor{darkslate!70}{\small Topic: Network Security Concepts}

Threat means that is a potential danger to an asset.

\begin{optlist}
\item True
\item False
\end{optlist}
\begin{explainbox}{}
\textbf{Correct Answer: True}\\[0.3em]
\textbf{Concept:} This statement correctly describes Threat.
\end{explainbox}
\end{questioncard}

\begin{questioncard}{36}{tfcol}
\qtypebadge{tfcol}{True / False} \hfill \textcolor{darkslate!70}{\small Topic: Network Security Concepts}

Threat is unrelated to is a potential danger to an asset.

\begin{optlist}
\item True
\item False
\end{optlist}
\begin{explainbox}{}
\textbf{Correct Answer: False}\\[0.3em]
\textbf{Concept:} Threat is directly related to is a potential danger to an asset.
\end{explainbox}
\end{questioncard}

\begin{questioncard}{37}{tfcol}
\qtypebadge{tfcol}{True / False} \hfill \textcolor{darkslate!70}{\small Topic: Network Security Concepts}

Exploit means that is a method that takes advantage of a vulnerability.

\begin{optlist}
\item True
\item False
\end{optlist}
\begin{explainbox}{}
\textbf{Correct Answer: True}\\[0.3em]
\textbf{Concept:} This statement correctly describes Exploit.
\end{explainbox}
\end{questioncard}

\begin{questioncard}{38}{tfcol}
\qtypebadge{tfcol}{True / False} \hfill \textcolor{darkslate!70}{\small Topic: Network Security Concepts}

Exploit is unrelated to is a method that takes advantage of a vulnerability.

\begin{optlist}
\item True
\item False
\end{optlist}
\begin{explainbox}{}
\textbf{Correct Answer: False}\\[0.3em]
\textbf{Concept:} Exploit is directly related to is a method that takes advantage of a vulnerability.
\end{explainbox}
\end{questioncard}

\begin{questioncard}{39}{tfcol}
\qtypebadge{tfcol}{True / False} \hfill \textcolor{darkslate!70}{\small Topic: Network Security Concepts}

Future security trend means that includes zero-trust architecture.

\begin{optlist}
\item True
\item False
\end{optlist}
\begin{explainbox}{}
\textbf{Correct Answer: True}\\[0.3em]
\textbf{Concept:} This statement correctly describes Future security trend.
\end{explainbox}
\end{questioncard}

\begin{questioncard}{40}{tfcol}
\qtypebadge{tfcol}{True / False} \hfill \textcolor{darkslate!70}{\small Topic: Network Security Concepts}

Future security trend is unrelated to includes zero-trust architecture.

\begin{optlist}
\item True
\item False
\end{optlist}
\begin{explainbox}{}
\textbf{Correct Answer: False}\\[0.3em]
\textbf{Concept:} Future security trend is directly related to includes zero-trust architecture.
\end{explainbox}
\end{questioncard}

\begin{questioncard}{41}{tfcol}
\qtypebadge{tfcol}{True / False} \hfill \textcolor{darkslate!70}{\small Topic: Network Security Concepts}

AI and machine learning means that help detect unknown threats and anomalies.

\begin{optlist}
\item True
\item False
\end{optlist}
\begin{explainbox}{}
\textbf{Correct Answer: True}\\[0.3em]
\textbf{Concept:} This statement correctly describes AI and machine learning.
\end{explainbox}
\end{questioncard}

\begin{questioncard}{42}{tfcol}
\qtypebadge{tfcol}{True / False} \hfill \textcolor{darkslate!70}{\small Topic: Network Security Concepts}

AI and machine learning is unrelated to help detect unknown threats and anomalies.

\begin{optlist}
\item True
\item False
\end{optlist}
\begin{explainbox}{}
\textbf{Correct Answer: False}\\[0.3em]
\textbf{Concept:} AI and machine learning is directly related to help detect unknown threats and anomalies.
\end{explainbox}
\end{questioncard}

\section{Network Security Standards}
\begin{questioncard}{1}{tfcol}
\qtypebadge{tfcol}{True / False} \hfill \textcolor{darkslate!70}{\small Topic: Network Security Standards}

ISO 27002 can be used to certify an enterprise information security system.

\begin{optlist}
\item True
\item False
\end{optlist}
\begin{explainbox}{}
\textbf{Correct Answer: False}\\[0.3em]
\textbf{Concept:} ISO 27001 is the certification standard; ISO 27002 is a code of practice with controls.
\end{explainbox}
\end{questioncard}

\begin{questioncard}{2}{tfcol}
\qtypebadge{tfcol}{True / False} \hfill \textcolor{darkslate!70}{\small Topic: Network Security Standards}

Classified Protection 2.0 defines protection levels from Level 1 to Level 5.

\begin{optlist}
\item True
\item False
\end{optlist}
\begin{explainbox}{}
\textbf{Correct Answer: True}\\[0.3em]
\textbf{Concept:} Classified Protection 2.0 uses five levels, with higher levels requiring stronger controls and oversight.
\end{explainbox}
\end{questioncard}

\begin{questioncard}{3}{tfcol}
\qtypebadge{tfcol}{True / False} \hfill \textcolor{darkslate!70}{\small Topic: Network Security Standards}

Network operators can always determine the final protection level on their own without review.

\begin{optlist}
\item True
\item False
\end{optlist}
\begin{explainbox}{}
\textbf{Correct Answer: False}\\[0.3em]
\textbf{Concept:} For Level 2 and above, expert review and approval are required; only Level 1 can be determined by the operator.
\end{explainbox}
\end{questioncard}

\begin{questioncard}{4}{tfcol}
\qtypebadge{tfcol}{True / False} \hfill \textcolor{darkslate!70}{\small Topic: Network Security Standards}

ISO 27001 provides a list of security controls but is not certifiable.

\begin{optlist}
\item True
\item False
\end{optlist}
\begin{explainbox}{}
\textbf{Correct Answer: False}\\[0.3em]
\textbf{Concept:} ISO 27001 is the certifiable standard specifying ISMS requirements; ISO 27002 provides controls.
\end{explainbox}
\end{questioncard}

\begin{questioncard}{5}{tfcol}
\qtypebadge{tfcol}{True / False} \hfill \textcolor{darkslate!70}{\small Topic: Network Security Standards}

Classified Protection 2.0 requires security and informatization to be planned together.

\begin{optlist}
\item True
\item False
\end{optlist}
\begin{explainbox}{}
\textbf{Correct Answer: True}\\[0.3em]
\textbf{Concept:} The framework emphasizes synchronized planning, construction, and use of security.
\end{explainbox}
\end{questioncard}

\begin{questioncard}{6}{tfcol}
\qtypebadge{tfcol}{True / False} \hfill \textcolor{darkslate!70}{\small Topic: Network Security Standards}

Level assessment is part of the Classified Protection lifecycle.

\begin{optlist}
\item True
\item False
\end{optlist}
\begin{explainbox}{}
\textbf{Correct Answer: True}\\[0.3em]
\textbf{Concept:} After construction, a level assessment verifies that protection measures meet the required level.
\end{explainbox}
\end{questioncard}

\begin{questioncard}{7}{tfcol}
\qtypebadge{tfcol}{True / False} \hfill \textcolor{darkslate!70}{\small Topic: Network Security Standards}

ISO 27001 is certifiable.

\begin{optlist}
\item True
\item False
\end{optlist}
\begin{explainbox}{}
\textbf{Correct Answer: True}\\[0.3em]
\textbf{Concept:} Organizations can be certified against ISO 27001 requirements.
\end{explainbox}
\end{questioncard}

\begin{questioncard}{8}{tfcol}
\qtypebadge{tfcol}{True / False} \hfill \textcolor{darkslate!70}{\small Topic: Network Security Standards}

Classified Protection 2.0 has only three levels.

\begin{optlist}
\item True
\item False
\end{optlist}
\begin{explainbox}{}
\textbf{Correct Answer: False}\\[0.3em]
\textbf{Concept:} Classified Protection 2.0 defines five levels.
\end{explainbox}
\end{questioncard}

\begin{questioncard}{9}{tfcol}
\qtypebadge{tfcol}{True / False} \hfill \textcolor{darkslate!70}{\small Topic: Network Security Standards}

ISO 27001 means that specifies certifiable ISMS requirements.

\begin{optlist}
\item True
\item False
\end{optlist}
\begin{explainbox}{}
\textbf{Correct Answer: True}\\[0.3em]
\textbf{Concept:} This statement correctly describes ISO 27001.
\end{explainbox}
\end{questioncard}

\begin{questioncard}{10}{tfcol}
\qtypebadge{tfcol}{True / False} \hfill \textcolor{darkslate!70}{\small Topic: Network Security Standards}

ISO 27001 is unrelated to specifies certifiable ISMS requirements.

\begin{optlist}
\item True
\item False
\end{optlist}
\begin{explainbox}{}
\textbf{Correct Answer: False}\\[0.3em]
\textbf{Concept:} ISO 27001 is directly related to specifies certifiable ISMS requirements.
\end{explainbox}
\end{questioncard}

\begin{questioncard}{11}{tfcol}
\qtypebadge{tfcol}{True / False} \hfill \textcolor{darkslate!70}{\small Topic: Network Security Standards}

ISO 27002 means that provides a code of practice for security controls.

\begin{optlist}
\item True
\item False
\end{optlist}
\begin{explainbox}{}
\textbf{Correct Answer: True}\\[0.3em]
\textbf{Concept:} This statement correctly describes ISO 27002.
\end{explainbox}
\end{questioncard}

\begin{questioncard}{12}{tfcol}
\qtypebadge{tfcol}{True / False} \hfill \textcolor{darkslate!70}{\small Topic: Network Security Standards}

ISO 27002 is unrelated to provides a code of practice for security controls.

\begin{optlist}
\item True
\item False
\end{optlist}
\begin{explainbox}{}
\textbf{Correct Answer: False}\\[0.3em]
\textbf{Concept:} ISO 27002 is directly related to provides a code of practice for security controls.
\end{explainbox}
\end{questioncard}

\begin{questioncard}{13}{tfcol}
\qtypebadge{tfcol}{True / False} \hfill \textcolor{darkslate!70}{\small Topic: Network Security Standards}

PDCA means that is the continuous improvement methodology used in ISO 27001.

\begin{optlist}
\item True
\item False
\end{optlist}
\begin{explainbox}{}
\textbf{Correct Answer: True}\\[0.3em]
\textbf{Concept:} This statement correctly describes PDCA.
\end{explainbox}
\end{questioncard}

\begin{questioncard}{14}{tfcol}
\qtypebadge{tfcol}{True / False} \hfill \textcolor{darkslate!70}{\small Topic: Network Security Standards}

PDCA is unrelated to is the continuous improvement methodology used in ISO 27001.

\begin{optlist}
\item True
\item False
\end{optlist}
\begin{explainbox}{}
\textbf{Correct Answer: False}\\[0.3em]
\textbf{Concept:} PDCA is directly related to is the continuous improvement methodology used in ISO 27001.
\end{explainbox}
\end{questioncard}

\begin{questioncard}{15}{tfcol}
\qtypebadge{tfcol}{True / False} \hfill \textcolor{darkslate!70}{\small Topic: Network Security Standards}

Classified Protection 2.0 means that is a Chinese national cybersecurity framework.

\begin{optlist}
\item True
\item False
\end{optlist}
\begin{explainbox}{}
\textbf{Correct Answer: True}\\[0.3em]
\textbf{Concept:} This statement correctly describes Classified Protection 2.0.
\end{explainbox}
\end{questioncard}

\begin{questioncard}{16}{tfcol}
\qtypebadge{tfcol}{True / False} \hfill \textcolor{darkslate!70}{\small Topic: Network Security Standards}

Classified Protection 2.0 is unrelated to is a Chinese national cybersecurity framework.

\begin{optlist}
\item True
\item False
\end{optlist}
\begin{explainbox}{}
\textbf{Correct Answer: False}\\[0.3em]
\textbf{Concept:} Classified Protection 2.0 is directly related to is a Chinese national cybersecurity framework.
\end{explainbox}
\end{questioncard}

\begin{questioncard}{17}{tfcol}
\qtypebadge{tfcol}{True / False} \hfill \textcolor{darkslate!70}{\small Topic: Network Security Standards}

Classified Protection levels means that range from Level 1 to Level 5.

\begin{optlist}
\item True
\item False
\end{optlist}
\begin{explainbox}{}
\textbf{Correct Answer: True}\\[0.3em]
\textbf{Concept:} This statement correctly describes Classified Protection levels.
\end{explainbox}
\end{questioncard}

\begin{questioncard}{18}{tfcol}
\qtypebadge{tfcol}{True / False} \hfill \textcolor{darkslate!70}{\small Topic: Network Security Standards}

Classified Protection levels is unrelated to range from Level 1 to Level 5.

\begin{optlist}
\item True
\item False
\end{optlist}
\begin{explainbox}{}
\textbf{Correct Answer: False}\\[0.3em]
\textbf{Concept:} Classified Protection levels is directly related to range from Level 1 to Level 5.
\end{explainbox}
\end{questioncard}

\begin{questioncard}{19}{tfcol}
\qtypebadge{tfcol}{True / False} \hfill \textcolor{darkslate!70}{\small Topic: Network Security Standards}

Level 2 and above means that require expert review and approval.

\begin{optlist}
\item True
\item False
\end{optlist}
\begin{explainbox}{}
\textbf{Correct Answer: True}\\[0.3em]
\textbf{Concept:} This statement correctly describes Level 2 and above.
\end{explainbox}
\end{questioncard}

\begin{questioncard}{20}{tfcol}
\qtypebadge{tfcol}{True / False} \hfill \textcolor{darkslate!70}{\small Topic: Network Security Standards}

Level 2 and above is unrelated to require expert review and approval.

\begin{optlist}
\item True
\item False
\end{optlist}
\begin{explainbox}{}
\textbf{Correct Answer: False}\\[0.3em]
\textbf{Concept:} Level 2 and above is directly related to require expert review and approval.
\end{explainbox}
\end{questioncard}

\begin{questioncard}{21}{tfcol}
\qtypebadge{tfcol}{True / False} \hfill \textcolor{darkslate!70}{\small Topic: Network Security Standards}

Classified Protection lifecycle means that includes classification, filing, construction, assessment, and supervision.

\begin{optlist}
\item True
\item False
\end{optlist}
\begin{explainbox}{}
\textbf{Correct Answer: True}\\[0.3em]
\textbf{Concept:} This statement correctly describes Classified Protection lifecycle.
\end{explainbox}
\end{questioncard}

\begin{questioncard}{22}{tfcol}
\qtypebadge{tfcol}{True / False} \hfill \textcolor{darkslate!70}{\small Topic: Network Security Standards}

Classified Protection lifecycle is unrelated to includes classification, filing, construction, assessment, and supervision.

\begin{optlist}
\item True
\item False
\end{optlist}
\begin{explainbox}{}
\textbf{Correct Answer: False}\\[0.3em]
\textbf{Concept:} Classified Protection lifecycle is directly related to includes classification, filing, construction, assessment, and supervision.
\end{explainbox}
\end{questioncard}

\begin{questioncard}{23}{tfcol}
\qtypebadge{tfcol}{True / False} \hfill \textcolor{darkslate!70}{\small Topic: Network Security Standards}

Security synchronization means that means planning security with information system construction.

\begin{optlist}
\item True
\item False
\end{optlist}
\begin{explainbox}{}
\textbf{Correct Answer: True}\\[0.3em]
\textbf{Concept:} This statement correctly describes Security synchronization.
\end{explainbox}
\end{questioncard}

\begin{questioncard}{24}{tfcol}
\qtypebadge{tfcol}{True / False} \hfill \textcolor{darkslate!70}{\small Topic: Network Security Standards}

Security synchronization is unrelated to means planning security with information system construction.

\begin{optlist}
\item True
\item False
\end{optlist}
\begin{explainbox}{}
\textbf{Correct Answer: False}\\[0.3em]
\textbf{Concept:} Security synchronization is directly related to means planning security with information system construction.
\end{explainbox}
\end{questioncard}

\begin{questioncard}{25}{tfcol}
\qtypebadge{tfcol}{True / False} \hfill \textcolor{darkslate!70}{\small Topic: Network Security Standards}

ISO 27001:2022 means that organizes controls into four themes.

\begin{optlist}
\item True
\item False
\end{optlist}
\begin{explainbox}{}
\textbf{Correct Answer: True}\\[0.3em]
\textbf{Concept:} This statement correctly describes ISO 27001:2022.
\end{explainbox}
\end{questioncard}

\begin{questioncard}{26}{tfcol}
\qtypebadge{tfcol}{True / False} \hfill \textcolor{darkslate!70}{\small Topic: Network Security Standards}

ISO 27001:2022 is unrelated to organizes controls into four themes.

\begin{optlist}
\item True
\item False
\end{optlist}
\begin{explainbox}{}
\textbf{Correct Answer: False}\\[0.3em]
\textbf{Concept:} ISO 27001:2022 is directly related to organizes controls into four themes.
\end{explainbox}
\end{questioncard}

\begin{questioncard}{27}{tfcol}
\qtypebadge{tfcol}{True / False} \hfill \textcolor{darkslate!70}{\small Topic: Network Security Standards}

Cybersecurity Law means that underpins Classified Protection in China.

\begin{optlist}
\item True
\item False
\end{optlist}
\begin{explainbox}{}
\textbf{Correct Answer: True}\\[0.3em]
\textbf{Concept:} This statement correctly describes Cybersecurity Law.
\end{explainbox}
\end{questioncard}

\begin{questioncard}{28}{tfcol}
\qtypebadge{tfcol}{True / False} \hfill \textcolor{darkslate!70}{\small Topic: Network Security Standards}

Cybersecurity Law is unrelated to underpins Classified Protection in China.

\begin{optlist}
\item True
\item False
\end{optlist}
\begin{explainbox}{}
\textbf{Correct Answer: False}\\[0.3em]
\textbf{Concept:} Cybersecurity Law is directly related to underpins Classified Protection in China.
\end{explainbox}
\end{questioncard}

\section{Network Fundamentals}
\begin{questioncard}{1}{tfcol}
\qtypebadge{tfcol}{True / False} \hfill \textcolor{darkslate!70}{\small Topic: Network Fundamentals}

The OSI model was developed before the TCP/IP model.

\begin{optlist}
\item True
\item False
\end{optlist}
\begin{explainbox}{}
\textbf{Correct Answer: False}\\[0.3em]
\textbf{Concept:} TCP/IP was developed earlier and became the de facto standard; OSI was later as a reference model.
\end{explainbox}
\end{questioncard}

\begin{questioncard}{2}{tfcol}
\qtypebadge{tfcol}{True / False} \hfill \textcolor{darkslate!70}{\small Topic: Network Fundamentals}

DNS always uses TCP port 53.

\begin{optlist}
\item True
\item False
\end{optlist}
\begin{explainbox}{}
\textbf{Correct Answer: False}\\[0.3em]
\textbf{Concept:} DNS primarily uses UDP port 53; TCP port 53 is used for zone transfers and large responses.
\end{explainbox}
\end{questioncard}

\begin{questioncard}{3}{tfcol}
\qtypebadge{tfcol}{True / False} \hfill \textcolor{darkslate!70}{\small Topic: Network Fundamentals}

A switch forwards packets based on IP addresses.

\begin{optlist}
\item True
\item False
\end{optlist}
\begin{explainbox}{}
\textbf{Correct Answer: False}\\[0.3em]
\textbf{Concept:} A switch forwards frames based on MAC addresses; routers forward packets based on IP addresses.
\end{explainbox}
\end{questioncard}

\begin{questioncard}{4}{tfcol}
\qtypebadge{tfcol}{True / False} \hfill \textcolor{darkslate!70}{\small Topic: Network Fundamentals}

IPv4 addresses are 32 bits long.

\begin{optlist}
\item True
\item False
\end{optlist}
\begin{explainbox}{}
\textbf{Correct Answer: True}\\[0.3em]
\textbf{Concept:} IPv4 addresses are 32-bit binary numbers, usually written in dotted-decimal notation.
\end{explainbox}
\end{questioncard}

\begin{questioncard}{5}{tfcol}
\qtypebadge{tfcol}{True / False} \hfill \textcolor{darkslate!70}{\small Topic: Network Fundamentals}

ICMP is a transport-layer protocol.

\begin{optlist}
\item True
\item False
\end{optlist}
\begin{explainbox}{}
\textbf{Correct Answer: False}\\[0.3em]
\textbf{Concept:} ICMP is a network-layer protocol used for diagnostics and error reporting.
\end{explainbox}
\end{questioncard}

\begin{questioncard}{6}{tfcol}
\qtypebadge{tfcol}{True / False} \hfill \textcolor{darkslate!70}{\small Topic: Network Fundamentals}

TCP provides connection-oriented reliable delivery.

\begin{optlist}
\item True
\item False
\end{optlist}
\begin{explainbox}{}
\textbf{Correct Answer: True}\\[0.3em]
\textbf{Concept:} TCP establishes connections and ensures reliable, ordered delivery.
\end{explainbox}
\end{questioncard}

\begin{questioncard}{7}{tfcol}
\qtypebadge{tfcol}{True / False} \hfill \textcolor{darkslate!70}{\small Topic: Network Fundamentals}

A hub creates separate collision domains per port.

\begin{optlist}
\item True
\item False
\end{optlist}
\begin{explainbox}{}
\textbf{Correct Answer: False}\\[0.3em]
\textbf{Concept:} A hub shares a single collision domain; switches create separate collision domains.
\end{explainbox}
\end{questioncard}

\begin{questioncard}{8}{tfcol}
\qtypebadge{tfcol}{True / False} \hfill \textcolor{darkslate!70}{\small Topic: Network Fundamentals}

IPv6 addresses are 128 bits long.

\begin{optlist}
\item True
\item False
\end{optlist}
\begin{explainbox}{}
\textbf{Correct Answer: True}\\[0.3em]
\textbf{Concept:} IPv6 addresses are 128-bit binary numbers.
\end{explainbox}
\end{questioncard}

\begin{questioncard}{9}{tfcol}
\qtypebadge{tfcol}{True / False} \hfill \textcolor{darkslate!70}{\small Topic: Network Fundamentals}

OSI model means that has seven layers.

\begin{optlist}
\item True
\item False
\end{optlist}
\begin{explainbox}{}
\textbf{Correct Answer: True}\\[0.3em]
\textbf{Concept:} This statement correctly describes OSI model.
\end{explainbox}
\end{questioncard}

\begin{questioncard}{10}{tfcol}
\qtypebadge{tfcol}{True / False} \hfill \textcolor{darkslate!70}{\small Topic: Network Fundamentals}

OSI model is unrelated to has seven layers.

\begin{optlist}
\item True
\item False
\end{optlist}
\begin{explainbox}{}
\textbf{Correct Answer: False}\\[0.3em]
\textbf{Concept:} OSI model is directly related to has seven layers.
\end{explainbox}
\end{questioncard}

\begin{questioncard}{11}{tfcol}
\qtypebadge{tfcol}{True / False} \hfill \textcolor{darkslate!70}{\small Topic: Network Fundamentals}

TCP/IP standard model means that has four layers.

\begin{optlist}
\item True
\item False
\end{optlist}
\begin{explainbox}{}
\textbf{Correct Answer: True}\\[0.3em]
\textbf{Concept:} This statement correctly describes TCP/IP standard model.
\end{explainbox}
\end{questioncard}

\begin{questioncard}{12}{tfcol}
\qtypebadge{tfcol}{True / False} \hfill \textcolor{darkslate!70}{\small Topic: Network Fundamentals}

TCP/IP standard model is unrelated to has four layers.

\begin{optlist}
\item True
\item False
\end{optlist}
\begin{explainbox}{}
\textbf{Correct Answer: False}\\[0.3em]
\textbf{Concept:} TCP/IP standard model is directly related to has four layers.
\end{explainbox}
\end{questioncard}

\begin{questioncard}{13}{tfcol}
\qtypebadge{tfcol}{True / False} \hfill \textcolor{darkslate!70}{\small Topic: Network Fundamentals}

TCP/IP equivalent model means that has five layers.

\begin{optlist}
\item True
\item False
\end{optlist}
\begin{explainbox}{}
\textbf{Correct Answer: True}\\[0.3em]
\textbf{Concept:} This statement correctly describes TCP/IP equivalent model.
\end{explainbox}
\end{questioncard}

\begin{questioncard}{14}{tfcol}
\qtypebadge{tfcol}{True / False} \hfill \textcolor{darkslate!70}{\small Topic: Network Fundamentals}

TCP/IP equivalent model is unrelated to has five layers.

\begin{optlist}
\item True
\item False
\end{optlist}
\begin{explainbox}{}
\textbf{Correct Answer: False}\\[0.3em]
\textbf{Concept:} TCP/IP equivalent model is directly related to has five layers.
\end{explainbox}
\end{questioncard}

\begin{questioncard}{15}{tfcol}
\qtypebadge{tfcol}{True / False} \hfill \textcolor{darkslate!70}{\small Topic: Network Fundamentals}

Physical layer means that transmits raw bits over a medium.

\begin{optlist}
\item True
\item False
\end{optlist}
\begin{explainbox}{}
\textbf{Correct Answer: True}\\[0.3em]
\textbf{Concept:} This statement correctly describes Physical layer.
\end{explainbox}
\end{questioncard}

\begin{questioncard}{16}{tfcol}
\qtypebadge{tfcol}{True / False} \hfill \textcolor{darkslate!70}{\small Topic: Network Fundamentals}

Physical layer is unrelated to transmits raw bits over a medium.

\begin{optlist}
\item True
\item False
\end{optlist}
\begin{explainbox}{}
\textbf{Correct Answer: False}\\[0.3em]
\textbf{Concept:} Physical layer is directly related to transmits raw bits over a medium.
\end{explainbox}
\end{questioncard}

\begin{questioncard}{17}{tfcol}
\qtypebadge{tfcol}{True / False} \hfill \textcolor{darkslate!70}{\small Topic: Network Fundamentals}

Data link layer means that handles framing and MAC addressing.

\begin{optlist}
\item True
\item False
\end{optlist}
\begin{explainbox}{}
\textbf{Correct Answer: True}\\[0.3em]
\textbf{Concept:} This statement correctly describes Data link layer.
\end{explainbox}
\end{questioncard}

\begin{questioncard}{18}{tfcol}
\qtypebadge{tfcol}{True / False} \hfill \textcolor{darkslate!70}{\small Topic: Network Fundamentals}

Data link layer is unrelated to handles framing and MAC addressing.

\begin{optlist}
\item True
\item False
\end{optlist}
\begin{explainbox}{}
\textbf{Correct Answer: False}\\[0.3em]
\textbf{Concept:} Data link layer is directly related to handles framing and MAC addressing.
\end{explainbox}
\end{questioncard}

\begin{questioncard}{19}{tfcol}
\qtypebadge{tfcol}{True / False} \hfill \textcolor{darkslate!70}{\small Topic: Network Fundamentals}

Network layer means that handles logical addressing and routing.

\begin{optlist}
\item True
\item False
\end{optlist}
\begin{explainbox}{}
\textbf{Correct Answer: True}\\[0.3em]
\textbf{Concept:} This statement correctly describes Network layer.
\end{explainbox}
\end{questioncard}

\begin{questioncard}{20}{tfcol}
\qtypebadge{tfcol}{True / False} \hfill \textcolor{darkslate!70}{\small Topic: Network Fundamentals}

Network layer is unrelated to handles logical addressing and routing.

\begin{optlist}
\item True
\item False
\end{optlist}
\begin{explainbox}{}
\textbf{Correct Answer: False}\\[0.3em]
\textbf{Concept:} Network layer is directly related to handles logical addressing and routing.
\end{explainbox}
\end{questioncard}

\begin{questioncard}{21}{tfcol}
\qtypebadge{tfcol}{True / False} \hfill \textcolor{darkslate!70}{\small Topic: Network Fundamentals}

Transport layer means that provides end-to-end communication and reliability.

\begin{optlist}
\item True
\item False
\end{optlist}
\begin{explainbox}{}
\textbf{Correct Answer: True}\\[0.3em]
\textbf{Concept:} This statement correctly describes Transport layer.
\end{explainbox}
\end{questioncard}

\begin{questioncard}{22}{tfcol}
\qtypebadge{tfcol}{True / False} \hfill \textcolor{darkslate!70}{\small Topic: Network Fundamentals}

Transport layer is unrelated to provides end-to-end communication and reliability.

\begin{optlist}
\item True
\item False
\end{optlist}
\begin{explainbox}{}
\textbf{Correct Answer: False}\\[0.3em]
\textbf{Concept:} Transport layer is directly related to provides end-to-end communication and reliability.
\end{explainbox}
\end{questioncard}

\begin{questioncard}{23}{tfcol}
\qtypebadge{tfcol}{True / False} \hfill \textcolor{darkslate!70}{\small Topic: Network Fundamentals}

Session layer means that manages application dialogs.

\begin{optlist}
\item True
\item False
\end{optlist}
\begin{explainbox}{}
\textbf{Correct Answer: True}\\[0.3em]
\textbf{Concept:} This statement correctly describes Session layer.
\end{explainbox}
\end{questioncard}

\begin{questioncard}{24}{tfcol}
\qtypebadge{tfcol}{True / False} \hfill \textcolor{darkslate!70}{\small Topic: Network Fundamentals}

Session layer is unrelated to manages application dialogs.

\begin{optlist}
\item True
\item False
\end{optlist}
\begin{explainbox}{}
\textbf{Correct Answer: False}\\[0.3em]
\textbf{Concept:} Session layer is directly related to manages application dialogs.
\end{explainbox}
\end{questioncard}

\begin{questioncard}{25}{tfcol}
\qtypebadge{tfcol}{True / False} \hfill \textcolor{darkslate!70}{\small Topic: Network Fundamentals}

Presentation layer means that handles encryption and data format conversion.

\begin{optlist}
\item True
\item False
\end{optlist}
\begin{explainbox}{}
\textbf{Correct Answer: True}\\[0.3em]
\textbf{Concept:} This statement correctly describes Presentation layer.
\end{explainbox}
\end{questioncard}

\begin{questioncard}{26}{tfcol}
\qtypebadge{tfcol}{True / False} \hfill \textcolor{darkslate!70}{\small Topic: Network Fundamentals}

Presentation layer is unrelated to handles encryption and data format conversion.

\begin{optlist}
\item True
\item False
\end{optlist}
\begin{explainbox}{}
\textbf{Correct Answer: False}\\[0.3em]
\textbf{Concept:} Presentation layer is directly related to handles encryption and data format conversion.
\end{explainbox}
\end{questioncard}

\begin{questioncard}{27}{tfcol}
\qtypebadge{tfcol}{True / False} \hfill \textcolor{darkslate!70}{\small Topic: Network Fundamentals}

Application layer means that provides network services to applications.

\begin{optlist}
\item True
\item False
\end{optlist}
\begin{explainbox}{}
\textbf{Correct Answer: True}\\[0.3em]
\textbf{Concept:} This statement correctly describes Application layer.
\end{explainbox}
\end{questioncard}

\begin{questioncard}{28}{tfcol}
\qtypebadge{tfcol}{True / False} \hfill \textcolor{darkslate!70}{\small Topic: Network Fundamentals}

Application layer is unrelated to provides network services to applications.

\begin{optlist}
\item True
\item False
\end{optlist}
\begin{explainbox}{}
\textbf{Correct Answer: False}\\[0.3em]
\textbf{Concept:} Application layer is directly related to provides network services to applications.
\end{explainbox}
\end{questioncard}

\begin{questioncard}{29}{tfcol}
\qtypebadge{tfcol}{True / False} \hfill \textcolor{darkslate!70}{\small Topic: Network Fundamentals}

TCP means that is connection-oriented and reliable.

\begin{optlist}
\item True
\item False
\end{optlist}
\begin{explainbox}{}
\textbf{Correct Answer: True}\\[0.3em]
\textbf{Concept:} This statement correctly describes TCP.
\end{explainbox}
\end{questioncard}

\begin{questioncard}{30}{tfcol}
\qtypebadge{tfcol}{True / False} \hfill \textcolor{darkslate!70}{\small Topic: Network Fundamentals}

TCP is unrelated to is connection-oriented and reliable.

\begin{optlist}
\item True
\item False
\end{optlist}
\begin{explainbox}{}
\textbf{Correct Answer: False}\\[0.3em]
\textbf{Concept:} TCP is directly related to is connection-oriented and reliable.
\end{explainbox}
\end{questioncard}

\begin{questioncard}{31}{tfcol}
\qtypebadge{tfcol}{True / False} \hfill \textcolor{darkslate!70}{\small Topic: Network Fundamentals}

UDP means that is connectionless and has low overhead.

\begin{optlist}
\item True
\item False
\end{optlist}
\begin{explainbox}{}
\textbf{Correct Answer: True}\\[0.3em]
\textbf{Concept:} This statement correctly describes UDP.
\end{explainbox}
\end{questioncard}

\begin{questioncard}{32}{tfcol}
\qtypebadge{tfcol}{True / False} \hfill \textcolor{darkslate!70}{\small Topic: Network Fundamentals}

UDP is unrelated to is connectionless and has low overhead.

\begin{optlist}
\item True
\item False
\end{optlist}
\begin{explainbox}{}
\textbf{Correct Answer: False}\\[0.3em]
\textbf{Concept:} UDP is directly related to is connectionless and has low overhead.
\end{explainbox}
\end{questioncard}

\begin{questioncard}{33}{tfcol}
\qtypebadge{tfcol}{True / False} \hfill \textcolor{darkslate!70}{\small Topic: Network Fundamentals}

FTP means that uses TCP ports 20 and 21.

\begin{optlist}
\item True
\item False
\end{optlist}
\begin{explainbox}{}
\textbf{Correct Answer: True}\\[0.3em]
\textbf{Concept:} This statement correctly describes FTP.
\end{explainbox}
\end{questioncard}

\begin{questioncard}{34}{tfcol}
\qtypebadge{tfcol}{True / False} \hfill \textcolor{darkslate!70}{\small Topic: Network Fundamentals}

FTP is unrelated to uses TCP ports 20 and 21.

\begin{optlist}
\item True
\item False
\end{optlist}
\begin{explainbox}{}
\textbf{Correct Answer: False}\\[0.3em]
\textbf{Concept:} FTP is directly related to uses TCP ports 20 and 21.
\end{explainbox}
\end{questioncard}

\begin{questioncard}{35}{tfcol}
\qtypebadge{tfcol}{True / False} \hfill \textcolor{darkslate!70}{\small Topic: Network Fundamentals}

HTTP means that uses TCP port 80.

\begin{optlist}
\item True
\item False
\end{optlist}
\begin{explainbox}{}
\textbf{Correct Answer: True}\\[0.3em]
\textbf{Concept:} This statement correctly describes HTTP.
\end{explainbox}
\end{questioncard}

\begin{questioncard}{36}{tfcol}
\qtypebadge{tfcol}{True / False} \hfill \textcolor{darkslate!70}{\small Topic: Network Fundamentals}

HTTP is unrelated to uses TCP port 80.

\begin{optlist}
\item True
\item False
\end{optlist}
\begin{explainbox}{}
\textbf{Correct Answer: False}\\[0.3em]
\textbf{Concept:} HTTP is directly related to uses TCP port 80.
\end{explainbox}
\end{questioncard}

\begin{questioncard}{37}{tfcol}
\qtypebadge{tfcol}{True / False} \hfill \textcolor{darkslate!70}{\small Topic: Network Fundamentals}

HTTPS means that uses TCP port 443.

\begin{optlist}
\item True
\item False
\end{optlist}
\begin{explainbox}{}
\textbf{Correct Answer: True}\\[0.3em]
\textbf{Concept:} This statement correctly describes HTTPS.
\end{explainbox}
\end{questioncard}

\begin{questioncard}{38}{tfcol}
\qtypebadge{tfcol}{True / False} \hfill \textcolor{darkslate!70}{\small Topic: Network Fundamentals}

HTTPS is unrelated to uses TCP port 443.

\begin{optlist}
\item True
\item False
\end{optlist}
\begin{explainbox}{}
\textbf{Correct Answer: False}\\[0.3em]
\textbf{Concept:} HTTPS is directly related to uses TCP port 443.
\end{explainbox}
\end{questioncard}

\begin{questioncard}{39}{tfcol}
\qtypebadge{tfcol}{True / False} \hfill \textcolor{darkslate!70}{\small Topic: Network Fundamentals}

SSH means that uses TCP port 22.

\begin{optlist}
\item True
\item False
\end{optlist}
\begin{explainbox}{}
\textbf{Correct Answer: True}\\[0.3em]
\textbf{Concept:} This statement correctly describes SSH.
\end{explainbox}
\end{questioncard}

\begin{questioncard}{40}{tfcol}
\qtypebadge{tfcol}{True / False} \hfill \textcolor{darkslate!70}{\small Topic: Network Fundamentals}

SSH is unrelated to uses TCP port 22.

\begin{optlist}
\item True
\item False
\end{optlist}
\begin{explainbox}{}
\textbf{Correct Answer: False}\\[0.3em]
\textbf{Concept:} SSH is directly related to uses TCP port 22.
\end{explainbox}
\end{questioncard}

\begin{questioncard}{41}{tfcol}
\qtypebadge{tfcol}{True / False} \hfill \textcolor{darkslate!70}{\small Topic: Network Fundamentals}

Telnet means that uses TCP port 23 and is insecure.

\begin{optlist}
\item True
\item False
\end{optlist}
\begin{explainbox}{}
\textbf{Correct Answer: True}\\[0.3em]
\textbf{Concept:} This statement correctly describes Telnet.
\end{explainbox}
\end{questioncard}

\begin{questioncard}{42}{tfcol}
\qtypebadge{tfcol}{True / False} \hfill \textcolor{darkslate!70}{\small Topic: Network Fundamentals}

Telnet is unrelated to uses TCP port 23 and is insecure.

\begin{optlist}
\item True
\item False
\end{optlist}
\begin{explainbox}{}
\textbf{Correct Answer: False}\\[0.3em]
\textbf{Concept:} Telnet is directly related to uses TCP port 23 and is insecure.
\end{explainbox}
\end{questioncard}

\begin{questioncard}{43}{tfcol}
\qtypebadge{tfcol}{True / False} \hfill \textcolor{darkslate!70}{\small Topic: Network Fundamentals}

DNS means that primarily uses UDP port 53.

\begin{optlist}
\item True
\item False
\end{optlist}
\begin{explainbox}{}
\textbf{Correct Answer: True}\\[0.3em]
\textbf{Concept:} This statement correctly describes DNS.
\end{explainbox}
\end{questioncard}

\begin{questioncard}{44}{tfcol}
\qtypebadge{tfcol}{True / False} \hfill \textcolor{darkslate!70}{\small Topic: Network Fundamentals}

DNS is unrelated to primarily uses UDP port 53.

\begin{optlist}
\item True
\item False
\end{optlist}
\begin{explainbox}{}
\textbf{Correct Answer: False}\\[0.3em]
\textbf{Concept:} DNS is directly related to primarily uses UDP port 53.
\end{explainbox}
\end{questioncard}

\begin{questioncard}{45}{tfcol}
\qtypebadge{tfcol}{True / False} \hfill \textcolor{darkslate!70}{\small Topic: Network Fundamentals}

SMTP means that uses TCP port 25.

\begin{optlist}
\item True
\item False
\end{optlist}
\begin{explainbox}{}
\textbf{Correct Answer: True}\\[0.3em]
\textbf{Concept:} This statement correctly describes SMTP.
\end{explainbox}
\end{questioncard}

\begin{questioncard}{46}{tfcol}
\qtypebadge{tfcol}{True / False} \hfill \textcolor{darkslate!70}{\small Topic: Network Fundamentals}

SMTP is unrelated to uses TCP port 25.

\begin{optlist}
\item True
\item False
\end{optlist}
\begin{explainbox}{}
\textbf{Correct Answer: False}\\[0.3em]
\textbf{Concept:} SMTP is directly related to uses TCP port 25.
\end{explainbox}
\end{questioncard}

\begin{questioncard}{47}{tfcol}
\qtypebadge{tfcol}{True / False} \hfill \textcolor{darkslate!70}{\small Topic: Network Fundamentals}

DHCP means that dynamically assigns IP addresses.

\begin{optlist}
\item True
\item False
\end{optlist}
\begin{explainbox}{}
\textbf{Correct Answer: True}\\[0.3em]
\textbf{Concept:} This statement correctly describes DHCP.
\end{explainbox}
\end{questioncard}

\begin{questioncard}{48}{tfcol}
\qtypebadge{tfcol}{True / False} \hfill \textcolor{darkslate!70}{\small Topic: Network Fundamentals}

DHCP is unrelated to dynamically assigns IP addresses.

\begin{optlist}
\item True
\item False
\end{optlist}
\begin{explainbox}{}
\textbf{Correct Answer: False}\\[0.3em]
\textbf{Concept:} DHCP is directly related to dynamically assigns IP addresses.
\end{explainbox}
\end{questioncard}

\begin{questioncard}{49}{tfcol}
\qtypebadge{tfcol}{True / False} \hfill \textcolor{darkslate!70}{\small Topic: Network Fundamentals}

ARP means that maps IP addresses to MAC addresses.

\begin{optlist}
\item True
\item False
\end{optlist}
\begin{explainbox}{}
\textbf{Correct Answer: True}\\[0.3em]
\textbf{Concept:} This statement correctly describes ARP.
\end{explainbox}
\end{questioncard}

\begin{questioncard}{50}{tfcol}
\qtypebadge{tfcol}{True / False} \hfill \textcolor{darkslate!70}{\small Topic: Network Fundamentals}

ARP is unrelated to maps IP addresses to MAC addresses.

\begin{optlist}
\item True
\item False
\end{optlist}
\begin{explainbox}{}
\textbf{Correct Answer: False}\\[0.3em]
\textbf{Concept:} ARP is directly related to maps IP addresses to MAC addresses.
\end{explainbox}
\end{questioncard}

\begin{questioncard}{51}{tfcol}
\qtypebadge{tfcol}{True / False} \hfill \textcolor{darkslate!70}{\small Topic: Network Fundamentals}

ICMP means that provides error reporting and diagnostics.

\begin{optlist}
\item True
\item False
\end{optlist}
\begin{explainbox}{}
\textbf{Correct Answer: True}\\[0.3em]
\textbf{Concept:} This statement correctly describes ICMP.
\end{explainbox}
\end{questioncard}

\begin{questioncard}{52}{tfcol}
\qtypebadge{tfcol}{True / False} \hfill \textcolor{darkslate!70}{\small Topic: Network Fundamentals}

ICMP is unrelated to provides error reporting and diagnostics.

\begin{optlist}
\item True
\item False
\end{optlist}
\begin{explainbox}{}
\textbf{Correct Answer: False}\\[0.3em]
\textbf{Concept:} ICMP is directly related to provides error reporting and diagnostics.
\end{explainbox}
\end{questioncard}

\begin{questioncard}{53}{tfcol}
\qtypebadge{tfcol}{True / False} \hfill \textcolor{darkslate!70}{\small Topic: Network Fundamentals}

IPv4 address means that is 32 bits long.

\begin{optlist}
\item True
\item False
\end{optlist}
\begin{explainbox}{}
\textbf{Correct Answer: True}\\[0.3em]
\textbf{Concept:} This statement correctly describes IPv4 address.
\end{explainbox}
\end{questioncard}

\begin{questioncard}{54}{tfcol}
\qtypebadge{tfcol}{True / False} \hfill \textcolor{darkslate!70}{\small Topic: Network Fundamentals}

IPv4 address is unrelated to is 32 bits long.

\begin{optlist}
\item True
\item False
\end{optlist}
\begin{explainbox}{}
\textbf{Correct Answer: False}\\[0.3em]
\textbf{Concept:} IPv4 address is directly related to is 32 bits long.
\end{explainbox}
\end{questioncard}

\begin{questioncard}{55}{tfcol}
\qtypebadge{tfcol}{True / False} \hfill \textcolor{darkslate!70}{\small Topic: Network Fundamentals}

IPv6 address means that is 128 bits long.

\begin{optlist}
\item True
\item False
\end{optlist}
\begin{explainbox}{}
\textbf{Correct Answer: True}\\[0.3em]
\textbf{Concept:} This statement correctly describes IPv6 address.
\end{explainbox}
\end{questioncard}

\begin{questioncard}{56}{tfcol}
\qtypebadge{tfcol}{True / False} \hfill \textcolor{darkslate!70}{\small Topic: Network Fundamentals}

IPv6 address is unrelated to is 128 bits long.

\begin{optlist}
\item True
\item False
\end{optlist}
\begin{explainbox}{}
\textbf{Correct Answer: False}\\[0.3em]
\textbf{Concept:} IPv6 address is directly related to is 128 bits long.
\end{explainbox}
\end{questioncard}

\begin{questioncard}{57}{tfcol}
\qtypebadge{tfcol}{True / False} \hfill \textcolor{darkslate!70}{\small Topic: Network Fundamentals}

MAC address means that is a Layer 2 hardware address.

\begin{optlist}
\item True
\item False
\end{optlist}
\begin{explainbox}{}
\textbf{Correct Answer: True}\\[0.3em]
\textbf{Concept:} This statement correctly describes MAC address.
\end{explainbox}
\end{questioncard}

\begin{questioncard}{58}{tfcol}
\qtypebadge{tfcol}{True / False} \hfill \textcolor{darkslate!70}{\small Topic: Network Fundamentals}

MAC address is unrelated to is a Layer 2 hardware address.

\begin{optlist}
\item True
\item False
\end{optlist}
\begin{explainbox}{}
\textbf{Correct Answer: False}\\[0.3em]
\textbf{Concept:} MAC address is directly related to is a Layer 2 hardware address.
\end{explainbox}
\end{questioncard}

\begin{questioncard}{59}{tfcol}
\qtypebadge{tfcol}{True / False} \hfill \textcolor{darkslate!70}{\small Topic: Network Fundamentals}

Switch means that operates at Layer 2 and forwards frames by MAC.

\begin{optlist}
\item True
\item False
\end{optlist}
\begin{explainbox}{}
\textbf{Correct Answer: True}\\[0.3em]
\textbf{Concept:} This statement correctly describes Switch.
\end{explainbox}
\end{questioncard}

\begin{questioncard}{60}{tfcol}
\qtypebadge{tfcol}{True / False} \hfill \textcolor{darkslate!70}{\small Topic: Network Fundamentals}

Switch is unrelated to operates at Layer 2 and forwards frames by MAC.

\begin{optlist}
\item True
\item False
\end{optlist}
\begin{explainbox}{}
\textbf{Correct Answer: False}\\[0.3em]
\textbf{Concept:} Switch is directly related to operates at Layer 2 and forwards frames by MAC.
\end{explainbox}
\end{questioncard}

\begin{questioncard}{61}{tfcol}
\qtypebadge{tfcol}{True / False} \hfill \textcolor{darkslate!70}{\small Topic: Network Fundamentals}

Router means that operates at Layer 3 and forwards packets by IP.

\begin{optlist}
\item True
\item False
\end{optlist}
\begin{explainbox}{}
\textbf{Correct Answer: True}\\[0.3em]
\textbf{Concept:} This statement correctly describes Router.
\end{explainbox}
\end{questioncard}

\begin{questioncard}{62}{tfcol}
\qtypebadge{tfcol}{True / False} \hfill \textcolor{darkslate!70}{\small Topic: Network Fundamentals}

Router is unrelated to operates at Layer 3 and forwards packets by IP.

\begin{optlist}
\item True
\item False
\end{optlist}
\begin{explainbox}{}
\textbf{Correct Answer: False}\\[0.3em]
\textbf{Concept:} Router is directly related to operates at Layer 3 and forwards packets by IP.
\end{explainbox}
\end{questioncard}

\begin{questioncard}{63}{tfcol}
\qtypebadge{tfcol}{True / False} \hfill \textcolor{darkslate!70}{\small Topic: Network Fundamentals}

Hub means that operates at Layer 1 and repeats signals.

\begin{optlist}
\item True
\item False
\end{optlist}
\begin{explainbox}{}
\textbf{Correct Answer: True}\\[0.3em]
\textbf{Concept:} This statement correctly describes Hub.
\end{explainbox}
\end{questioncard}

\begin{questioncard}{64}{tfcol}
\qtypebadge{tfcol}{True / False} \hfill \textcolor{darkslate!70}{\small Topic: Network Fundamentals}

Hub is unrelated to operates at Layer 1 and repeats signals.

\begin{optlist}
\item True
\item False
\end{optlist}
\begin{explainbox}{}
\textbf{Correct Answer: False}\\[0.3em]
\textbf{Concept:} Hub is directly related to operates at Layer 1 and repeats signals.
\end{explainbox}
\end{questioncard}

\begin{questioncard}{65}{tfcol}
\qtypebadge{tfcol}{True / False} \hfill \textcolor{darkslate!70}{\small Topic: Network Fundamentals}

VLAN means that logically segments a network into broadcast domains.

\begin{optlist}
\item True
\item False
\end{optlist}
\begin{explainbox}{}
\textbf{Correct Answer: True}\\[0.3em]
\textbf{Concept:} This statement correctly describes VLAN.
\end{explainbox}
\end{questioncard}

\begin{questioncard}{66}{tfcol}
\qtypebadge{tfcol}{True / False} \hfill \textcolor{darkslate!70}{\small Topic: Network Fundamentals}

VLAN is unrelated to logically segments a network into broadcast domains.

\begin{optlist}
\item True
\item False
\end{optlist}
\begin{explainbox}{}
\textbf{Correct Answer: False}\\[0.3em]
\textbf{Concept:} VLAN is directly related to logically segments a network into broadcast domains.
\end{explainbox}
\end{questioncard}

\begin{questioncard}{67}{tfcol}
\qtypebadge{tfcol}{True / False} \hfill \textcolor{darkslate!70}{\small Topic: Network Fundamentals}

STP means that prevents Layer 2 loops.

\begin{optlist}
\item True
\item False
\end{optlist}
\begin{explainbox}{}
\textbf{Correct Answer: True}\\[0.3em]
\textbf{Concept:} This statement correctly describes STP.
\end{explainbox}
\end{questioncard}

\begin{questioncard}{68}{tfcol}
\qtypebadge{tfcol}{True / False} \hfill \textcolor{darkslate!70}{\small Topic: Network Fundamentals}

STP is unrelated to prevents Layer 2 loops.

\begin{optlist}
\item True
\item False
\end{optlist}
\begin{explainbox}{}
\textbf{Correct Answer: False}\\[0.3em]
\textbf{Concept:} STP is directly related to prevents Layer 2 loops.
\end{explainbox}
\end{questioncard}

\begin{questioncard}{69}{tfcol}
\qtypebadge{tfcol}{True / False} \hfill \textcolor{darkslate!70}{\small Topic: Network Fundamentals}

RFC 1918 means that defines private IPv4 address ranges.

\begin{optlist}
\item True
\item False
\end{optlist}
\begin{explainbox}{}
\textbf{Correct Answer: True}\\[0.3em]
\textbf{Concept:} This statement correctly describes RFC 1918.
\end{explainbox}
\end{questioncard}

\begin{questioncard}{70}{tfcol}
\qtypebadge{tfcol}{True / False} \hfill \textcolor{darkslate!70}{\small Topic: Network Fundamentals}

RFC 1918 is unrelated to defines private IPv4 address ranges.

\begin{optlist}
\item True
\item False
\end{optlist}
\begin{explainbox}{}
\textbf{Correct Answer: False}\\[0.3em]
\textbf{Concept:} RFC 1918 is directly related to defines private IPv4 address ranges.
\end{explainbox}
\end{questioncard}

\section{OSI and TCP/IP Models}
\begin{questioncard}{1}{tfcol}
\qtypebadge{tfcol}{True / False} \hfill \textcolor{darkslate!70}{\small Topic: OSI and TCP/IP Models}

The TCP/IP model combines the OSI session and presentation layers into the application layer.

\begin{optlist}
\item True
\item False
\end{optlist}
\begin{explainbox}{}
\textbf{Correct Answer: True}\\[0.3em]
\textbf{Concept:} The TCP/IP model does not have separate session and presentation layers; their functions are included in the application layer.
\end{explainbox}
\end{questioncard}

\section{Network Protocols}
\begin{questioncard}{1}{tfcol}
\qtypebadge{tfcol}{True / False} \hfill \textcolor{darkslate!70}{\small Topic: Network Protocols}

UDP guarantees ordered delivery of packets.

\begin{optlist}
\item True
\item False
\end{optlist}
\begin{explainbox}{}
\textbf{Correct Answer: False}\\[0.3em]
\textbf{Concept:} UDP is connectionless and does not guarantee delivery, ordering, or duplicate protection.
\end{explainbox}
\end{questioncard}

\section{Network Devices}
\begin{questioncard}{1}{tfcol}
\qtypebadge{tfcol}{True / False} \hfill \textcolor{darkslate!70}{\small Topic: Network Devices}

A router forwards frames based on MAC addresses.

\begin{optlist}
\item True
\item False
\end{optlist}
\begin{explainbox}{}
\textbf{Correct Answer: False}\\[0.3em]
\textbf{Concept:} Routers forward packets based on IP addresses; switches forward frames based on MAC addresses.
\end{explainbox}
\end{questioncard}

\begin{questioncard}{2}{tfcol}
\qtypebadge{tfcol}{True / False} \hfill \textcolor{darkslate!70}{\small Topic: Network Devices}

A switch creates a separate collision domain for each port.

\begin{optlist}
\item True
\item False
\end{optlist}
\begin{explainbox}{}
\textbf{Correct Answer: True}\\[0.3em]
\textbf{Concept:} Each switch port is its own collision domain, reducing collisions compared to a hub.
\end{explainbox}
\end{questioncard}

\section{Enterprise Security Threats}
\begin{questioncard}{1}{tfcol}
\qtypebadge{tfcol}{True / False} \hfill \textcolor{darkslate!70}{\small Topic: Enterprise Security Threats}

A worm requires a host program to spread.

\begin{optlist}
\item True
\item False
\end{optlist}
\begin{explainbox}{}
\textbf{Correct Answer: False}\\[0.3em]
\textbf{Concept:} Worms are standalone malware that can self-replicate; viruses require host programs.
\end{explainbox}
\end{questioncard}

\begin{questioncard}{2}{tfcol}
\qtypebadge{tfcol}{True / False} \hfill \textcolor{darkslate!70}{\small Topic: Enterprise Security Threats}

Phishing attacks target human psychology rather than software vulnerabilities.

\begin{optlist}
\item True
\item False
\end{optlist}
\begin{explainbox}{}
\textbf{Correct Answer: True}\\[0.3em]
\textbf{Concept:} Phishing uses deception and social engineering to trick users.
\end{explainbox}
\end{questioncard}

\begin{questioncard}{3}{tfcol}
\qtypebadge{tfcol}{True / False} \hfill \textcolor{darkslate!70}{\small Topic: Enterprise Security Threats}

Using HTTP instead of HTTPS improves confidentiality.

\begin{optlist}
\item True
\item False
\end{optlist}
\begin{explainbox}{}
\textbf{Correct Answer: False}\\[0.3em]
\textbf{Concept:} HTTP sends data in plaintext; HTTPS encrypts data to protect confidentiality.
\end{explainbox}
\end{questioncard}

\begin{questioncard}{4}{tfcol}
\qtypebadge{tfcol}{True / False} \hfill \textcolor{darkslate!70}{\small Topic: Enterprise Security Threats}

The DMZ zone is typically more trusted than the Untrust zone but less trusted than the Trust zone.

\begin{optlist}
\item True
\item False
\end{optlist}
\begin{explainbox}{}
\textbf{Correct Answer: True}\\[0.3em]
\textbf{Concept:} DMZ has an intermediate trust level for public-facing servers.
\end{explainbox}
\end{questioncard}

\begin{questioncard}{5}{tfcol}
\qtypebadge{tfcol}{True / False} \hfill \textcolor{darkslate!70}{\small Topic: Enterprise Security Threats}

Patch management helps eliminate known vulnerabilities on hosts.

\begin{optlist}
\item True
\item False
\end{optlist}
\begin{explainbox}{}
\textbf{Correct Answer: True}\\[0.3em]
\textbf{Concept:} Applying patches closes security holes that attackers could exploit.
\end{explainbox}
\end{questioncard}

\begin{questioncard}{6}{tfcol}
\qtypebadge{tfcol}{True / False} \hfill \textcolor{darkslate!70}{\small Topic: Enterprise Security Threats}

All security threats originate from external attackers.

\begin{optlist}
\item True
\item False
\end{optlist}
\begin{explainbox}{}
\textbf{Correct Answer: False}\\[0.3em]
\textbf{Concept:} Threats can be external or internal, accidental or malicious.
\end{explainbox}
\end{questioncard}

\begin{questioncard}{7}{tfcol}
\qtypebadge{tfcol}{True / False} \hfill \textcolor{darkslate!70}{\small Topic: Enterprise Security Threats}

Phishing is a type of social engineering.

\begin{optlist}
\item True
\item False
\end{optlist}
\begin{explainbox}{}
\textbf{Correct Answer: True}\\[0.3em]
\textbf{Concept:} Phishing manipulates users into revealing information.
\end{explainbox}
\end{questioncard}

\begin{questioncard}{8}{tfcol}
\qtypebadge{tfcol}{True / False} \hfill \textcolor{darkslate!70}{\small Topic: Enterprise Security Threats}

A firewall at the network border can enforce zone-based policies.

\begin{optlist}
\item True
\item False
\end{optlist}
\begin{explainbox}{}
\textbf{Correct Answer: True}\\[0.3em]
\textbf{Concept:} Firewalls enforce policies between security zones.
\end{explainbox}
\end{questioncard}

\begin{questioncard}{9}{tfcol}
\qtypebadge{tfcol}{True / False} \hfill \textcolor{darkslate!70}{\small Topic: Enterprise Security Threats}

Least privilege reduces the impact of compromised accounts.

\begin{optlist}
\item True
\item False
\end{optlist}
\begin{explainbox}{}
\textbf{Correct Answer: True}\\[0.3em]
\textbf{Concept:} Limiting access limits damage from compromised credentials.
\end{explainbox}
\end{questioncard}

\begin{questioncard}{10}{tfcol}
\qtypebadge{tfcol}{True / False} \hfill \textcolor{darkslate!70}{\small Topic: Enterprise Security Threats}

Passive attack means that monitors traffic without altering it.

\begin{optlist}
\item True
\item False
\end{optlist}
\begin{explainbox}{}
\textbf{Correct Answer: True}\\[0.3em]
\textbf{Concept:} This statement correctly describes Passive attack.
\end{explainbox}
\end{questioncard}

\begin{questioncard}{11}{tfcol}
\qtypebadge{tfcol}{True / False} \hfill \textcolor{darkslate!70}{\small Topic: Enterprise Security Threats}

Passive attack is unrelated to monitors traffic without altering it.

\begin{optlist}
\item True
\item False
\end{optlist}
\begin{explainbox}{}
\textbf{Correct Answer: False}\\[0.3em]
\textbf{Concept:} Passive attack is directly related to monitors traffic without altering it.
\end{explainbox}
\end{questioncard}

\begin{questioncard}{12}{tfcol}
\qtypebadge{tfcol}{True / False} \hfill \textcolor{darkslate!70}{\small Topic: Enterprise Security Threats}

Active attack means that modifies or disrupts systems or data.

\begin{optlist}
\item True
\item False
\end{optlist}
\begin{explainbox}{}
\textbf{Correct Answer: True}\\[0.3em]
\textbf{Concept:} This statement correctly describes Active attack.
\end{explainbox}
\end{questioncard}

\begin{questioncard}{13}{tfcol}
\qtypebadge{tfcol}{True / False} \hfill \textcolor{darkslate!70}{\small Topic: Enterprise Security Threats}

Active attack is unrelated to modifies or disrupts systems or data.

\begin{optlist}
\item True
\item False
\end{optlist}
\begin{explainbox}{}
\textbf{Correct Answer: False}\\[0.3em]
\textbf{Concept:} Active attack is directly related to modifies or disrupts systems or data.
\end{explainbox}
\end{questioncard}

\begin{questioncard}{14}{tfcol}
\qtypebadge{tfcol}{True / False} \hfill \textcolor{darkslate!70}{\small Topic: Enterprise Security Threats}

Malware means that includes viruses, worms, Trojans, and ransomware.

\begin{optlist}
\item True
\item False
\end{optlist}
\begin{explainbox}{}
\textbf{Correct Answer: True}\\[0.3em]
\textbf{Concept:} This statement correctly describes Malware.
\end{explainbox}
\end{questioncard}

\begin{questioncard}{15}{tfcol}
\qtypebadge{tfcol}{True / False} \hfill \textcolor{darkslate!70}{\small Topic: Enterprise Security Threats}

Malware is unrelated to includes viruses, worms, Trojans, and ransomware.

\begin{optlist}
\item True
\item False
\end{optlist}
\begin{explainbox}{}
\textbf{Correct Answer: False}\\[0.3em]
\textbf{Concept:} Malware is directly related to includes viruses, worms, Trojans, and ransomware.
\end{explainbox}
\end{questioncard}

\begin{questioncard}{16}{tfcol}
\qtypebadge{tfcol}{True / False} \hfill \textcolor{darkslate!70}{\small Topic: Enterprise Security Threats}

Virus means that requires a host program to spread.

\begin{optlist}
\item True
\item False
\end{optlist}
\begin{explainbox}{}
\textbf{Correct Answer: True}\\[0.3em]
\textbf{Concept:} This statement correctly describes Virus.
\end{explainbox}
\end{questioncard}

\begin{questioncard}{17}{tfcol}
\qtypebadge{tfcol}{True / False} \hfill \textcolor{darkslate!70}{\small Topic: Enterprise Security Threats}

Virus is unrelated to requires a host program to spread.

\begin{optlist}
\item True
\item False
\end{optlist}
\begin{explainbox}{}
\textbf{Correct Answer: False}\\[0.3em]
\textbf{Concept:} Virus is directly related to requires a host program to spread.
\end{explainbox}
\end{questioncard}

\begin{questioncard}{18}{tfcol}
\qtypebadge{tfcol}{True / False} \hfill \textcolor{darkslate!70}{\small Topic: Enterprise Security Threats}

Worm means that can self-replicate across networks.

\begin{optlist}
\item True
\item False
\end{optlist}
\begin{explainbox}{}
\textbf{Correct Answer: True}\\[0.3em]
\textbf{Concept:} This statement correctly describes Worm.
\end{explainbox}
\end{questioncard}

\begin{questioncard}{19}{tfcol}
\qtypebadge{tfcol}{True / False} \hfill \textcolor{darkslate!70}{\small Topic: Enterprise Security Threats}

Worm is unrelated to can self-replicate across networks.

\begin{optlist}
\item True
\item False
\end{optlist}
\begin{explainbox}{}
\textbf{Correct Answer: False}\\[0.3em]
\textbf{Concept:} Worm is directly related to can self-replicate across networks.
\end{explainbox}
\end{questioncard}

\begin{questioncard}{20}{tfcol}
\qtypebadge{tfcol}{True / False} \hfill \textcolor{darkslate!70}{\small Topic: Enterprise Security Threats}

Trojan means that disguises itself as legitimate software.

\begin{optlist}
\item True
\item False
\end{optlist}
\begin{explainbox}{}
\textbf{Correct Answer: True}\\[0.3em]
\textbf{Concept:} This statement correctly describes Trojan.
\end{explainbox}
\end{questioncard}

\begin{questioncard}{21}{tfcol}
\qtypebadge{tfcol}{True / False} \hfill \textcolor{darkslate!70}{\small Topic: Enterprise Security Threats}

Trojan is unrelated to disguises itself as legitimate software.

\begin{optlist}
\item True
\item False
\end{optlist}
\begin{explainbox}{}
\textbf{Correct Answer: False}\\[0.3em]
\textbf{Concept:} Trojan is directly related to disguises itself as legitimate software.
\end{explainbox}
\end{questioncard}

\begin{questioncard}{22}{tfcol}
\qtypebadge{tfcol}{True / False} \hfill \textcolor{darkslate!70}{\small Topic: Enterprise Security Threats}

Ransomware means that encrypts data and demands payment.

\begin{optlist}
\item True
\item False
\end{optlist}
\begin{explainbox}{}
\textbf{Correct Answer: True}\\[0.3em]
\textbf{Concept:} This statement correctly describes Ransomware.
\end{explainbox}
\end{questioncard}

\begin{questioncard}{23}{tfcol}
\qtypebadge{tfcol}{True / False} \hfill \textcolor{darkslate!70}{\small Topic: Enterprise Security Threats}

Ransomware is unrelated to encrypts data and demands payment.

\begin{optlist}
\item True
\item False
\end{optlist}
\begin{explainbox}{}
\textbf{Correct Answer: False}\\[0.3em]
\textbf{Concept:} Ransomware is directly related to encrypts data and demands payment.
\end{explainbox}
\end{questioncard}

\begin{questioncard}{24}{tfcol}
\qtypebadge{tfcol}{True / False} \hfill \textcolor{darkslate!70}{\small Topic: Enterprise Security Threats}

Phishing means that uses deception to steal information.

\begin{optlist}
\item True
\item False
\end{optlist}
\begin{explainbox}{}
\textbf{Correct Answer: True}\\[0.3em]
\textbf{Concept:} This statement correctly describes Phishing.
\end{explainbox}
\end{questioncard}

\begin{questioncard}{25}{tfcol}
\qtypebadge{tfcol}{True / False} \hfill \textcolor{darkslate!70}{\small Topic: Enterprise Security Threats}

Phishing is unrelated to uses deception to steal information.

\begin{optlist}
\item True
\item False
\end{optlist}
\begin{explainbox}{}
\textbf{Correct Answer: False}\\[0.3em]
\textbf{Concept:} Phishing is directly related to uses deception to steal information.
\end{explainbox}
\end{questioncard}

\begin{questioncard}{26}{tfcol}
\qtypebadge{tfcol}{True / False} \hfill \textcolor{darkslate!70}{\small Topic: Enterprise Security Threats}

Social engineering means that manipulates people into breaking security rules.

\begin{optlist}
\item True
\item False
\end{optlist}
\begin{explainbox}{}
\textbf{Correct Answer: True}\\[0.3em]
\textbf{Concept:} This statement correctly describes Social engineering.
\end{explainbox}
\end{questioncard}

\begin{questioncard}{27}{tfcol}
\qtypebadge{tfcol}{True / False} \hfill \textcolor{darkslate!70}{\small Topic: Enterprise Security Threats}

Social engineering is unrelated to manipulates people into breaking security rules.

\begin{optlist}
\item True
\item False
\end{optlist}
\begin{explainbox}{}
\textbf{Correct Answer: False}\\[0.3em]
\textbf{Concept:} Social engineering is directly related to manipulates people into breaking security rules.
\end{explainbox}
\end{questioncard}

\begin{questioncard}{28}{tfcol}
\qtypebadge{tfcol}{True / False} \hfill \textcolor{darkslate!70}{\small Topic: Enterprise Security Threats}

DoS attack means that aims to make a service unavailable.

\begin{optlist}
\item True
\item False
\end{optlist}
\begin{explainbox}{}
\textbf{Correct Answer: True}\\[0.3em]
\textbf{Concept:} This statement correctly describes DoS attack.
\end{explainbox}
\end{questioncard}

\begin{questioncard}{29}{tfcol}
\qtypebadge{tfcol}{True / False} \hfill \textcolor{darkslate!70}{\small Topic: Enterprise Security Threats}

DoS attack is unrelated to aims to make a service unavailable.

\begin{optlist}
\item True
\item False
\end{optlist}
\begin{explainbox}{}
\textbf{Correct Answer: False}\\[0.3em]
\textbf{Concept:} DoS attack is directly related to aims to make a service unavailable.
\end{explainbox}
\end{questioncard}

\begin{questioncard}{30}{tfcol}
\qtypebadge{tfcol}{True / False} \hfill \textcolor{darkslate!70}{\small Topic: Enterprise Security Threats}

DDoS attack means that uses many distributed sources to flood a target.

\begin{optlist}
\item True
\item False
\end{optlist}
\begin{explainbox}{}
\textbf{Correct Answer: True}\\[0.3em]
\textbf{Concept:} This statement correctly describes DDoS attack.
\end{explainbox}
\end{questioncard}

\begin{questioncard}{31}{tfcol}
\qtypebadge{tfcol}{True / False} \hfill \textcolor{darkslate!70}{\small Topic: Enterprise Security Threats}

DDoS attack is unrelated to uses many distributed sources to flood a target.

\begin{optlist}
\item True
\item False
\end{optlist}
\begin{explainbox}{}
\textbf{Correct Answer: False}\\[0.3em]
\textbf{Concept:} DDoS attack is directly related to uses many distributed sources to flood a target.
\end{explainbox}
\end{questioncard}

\begin{questioncard}{32}{tfcol}
\qtypebadge{tfcol}{True / False} \hfill \textcolor{darkslate!70}{\small Topic: Enterprise Security Threats}

Man-in-the-middle attack means that intercepts communication between two parties.

\begin{optlist}
\item True
\item False
\end{optlist}
\begin{explainbox}{}
\textbf{Correct Answer: True}\\[0.3em]
\textbf{Concept:} This statement correctly describes Man-in-the-middle attack.
\end{explainbox}
\end{questioncard}

\begin{questioncard}{33}{tfcol}
\qtypebadge{tfcol}{True / False} \hfill \textcolor{darkslate!70}{\small Topic: Enterprise Security Threats}

Man-in-the-middle attack is unrelated to intercepts communication between two parties.

\begin{optlist}
\item True
\item False
\end{optlist}
\begin{explainbox}{}
\textbf{Correct Answer: False}\\[0.3em]
\textbf{Concept:} Man-in-the-middle attack is directly related to intercepts communication between two parties.
\end{explainbox}
\end{questioncard}

\begin{questioncard}{34}{tfcol}
\qtypebadge{tfcol}{True / False} \hfill \textcolor{darkslate!70}{\small Topic: Enterprise Security Threats}

ARP spoofing means that forges ARP replies to redirect traffic.

\begin{optlist}
\item True
\item False
\end{optlist}
\begin{explainbox}{}
\textbf{Correct Answer: True}\\[0.3em]
\textbf{Concept:} This statement correctly describes ARP spoofing.
\end{explainbox}
\end{questioncard}

\begin{questioncard}{35}{tfcol}
\qtypebadge{tfcol}{True / False} \hfill \textcolor{darkslate!70}{\small Topic: Enterprise Security Threats}

ARP spoofing is unrelated to forges ARP replies to redirect traffic.

\begin{optlist}
\item True
\item False
\end{optlist}
\begin{explainbox}{}
\textbf{Correct Answer: False}\\[0.3em]
\textbf{Concept:} ARP spoofing is directly related to forges ARP replies to redirect traffic.
\end{explainbox}
\end{questioncard}

\begin{questioncard}{36}{tfcol}
\qtypebadge{tfcol}{True / False} \hfill \textcolor{darkslate!70}{\small Topic: Enterprise Security Threats}

SYN flood means that sends many half-open TCP connections.

\begin{optlist}
\item True
\item False
\end{optlist}
\begin{explainbox}{}
\textbf{Correct Answer: True}\\[0.3em]
\textbf{Concept:} This statement correctly describes SYN flood.
\end{explainbox}
\end{questioncard}

\begin{questioncard}{37}{tfcol}
\qtypebadge{tfcol}{True / False} \hfill \textcolor{darkslate!70}{\small Topic: Enterprise Security Threats}

SYN flood is unrelated to sends many half-open TCP connections.

\begin{optlist}
\item True
\item False
\end{optlist}
\begin{explainbox}{}
\textbf{Correct Answer: False}\\[0.3em]
\textbf{Concept:} SYN flood is directly related to sends many half-open TCP connections.
\end{explainbox}
\end{questioncard}

\begin{questioncard}{38}{tfcol}
\qtypebadge{tfcol}{True / False} \hfill \textcolor{darkslate!70}{\small Topic: Enterprise Security Threats}

Ping of death means that sends oversized ICMP packets.

\begin{optlist}
\item True
\item False
\end{optlist}
\begin{explainbox}{}
\textbf{Correct Answer: True}\\[0.3em]
\textbf{Concept:} This statement correctly describes Ping of death.
\end{explainbox}
\end{questioncard}

\begin{questioncard}{39}{tfcol}
\qtypebadge{tfcol}{True / False} \hfill \textcolor{darkslate!70}{\small Topic: Enterprise Security Threats}

Ping of death is unrelated to sends oversized ICMP packets.

\begin{optlist}
\item True
\item False
\end{optlist}
\begin{explainbox}{}
\textbf{Correct Answer: False}\\[0.3em]
\textbf{Concept:} Ping of death is directly related to sends oversized ICMP packets.
\end{explainbox}
\end{questioncard}

\begin{questioncard}{40}{tfcol}
\qtypebadge{tfcol}{True / False} \hfill \textcolor{darkslate!70}{\small Topic: Enterprise Security Threats}

Smurf attack means that uses broadcast pings to amplify traffic.

\begin{optlist}
\item True
\item False
\end{optlist}
\begin{explainbox}{}
\textbf{Correct Answer: True}\\[0.3em]
\textbf{Concept:} This statement correctly describes Smurf attack.
\end{explainbox}
\end{questioncard}

\begin{questioncard}{41}{tfcol}
\qtypebadge{tfcol}{True / False} \hfill \textcolor{darkslate!70}{\small Topic: Enterprise Security Threats}

Smurf attack is unrelated to uses broadcast pings to amplify traffic.

\begin{optlist}
\item True
\item False
\end{optlist}
\begin{explainbox}{}
\textbf{Correct Answer: False}\\[0.3em]
\textbf{Concept:} Smurf attack is directly related to uses broadcast pings to amplify traffic.
\end{explainbox}
\end{questioncard}

\begin{questioncard}{42}{tfcol}
\qtypebadge{tfcol}{True / False} \hfill \textcolor{darkslate!70}{\small Topic: Enterprise Security Threats}

Least privilege means that limits user access to the minimum necessary.

\begin{optlist}
\item True
\item False
\end{optlist}
\begin{explainbox}{}
\textbf{Correct Answer: True}\\[0.3em]
\textbf{Concept:} This statement correctly describes Least privilege.
\end{explainbox}
\end{questioncard}

\begin{questioncard}{43}{tfcol}
\qtypebadge{tfcol}{True / False} \hfill \textcolor{darkslate!70}{\small Topic: Enterprise Security Threats}

Least privilege is unrelated to limits user access to the minimum necessary.

\begin{optlist}
\item True
\item False
\end{optlist}
\begin{explainbox}{}
\textbf{Correct Answer: False}\\[0.3em]
\textbf{Concept:} Least privilege is directly related to limits user access to the minimum necessary.
\end{explainbox}
\end{questioncard}

\begin{questioncard}{44}{tfcol}
\qtypebadge{tfcol}{True / False} \hfill \textcolor{darkslate!70}{\small Topic: Enterprise Security Threats}

Patch management means that eliminates known vulnerabilities.

\begin{optlist}
\item True
\item False
\end{optlist}
\begin{explainbox}{}
\textbf{Correct Answer: True}\\[0.3em]
\textbf{Concept:} This statement correctly describes Patch management.
\end{explainbox}
\end{questioncard}

\begin{questioncard}{45}{tfcol}
\qtypebadge{tfcol}{True / False} \hfill \textcolor{darkslate!70}{\small Topic: Enterprise Security Threats}

Patch management is unrelated to eliminates known vulnerabilities.

\begin{optlist}
\item True
\item False
\end{optlist}
\begin{explainbox}{}
\textbf{Correct Answer: False}\\[0.3em]
\textbf{Concept:} Patch management is directly related to eliminates known vulnerabilities.
\end{explainbox}
\end{questioncard}

\begin{questioncard}{46}{tfcol}
\qtypebadge{tfcol}{True / False} \hfill \textcolor{darkslate!70}{\small Topic: Enterprise Security Threats}

SIEM means that centralizes log collection and analysis.

\begin{optlist}
\item True
\item False
\end{optlist}
\begin{explainbox}{}
\textbf{Correct Answer: True}\\[0.3em]
\textbf{Concept:} This statement correctly describes SIEM.
\end{explainbox}
\end{questioncard}

\begin{questioncard}{47}{tfcol}
\qtypebadge{tfcol}{True / False} \hfill \textcolor{darkslate!70}{\small Topic: Enterprise Security Threats}

SIEM is unrelated to centralizes log collection and analysis.

\begin{optlist}
\item True
\item False
\end{optlist}
\begin{explainbox}{}
\textbf{Correct Answer: False}\\[0.3em]
\textbf{Concept:} SIEM is directly related to centralizes log collection and analysis.
\end{explainbox}
\end{questioncard}

\begin{questioncard}{48}{tfcol}
\qtypebadge{tfcol}{True / False} \hfill \textcolor{darkslate!70}{\small Topic: Enterprise Security Threats}

DMZ means that hosts public-facing servers between Internet and internal network.

\begin{optlist}
\item True
\item False
\end{optlist}
\begin{explainbox}{}
\textbf{Correct Answer: True}\\[0.3em]
\textbf{Concept:} This statement correctly describes DMZ.
\end{explainbox}
\end{questioncard}

\begin{questioncard}{49}{tfcol}
\qtypebadge{tfcol}{True / False} \hfill \textcolor{darkslate!70}{\small Topic: Enterprise Security Threats}

DMZ is unrelated to hosts public-facing servers between Internet and internal network.

\begin{optlist}
\item True
\item False
\end{optlist}
\begin{explainbox}{}
\textbf{Correct Answer: False}\\[0.3em]
\textbf{Concept:} DMZ is directly related to hosts public-facing servers between Internet and internal network.
\end{explainbox}
\end{questioncard}

\begin{questioncard}{50}{tfcol}
\qtypebadge{tfcol}{True / False} \hfill \textcolor{darkslate!70}{\small Topic: Enterprise Security Threats}

Insider threat means that originates from within the organization.

\begin{optlist}
\item True
\item False
\end{optlist}
\begin{explainbox}{}
\textbf{Correct Answer: True}\\[0.3em]
\textbf{Concept:} This statement correctly describes Insider threat.
\end{explainbox}
\end{questioncard}

\begin{questioncard}{51}{tfcol}
\qtypebadge{tfcol}{True / False} \hfill \textcolor{darkslate!70}{\small Topic: Enterprise Security Threats}

Insider threat is unrelated to originates from within the organization.

\begin{optlist}
\item True
\item False
\end{optlist}
\begin{explainbox}{}
\textbf{Correct Answer: False}\\[0.3em]
\textbf{Concept:} Insider threat is directly related to originates from within the organization.
\end{explainbox}
\end{questioncard}

\section{Communication Network Security}
\begin{questioncard}{1}{tfcol}
\qtypebadge{tfcol}{True / False} \hfill \textcolor{darkslate!70}{\small Topic: Communication Network Security}

Encryption is an effective way to protect data confidentiality over public networks.

\begin{optlist}
\item True
\item False
\end{optlist}
\begin{explainbox}{}
\textbf{Correct Answer: True}\\[0.3em]
\textbf{Concept:} Encryption ensures that intercepted data cannot be read without the proper key.
\end{explainbox}
\end{questioncard}

\section{Zone Border Security}
\begin{questioncard}{1}{tfcol}
\qtypebadge{tfcol}{True / False} \hfill \textcolor{darkslate!70}{\small Topic: Zone Border Security}

A firewall at the zone border can enforce security policies between trusted and untrusted networks.

\begin{optlist}
\item True
\item False
\end{optlist}
\begin{explainbox}{}
\textbf{Correct Answer: True}\\[0.3em]
\textbf{Concept:} Firewalls inspect and control traffic crossing security zone boundaries based on policy.
\end{explainbox}
\end{questioncard}

\section{Firewall Basics}
\begin{questioncard}{1}{tfcol}
\qtypebadge{tfcol}{True / False} \hfill \textcolor{darkslate!70}{\small Topic: Firewall Basics}

Security policies on Huawei firewalls are matched from bottom to top by default.

\begin{optlist}
\item True
\item False
\end{optlist}
\begin{explainbox}{}
\textbf{Correct Answer: False}\\[0.3em]
\textbf{Concept:} Policies are matched from top to bottom based on configured order/priority.
\end{explainbox}
\end{questioncard}

\begin{questioncard}{2}{tfcol}
\qtypebadge{tfcol}{True / False} \hfill \textcolor{darkslate!70}{\small Topic: Firewall Basics}

A packet-filtering firewall tracks the state of every connection.

\begin{optlist}
\item True
\item False
\end{optlist}
\begin{explainbox}{}
\textbf{Correct Answer: False}\\[0.3em]
\textbf{Concept:} Only stateful firewalls maintain connection state; packet filters inspect packets individually.
\end{explainbox}
\end{questioncard}

\begin{questioncard}{3}{tfcol}
\qtypebadge{tfcol}{True / False} \hfill \textcolor{darkslate!70}{\small Topic: Firewall Basics}

ASPF can generate temporary server-map entries for dynamic protocols.

\begin{optlist}
\item True
\item False
\end{optlist}
\begin{explainbox}{}
\textbf{Correct Answer: True}\\[0.3em]
\textbf{Concept:} ASPF creates server-map entries to permit negotiated secondary connections.
\end{explainbox}
\end{questioncard}

\begin{questioncard}{4}{tfcol}
\qtypebadge{tfcol}{True / False} \hfill \textcolor{darkslate!70}{\small Topic: Firewall Basics}

Next-generation firewalls can identify applications regardless of port.

\begin{optlist}
\item True
\item False
\end{optlist}
\begin{explainbox}{}
\textbf{Correct Answer: True}\\[0.3em]
\textbf{Concept:} NGFWs use application-layer inspection to identify applications beyond simple port numbers.
\end{explainbox}
\end{questioncard}

\begin{questioncard}{5}{tfcol}
\qtypebadge{tfcol}{True / False} \hfill \textcolor{darkslate!70}{\small Topic: Firewall Basics}

Traffic between interfaces in the same security zone is always allowed.

\begin{optlist}
\item True
\item False
\end{optlist}
\begin{explainbox}{}
\textbf{Correct Answer: False}\\[0.3em]
\textbf{Concept:} Even intra-zone traffic can be controlled by policies, though default behavior may vary by configuration.
\end{explainbox}
\end{questioncard}

\begin{questioncard}{6}{tfcol}
\qtypebadge{tfcol}{True / False} \hfill \textcolor{darkslate!70}{\small Topic: Firewall Basics}

A deny action in a security policy can generate logs.

\begin{optlist}
\item True
\item False
\end{optlist}
\begin{explainbox}{}
\textbf{Correct Answer: True}\\[0.3em]
\textbf{Concept:} Both permit and deny actions can be configured to log traffic.
\end{explainbox}
\end{questioncard}

\begin{questioncard}{7}{tfcol}
\qtypebadge{tfcol}{True / False} \hfill \textcolor{darkslate!70}{\small Topic: Firewall Basics}

The default Huawei firewall action is permit.

\begin{optlist}
\item True
\item False
\end{optlist}
\begin{explainbox}{}
\textbf{Correct Answer: False}\\[0.3em]
\textbf{Concept:} The default action is deny unless explicitly permitted.
\end{explainbox}
\end{questioncard}

\begin{questioncard}{8}{tfcol}
\qtypebadge{tfcol}{True / False} \hfill \textcolor{darkslate!70}{\small Topic: Firewall Basics}

ASPF inspects application-layer negotiation.

\begin{optlist}
\item True
\item False
\end{optlist}
\begin{explainbox}{}
\textbf{Correct Answer: True}\\[0.3em]
\textbf{Concept:} ASPF handles dynamic protocols such as FTP.
\end{explainbox}
\end{questioncard}

\begin{questioncard}{9}{tfcol}
\qtypebadge{tfcol}{True / False} \hfill \textcolor{darkslate!70}{\small Topic: Firewall Basics}

Stateful firewalls use session tables to track connections.

\begin{optlist}
\item True
\item False
\end{optlist}
\begin{explainbox}{}
\textbf{Correct Answer: True}\\[0.3em]
\textbf{Concept:} Session tables store connection state for return traffic.
\end{explainbox}
\end{questioncard}

\begin{questioncard}{10}{tfcol}
\qtypebadge{tfcol}{True / False} \hfill \textcolor{darkslate!70}{\small Topic: Firewall Basics}

Trust zone priority means that is 85 by default.

\begin{optlist}
\item True
\item False
\end{optlist}
\begin{explainbox}{}
\textbf{Correct Answer: True}\\[0.3em]
\textbf{Concept:} This statement correctly describes Trust zone priority.
\end{explainbox}
\end{questioncard}

\begin{questioncard}{11}{tfcol}
\qtypebadge{tfcol}{True / False} \hfill \textcolor{darkslate!70}{\small Topic: Firewall Basics}

Trust zone priority is unrelated to is 85 by default.

\begin{optlist}
\item True
\item False
\end{optlist}
\begin{explainbox}{}
\textbf{Correct Answer: False}\\[0.3em]
\textbf{Concept:} Trust zone priority is directly related to is 85 by default.
\end{explainbox}
\end{questioncard}

\begin{questioncard}{12}{tfcol}
\qtypebadge{tfcol}{True / False} \hfill \textcolor{darkslate!70}{\small Topic: Firewall Basics}

DMZ zone priority means that is 50 by default.

\begin{optlist}
\item True
\item False
\end{optlist}
\begin{explainbox}{}
\textbf{Correct Answer: True}\\[0.3em]
\textbf{Concept:} This statement correctly describes DMZ zone priority.
\end{explainbox}
\end{questioncard}

\begin{questioncard}{13}{tfcol}
\qtypebadge{tfcol}{True / False} \hfill \textcolor{darkslate!70}{\small Topic: Firewall Basics}

DMZ zone priority is unrelated to is 50 by default.

\begin{optlist}
\item True
\item False
\end{optlist}
\begin{explainbox}{}
\textbf{Correct Answer: False}\\[0.3em]
\textbf{Concept:} DMZ zone priority is directly related to is 50 by default.
\end{explainbox}
\end{questioncard}

\begin{questioncard}{14}{tfcol}
\qtypebadge{tfcol}{True / False} \hfill \textcolor{darkslate!70}{\small Topic: Firewall Basics}

Untrust zone priority means that is 5 by default.

\begin{optlist}
\item True
\item False
\end{optlist}
\begin{explainbox}{}
\textbf{Correct Answer: True}\\[0.3em]
\textbf{Concept:} This statement correctly describes Untrust zone priority.
\end{explainbox}
\end{questioncard}

\begin{questioncard}{15}{tfcol}
\qtypebadge{tfcol}{True / False} \hfill \textcolor{darkslate!70}{\small Topic: Firewall Basics}

Untrust zone priority is unrelated to is 5 by default.

\begin{optlist}
\item True
\item False
\end{optlist}
\begin{explainbox}{}
\textbf{Correct Answer: False}\\[0.3em]
\textbf{Concept:} Untrust zone priority is directly related to is 5 by default.
\end{explainbox}
\end{questioncard}

\begin{questioncard}{16}{tfcol}
\qtypebadge{tfcol}{True / False} \hfill \textcolor{darkslate!70}{\small Topic: Firewall Basics}

Local zone priority means that is 100 by default.

\begin{optlist}
\item True
\item False
\end{optlist}
\begin{explainbox}{}
\textbf{Correct Answer: True}\\[0.3em]
\textbf{Concept:} This statement correctly describes Local zone priority.
\end{explainbox}
\end{questioncard}

\begin{questioncard}{17}{tfcol}
\qtypebadge{tfcol}{True / False} \hfill \textcolor{darkslate!70}{\small Topic: Firewall Basics}

Local zone priority is unrelated to is 100 by default.

\begin{optlist}
\item True
\item False
\end{optlist}
\begin{explainbox}{}
\textbf{Correct Answer: False}\\[0.3em]
\textbf{Concept:} Local zone priority is directly related to is 100 by default.
\end{explainbox}
\end{questioncard}

\begin{questioncard}{18}{tfcol}
\qtypebadge{tfcol}{True / False} \hfill \textcolor{darkslate!70}{\small Topic: Firewall Basics}

Default security policy action means that is deny.

\begin{optlist}
\item True
\item False
\end{optlist}
\begin{explainbox}{}
\textbf{Correct Answer: True}\\[0.3em]
\textbf{Concept:} This statement correctly describes Default security policy action.
\end{explainbox}
\end{questioncard}

\begin{questioncard}{19}{tfcol}
\qtypebadge{tfcol}{True / False} \hfill \textcolor{darkslate!70}{\small Topic: Firewall Basics}

Default security policy action is unrelated to is deny.

\begin{optlist}
\item True
\item False
\end{optlist}
\begin{explainbox}{}
\textbf{Correct Answer: False}\\[0.3em]
\textbf{Concept:} Default security policy action is directly related to is deny.
\end{explainbox}
\end{questioncard}

\begin{questioncard}{20}{tfcol}
\qtypebadge{tfcol}{True / False} \hfill \textcolor{darkslate!70}{\small Topic: Firewall Basics}

Stateful firewall means that tracks connections in a session table.

\begin{optlist}
\item True
\item False
\end{optlist}
\begin{explainbox}{}
\textbf{Correct Answer: True}\\[0.3em]
\textbf{Concept:} This statement correctly describes Stateful firewall.
\end{explainbox}
\end{questioncard}

\begin{questioncard}{21}{tfcol}
\qtypebadge{tfcol}{True / False} \hfill \textcolor{darkslate!70}{\small Topic: Firewall Basics}

Stateful firewall is unrelated to tracks connections in a session table.

\begin{optlist}
\item True
\item False
\end{optlist}
\begin{explainbox}{}
\textbf{Correct Answer: False}\\[0.3em]
\textbf{Concept:} Stateful firewall is directly related to tracks connections in a session table.
\end{explainbox}
\end{questioncard}

\begin{questioncard}{22}{tfcol}
\qtypebadge{tfcol}{True / False} \hfill \textcolor{darkslate!70}{\small Topic: Firewall Basics}

Packet-filtering firewall means that inspects packets individually without state.

\begin{optlist}
\item True
\item False
\end{optlist}
\begin{explainbox}{}
\textbf{Correct Answer: True}\\[0.3em]
\textbf{Concept:} This statement correctly describes Packet-filtering firewall.
\end{explainbox}
\end{questioncard}

\begin{questioncard}{23}{tfcol}
\qtypebadge{tfcol}{True / False} \hfill \textcolor{darkslate!70}{\small Topic: Firewall Basics}

Packet-filtering firewall is unrelated to inspects packets individually without state.

\begin{optlist}
\item True
\item False
\end{optlist}
\begin{explainbox}{}
\textbf{Correct Answer: False}\\[0.3em]
\textbf{Concept:} Packet-filtering firewall is directly related to inspects packets individually without state.
\end{explainbox}
\end{questioncard}

\begin{questioncard}{24}{tfcol}
\qtypebadge{tfcol}{True / False} \hfill \textcolor{darkslate!70}{\small Topic: Firewall Basics}

ASPF means that inspects multi-channel protocols like FTP.

\begin{optlist}
\item True
\item False
\end{optlist}
\begin{explainbox}{}
\textbf{Correct Answer: True}\\[0.3em]
\textbf{Concept:} This statement correctly describes ASPF.
\end{explainbox}
\end{questioncard}

\begin{questioncard}{25}{tfcol}
\qtypebadge{tfcol}{True / False} \hfill \textcolor{darkslate!70}{\small Topic: Firewall Basics}

ASPF is unrelated to inspects multi-channel protocols like FTP.

\begin{optlist}
\item True
\item False
\end{optlist}
\begin{explainbox}{}
\textbf{Correct Answer: False}\\[0.3em]
\textbf{Concept:} ASPF is directly related to inspects multi-channel protocols like FTP.
\end{explainbox}
\end{questioncard}

\begin{questioncard}{26}{tfcol}
\qtypebadge{tfcol}{True / False} \hfill \textcolor{darkslate!70}{\small Topic: Firewall Basics}

Next-generation firewall means that includes application awareness and IPS.

\begin{optlist}
\item True
\item False
\end{optlist}
\begin{explainbox}{}
\textbf{Correct Answer: True}\\[0.3em]
\textbf{Concept:} This statement correctly describes Next-generation firewall.
\end{explainbox}
\end{questioncard}

\begin{questioncard}{27}{tfcol}
\qtypebadge{tfcol}{True / False} \hfill \textcolor{darkslate!70}{\small Topic: Firewall Basics}

Next-generation firewall is unrelated to includes application awareness and IPS.

\begin{optlist}
\item True
\item False
\end{optlist}
\begin{explainbox}{}
\textbf{Correct Answer: False}\\[0.3em]
\textbf{Concept:} Next-generation firewall is directly related to includes application awareness and IPS.
\end{explainbox}
\end{questioncard}

\begin{questioncard}{28}{tfcol}
\qtypebadge{tfcol}{True / False} \hfill \textcolor{darkslate!70}{\small Topic: Firewall Basics}

Security zones means that group interfaces with similar trust levels.

\begin{optlist}
\item True
\item False
\end{optlist}
\begin{explainbox}{}
\textbf{Correct Answer: True}\\[0.3em]
\textbf{Concept:} This statement correctly describes Security zones.
\end{explainbox}
\end{questioncard}

\begin{questioncard}{29}{tfcol}
\qtypebadge{tfcol}{True / False} \hfill \textcolor{darkslate!70}{\small Topic: Firewall Basics}

Security zones is unrelated to group interfaces with similar trust levels.

\begin{optlist}
\item True
\item False
\end{optlist}
\begin{explainbox}{}
\textbf{Correct Answer: False}\\[0.3em]
\textbf{Concept:} Security zones is directly related to group interfaces with similar trust levels.
\end{explainbox}
\end{questioncard}

\begin{questioncard}{30}{tfcol}
\qtypebadge{tfcol}{True / False} \hfill \textcolor{darkslate!70}{\small Topic: Firewall Basics}

Intra-zone traffic means that can be controlled by policies.

\begin{optlist}
\item True
\item False
\end{optlist}
\begin{explainbox}{}
\textbf{Correct Answer: True}\\[0.3em]
\textbf{Concept:} This statement correctly describes Intra-zone traffic.
\end{explainbox}
\end{questioncard}

\begin{questioncard}{31}{tfcol}
\qtypebadge{tfcol}{True / False} \hfill \textcolor{darkslate!70}{\small Topic: Firewall Basics}

Intra-zone traffic is unrelated to can be controlled by policies.

\begin{optlist}
\item True
\item False
\end{optlist}
\begin{explainbox}{}
\textbf{Correct Answer: False}\\[0.3em]
\textbf{Concept:} Intra-zone traffic is directly related to can be controlled by policies.
\end{explainbox}
\end{questioncard}

\begin{questioncard}{32}{tfcol}
\qtypebadge{tfcol}{True / False} \hfill \textcolor{darkslate!70}{\small Topic: Firewall Basics}

Policy hit count means that helps verify which rules are used.

\begin{optlist}
\item True
\item False
\end{optlist}
\begin{explainbox}{}
\textbf{Correct Answer: True}\\[0.3em]
\textbf{Concept:} This statement correctly describes Policy hit count.
\end{explainbox}
\end{questioncard}

\begin{questioncard}{33}{tfcol}
\qtypebadge{tfcol}{True / False} \hfill \textcolor{darkslate!70}{\small Topic: Firewall Basics}

Policy hit count is unrelated to helps verify which rules are used.

\begin{optlist}
\item True
\item False
\end{optlist}
\begin{explainbox}{}
\textbf{Correct Answer: False}\\[0.3em]
\textbf{Concept:} Policy hit count is directly related to helps verify which rules are used.
\end{explainbox}
\end{questioncard}

\begin{questioncard}{34}{tfcol}
\qtypebadge{tfcol}{True / False} \hfill \textcolor{darkslate!70}{\small Topic: Firewall Basics}

Security policy logging means that can be enabled for permit and deny actions.

\begin{optlist}
\item True
\item False
\end{optlist}
\begin{explainbox}{}
\textbf{Correct Answer: True}\\[0.3em]
\textbf{Concept:} This statement correctly describes Security policy logging.
\end{explainbox}
\end{questioncard}

\begin{questioncard}{35}{tfcol}
\qtypebadge{tfcol}{True / False} \hfill \textcolor{darkslate!70}{\small Topic: Firewall Basics}

Security policy logging is unrelated to can be enabled for permit and deny actions.

\begin{optlist}
\item True
\item False
\end{optlist}
\begin{explainbox}{}
\textbf{Correct Answer: False}\\[0.3em]
\textbf{Concept:} Security policy logging is directly related to can be enabled for permit and deny actions.
\end{explainbox}
\end{questioncard}

\begin{questioncard}{36}{tfcol}
\qtypebadge{tfcol}{True / False} \hfill \textcolor{darkslate!70}{\small Topic: Firewall Basics}

Routing mode means that is a firewall deployment mode.

\begin{optlist}
\item True
\item False
\end{optlist}
\begin{explainbox}{}
\textbf{Correct Answer: True}\\[0.3em]
\textbf{Concept:} This statement correctly describes Routing mode.
\end{explainbox}
\end{questioncard}

\begin{questioncard}{37}{tfcol}
\qtypebadge{tfcol}{True / False} \hfill \textcolor{darkslate!70}{\small Topic: Firewall Basics}

Routing mode is unrelated to is a firewall deployment mode.

\begin{optlist}
\item True
\item False
\end{optlist}
\begin{explainbox}{}
\textbf{Correct Answer: False}\\[0.3em]
\textbf{Concept:} Routing mode is directly related to is a firewall deployment mode.
\end{explainbox}
\end{questioncard}

\begin{questioncard}{38}{tfcol}
\qtypebadge{tfcol}{True / False} \hfill \textcolor{darkslate!70}{\small Topic: Firewall Basics}

Transparent mode means that is a firewall deployment mode.

\begin{optlist}
\item True
\item False
\end{optlist}
\begin{explainbox}{}
\textbf{Correct Answer: True}\\[0.3em]
\textbf{Concept:} This statement correctly describes Transparent mode.
\end{explainbox}
\end{questioncard}

\begin{questioncard}{39}{tfcol}
\qtypebadge{tfcol}{True / False} \hfill \textcolor{darkslate!70}{\small Topic: Firewall Basics}

Transparent mode is unrelated to is a firewall deployment mode.

\begin{optlist}
\item True
\item False
\end{optlist}
\begin{explainbox}{}
\textbf{Correct Answer: False}\\[0.3em]
\textbf{Concept:} Transparent mode is directly related to is a firewall deployment mode.
\end{explainbox}
\end{questioncard}

\begin{questioncard}{40}{tfcol}
\qtypebadge{tfcol}{True / False} \hfill \textcolor{darkslate!70}{\small Topic: Firewall Basics}

Hybrid mode means that is a firewall deployment mode.

\begin{optlist}
\item True
\item False
\end{optlist}
\begin{explainbox}{}
\textbf{Correct Answer: True}\\[0.3em]
\textbf{Concept:} This statement correctly describes Hybrid mode.
\end{explainbox}
\end{questioncard}

\begin{questioncard}{41}{tfcol}
\qtypebadge{tfcol}{True / False} \hfill \textcolor{darkslate!70}{\small Topic: Firewall Basics}

Hybrid mode is unrelated to is a firewall deployment mode.

\begin{optlist}
\item True
\item False
\end{optlist}
\begin{explainbox}{}
\textbf{Correct Answer: False}\\[0.3em]
\textbf{Concept:} Hybrid mode is directly related to is a firewall deployment mode.
\end{explainbox}
\end{questioncard}

\section{Security Zones}
\begin{questioncard}{1}{tfcol}
\qtypebadge{tfcol}{True / False} \hfill \textcolor{darkslate!70}{\small Topic: Security Zones}

The Local zone represents traffic destined to or originated from the firewall itself.

\begin{optlist}
\item True
\item False
\end{optlist}
\begin{explainbox}{}
\textbf{Correct Answer: True}\\[0.3em]
\textbf{Concept:} Local zone traffic involves the firewall services and management interfaces.
\end{explainbox}
\end{questioncard}

\section{Security Policies}
\begin{questioncard}{1}{tfcol}
\qtypebadge{tfcol}{True / False} \hfill \textcolor{darkslate!70}{\small Topic: Security Policies}

By default, Huawei firewalls permit all traffic unless a policy explicitly denies it.

\begin{optlist}
\item True
\item False
\end{optlist}
\begin{explainbox}{}
\textbf{Correct Answer: False}\\[0.3em]
\textbf{Concept:} The default action is to deny traffic unless a security policy explicitly permits it.
\end{explainbox}
\end{questioncard}

\section{Stateful Inspection}
\begin{questioncard}{1}{tfcol}
\qtypebadge{tfcol}{True / False} \hfill \textcolor{darkslate!70}{\small Topic: Stateful Inspection}

A stateful firewall can allow return traffic for an established session without a separate policy.

\begin{optlist}
\item True
\item False
\end{optlist}
\begin{explainbox}{}
\textbf{Correct Answer: True}\\[0.3em]
\textbf{Concept:} Return packets matching an existing session table entry are permitted automatically.
\end{explainbox}
\end{questioncard}

\section{ASPF}
\begin{questioncard}{1}{tfcol}
\qtypebadge{tfcol}{True / False} \hfill \textcolor{darkslate!70}{\small Topic: ASPF}

ASPF is used to support protocols that open secondary or dynamic connections.

\begin{optlist}
\item True
\item False
\end{optlist}
\begin{explainbox}{}
\textbf{Correct Answer: True}\\[0.3em]
\textbf{Concept:} ASPF inspects application-layer negotiation to dynamically permit related channels.
\end{explainbox}
\end{questioncard}

\section{NAT}
\begin{questioncard}{1}{tfcol}
\qtypebadge{tfcol}{True / False} \hfill \textcolor{darkslate!70}{\small Topic: NAT}

NAT can hide internal network topology from external observers.

\begin{optlist}
\item True
\item False
\end{optlist}
\begin{explainbox}{}
\textbf{Correct Answer: True}\\[0.3em]
\textbf{Concept:} NAT translates private addresses, making internal topology invisible from outside.
\end{explainbox}
\end{questioncard}

\begin{questioncard}{2}{tfcol}
\qtypebadge{tfcol}{True / False} \hfill \textcolor{darkslate!70}{\small Topic: NAT}

Easy-IP uses a fixed public IP address pool.

\begin{optlist}
\item True
\item False
\end{optlist}
\begin{explainbox}{}
\textbf{Correct Answer: False}\\[0.3em]
\textbf{Concept:} Easy-IP uses the firewall egress interface IP address dynamically.
\end{explainbox}
\end{questioncard}

\begin{questioncard}{3}{tfcol}
\qtypebadge{tfcol}{True / False} \hfill \textcolor{darkslate!70}{\small Topic: NAT}

NAT ALG is required for HTTP because HTTP embeds IP addresses in payloads.

\begin{optlist}
\item True
\item False
\end{optlist}
\begin{explainbox}{}
\textbf{Correct Answer: False}\\[0.3em]
\textbf{Concept:} HTTP generally does not embed IP addresses in payloads, so ALG is not required.
\end{explainbox}
\end{questioncard}

\begin{questioncard}{4}{tfcol}
\qtypebadge{tfcol}{True / False} \hfill \textcolor{darkslate!70}{\small Topic: NAT}

NAPT conserves public IPv4 addresses by using port multiplexing.

\begin{optlist}
\item True
\item False
\end{optlist}
\begin{explainbox}{}
\textbf{Correct Answer: True}\\[0.3em]
\textbf{Concept:} NAPT maps many private addresses to one public IP using unique port numbers.
\end{explainbox}
\end{questioncard}

\begin{questioncard}{5}{tfcol}
\qtypebadge{tfcol}{True / False} \hfill \textcolor{darkslate!70}{\small Topic: NAT}

FTP data channel failure can be caused by disabled NAT ALG.

\begin{optlist}
\item True
\item False
\end{optlist}
\begin{explainbox}{}
\textbf{Correct Answer: True}\\[0.3em]
\textbf{Concept:} Without ALG, the firewall cannot identify and permit the negotiated FTP data channel.
\end{explainbox}
\end{questioncard}

\begin{questioncard}{6}{tfcol}
\qtypebadge{tfcol}{True / False} \hfill \textcolor{darkslate!70}{\small Topic: NAT}

NAPT translates both IP addresses and port numbers.

\begin{optlist}
\item True
\item False
\end{optlist}
\begin{explainbox}{}
\textbf{Correct Answer: True}\\[0.3em]
\textbf{Concept:} NAPT uses port multiplexing for many-to-one translation.
\end{explainbox}
\end{questioncard}

\begin{questioncard}{7}{tfcol}
\qtypebadge{tfcol}{True / False} \hfill \textcolor{darkslate!70}{\small Topic: NAT}

NAT Server is used for outbound Internet access.

\begin{optlist}
\item True
\item False
\end{optlist}
\begin{explainbox}{}
\textbf{Correct Answer: False}\\[0.3em]
\textbf{Concept:} NAT Server publishes inbound services.
\end{explainbox}
\end{questioncard}

\begin{questioncard}{8}{tfcol}
\qtypebadge{tfcol}{True / False} \hfill \textcolor{darkslate!70}{\small Topic: NAT}

Static NAT preserves a one-to-one public-private mapping.

\begin{optlist}
\item True
\item False
\end{optlist}
\begin{explainbox}{}
\textbf{Correct Answer: True}\\[0.3em]
\textbf{Concept:} Static NAT maps one public IP to one private IP consistently.
\end{explainbox}
\end{questioncard}

\begin{questioncard}{9}{tfcol}
\qtypebadge{tfcol}{True / False} \hfill \textcolor{darkslate!70}{\small Topic: NAT}

NAT means that translates IP addresses between private and public.

\begin{optlist}
\item True
\item False
\end{optlist}
\begin{explainbox}{}
\textbf{Correct Answer: True}\\[0.3em]
\textbf{Concept:} This statement correctly describes NAT.
\end{explainbox}
\end{questioncard}

\begin{questioncard}{10}{tfcol}
\qtypebadge{tfcol}{True / False} \hfill \textcolor{darkslate!70}{\small Topic: NAT}

NAT is unrelated to translates IP addresses between private and public.

\begin{optlist}
\item True
\item False
\end{optlist}
\begin{explainbox}{}
\textbf{Correct Answer: False}\\[0.3em]
\textbf{Concept:} NAT is directly related to translates IP addresses between private and public.
\end{explainbox}
\end{questioncard}

\begin{questioncard}{11}{tfcol}
\qtypebadge{tfcol}{True / False} \hfill \textcolor{darkslate!70}{\small Topic: NAT}

Source NAT means that translates the source IP of outbound packets.

\begin{optlist}
\item True
\item False
\end{optlist}
\begin{explainbox}{}
\textbf{Correct Answer: True}\\[0.3em]
\textbf{Concept:} This statement correctly describes Source NAT.
\end{explainbox}
\end{questioncard}

\begin{questioncard}{12}{tfcol}
\qtypebadge{tfcol}{True / False} \hfill \textcolor{darkslate!70}{\small Topic: NAT}

Source NAT is unrelated to translates the source IP of outbound packets.

\begin{optlist}
\item True
\item False
\end{optlist}
\begin{explainbox}{}
\textbf{Correct Answer: False}\\[0.3em]
\textbf{Concept:} Source NAT is directly related to translates the source IP of outbound packets.
\end{explainbox}
\end{questioncard}

\begin{questioncard}{13}{tfcol}
\qtypebadge{tfcol}{True / False} \hfill \textcolor{darkslate!70}{\small Topic: NAT}

Destination NAT means that translates the destination IP of inbound packets.

\begin{optlist}
\item True
\item False
\end{optlist}
\begin{explainbox}{}
\textbf{Correct Answer: True}\\[0.3em]
\textbf{Concept:} This statement correctly describes Destination NAT.
\end{explainbox}
\end{questioncard}

\begin{questioncard}{14}{tfcol}
\qtypebadge{tfcol}{True / False} \hfill \textcolor{darkslate!70}{\small Topic: NAT}

Destination NAT is unrelated to translates the destination IP of inbound packets.

\begin{optlist}
\item True
\item False
\end{optlist}
\begin{explainbox}{}
\textbf{Correct Answer: False}\\[0.3em]
\textbf{Concept:} Destination NAT is directly related to translates the destination IP of inbound packets.
\end{explainbox}
\end{questioncard}

\begin{questioncard}{15}{tfcol}
\qtypebadge{tfcol}{True / False} \hfill \textcolor{darkslate!70}{\small Topic: NAT}

NAT Server means that publishes internal services to the Internet.

\begin{optlist}
\item True
\item False
\end{optlist}
\begin{explainbox}{}
\textbf{Correct Answer: True}\\[0.3em]
\textbf{Concept:} This statement correctly describes NAT Server.
\end{explainbox}
\end{questioncard}

\begin{questioncard}{16}{tfcol}
\qtypebadge{tfcol}{True / False} \hfill \textcolor{darkslate!70}{\small Topic: NAT}

NAT Server is unrelated to publishes internal services to the Internet.

\begin{optlist}
\item True
\item False
\end{optlist}
\begin{explainbox}{}
\textbf{Correct Answer: False}\\[0.3em]
\textbf{Concept:} NAT Server is directly related to publishes internal services to the Internet.
\end{explainbox}
\end{questioncard}

\begin{questioncard}{17}{tfcol}
\qtypebadge{tfcol}{True / False} \hfill \textcolor{darkslate!70}{\small Topic: NAT}

Bidirectional NAT means that translates both source and destination addresses.

\begin{optlist}
\item True
\item False
\end{optlist}
\begin{explainbox}{}
\textbf{Correct Answer: True}\\[0.3em]
\textbf{Concept:} This statement correctly describes Bidirectional NAT.
\end{explainbox}
\end{questioncard}

\begin{questioncard}{18}{tfcol}
\qtypebadge{tfcol}{True / False} \hfill \textcolor{darkslate!70}{\small Topic: NAT}

Bidirectional NAT is unrelated to translates both source and destination addresses.

\begin{optlist}
\item True
\item False
\end{optlist}
\begin{explainbox}{}
\textbf{Correct Answer: False}\\[0.3em]
\textbf{Concept:} Bidirectional NAT is directly related to translates both source and destination addresses.
\end{explainbox}
\end{questioncard}

\begin{questioncard}{19}{tfcol}
\qtypebadge{tfcol}{True / False} \hfill \textcolor{darkslate!70}{\small Topic: NAT}

NAPT means that translates many private addresses to one public IP using ports.

\begin{optlist}
\item True
\item False
\end{optlist}
\begin{explainbox}{}
\textbf{Correct Answer: True}\\[0.3em]
\textbf{Concept:} This statement correctly describes NAPT.
\end{explainbox}
\end{questioncard}

\begin{questioncard}{20}{tfcol}
\qtypebadge{tfcol}{True / False} \hfill \textcolor{darkslate!70}{\small Topic: NAT}

NAPT is unrelated to translates many private addresses to one public IP using ports.

\begin{optlist}
\item True
\item False
\end{optlist}
\begin{explainbox}{}
\textbf{Correct Answer: False}\\[0.3em]
\textbf{Concept:} NAPT is directly related to translates many private addresses to one public IP using ports.
\end{explainbox}
\end{questioncard}

\begin{questioncard}{21}{tfcol}
\qtypebadge{tfcol}{True / False} \hfill \textcolor{darkslate!70}{\small Topic: NAT}

NAT No-PAT means that performs one-to-one address translation preserving ports.

\begin{optlist}
\item True
\item False
\end{optlist}
\begin{explainbox}{}
\textbf{Correct Answer: True}\\[0.3em]
\textbf{Concept:} This statement correctly describes NAT No-PAT.
\end{explainbox}
\end{questioncard}

\begin{questioncard}{22}{tfcol}
\qtypebadge{tfcol}{True / False} \hfill \textcolor{darkslate!70}{\small Topic: NAT}

NAT No-PAT is unrelated to performs one-to-one address translation preserving ports.

\begin{optlist}
\item True
\item False
\end{optlist}
\begin{explainbox}{}
\textbf{Correct Answer: False}\\[0.3em]
\textbf{Concept:} NAT No-PAT is directly related to performs one-to-one address translation preserving ports.
\end{explainbox}
\end{questioncard}

\begin{questioncard}{23}{tfcol}
\qtypebadge{tfcol}{True / False} \hfill \textcolor{darkslate!70}{\small Topic: NAT}

Easy-IP means that uses the outbound interface IP for translation.

\begin{optlist}
\item True
\item False
\end{optlist}
\begin{explainbox}{}
\textbf{Correct Answer: True}\\[0.3em]
\textbf{Concept:} This statement correctly describes Easy-IP.
\end{explainbox}
\end{questioncard}

\begin{questioncard}{24}{tfcol}
\qtypebadge{tfcol}{True / False} \hfill \textcolor{darkslate!70}{\small Topic: NAT}

Easy-IP is unrelated to uses the outbound interface IP for translation.

\begin{optlist}
\item True
\item False
\end{optlist}
\begin{explainbox}{}
\textbf{Correct Answer: False}\\[0.3em]
\textbf{Concept:} Easy-IP is directly related to uses the outbound interface IP for translation.
\end{explainbox}
\end{questioncard}

\begin{questioncard}{25}{tfcol}
\qtypebadge{tfcol}{True / False} \hfill \textcolor{darkslate!70}{\small Topic: NAT}

NAT ALG means that rewrites addresses inside application payloads.

\begin{optlist}
\item True
\item False
\end{optlist}
\begin{explainbox}{}
\textbf{Correct Answer: True}\\[0.3em]
\textbf{Concept:} This statement correctly describes NAT ALG.
\end{explainbox}
\end{questioncard}

\begin{questioncard}{26}{tfcol}
\qtypebadge{tfcol}{True / False} \hfill \textcolor{darkslate!70}{\small Topic: NAT}

NAT ALG is unrelated to rewrites addresses inside application payloads.

\begin{optlist}
\item True
\item False
\end{optlist}
\begin{explainbox}{}
\textbf{Correct Answer: False}\\[0.3em]
\textbf{Concept:} NAT ALG is directly related to rewrites addresses inside application payloads.
\end{explainbox}
\end{questioncard}

\begin{questioncard}{27}{tfcol}
\qtypebadge{tfcol}{True / False} \hfill \textcolor{darkslate!70}{\small Topic: NAT}

RFC 1918 means that defines private IPv4 address ranges.

\begin{optlist}
\item True
\item False
\end{optlist}
\begin{explainbox}{}
\textbf{Correct Answer: True}\\[0.3em]
\textbf{Concept:} This statement correctly describes RFC 1918.
\end{explainbox}
\end{questioncard}

\begin{questioncard}{28}{tfcol}
\qtypebadge{tfcol}{True / False} \hfill \textcolor{darkslate!70}{\small Topic: NAT}

RFC 1918 is unrelated to defines private IPv4 address ranges.

\begin{optlist}
\item True
\item False
\end{optlist}
\begin{explainbox}{}
\textbf{Correct Answer: False}\\[0.3em]
\textbf{Concept:} RFC 1918 is directly related to defines private IPv4 address ranges.
\end{explainbox}
\end{questioncard}

\begin{questioncard}{29}{tfcol}
\qtypebadge{tfcol}{True / False} \hfill \textcolor{darkslate!70}{\small Topic: NAT}

Hairpin NAT means that lets internal users access servers via public NAT address.

\begin{optlist}
\item True
\item False
\end{optlist}
\begin{explainbox}{}
\textbf{Correct Answer: True}\\[0.3em]
\textbf{Concept:} This statement correctly describes Hairpin NAT.
\end{explainbox}
\end{questioncard}

\begin{questioncard}{30}{tfcol}
\qtypebadge{tfcol}{True / False} \hfill \textcolor{darkslate!70}{\small Topic: NAT}

Hairpin NAT is unrelated to lets internal users access servers via public NAT address.

\begin{optlist}
\item True
\item False
\end{optlist}
\begin{explainbox}{}
\textbf{Correct Answer: False}\\[0.3em]
\textbf{Concept:} Hairpin NAT is directly related to lets internal users access servers via public NAT address.
\end{explainbox}
\end{questioncard}

\begin{questioncard}{31}{tfcol}
\qtypebadge{tfcol}{True / False} \hfill \textcolor{darkslate!70}{\small Topic: NAT}

Static NAT means that provides a persistent one-to-one mapping.

\begin{optlist}
\item True
\item False
\end{optlist}
\begin{explainbox}{}
\textbf{Correct Answer: True}\\[0.3em]
\textbf{Concept:} This statement correctly describes Static NAT.
\end{explainbox}
\end{questioncard}

\begin{questioncard}{32}{tfcol}
\qtypebadge{tfcol}{True / False} \hfill \textcolor{darkslate!70}{\small Topic: NAT}

Static NAT is unrelated to provides a persistent one-to-one mapping.

\begin{optlist}
\item True
\item False
\end{optlist}
\begin{explainbox}{}
\textbf{Correct Answer: False}\\[0.3em]
\textbf{Concept:} Static NAT is directly related to provides a persistent one-to-one mapping.
\end{explainbox}
\end{questioncard}

\section{NAT Basics}
\begin{questioncard}{1}{tfcol}
\qtypebadge{tfcol}{True / False} \hfill \textcolor{darkslate!70}{\small Topic: NAT Basics}

NAT can conserve public IPv4 addresses.

\begin{optlist}
\item True
\item False
\end{optlist}
\begin{explainbox}{}
\textbf{Correct Answer: True}\\[0.3em]
\textbf{Concept:} NAT allows many private hosts to share a smaller number of public IP addresses.
\end{explainbox}
\end{questioncard}

\section{Source NAT}
\begin{questioncard}{1}{tfcol}
\qtypebadge{tfcol}{True / False} \hfill \textcolor{darkslate!70}{\small Topic: Source NAT}

NAPT translates both IP addresses and port numbers.

\begin{optlist}
\item True
\item False
\end{optlist}
\begin{explainbox}{}
\textbf{Correct Answer: True}\\[0.3em]
\textbf{Concept:} NAPT multiplexes many private addresses onto one public IP by translating ports.
\end{explainbox}
\end{questioncard}

\section{Destination NAT}
\begin{questioncard}{1}{tfcol}
\qtypebadge{tfcol}{True / False} \hfill \textcolor{darkslate!70}{\small Topic: Destination NAT}

Destination NAT is typically used for outbound Internet access.

\begin{optlist}
\item True
\item False
\end{optlist}
\begin{explainbox}{}
\textbf{Correct Answer: False}\\[0.3em]
\textbf{Concept:} Destination NAT is used for inbound traffic to expose internal services.
\end{explainbox}
\end{questioncard}

\section{NAT Server}
\begin{questioncard}{1}{tfcol}
\qtypebadge{tfcol}{True / False} \hfill \textcolor{darkslate!70}{\small Topic: NAT Server}

NAT Server maps public IP addresses to private server addresses.

\begin{optlist}
\item True
\item False
\end{optlist}
\begin{explainbox}{}
\textbf{Correct Answer: True}\\[0.3em]
\textbf{Concept:} NAT Server publishes internal services by mapping public IPs/ports to private servers.
\end{explainbox}
\end{questioncard}

\section{NAT ALG}
\begin{questioncard}{1}{tfcol}
\qtypebadge{tfcol}{True / False} \hfill \textcolor{darkslate!70}{\small Topic: NAT ALG}

NAT ALG is needed when protocols embed IP addresses inside application payloads.

\begin{optlist}
\item True
\item False
\end{optlist}
\begin{explainbox}{}
\textbf{Correct Answer: True}\\[0.3em]
\textbf{Concept:} ALG inspects application-layer data to update embedded addresses/ports.
\end{explainbox}
\end{questioncard}

\section{Firewall Hot Standby}
\begin{questioncard}{1}{tfcol}
\qtypebadge{tfcol}{True / False} \hfill \textcolor{darkslate!70}{\small Topic: Firewall Hot Standby}

In active/standby mode, both firewalls actively forward the same traffic.

\begin{optlist}
\item True
\item False
\end{optlist}
\begin{explainbox}{}
\textbf{Correct Answer: False}\\[0.3em]
\textbf{Concept:} Only the active firewall forwards traffic; the standby takes over upon failure.
\end{explainbox}
\end{questioncard}

\begin{questioncard}{2}{tfcol}
\qtypebadge{tfcol}{True / False} \hfill \textcolor{darkslate!70}{\small Topic: Firewall Hot Standby}

VRRP provides default gateway redundancy for end hosts.

\begin{optlist}
\item True
\item False
\end{optlist}
\begin{explainbox}{}
\textbf{Correct Answer: True}\\[0.3em]
\textbf{Concept:} VRRP allows multiple routers to act as a single virtual gateway.
\end{explainbox}
\end{questioncard}

\begin{questioncard}{3}{tfcol}
\qtypebadge{tfcol}{True / False} \hfill \textcolor{darkslate!70}{\small Topic: Firewall Hot Standby}

In VRRP, only the router with the physical IP matching the virtual IP can become master.

\begin{optlist}
\item True
\item False
\end{optlist}
\begin{explainbox}{}
\textbf{Correct Answer: False}\\[0.3em]
\textbf{Concept:} Any router with sufficient priority can become master, not just the IP owner.
\end{explainbox}
\end{questioncard}

\begin{questioncard}{4}{tfcol}
\qtypebadge{tfcol}{True / False} \hfill \textcolor{darkslate!70}{\small Topic: Firewall Hot Standby}

VGMP prevents individual VRRP groups from becoming inconsistent during failover.

\begin{optlist}
\item True
\item False
\end{optlist}
\begin{explainbox}{}
\textbf{Correct Answer: True}\\[0.3em]
\textbf{Concept:} VGMP manages multiple VRRP groups as a single entity.
\end{explainbox}
\end{questioncard}

\begin{questioncard}{5}{tfcol}
\qtypebadge{tfcol}{True / False} \hfill \textcolor{darkslate!70}{\small Topic: Firewall Hot Standby}

HRP can synchronize NAT sessions between active and standby firewalls.

\begin{optlist}
\item True
\item False
\end{optlist}
\begin{explainbox}{}
\textbf{Correct Answer: True}\\[0.3em]
\textbf{Concept:} HRP keeps NAT/session state consistent across hot-standby peers.
\end{explainbox}
\end{questioncard}

\begin{questioncard}{6}{tfcol}
\qtypebadge{tfcol}{True / False} \hfill \textcolor{darkslate!70}{\small Topic: Firewall Hot Standby}

Load-sharing mode improves resource utilization compared to active/standby mode.

\begin{optlist}
\item True
\item False
\end{optlist}
\begin{explainbox}{}
\textbf{Correct Answer: True}\\[0.3em]
\textbf{Concept:} Multiple firewalls forward traffic simultaneously in load-sharing mode.
\end{explainbox}
\end{questioncard}

\begin{questioncard}{7}{tfcol}
\qtypebadge{tfcol}{True / False} \hfill \textcolor{darkslate!70}{\small Topic: Firewall Hot Standby}

The HRP heartbeat link can share bandwidth with regular user traffic without risk.

\begin{optlist}
\item True
\item False
\end{optlist}
\begin{explainbox}{}
\textbf{Correct Answer: False}\\[0.3em]
\textbf{Concept:} Sharing heartbeat and user traffic can cause false failover due to congestion.
\end{explainbox}
\end{questioncard}

\begin{questioncard}{8}{tfcol}
\qtypebadge{tfcol}{True / False} \hfill \textcolor{darkslate!70}{\small Topic: Firewall Hot Standby}

VRRP provides a virtual gateway IP.

\begin{optlist}
\item True
\item False
\end{optlist}
\begin{explainbox}{}
\textbf{Correct Answer: True}\\[0.3em]
\textbf{Concept:} VRRP allows multiple routers to share a virtual IP.
\end{explainbox}
\end{questioncard}

\begin{questioncard}{9}{tfcol}
\qtypebadge{tfcol}{True / False} \hfill \textcolor{darkslate!70}{\small Topic: Firewall Hot Standby}

HRP synchronizes only configurations.

\begin{optlist}
\item True
\item False
\end{optlist}
\begin{explainbox}{}
\textbf{Correct Answer: False}\\[0.3em]
\textbf{Concept:} HRP synchronizes configurations and runtime state.
\end{explainbox}
\end{questioncard}

\begin{questioncard}{10}{tfcol}
\qtypebadge{tfcol}{True / False} \hfill \textcolor{darkslate!70}{\small Topic: Firewall Hot Standby}

Load-sharing mode allows multiple firewalls to forward traffic.

\begin{optlist}
\item True
\item False
\end{optlist}
\begin{explainbox}{}
\textbf{Correct Answer: True}\\[0.3em]
\textbf{Concept:} Multiple active firewalls share the traffic load.
\end{explainbox}
\end{questioncard}

\begin{questioncard}{11}{tfcol}
\qtypebadge{tfcol}{True / False} \hfill \textcolor{darkslate!70}{\small Topic: Firewall Hot Standby}

VRRP means that provides gateway redundancy with a virtual IP.

\begin{optlist}
\item True
\item False
\end{optlist}
\begin{explainbox}{}
\textbf{Correct Answer: True}\\[0.3em]
\textbf{Concept:} This statement correctly describes VRRP.
\end{explainbox}
\end{questioncard}

\begin{questioncard}{12}{tfcol}
\qtypebadge{tfcol}{True / False} \hfill \textcolor{darkslate!70}{\small Topic: Firewall Hot Standby}

VRRP is unrelated to provides gateway redundancy with a virtual IP.

\begin{optlist}
\item True
\item False
\end{optlist}
\begin{explainbox}{}
\textbf{Correct Answer: False}\\[0.3em]
\textbf{Concept:} VRRP is directly related to provides gateway redundancy with a virtual IP.
\end{explainbox}
\end{questioncard}

\begin{questioncard}{13}{tfcol}
\qtypebadge{tfcol}{True / False} \hfill \textcolor{darkslate!70}{\small Topic: Firewall Hot Standby}

VRRP master means that forwards traffic for the virtual IP.

\begin{optlist}
\item True
\item False
\end{optlist}
\begin{explainbox}{}
\textbf{Correct Answer: True}\\[0.3em]
\textbf{Concept:} This statement correctly describes VRRP master.
\end{explainbox}
\end{questioncard}

\begin{questioncard}{14}{tfcol}
\qtypebadge{tfcol}{True / False} \hfill \textcolor{darkslate!70}{\small Topic: Firewall Hot Standby}

VRRP master is unrelated to forwards traffic for the virtual IP.

\begin{optlist}
\item True
\item False
\end{optlist}
\begin{explainbox}{}
\textbf{Correct Answer: False}\\[0.3em]
\textbf{Concept:} VRRP master is directly related to forwards traffic for the virtual IP.
\end{explainbox}
\end{questioncard}

\begin{questioncard}{15}{tfcol}
\qtypebadge{tfcol}{True / False} \hfill \textcolor{darkslate!70}{\small Topic: Firewall Hot Standby}

VRRP backup means that takes over if the master fails.

\begin{optlist}
\item True
\item False
\end{optlist}
\begin{explainbox}{}
\textbf{Correct Answer: True}\\[0.3em]
\textbf{Concept:} This statement correctly describes VRRP backup.
\end{explainbox}
\end{questioncard}

\begin{questioncard}{16}{tfcol}
\qtypebadge{tfcol}{True / False} \hfill \textcolor{darkslate!70}{\small Topic: Firewall Hot Standby}

VRRP backup is unrelated to takes over if the master fails.

\begin{optlist}
\item True
\item False
\end{optlist}
\begin{explainbox}{}
\textbf{Correct Answer: False}\\[0.3em]
\textbf{Concept:} VRRP backup is directly related to takes over if the master fails.
\end{explainbox}
\end{questioncard}

\begin{questioncard}{17}{tfcol}
\qtypebadge{tfcol}{True / False} \hfill \textcolor{darkslate!70}{\small Topic: Firewall Hot Standby}

VGMP means that manages VRRP group states consistently.

\begin{optlist}
\item True
\item False
\end{optlist}
\begin{explainbox}{}
\textbf{Correct Answer: True}\\[0.3em]
\textbf{Concept:} This statement correctly describes VGMP.
\end{explainbox}
\end{questioncard}

\begin{questioncard}{18}{tfcol}
\qtypebadge{tfcol}{True / False} \hfill \textcolor{darkslate!70}{\small Topic: Firewall Hot Standby}

VGMP is unrelated to manages VRRP group states consistently.

\begin{optlist}
\item True
\item False
\end{optlist}
\begin{explainbox}{}
\textbf{Correct Answer: False}\\[0.3em]
\textbf{Concept:} VGMP is directly related to manages VRRP group states consistently.
\end{explainbox}
\end{questioncard}

\begin{questioncard}{19}{tfcol}
\qtypebadge{tfcol}{True / False} \hfill \textcolor{darkslate!70}{\small Topic: Firewall Hot Standby}

HRP means that synchronizes sessions and configurations.

\begin{optlist}
\item True
\item False
\end{optlist}
\begin{explainbox}{}
\textbf{Correct Answer: True}\\[0.3em]
\textbf{Concept:} This statement correctly describes HRP.
\end{explainbox}
\end{questioncard}

\begin{questioncard}{20}{tfcol}
\qtypebadge{tfcol}{True / False} \hfill \textcolor{darkslate!70}{\small Topic: Firewall Hot Standby}

HRP is unrelated to synchronizes sessions and configurations.

\begin{optlist}
\item True
\item False
\end{optlist}
\begin{explainbox}{}
\textbf{Correct Answer: False}\\[0.3em]
\textbf{Concept:} HRP is directly related to synchronizes sessions and configurations.
\end{explainbox}
\end{questioncard}

\begin{questioncard}{21}{tfcol}
\qtypebadge{tfcol}{True / False} \hfill \textcolor{darkslate!70}{\small Topic: Firewall Hot Standby}

Active/standby mode means that has one firewall forwarding traffic.

\begin{optlist}
\item True
\item False
\end{optlist}
\begin{explainbox}{}
\textbf{Correct Answer: True}\\[0.3em]
\textbf{Concept:} This statement correctly describes Active/standby mode.
\end{explainbox}
\end{questioncard}

\begin{questioncard}{22}{tfcol}
\qtypebadge{tfcol}{True / False} \hfill \textcolor{darkslate!70}{\small Topic: Firewall Hot Standby}

Active/standby mode is unrelated to has one firewall forwarding traffic.

\begin{optlist}
\item True
\item False
\end{optlist}
\begin{explainbox}{}
\textbf{Correct Answer: False}\\[0.3em]
\textbf{Concept:} Active/standby mode is directly related to has one firewall forwarding traffic.
\end{explainbox}
\end{questioncard}

\begin{questioncard}{23}{tfcol}
\qtypebadge{tfcol}{True / False} \hfill \textcolor{darkslate!70}{\small Topic: Firewall Hot Standby}

Load-sharing mode means that has multiple firewalls forwarding traffic.

\begin{optlist}
\item True
\item False
\end{optlist}
\begin{explainbox}{}
\textbf{Correct Answer: True}\\[0.3em]
\textbf{Concept:} This statement correctly describes Load-sharing mode.
\end{explainbox}
\end{questioncard}

\begin{questioncard}{24}{tfcol}
\qtypebadge{tfcol}{True / False} \hfill \textcolor{darkslate!70}{\small Topic: Firewall Hot Standby}

Load-sharing mode is unrelated to has multiple firewalls forwarding traffic.

\begin{optlist}
\item True
\item False
\end{optlist}
\begin{explainbox}{}
\textbf{Correct Answer: False}\\[0.3em]
\textbf{Concept:} Load-sharing mode is directly related to has multiple firewalls forwarding traffic.
\end{explainbox}
\end{questioncard}

\begin{questioncard}{25}{tfcol}
\qtypebadge{tfcol}{True / False} \hfill \textcolor{darkslate!70}{\small Topic: Firewall Hot Standby}

Heartbeat link means that should be dedicated and reliable.

\begin{optlist}
\item True
\item False
\end{optlist}
\begin{explainbox}{}
\textbf{Correct Answer: True}\\[0.3em]
\textbf{Concept:} This statement correctly describes Heartbeat link.
\end{explainbox}
\end{questioncard}

\begin{questioncard}{26}{tfcol}
\qtypebadge{tfcol}{True / False} \hfill \textcolor{darkslate!70}{\small Topic: Firewall Hot Standby}

Heartbeat link is unrelated to should be dedicated and reliable.

\begin{optlist}
\item True
\item False
\end{optlist}
\begin{explainbox}{}
\textbf{Correct Answer: False}\\[0.3em]
\textbf{Concept:} Heartbeat link is directly related to should be dedicated and reliable.
\end{explainbox}
\end{questioncard}

\begin{questioncard}{27}{tfcol}
\qtypebadge{tfcol}{True / False} \hfill \textcolor{darkslate!70}{\small Topic: Firewall Hot Standby}

HRP session backup means that ensures connections survive failover.

\begin{optlist}
\item True
\item False
\end{optlist}
\begin{explainbox}{}
\textbf{Correct Answer: True}\\[0.3em]
\textbf{Concept:} This statement correctly describes HRP session backup.
\end{explainbox}
\end{questioncard}

\begin{questioncard}{28}{tfcol}
\qtypebadge{tfcol}{True / False} \hfill \textcolor{darkslate!70}{\small Topic: Firewall Hot Standby}

HRP session backup is unrelated to ensures connections survive failover.

\begin{optlist}
\item True
\item False
\end{optlist}
\begin{explainbox}{}
\textbf{Correct Answer: False}\\[0.3em]
\textbf{Concept:} HRP session backup is directly related to ensures connections survive failover.
\end{explainbox}
\end{questioncard}

\begin{questioncard}{29}{tfcol}
\qtypebadge{tfcol}{True / False} \hfill \textcolor{darkslate!70}{\small Topic: Firewall Hot Standby}

Preempt delay means that reduces rapid master-backup flapping.

\begin{optlist}
\item True
\item False
\end{optlist}
\begin{explainbox}{}
\textbf{Correct Answer: True}\\[0.3em]
\textbf{Concept:} This statement correctly describes Preempt delay.
\end{explainbox}
\end{questioncard}

\begin{questioncard}{30}{tfcol}
\qtypebadge{tfcol}{True / False} \hfill \textcolor{darkslate!70}{\small Topic: Firewall Hot Standby}

Preempt delay is unrelated to reduces rapid master-backup flapping.

\begin{optlist}
\item True
\item False
\end{optlist}
\begin{explainbox}{}
\textbf{Correct Answer: False}\\[0.3em]
\textbf{Concept:} Preempt delay is directly related to reduces rapid master-backup flapping.
\end{explainbox}
\end{questioncard}

\section{VRRP}
\begin{questioncard}{1}{tfcol}
\qtypebadge{tfcol}{True / False} \hfill \textcolor{darkslate!70}{\small Topic: VRRP}

VRRP uses a virtual IP address shared between multiple physical routers.

\begin{optlist}
\item True
\item False
\end{optlist}
\begin{explainbox}{}
\textbf{Correct Answer: True}\\[0.3em]
\textbf{Concept:} VRRP allows multiple routers to present a single virtual gateway IP to hosts.
\end{explainbox}
\end{questioncard}

\section{VGMP}
\begin{questioncard}{1}{tfcol}
\qtypebadge{tfcol}{True / False} \hfill \textcolor{darkslate!70}{\small Topic: VGMP}

VGMP ensures that VRRP groups on a firewall switch state consistently.

\begin{optlist}
\item True
\item False
\end{optlist}
\begin{explainbox}{}
\textbf{Correct Answer: True}\\[0.3em]
\textbf{Concept:} VGMP prevents split-brain scenarios by managing all VRRP groups as a single entity.
\end{explainbox}
\end{questioncard}

\section{HRP}
\begin{questioncard}{1}{tfcol}
\qtypebadge{tfcol}{True / False} \hfill \textcolor{darkslate!70}{\small Topic: HRP}

HRP only synchronizes configuration files, not session tables.

\begin{optlist}
\item True
\item False
\end{optlist}
\begin{explainbox}{}
\textbf{Correct Answer: False}\\[0.3em]
\textbf{Concept:} HRP synchronizes both configurations and runtime state such as session tables.
\end{explainbox}
\end{questioncard}

\section{Intrusion Prevention and Antivirus}
\begin{questioncard}{1}{tfcol}
\qtypebadge{tfcol}{True / False} \hfill \textcolor{darkslate!70}{\small Topic: Intrusion Prevention and Antivirus}

IPS is typically deployed inline to block threats.

\begin{optlist}
\item True
\item False
\end{optlist}
\begin{explainbox}{}
\textbf{Correct Answer: True}\\[0.3em]
\textbf{Concept:} Inline deployment allows IPS to actively block malicious traffic.
\end{explainbox}
\end{questioncard}

\begin{questioncard}{2}{tfcol}
\qtypebadge{tfcol}{True / False} \hfill \textcolor{darkslate!70}{\small Topic: Intrusion Prevention and Antivirus}

Anomaly-based detection requires a database of known attack signatures.

\begin{optlist}
\item True
\item False
\end{optlist}
\begin{explainbox}{}
\textbf{Correct Answer: False}\\[0.3em]
\textbf{Concept:} Anomaly detection uses baselines of normal behavior, not attack signatures.
\end{explainbox}
\end{questioncard}

\begin{questioncard}{3}{tfcol}
\qtypebadge{tfcol}{True / False} \hfill \textcolor{darkslate!70}{\small Topic: Intrusion Prevention and Antivirus}

Heuristic analysis can detect previously unknown malware.

\begin{optlist}
\item True
\item False
\end{optlist}
\begin{explainbox}{}
\textbf{Correct Answer: True}\\[0.3em]
\textbf{Concept:} Heuristics identify suspicious behavior even without a known signature.
\end{explainbox}
\end{questioncard}

\begin{questioncard}{4}{tfcol}
\qtypebadge{tfcol}{True / False} \hfill \textcolor{darkslate!70}{\small Topic: Intrusion Prevention and Antivirus}

Signature-based antivirus is effective against all zero-day threats.

\begin{optlist}
\item True
\item False
\end{optlist}
\begin{explainbox}{}
\textbf{Correct Answer: False}\\[0.3em]
\textbf{Concept:} Zero-day threats lack signatures, so signature-only detection cannot catch them.
\end{explainbox}
\end{questioncard}

\begin{questioncard}{5}{tfcol}
\qtypebadge{tfcol}{True / False} \hfill \textcolor{darkslate!70}{\small Topic: Intrusion Prevention and Antivirus}

An IDS can block malicious traffic inline.

\begin{optlist}
\item True
\item False
\end{optlist}
\begin{explainbox}{}
\textbf{Correct Answer: False}\\[0.3em]
\textbf{Concept:} IDS typically detects and alerts; IPS blocks inline.
\end{explainbox}
\end{questioncard}

\begin{questioncard}{6}{tfcol}
\qtypebadge{tfcol}{True / False} \hfill \textcolor{darkslate!70}{\small Topic: Intrusion Prevention and Antivirus}

Heuristic detection can identify unknown threats.

\begin{optlist}
\item True
\item False
\end{optlist}
\begin{explainbox}{}
\textbf{Correct Answer: True}\\[0.3em]
\textbf{Concept:} Heuristics analyze behavior to find unknown malware.
\end{explainbox}
\end{questioncard}

\begin{questioncard}{7}{tfcol}
\qtypebadge{tfcol}{True / False} \hfill \textcolor{darkslate!70}{\small Topic: Intrusion Prevention and Antivirus}

Sandboxing executes files in an isolated environment.

\begin{optlist}
\item True
\item False
\end{optlist}
\begin{explainbox}{}
\textbf{Correct Answer: True}\\[0.3em]
\textbf{Concept:} Sandboxing observes behavior without risking production systems.
\end{explainbox}
\end{questioncard}

\begin{questioncard}{8}{tfcol}
\qtypebadge{tfcol}{True / False} \hfill \textcolor{darkslate!70}{\small Topic: Intrusion Prevention and Antivirus}

IDS means that detects intrusions and generates alerts.

\begin{optlist}
\item True
\item False
\end{optlist}
\begin{explainbox}{}
\textbf{Correct Answer: True}\\[0.3em]
\textbf{Concept:} This statement correctly describes IDS.
\end{explainbox}
\end{questioncard}

\begin{questioncard}{9}{tfcol}
\qtypebadge{tfcol}{True / False} \hfill \textcolor{darkslate!70}{\small Topic: Intrusion Prevention and Antivirus}

IDS is unrelated to detects intrusions and generates alerts.

\begin{optlist}
\item True
\item False
\end{optlist}
\begin{explainbox}{}
\textbf{Correct Answer: False}\\[0.3em]
\textbf{Concept:} IDS is directly related to detects intrusions and generates alerts.
\end{explainbox}
\end{questioncard}

\begin{questioncard}{10}{tfcol}
\qtypebadge{tfcol}{True / False} \hfill \textcolor{darkslate!70}{\small Topic: Intrusion Prevention and Antivirus}

IPS means that detects and blocks intrusions inline.

\begin{optlist}
\item True
\item False
\end{optlist}
\begin{explainbox}{}
\textbf{Correct Answer: True}\\[0.3em]
\textbf{Concept:} This statement correctly describes IPS.
\end{explainbox}
\end{questioncard}

\begin{questioncard}{11}{tfcol}
\qtypebadge{tfcol}{True / False} \hfill \textcolor{darkslate!70}{\small Topic: Intrusion Prevention and Antivirus}

IPS is unrelated to detects and blocks intrusions inline.

\begin{optlist}
\item True
\item False
\end{optlist}
\begin{explainbox}{}
\textbf{Correct Answer: False}\\[0.3em]
\textbf{Concept:} IPS is directly related to detects and blocks intrusions inline.
\end{explainbox}
\end{questioncard}

\begin{questioncard}{12}{tfcol}
\qtypebadge{tfcol}{True / False} \hfill \textcolor{darkslate!70}{\small Topic: Intrusion Prevention and Antivirus}

Signature detection means that matches known attack patterns.

\begin{optlist}
\item True
\item False
\end{optlist}
\begin{explainbox}{}
\textbf{Correct Answer: True}\\[0.3em]
\textbf{Concept:} This statement correctly describes Signature detection.
\end{explainbox}
\end{questioncard}

\begin{questioncard}{13}{tfcol}
\qtypebadge{tfcol}{True / False} \hfill \textcolor{darkslate!70}{\small Topic: Intrusion Prevention and Antivirus}

Signature detection is unrelated to matches known attack patterns.

\begin{optlist}
\item True
\item False
\end{optlist}
\begin{explainbox}{}
\textbf{Correct Answer: False}\\[0.3em]
\textbf{Concept:} Signature detection is directly related to matches known attack patterns.
\end{explainbox}
\end{questioncard}

\begin{questioncard}{14}{tfcol}
\qtypebadge{tfcol}{True / False} \hfill \textcolor{darkslate!70}{\small Topic: Intrusion Prevention and Antivirus}

Anomaly detection means that flags deviations from normal behavior.

\begin{optlist}
\item True
\item False
\end{optlist}
\begin{explainbox}{}
\textbf{Correct Answer: True}\\[0.3em]
\textbf{Concept:} This statement correctly describes Anomaly detection.
\end{explainbox}
\end{questioncard}

\begin{questioncard}{15}{tfcol}
\qtypebadge{tfcol}{True / False} \hfill \textcolor{darkslate!70}{\small Topic: Intrusion Prevention and Antivirus}

Anomaly detection is unrelated to flags deviations from normal behavior.

\begin{optlist}
\item True
\item False
\end{optlist}
\begin{explainbox}{}
\textbf{Correct Answer: False}\\[0.3em]
\textbf{Concept:} Anomaly detection is directly related to flags deviations from normal behavior.
\end{explainbox}
\end{questioncard}

\begin{questioncard}{16}{tfcol}
\qtypebadge{tfcol}{True / False} \hfill \textcolor{darkslate!70}{\small Topic: Intrusion Prevention and Antivirus}

False positive means that incorrectly blocks legitimate traffic.

\begin{optlist}
\item True
\item False
\end{optlist}
\begin{explainbox}{}
\textbf{Correct Answer: True}\\[0.3em]
\textbf{Concept:} This statement correctly describes False positive.
\end{explainbox}
\end{questioncard}

\begin{questioncard}{17}{tfcol}
\qtypebadge{tfcol}{True / False} \hfill \textcolor{darkslate!70}{\small Topic: Intrusion Prevention and Antivirus}

False positive is unrelated to incorrectly blocks legitimate traffic.

\begin{optlist}
\item True
\item False
\end{optlist}
\begin{explainbox}{}
\textbf{Correct Answer: False}\\[0.3em]
\textbf{Concept:} False positive is directly related to incorrectly blocks legitimate traffic.
\end{explainbox}
\end{questioncard}

\begin{questioncard}{18}{tfcol}
\qtypebadge{tfcol}{True / False} \hfill \textcolor{darkslate!70}{\small Topic: Intrusion Prevention and Antivirus}

False negative means that misses an actual attack.

\begin{optlist}
\item True
\item False
\end{optlist}
\begin{explainbox}{}
\textbf{Correct Answer: True}\\[0.3em]
\textbf{Concept:} This statement correctly describes False negative.
\end{explainbox}
\end{questioncard}

\begin{questioncard}{19}{tfcol}
\qtypebadge{tfcol}{True / False} \hfill \textcolor{darkslate!70}{\small Topic: Intrusion Prevention and Antivirus}

False negative is unrelated to misses an actual attack.

\begin{optlist}
\item True
\item False
\end{optlist}
\begin{explainbox}{}
\textbf{Correct Answer: False}\\[0.3em]
\textbf{Concept:} False negative is directly related to misses an actual attack.
\end{explainbox}
\end{questioncard}

\begin{questioncard}{20}{tfcol}
\qtypebadge{tfcol}{True / False} \hfill \textcolor{darkslate!70}{\small Topic: Intrusion Prevention and Antivirus}

Sandboxing means that observes file behavior in an isolated environment.

\begin{optlist}
\item True
\item False
\end{optlist}
\begin{explainbox}{}
\textbf{Correct Answer: True}\\[0.3em]
\textbf{Concept:} This statement correctly describes Sandboxing.
\end{explainbox}
\end{questioncard}

\begin{questioncard}{21}{tfcol}
\qtypebadge{tfcol}{True / False} \hfill \textcolor{darkslate!70}{\small Topic: Intrusion Prevention and Antivirus}

Sandboxing is unrelated to observes file behavior in an isolated environment.

\begin{optlist}
\item True
\item False
\end{optlist}
\begin{explainbox}{}
\textbf{Correct Answer: False}\\[0.3em]
\textbf{Concept:} Sandboxing is directly related to observes file behavior in an isolated environment.
\end{explainbox}
\end{questioncard}

\begin{questioncard}{22}{tfcol}
\qtypebadge{tfcol}{True / False} \hfill \textcolor{darkslate!70}{\small Topic: Intrusion Prevention and Antivirus}

Heuristic analysis means that detects malware by examining behavior.

\begin{optlist}
\item True
\item False
\end{optlist}
\begin{explainbox}{}
\textbf{Correct Answer: True}\\[0.3em]
\textbf{Concept:} This statement correctly describes Heuristic analysis.
\end{explainbox}
\end{questioncard}

\begin{questioncard}{23}{tfcol}
\qtypebadge{tfcol}{True / False} \hfill \textcolor{darkslate!70}{\small Topic: Intrusion Prevention and Antivirus}

Heuristic analysis is unrelated to detects malware by examining behavior.

\begin{optlist}
\item True
\item False
\end{optlist}
\begin{explainbox}{}
\textbf{Correct Answer: False}\\[0.3em]
\textbf{Concept:} Heuristic analysis is directly related to detects malware by examining behavior.
\end{explainbox}
\end{questioncard}

\begin{questioncard}{24}{tfcol}
\qtypebadge{tfcol}{True / False} \hfill \textcolor{darkslate!70}{\small Topic: Intrusion Prevention and Antivirus}

Whitelisting means that allows only approved applications.

\begin{optlist}
\item True
\item False
\end{optlist}
\begin{explainbox}{}
\textbf{Correct Answer: True}\\[0.3em]
\textbf{Concept:} This statement correctly describes Whitelisting.
\end{explainbox}
\end{questioncard}

\begin{questioncard}{25}{tfcol}
\qtypebadge{tfcol}{True / False} \hfill \textcolor{darkslate!70}{\small Topic: Intrusion Prevention and Antivirus}

Whitelisting is unrelated to allows only approved applications.

\begin{optlist}
\item True
\item False
\end{optlist}
\begin{explainbox}{}
\textbf{Correct Answer: False}\\[0.3em]
\textbf{Concept:} Whitelisting is directly related to allows only approved applications.
\end{explainbox}
\end{questioncard}

\begin{questioncard}{26}{tfcol}
\qtypebadge{tfcol}{True / False} \hfill \textcolor{darkslate!70}{\small Topic: Intrusion Prevention and Antivirus}

Antivirus signatures means that must be updated regularly.

\begin{optlist}
\item True
\item False
\end{optlist}
\begin{explainbox}{}
\textbf{Correct Answer: True}\\[0.3em]
\textbf{Concept:} This statement correctly describes Antivirus signatures.
\end{explainbox}
\end{questioncard}

\begin{questioncard}{27}{tfcol}
\qtypebadge{tfcol}{True / False} \hfill \textcolor{darkslate!70}{\small Topic: Intrusion Prevention and Antivirus}

Antivirus signatures is unrelated to must be updated regularly.

\begin{optlist}
\item True
\item False
\end{optlist}
\begin{explainbox}{}
\textbf{Correct Answer: False}\\[0.3em]
\textbf{Concept:} Antivirus signatures is directly related to must be updated regularly.
\end{explainbox}
\end{questioncard}

\begin{questioncard}{28}{tfcol}
\qtypebadge{tfcol}{True / False} \hfill \textcolor{darkslate!70}{\small Topic: Intrusion Prevention and Antivirus}

NIDS means that monitors network traffic for intrusions.

\begin{optlist}
\item True
\item False
\end{optlist}
\begin{explainbox}{}
\textbf{Correct Answer: True}\\[0.3em]
\textbf{Concept:} This statement correctly describes NIDS.
\end{explainbox}
\end{questioncard}

\begin{questioncard}{29}{tfcol}
\qtypebadge{tfcol}{True / False} \hfill \textcolor{darkslate!70}{\small Topic: Intrusion Prevention and Antivirus}

NIDS is unrelated to monitors network traffic for intrusions.

\begin{optlist}
\item True
\item False
\end{optlist}
\begin{explainbox}{}
\textbf{Correct Answer: False}\\[0.3em]
\textbf{Concept:} NIDS is directly related to monitors network traffic for intrusions.
\end{explainbox}
\end{questioncard}

\begin{questioncard}{30}{tfcol}
\qtypebadge{tfcol}{True / False} \hfill \textcolor{darkslate!70}{\small Topic: Intrusion Prevention and Antivirus}

HIDS means that monitors host activity for intrusions.

\begin{optlist}
\item True
\item False
\end{optlist}
\begin{explainbox}{}
\textbf{Correct Answer: True}\\[0.3em]
\textbf{Concept:} This statement correctly describes HIDS.
\end{explainbox}
\end{questioncard}

\begin{questioncard}{31}{tfcol}
\qtypebadge{tfcol}{True / False} \hfill \textcolor{darkslate!70}{\small Topic: Intrusion Prevention and Antivirus}

HIDS is unrelated to monitors host activity for intrusions.

\begin{optlist}
\item True
\item False
\end{optlist}
\begin{explainbox}{}
\textbf{Correct Answer: False}\\[0.3em]
\textbf{Concept:} HIDS is directly related to monitors host activity for intrusions.
\end{explainbox}
\end{questioncard}

\begin{questioncard}{32}{tfcol}
\qtypebadge{tfcol}{True / False} \hfill \textcolor{darkslate!70}{\small Topic: Intrusion Prevention and Antivirus}

Reconnaissance means that is an early attack phase.

\begin{optlist}
\item True
\item False
\end{optlist}
\begin{explainbox}{}
\textbf{Correct Answer: True}\\[0.3em]
\textbf{Concept:} This statement correctly describes Reconnaissance.
\end{explainbox}
\end{questioncard}

\begin{questioncard}{33}{tfcol}
\qtypebadge{tfcol}{True / False} \hfill \textcolor{darkslate!70}{\small Topic: Intrusion Prevention and Antivirus}

Reconnaissance is unrelated to is an early attack phase.

\begin{optlist}
\item True
\item False
\end{optlist}
\begin{explainbox}{}
\textbf{Correct Answer: False}\\[0.3em]
\textbf{Concept:} Reconnaissance is directly related to is an early attack phase.
\end{explainbox}
\end{questioncard}

\begin{questioncard}{34}{tfcol}
\qtypebadge{tfcol}{True / False} \hfill \textcolor{darkslate!70}{\small Topic: Intrusion Prevention and Antivirus}

Exploitation means that gains unauthorized access using vulnerabilities.

\begin{optlist}
\item True
\item False
\end{optlist}
\begin{explainbox}{}
\textbf{Correct Answer: True}\\[0.3em]
\textbf{Concept:} This statement correctly describes Exploitation.
\end{explainbox}
\end{questioncard}

\begin{questioncard}{35}{tfcol}
\qtypebadge{tfcol}{True / False} \hfill \textcolor{darkslate!70}{\small Topic: Intrusion Prevention and Antivirus}

Exploitation is unrelated to gains unauthorized access using vulnerabilities.

\begin{optlist}
\item True
\item False
\end{optlist}
\begin{explainbox}{}
\textbf{Correct Answer: False}\\[0.3em]
\textbf{Concept:} Exploitation is directly related to gains unauthorized access using vulnerabilities.
\end{explainbox}
\end{questioncard}

\section{Intrusion Prevention}
\begin{questioncard}{1}{tfcol}
\qtypebadge{tfcol}{True / False} \hfill \textcolor{darkslate!70}{\small Topic: Intrusion Prevention}

An IDS is typically deployed inline and can block attacks.

\begin{optlist}
\item True
\item False
\end{optlist}
\begin{explainbox}{}
\textbf{Correct Answer: False}\\[0.3em]
\textbf{Concept:} IDS is usually passive/out-of-band and only detects and alerts; IPS blocks inline.
\end{explainbox}
\end{questioncard}

\begin{questioncard}{2}{tfcol}
\qtypebadge{tfcol}{True / False} \hfill \textcolor{darkslate!70}{\small Topic: Intrusion Prevention}

Signature-based IPS can detect zero-day attacks effectively.

\begin{optlist}
\item True
\item False
\end{optlist}
\begin{explainbox}{}
\textbf{Correct Answer: False}\\[0.3em]
\textbf{Concept:} Signatures require known attack patterns; zero-day attacks are unknown and often missed.
\end{explainbox}
\end{questioncard}

\section{Antivirus}
\begin{questioncard}{1}{tfcol}
\qtypebadge{tfcol}{True / False} \hfill \textcolor{darkslate!70}{\small Topic: Antivirus}

Antivirus signature databases must be updated regularly to detect new malware.

\begin{optlist}
\item True
\item False
\end{optlist}
\begin{explainbox}{}
\textbf{Correct Answer: True}\\[0.3em]
\textbf{Concept:} New malware variants require updated signatures and detection engines.
\end{explainbox}
\end{questioncard}

\begin{questioncard}{2}{tfcol}
\qtypebadge{tfcol}{True / False} \hfill \textcolor{darkslate!70}{\small Topic: Antivirus}

Sandboxing can detect unknown malware by observing its runtime behavior.

\begin{optlist}
\item True
\item False
\end{optlist}
\begin{explainbox}{}
\textbf{Correct Answer: True}\\[0.3em]
\textbf{Concept:} Sandboxing reveals malicious behavior of unknown samples in an isolated environment.
\end{explainbox}
\end{questioncard}

\section{AAA and User Authentication}
\begin{questioncard}{1}{tfcol}
\qtypebadge{tfcol}{True / False} \hfill \textcolor{darkslate!70}{\small Topic: AAA and User Authentication}

Authentication verifies who a user is.

\begin{optlist}
\item True
\item False
\end{optlist}
\begin{explainbox}{}
\textbf{Correct Answer: True}\\[0.3em]
\textbf{Concept:} Authentication confirms identity before granting access.
\end{explainbox}
\end{questioncard}

\begin{questioncard}{2}{tfcol}
\qtypebadge{tfcol}{True / False} \hfill \textcolor{darkslate!70}{\small Topic: AAA and User Authentication}

TACACS+ separates authentication, authorization, and accounting.

\begin{optlist}
\item True
\item False
\end{optlist}
\begin{explainbox}{}
\textbf{Correct Answer: True}\\[0.3em]
\textbf{Concept:} TACACS+ provides independent AAA functions.
\end{explainbox}
\end{questioncard}

\begin{questioncard}{3}{tfcol}
\qtypebadge{tfcol}{True / False} \hfill \textcolor{darkslate!70}{\small Topic: AAA and User Authentication}

RADIUS encrypts the entire authentication request.

\begin{optlist}
\item True
\item False
\end{optlist}
\begin{explainbox}{}
\textbf{Correct Answer: False}\\[0.3em]
\textbf{Concept:} RADIUS encrypts only the password attribute.
\end{explainbox}
\end{questioncard}

\begin{questioncard}{4}{tfcol}
\qtypebadge{tfcol}{True / False} \hfill \textcolor{darkslate!70}{\small Topic: AAA and User Authentication}

Local authentication does not require an external server.

\begin{optlist}
\item True
\item False
\end{optlist}
\begin{explainbox}{}
\textbf{Correct Answer: True}\\[0.3em]
\textbf{Concept:} Local authentication uses the firewall own user database.
\end{explainbox}
\end{questioncard}

\begin{questioncard}{5}{tfcol}
\qtypebadge{tfcol}{True / False} \hfill \textcolor{darkslate!70}{\small Topic: AAA and User Authentication}

Authorization happens after authentication.

\begin{optlist}
\item True
\item False
\end{optlist}
\begin{explainbox}{}
\textbf{Correct Answer: True}\\[0.3em]
\textbf{Concept:} Authorization defines what an authenticated user may do.
\end{explainbox}
\end{questioncard}

\begin{questioncard}{6}{tfcol}
\qtypebadge{tfcol}{True / False} \hfill \textcolor{darkslate!70}{\small Topic: AAA and User Authentication}

TACACS+ uses UDP.

\begin{optlist}
\item True
\item False
\end{optlist}
\begin{explainbox}{}
\textbf{Correct Answer: False}\\[0.3em]
\textbf{Concept:} TACACS+ uses TCP.
\end{explainbox}
\end{questioncard}

\begin{questioncard}{7}{tfcol}
\qtypebadge{tfcol}{True / False} \hfill \textcolor{darkslate!70}{\small Topic: AAA and User Authentication}

Portal authentication requires a web browser.

\begin{optlist}
\item True
\item False
\end{optlist}
\begin{explainbox}{}
\textbf{Correct Answer: True}\\[0.3em]
\textbf{Concept:} Portal presents a web login page to users.
\end{explainbox}
\end{questioncard}

\begin{questioncard}{8}{tfcol}
\qtypebadge{tfcol}{True / False} \hfill \textcolor{darkslate!70}{\small Topic: AAA and User Authentication}

Authentication means that verifies user identity.

\begin{optlist}
\item True
\item False
\end{optlist}
\begin{explainbox}{}
\textbf{Correct Answer: True}\\[0.3em]
\textbf{Concept:} This statement correctly describes Authentication.
\end{explainbox}
\end{questioncard}

\begin{questioncard}{9}{tfcol}
\qtypebadge{tfcol}{True / False} \hfill \textcolor{darkslate!70}{\small Topic: AAA and User Authentication}

Authentication is unrelated to verifies user identity.

\begin{optlist}
\item True
\item False
\end{optlist}
\begin{explainbox}{}
\textbf{Correct Answer: False}\\[0.3em]
\textbf{Concept:} Authentication is directly related to verifies user identity.
\end{explainbox}
\end{questioncard}

\begin{questioncard}{10}{tfcol}
\qtypebadge{tfcol}{True / False} \hfill \textcolor{darkslate!70}{\small Topic: AAA and User Authentication}

Authorization means that determines what a user can access.

\begin{optlist}
\item True
\item False
\end{optlist}
\begin{explainbox}{}
\textbf{Correct Answer: True}\\[0.3em]
\textbf{Concept:} This statement correctly describes Authorization.
\end{explainbox}
\end{questioncard}

\begin{questioncard}{11}{tfcol}
\qtypebadge{tfcol}{True / False} \hfill \textcolor{darkslate!70}{\small Topic: AAA and User Authentication}

Authorization is unrelated to determines what a user can access.

\begin{optlist}
\item True
\item False
\end{optlist}
\begin{explainbox}{}
\textbf{Correct Answer: False}\\[0.3em]
\textbf{Concept:} Authorization is directly related to determines what a user can access.
\end{explainbox}
\end{questioncard}

\begin{questioncard}{12}{tfcol}
\qtypebadge{tfcol}{True / False} \hfill \textcolor{darkslate!70}{\small Topic: AAA and User Authentication}

Accounting means that records user activities and resource usage.

\begin{optlist}
\item True
\item False
\end{optlist}
\begin{explainbox}{}
\textbf{Correct Answer: True}\\[0.3em]
\textbf{Concept:} This statement correctly describes Accounting.
\end{explainbox}
\end{questioncard}

\begin{questioncard}{13}{tfcol}
\qtypebadge{tfcol}{True / False} \hfill \textcolor{darkslate!70}{\small Topic: AAA and User Authentication}

Accounting is unrelated to records user activities and resource usage.

\begin{optlist}
\item True
\item False
\end{optlist}
\begin{explainbox}{}
\textbf{Correct Answer: False}\\[0.3em]
\textbf{Concept:} Accounting is directly related to records user activities and resource usage.
\end{explainbox}
\end{questioncard}

\begin{questioncard}{14}{tfcol}
\qtypebadge{tfcol}{True / False} \hfill \textcolor{darkslate!70}{\small Topic: AAA and User Authentication}

RADIUS means that uses UDP ports 1812 and 1813.

\begin{optlist}
\item True
\item False
\end{optlist}
\begin{explainbox}{}
\textbf{Correct Answer: True}\\[0.3em]
\textbf{Concept:} This statement correctly describes RADIUS.
\end{explainbox}
\end{questioncard}

\begin{questioncard}{15}{tfcol}
\qtypebadge{tfcol}{True / False} \hfill \textcolor{darkslate!70}{\small Topic: AAA and User Authentication}

RADIUS is unrelated to uses UDP ports 1812 and 1813.

\begin{optlist}
\item True
\item False
\end{optlist}
\begin{explainbox}{}
\textbf{Correct Answer: False}\\[0.3em]
\textbf{Concept:} RADIUS is directly related to uses UDP ports 1812 and 1813.
\end{explainbox}
\end{questioncard}

\begin{questioncard}{16}{tfcol}
\qtypebadge{tfcol}{True / False} \hfill \textcolor{darkslate!70}{\small Topic: AAA and User Authentication}

TACACS+ means that uses TCP port 49.

\begin{optlist}
\item True
\item False
\end{optlist}
\begin{explainbox}{}
\textbf{Correct Answer: True}\\[0.3em]
\textbf{Concept:} This statement correctly describes TACACS+.
\end{explainbox}
\end{questioncard}

\begin{questioncard}{17}{tfcol}
\qtypebadge{tfcol}{True / False} \hfill \textcolor{darkslate!70}{\small Topic: AAA and User Authentication}

TACACS+ is unrelated to uses TCP port 49.

\begin{optlist}
\item True
\item False
\end{optlist}
\begin{explainbox}{}
\textbf{Correct Answer: False}\\[0.3em]
\textbf{Concept:} TACACS+ is directly related to uses TCP port 49.
\end{explainbox}
\end{questioncard}

\begin{questioncard}{18}{tfcol}
\qtypebadge{tfcol}{True / False} \hfill \textcolor{darkslate!70}{\small Topic: AAA and User Authentication}

Portal authentication means that uses a web login page.

\begin{optlist}
\item True
\item False
\end{optlist}
\begin{explainbox}{}
\textbf{Correct Answer: True}\\[0.3em]
\textbf{Concept:} This statement correctly describes Portal authentication.
\end{explainbox}
\end{questioncard}

\begin{questioncard}{19}{tfcol}
\qtypebadge{tfcol}{True / False} \hfill \textcolor{darkslate!70}{\small Topic: AAA and User Authentication}

Portal authentication is unrelated to uses a web login page.

\begin{optlist}
\item True
\item False
\end{optlist}
\begin{explainbox}{}
\textbf{Correct Answer: False}\\[0.3em]
\textbf{Concept:} Portal authentication is directly related to uses a web login page.
\end{explainbox}
\end{questioncard}

\begin{questioncard}{20}{tfcol}
\qtypebadge{tfcol}{True / False} \hfill \textcolor{darkslate!70}{\small Topic: AAA and User Authentication}

802.1X means that provides port-based network access control.

\begin{optlist}
\item True
\item False
\end{optlist}
\begin{explainbox}{}
\textbf{Correct Answer: True}\\[0.3em]
\textbf{Concept:} This statement correctly describes 802.1X.
\end{explainbox}
\end{questioncard}

\begin{questioncard}{21}{tfcol}
\qtypebadge{tfcol}{True / False} \hfill \textcolor{darkslate!70}{\small Topic: AAA and User Authentication}

802.1X is unrelated to provides port-based network access control.

\begin{optlist}
\item True
\item False
\end{optlist}
\begin{explainbox}{}
\textbf{Correct Answer: False}\\[0.3em]
\textbf{Concept:} 802.1X is directly related to provides port-based network access control.
\end{explainbox}
\end{questioncard}

\begin{questioncard}{22}{tfcol}
\qtypebadge{tfcol}{True / False} \hfill \textcolor{darkslate!70}{\small Topic: AAA and User Authentication}

Local authentication means that uses the device user database.

\begin{optlist}
\item True
\item False
\end{optlist}
\begin{explainbox}{}
\textbf{Correct Answer: True}\\[0.3em]
\textbf{Concept:} This statement correctly describes Local authentication.
\end{explainbox}
\end{questioncard}

\begin{questioncard}{23}{tfcol}
\qtypebadge{tfcol}{True / False} \hfill \textcolor{darkslate!70}{\small Topic: AAA and User Authentication}

Local authentication is unrelated to uses the device user database.

\begin{optlist}
\item True
\item False
\end{optlist}
\begin{explainbox}{}
\textbf{Correct Answer: False}\\[0.3em]
\textbf{Concept:} Local authentication is directly related to uses the device user database.
\end{explainbox}
\end{questioncard}

\begin{questioncard}{24}{tfcol}
\qtypebadge{tfcol}{True / False} \hfill \textcolor{darkslate!70}{\small Topic: AAA and User Authentication}

LDAP means that queries a directory server for authentication.

\begin{optlist}
\item True
\item False
\end{optlist}
\begin{explainbox}{}
\textbf{Correct Answer: True}\\[0.3em]
\textbf{Concept:} This statement correctly describes LDAP.
\end{explainbox}
\end{questioncard}

\begin{questioncard}{25}{tfcol}
\qtypebadge{tfcol}{True / False} \hfill \textcolor{darkslate!70}{\small Topic: AAA and User Authentication}

LDAP is unrelated to queries a directory server for authentication.

\begin{optlist}
\item True
\item False
\end{optlist}
\begin{explainbox}{}
\textbf{Correct Answer: False}\\[0.3em]
\textbf{Concept:} LDAP is directly related to queries a directory server for authentication.
\end{explainbox}
\end{questioncard}

\begin{questioncard}{26}{tfcol}
\qtypebadge{tfcol}{True / False} \hfill \textcolor{darkslate!70}{\small Topic: AAA and User Authentication}

User group means that applies policies to multiple users collectively.

\begin{optlist}
\item True
\item False
\end{optlist}
\begin{explainbox}{}
\textbf{Correct Answer: True}\\[0.3em]
\textbf{Concept:} This statement correctly describes User group.
\end{explainbox}
\end{questioncard}

\begin{questioncard}{27}{tfcol}
\qtypebadge{tfcol}{True / False} \hfill \textcolor{darkslate!70}{\small Topic: AAA and User Authentication}

User group is unrelated to applies policies to multiple users collectively.

\begin{optlist}
\item True
\item False
\end{optlist}
\begin{explainbox}{}
\textbf{Correct Answer: False}\\[0.3em]
\textbf{Concept:} User group is directly related to applies policies to multiple users collectively.
\end{explainbox}
\end{questioncard}

\begin{questioncard}{28}{tfcol}
\qtypebadge{tfcol}{True / False} \hfill \textcolor{darkslate!70}{\small Topic: AAA and User Authentication}

SSO means that allows one authentication for multiple services.

\begin{optlist}
\item True
\item False
\end{optlist}
\begin{explainbox}{}
\textbf{Correct Answer: True}\\[0.3em]
\textbf{Concept:} This statement correctly describes SSO.
\end{explainbox}
\end{questioncard}

\begin{questioncard}{29}{tfcol}
\qtypebadge{tfcol}{True / False} \hfill \textcolor{darkslate!70}{\small Topic: AAA and User Authentication}

SSO is unrelated to allows one authentication for multiple services.

\begin{optlist}
\item True
\item False
\end{optlist}
\begin{explainbox}{}
\textbf{Correct Answer: False}\\[0.3em]
\textbf{Concept:} SSO is directly related to allows one authentication for multiple services.
\end{explainbox}
\end{questioncard}

\section{AAA Basics}
\begin{questioncard}{1}{tfcol}
\qtypebadge{tfcol}{True / False} \hfill \textcolor{darkslate!70}{\small Topic: AAA Basics}

Authorization determines whether an authenticated user is allowed to perform an action.

\begin{optlist}
\item True
\item False
\end{optlist}
\begin{explainbox}{}
\textbf{Correct Answer: True}\\[0.3em]
\textbf{Concept:} Authorization enforces access control decisions based on identity and policy.
\end{explainbox}
\end{questioncard}

\begin{questioncard}{2}{tfcol}
\qtypebadge{tfcol}{True / False} \hfill \textcolor{darkslate!70}{\small Topic: AAA Basics}

Accounting in AAA records user activities and resource usage.

\begin{optlist}
\item True
\item False
\end{optlist}
\begin{explainbox}{}
\textbf{Correct Answer: True}\\[0.3em]
\textbf{Concept:} Accounting provides logs for auditing, billing, and forensic analysis.
\end{explainbox}
\end{questioncard}

\section{Firewall User Authentication}
\begin{questioncard}{1}{tfcol}
\qtypebadge{tfcol}{True / False} \hfill \textcolor{darkslate!70}{\small Topic: Firewall User Authentication}

Portal authentication is commonly used for guest access and web-based login.

\begin{optlist}
\item True
\item False
\end{optlist}
\begin{explainbox}{}
\textbf{Correct Answer: True}\\[0.3em]
\textbf{Concept:} Portal authentication provides a convenient web-based login experience.
\end{explainbox}
\end{questioncard}

\begin{questioncard}{2}{tfcol}
\qtypebadge{tfcol}{True / False} \hfill \textcolor{darkslate!70}{\small Topic: Firewall User Authentication}

RADIUS encrypts the entire authentication packet including the username.

\begin{optlist}
\item True
\item False
\end{optlist}
\begin{explainbox}{}
\textbf{Correct Answer: False}\\[0.3em]
\textbf{Concept:} RADIUS encrypts only the password field; the rest of the packet is not encrypted.
\end{explainbox}
\end{questioncard}

\section{Cryptography and PKI}
\begin{questioncard}{1}{tfcol}
\qtypebadge{tfcol}{True / False} \hfill \textcolor{darkslate!70}{\small Topic: Cryptography and PKI}

AES is a symmetric encryption algorithm.

\begin{optlist}
\item True
\item False
\end{optlist}
\begin{explainbox}{}
\textbf{Correct Answer: True}\\[0.3em]
\textbf{Concept:} AES uses the same key for encryption and decryption.
\end{explainbox}
\end{questioncard}

\begin{questioncard}{2}{tfcol}
\qtypebadge{tfcol}{True / False} \hfill \textcolor{darkslate!70}{\small Topic: Cryptography and PKI}

RSA is faster than AES for bulk data encryption.

\begin{optlist}
\item True
\item False
\end{optlist}
\begin{explainbox}{}
\textbf{Correct Answer: False}\\[0.3em]
\textbf{Concept:} RSA is slower and typically used for key exchange or signatures, not bulk encryption.
\end{explainbox}
\end{questioncard}

\begin{questioncard}{3}{tfcol}
\qtypebadge{tfcol}{True / False} \hfill \textcolor{darkslate!70}{\small Topic: Cryptography and PKI}

SHA-256 produces a 256-bit digest.

\begin{optlist}
\item True
\item False
\end{optlist}
\begin{explainbox}{}
\textbf{Correct Answer: True}\\[0.3em]
\textbf{Concept:} SHA-256 generates a fixed 256-bit output.
\end{explainbox}
\end{questioncard}

\begin{questioncard}{4}{tfcol}
\qtypebadge{tfcol}{True / False} \hfill \textcolor{darkslate!70}{\small Topic: Cryptography and PKI}

Hash functions are reversible.

\begin{optlist}
\item True
\item False
\end{optlist}
\begin{explainbox}{}
\textbf{Correct Answer: False}\\[0.3em]
\textbf{Concept:} Hash functions are designed to be one-way.
\end{explainbox}
\end{questioncard}

\begin{questioncard}{5}{tfcol}
\qtypebadge{tfcol}{True / False} \hfill \textcolor{darkslate!70}{\small Topic: Cryptography and PKI}

A digital signature provides non-repudiation.

\begin{optlist}
\item True
\item False
\end{optlist}
\begin{explainbox}{}
\textbf{Correct Answer: True}\\[0.3em]
\textbf{Concept:} Digital signatures prove the signer created the message and cannot easily deny it.
\end{explainbox}
\end{questioncard}

\begin{questioncard}{6}{tfcol}
\qtypebadge{tfcol}{True / False} \hfill \textcolor{darkslate!70}{\small Topic: Cryptography and PKI}

A root CA certificate is typically signed by another CA.

\begin{optlist}
\item True
\item False
\end{optlist}
\begin{explainbox}{}
\textbf{Correct Answer: False}\\[0.3em]
\textbf{Concept:} Root CA certificates are self-signed and act as trust anchors.
\end{explainbox}
\end{questioncard}

\begin{questioncard}{7}{tfcol}
\qtypebadge{tfcol}{True / False} \hfill \textcolor{darkslate!70}{\small Topic: Cryptography and PKI}

OCSP provides real-time certificate revocation checking.

\begin{optlist}
\item True
\item False
\end{optlist}
\begin{explainbox}{}
\textbf{Correct Answer: True}\\[0.3em]
\textbf{Concept:} OCSP allows clients to query revocation status online.
\end{explainbox}
\end{questioncard}

\begin{questioncard}{8}{tfcol}
\qtypebadge{tfcol}{True / False} \hfill \textcolor{darkslate!70}{\small Topic: Cryptography and PKI}

A hash function output can be reversed to find the input.

\begin{optlist}
\item True
\item False
\end{optlist}
\begin{explainbox}{}
\textbf{Correct Answer: False}\\[0.3em]
\textbf{Concept:} Hash functions are one-way.
\end{explainbox}
\end{questioncard}

\begin{questioncard}{9}{tfcol}
\qtypebadge{tfcol}{True / False} \hfill \textcolor{darkslate!70}{\small Topic: Cryptography and PKI}

A digital signature verifies integrity and authenticity.

\begin{optlist}
\item True
\item False
\end{optlist}
\begin{explainbox}{}
\textbf{Correct Answer: True}\\[0.3em]
\textbf{Concept:} Signatures prove the message came from the signer and was not altered.
\end{explainbox}
\end{questioncard}

\begin{questioncard}{10}{tfcol}
\qtypebadge{tfcol}{True / False} \hfill \textcolor{darkslate!70}{\small Topic: Cryptography and PKI}

RSA is a symmetric encryption algorithm.

\begin{optlist}
\item True
\item False
\end{optlist}
\begin{explainbox}{}
\textbf{Correct Answer: False}\\[0.3em]
\textbf{Concept:} RSA is asymmetric.
\end{explainbox}
\end{questioncard}

\begin{questioncard}{11}{tfcol}
\qtypebadge{tfcol}{True / False} \hfill \textcolor{darkslate!70}{\small Topic: Cryptography and PKI}

Symmetric encryption means that uses the same key for encryption and decryption.

\begin{optlist}
\item True
\item False
\end{optlist}
\begin{explainbox}{}
\textbf{Correct Answer: True}\\[0.3em]
\textbf{Concept:} This statement correctly describes Symmetric encryption.
\end{explainbox}
\end{questioncard}

\begin{questioncard}{12}{tfcol}
\qtypebadge{tfcol}{True / False} \hfill \textcolor{darkslate!70}{\small Topic: Cryptography and PKI}

Symmetric encryption is unrelated to uses the same key for encryption and decryption.

\begin{optlist}
\item True
\item False
\end{optlist}
\begin{explainbox}{}
\textbf{Correct Answer: False}\\[0.3em]
\textbf{Concept:} Symmetric encryption is directly related to uses the same key for encryption and decryption.
\end{explainbox}
\end{questioncard}

\begin{questioncard}{13}{tfcol}
\qtypebadge{tfcol}{True / False} \hfill \textcolor{darkslate!70}{\small Topic: Cryptography and PKI}

Asymmetric encryption means that uses public and private key pairs.

\begin{optlist}
\item True
\item False
\end{optlist}
\begin{explainbox}{}
\textbf{Correct Answer: True}\\[0.3em]
\textbf{Concept:} This statement correctly describes Asymmetric encryption.
\end{explainbox}
\end{questioncard}

\begin{questioncard}{14}{tfcol}
\qtypebadge{tfcol}{True / False} \hfill \textcolor{darkslate!70}{\small Topic: Cryptography and PKI}

Asymmetric encryption is unrelated to uses public and private key pairs.

\begin{optlist}
\item True
\item False
\end{optlist}
\begin{explainbox}{}
\textbf{Correct Answer: False}\\[0.3em]
\textbf{Concept:} Asymmetric encryption is directly related to uses public and private key pairs.
\end{explainbox}
\end{questioncard}

\begin{questioncard}{15}{tfcol}
\qtypebadge{tfcol}{True / False} \hfill \textcolor{darkslate!70}{\small Topic: Cryptography and PKI}

AES means that is a symmetric encryption standard.

\begin{optlist}
\item True
\item False
\end{optlist}
\begin{explainbox}{}
\textbf{Correct Answer: True}\\[0.3em]
\textbf{Concept:} This statement correctly describes AES.
\end{explainbox}
\end{questioncard}

\begin{questioncard}{16}{tfcol}
\qtypebadge{tfcol}{True / False} \hfill \textcolor{darkslate!70}{\small Topic: Cryptography and PKI}

AES is unrelated to is a symmetric encryption standard.

\begin{optlist}
\item True
\item False
\end{optlist}
\begin{explainbox}{}
\textbf{Correct Answer: False}\\[0.3em]
\textbf{Concept:} AES is directly related to is a symmetric encryption standard.
\end{explainbox}
\end{questioncard}

\begin{questioncard}{17}{tfcol}
\qtypebadge{tfcol}{True / False} \hfill \textcolor{darkslate!70}{\small Topic: Cryptography and PKI}

RSA means that is an asymmetric encryption algorithm.

\begin{optlist}
\item True
\item False
\end{optlist}
\begin{explainbox}{}
\textbf{Correct Answer: True}\\[0.3em]
\textbf{Concept:} This statement correctly describes RSA.
\end{explainbox}
\end{questioncard}

\begin{questioncard}{18}{tfcol}
\qtypebadge{tfcol}{True / False} \hfill \textcolor{darkslate!70}{\small Topic: Cryptography and PKI}

RSA is unrelated to is an asymmetric encryption algorithm.

\begin{optlist}
\item True
\item False
\end{optlist}
\begin{explainbox}{}
\textbf{Correct Answer: False}\\[0.3em]
\textbf{Concept:} RSA is directly related to is an asymmetric encryption algorithm.
\end{explainbox}
\end{questioncard}

\begin{questioncard}{19}{tfcol}
\qtypebadge{tfcol}{True / False} \hfill \textcolor{darkslate!70}{\small Topic: Cryptography and PKI}

DES means that uses a 56-bit effective key.

\begin{optlist}
\item True
\item False
\end{optlist}
\begin{explainbox}{}
\textbf{Correct Answer: True}\\[0.3em]
\textbf{Concept:} This statement correctly describes DES.
\end{explainbox}
\end{questioncard}

\begin{questioncard}{20}{tfcol}
\qtypebadge{tfcol}{True / False} \hfill \textcolor{darkslate!70}{\small Topic: Cryptography and PKI}

DES is unrelated to uses a 56-bit effective key.

\begin{optlist}
\item True
\item False
\end{optlist}
\begin{explainbox}{}
\textbf{Correct Answer: False}\\[0.3em]
\textbf{Concept:} DES is directly related to uses a 56-bit effective key.
\end{explainbox}
\end{questioncard}

\begin{questioncard}{21}{tfcol}
\qtypebadge{tfcol}{True / False} \hfill \textcolor{darkslate!70}{\small Topic: Cryptography and PKI}

Hash function means that produces a fixed-length digest.

\begin{optlist}
\item True
\item False
\end{optlist}
\begin{explainbox}{}
\textbf{Correct Answer: True}\\[0.3em]
\textbf{Concept:} This statement correctly describes Hash function.
\end{explainbox}
\end{questioncard}

\begin{questioncard}{22}{tfcol}
\qtypebadge{tfcol}{True / False} \hfill \textcolor{darkslate!70}{\small Topic: Cryptography and PKI}

Hash function is unrelated to produces a fixed-length digest.

\begin{optlist}
\item True
\item False
\end{optlist}
\begin{explainbox}{}
\textbf{Correct Answer: False}\\[0.3em]
\textbf{Concept:} Hash function is directly related to produces a fixed-length digest.
\end{explainbox}
\end{questioncard}

\begin{questioncard}{23}{tfcol}
\qtypebadge{tfcol}{True / False} \hfill \textcolor{darkslate!70}{\small Topic: Cryptography and PKI}

MD5 means that is considered broken due to collision attacks.

\begin{optlist}
\item True
\item False
\end{optlist}
\begin{explainbox}{}
\textbf{Correct Answer: True}\\[0.3em]
\textbf{Concept:} This statement correctly describes MD5.
\end{explainbox}
\end{questioncard}

\begin{questioncard}{24}{tfcol}
\qtypebadge{tfcol}{True / False} \hfill \textcolor{darkslate!70}{\small Topic: Cryptography and PKI}

MD5 is unrelated to is considered broken due to collision attacks.

\begin{optlist}
\item True
\item False
\end{optlist}
\begin{explainbox}{}
\textbf{Correct Answer: False}\\[0.3em]
\textbf{Concept:} MD5 is directly related to is considered broken due to collision attacks.
\end{explainbox}
\end{questioncard}

\begin{questioncard}{25}{tfcol}
\qtypebadge{tfcol}{True / False} \hfill \textcolor{darkslate!70}{\small Topic: Cryptography and PKI}

SHA-256 means that produces a 256-bit digest.

\begin{optlist}
\item True
\item False
\end{optlist}
\begin{explainbox}{}
\textbf{Correct Answer: True}\\[0.3em]
\textbf{Concept:} This statement correctly describes SHA-256.
\end{explainbox}
\end{questioncard}

\begin{questioncard}{26}{tfcol}
\qtypebadge{tfcol}{True / False} \hfill \textcolor{darkslate!70}{\small Topic: Cryptography and PKI}

SHA-256 is unrelated to produces a 256-bit digest.

\begin{optlist}
\item True
\item False
\end{optlist}
\begin{explainbox}{}
\textbf{Correct Answer: False}\\[0.3em]
\textbf{Concept:} SHA-256 is directly related to produces a 256-bit digest.
\end{explainbox}
\end{questioncard}

\begin{questioncard}{27}{tfcol}
\qtypebadge{tfcol}{True / False} \hfill \textcolor{darkslate!70}{\small Topic: Cryptography and PKI}

HMAC means that provides message authentication using a hash and key.

\begin{optlist}
\item True
\item False
\end{optlist}
\begin{explainbox}{}
\textbf{Correct Answer: True}\\[0.3em]
\textbf{Concept:} This statement correctly describes HMAC.
\end{explainbox}
\end{questioncard}

\begin{questioncard}{28}{tfcol}
\qtypebadge{tfcol}{True / False} \hfill \textcolor{darkslate!70}{\small Topic: Cryptography and PKI}

HMAC is unrelated to provides message authentication using a hash and key.

\begin{optlist}
\item True
\item False
\end{optlist}
\begin{explainbox}{}
\textbf{Correct Answer: False}\\[0.3em]
\textbf{Concept:} HMAC is directly related to provides message authentication using a hash and key.
\end{explainbox}
\end{questioncard}

\begin{questioncard}{29}{tfcol}
\qtypebadge{tfcol}{True / False} \hfill \textcolor{darkslate!70}{\small Topic: Cryptography and PKI}

Digital signature means that verifies authenticity and integrity.

\begin{optlist}
\item True
\item False
\end{optlist}
\begin{explainbox}{}
\textbf{Correct Answer: True}\\[0.3em]
\textbf{Concept:} This statement correctly describes Digital signature.
\end{explainbox}
\end{questioncard}

\begin{questioncard}{30}{tfcol}
\qtypebadge{tfcol}{True / False} \hfill \textcolor{darkslate!70}{\small Topic: Cryptography and PKI}

Digital signature is unrelated to verifies authenticity and integrity.

\begin{optlist}
\item True
\item False
\end{optlist}
\begin{explainbox}{}
\textbf{Correct Answer: False}\\[0.3em]
\textbf{Concept:} Digital signature is directly related to verifies authenticity and integrity.
\end{explainbox}
\end{questioncard}

\begin{questioncard}{31}{tfcol}
\qtypebadge{tfcol}{True / False} \hfill \textcolor{darkslate!70}{\small Topic: Cryptography and PKI}

Public key means that is used to verify a digital signature.

\begin{optlist}
\item True
\item False
\end{optlist}
\begin{explainbox}{}
\textbf{Correct Answer: True}\\[0.3em]
\textbf{Concept:} This statement correctly describes Public key.
\end{explainbox}
\end{questioncard}

\begin{questioncard}{32}{tfcol}
\qtypebadge{tfcol}{True / False} \hfill \textcolor{darkslate!70}{\small Topic: Cryptography and PKI}

Public key is unrelated to is used to verify a digital signature.

\begin{optlist}
\item True
\item False
\end{optlist}
\begin{explainbox}{}
\textbf{Correct Answer: False}\\[0.3em]
\textbf{Concept:} Public key is directly related to is used to verify a digital signature.
\end{explainbox}
\end{questioncard}

\begin{questioncard}{33}{tfcol}
\qtypebadge{tfcol}{True / False} \hfill \textcolor{darkslate!70}{\small Topic: Cryptography and PKI}

Private key means that is used to create a digital signature.

\begin{optlist}
\item True
\item False
\end{optlist}
\begin{explainbox}{}
\textbf{Correct Answer: True}\\[0.3em]
\textbf{Concept:} This statement correctly describes Private key.
\end{explainbox}
\end{questioncard}

\begin{questioncard}{34}{tfcol}
\qtypebadge{tfcol}{True / False} \hfill \textcolor{darkslate!70}{\small Topic: Cryptography and PKI}

Private key is unrelated to is used to create a digital signature.

\begin{optlist}
\item True
\item False
\end{optlist}
\begin{explainbox}{}
\textbf{Correct Answer: False}\\[0.3em]
\textbf{Concept:} Private key is directly related to is used to create a digital signature.
\end{explainbox}
\end{questioncard}

\begin{questioncard}{35}{tfcol}
\qtypebadge{tfcol}{True / False} \hfill \textcolor{darkslate!70}{\small Topic: Cryptography and PKI}

PKI means that manages public keys and certificates.

\begin{optlist}
\item True
\item False
\end{optlist}
\begin{explainbox}{}
\textbf{Correct Answer: True}\\[0.3em]
\textbf{Concept:} This statement correctly describes PKI.
\end{explainbox}
\end{questioncard}

\begin{questioncard}{36}{tfcol}
\qtypebadge{tfcol}{True / False} \hfill \textcolor{darkslate!70}{\small Topic: Cryptography and PKI}

PKI is unrelated to manages public keys and certificates.

\begin{optlist}
\item True
\item False
\end{optlist}
\begin{explainbox}{}
\textbf{Correct Answer: False}\\[0.3em]
\textbf{Concept:} PKI is directly related to manages public keys and certificates.
\end{explainbox}
\end{questioncard}

\begin{questioncard}{37}{tfcol}
\qtypebadge{tfcol}{True / False} \hfill \textcolor{darkslate!70}{\small Topic: Cryptography and PKI}

CA means that issues and signs digital certificates.

\begin{optlist}
\item True
\item False
\end{optlist}
\begin{explainbox}{}
\textbf{Correct Answer: True}\\[0.3em]
\textbf{Concept:} This statement correctly describes CA.
\end{explainbox}
\end{questioncard}

\begin{questioncard}{38}{tfcol}
\qtypebadge{tfcol}{True / False} \hfill \textcolor{darkslate!70}{\small Topic: Cryptography and PKI}

CA is unrelated to issues and signs digital certificates.

\begin{optlist}
\item True
\item False
\end{optlist}
\begin{explainbox}{}
\textbf{Correct Answer: False}\\[0.3em]
\textbf{Concept:} CA is directly related to issues and signs digital certificates.
\end{explainbox}
\end{questioncard}

\begin{questioncard}{39}{tfcol}
\qtypebadge{tfcol}{True / False} \hfill \textcolor{darkslate!70}{\small Topic: Cryptography and PKI}

RA means that validates identity before certificate issuance.

\begin{optlist}
\item True
\item False
\end{optlist}
\begin{explainbox}{}
\textbf{Correct Answer: True}\\[0.3em]
\textbf{Concept:} This statement correctly describes RA.
\end{explainbox}
\end{questioncard}

\begin{questioncard}{40}{tfcol}
\qtypebadge{tfcol}{True / False} \hfill \textcolor{darkslate!70}{\small Topic: Cryptography and PKI}

RA is unrelated to validates identity before certificate issuance.

\begin{optlist}
\item True
\item False
\end{optlist}
\begin{explainbox}{}
\textbf{Correct Answer: False}\\[0.3em]
\textbf{Concept:} RA is directly related to validates identity before certificate issuance.
\end{explainbox}
\end{questioncard}

\begin{questioncard}{41}{tfcol}
\qtypebadge{tfcol}{True / False} \hfill \textcolor{darkslate!70}{\small Topic: Cryptography and PKI}

CRL means that lists revoked certificates.

\begin{optlist}
\item True
\item False
\end{optlist}
\begin{explainbox}{}
\textbf{Correct Answer: True}\\[0.3em]
\textbf{Concept:} This statement correctly describes CRL.
\end{explainbox}
\end{questioncard}

\begin{questioncard}{42}{tfcol}
\qtypebadge{tfcol}{True / False} \hfill \textcolor{darkslate!70}{\small Topic: Cryptography and PKI}

CRL is unrelated to lists revoked certificates.

\begin{optlist}
\item True
\item False
\end{optlist}
\begin{explainbox}{}
\textbf{Correct Answer: False}\\[0.3em]
\textbf{Concept:} CRL is directly related to lists revoked certificates.
\end{explainbox}
\end{questioncard}

\begin{questioncard}{43}{tfcol}
\qtypebadge{tfcol}{True / False} \hfill \textcolor{darkslate!70}{\small Topic: Cryptography and PKI}

OCSP means that provides online certificate status checking.

\begin{optlist}
\item True
\item False
\end{optlist}
\begin{explainbox}{}
\textbf{Correct Answer: True}\\[0.3em]
\textbf{Concept:} This statement correctly describes OCSP.
\end{explainbox}
\end{questioncard}

\begin{questioncard}{44}{tfcol}
\qtypebadge{tfcol}{True / False} \hfill \textcolor{darkslate!70}{\small Topic: Cryptography and PKI}

OCSP is unrelated to provides online certificate status checking.

\begin{optlist}
\item True
\item False
\end{optlist}
\begin{explainbox}{}
\textbf{Correct Answer: False}\\[0.3em]
\textbf{Concept:} OCSP is directly related to provides online certificate status checking.
\end{explainbox}
\end{questioncard}

\begin{questioncard}{45}{tfcol}
\qtypebadge{tfcol}{True / False} \hfill \textcolor{darkslate!70}{\small Topic: Cryptography and PKI}

X.509 means that is the standard certificate format.

\begin{optlist}
\item True
\item False
\end{optlist}
\begin{explainbox}{}
\textbf{Correct Answer: True}\\[0.3em]
\textbf{Concept:} This statement correctly describes X.509.
\end{explainbox}
\end{questioncard}

\begin{questioncard}{46}{tfcol}
\qtypebadge{tfcol}{True / False} \hfill \textcolor{darkslate!70}{\small Topic: Cryptography and PKI}

X.509 is unrelated to is the standard certificate format.

\begin{optlist}
\item True
\item False
\end{optlist}
\begin{explainbox}{}
\textbf{Correct Answer: False}\\[0.3em]
\textbf{Concept:} X.509 is directly related to is the standard certificate format.
\end{explainbox}
\end{questioncard}

\begin{questioncard}{47}{tfcol}
\qtypebadge{tfcol}{True / False} \hfill \textcolor{darkslate!70}{\small Topic: Cryptography and PKI}

Root CA certificate means that is self-signed and acts as a trust anchor.

\begin{optlist}
\item True
\item False
\end{optlist}
\begin{explainbox}{}
\textbf{Correct Answer: True}\\[0.3em]
\textbf{Concept:} This statement correctly describes Root CA certificate.
\end{explainbox}
\end{questioncard}

\begin{questioncard}{48}{tfcol}
\qtypebadge{tfcol}{True / False} \hfill \textcolor{darkslate!70}{\small Topic: Cryptography and PKI}

Root CA certificate is unrelated to is self-signed and acts as a trust anchor.

\begin{optlist}
\item True
\item False
\end{optlist}
\begin{explainbox}{}
\textbf{Correct Answer: False}\\[0.3em]
\textbf{Concept:} Root CA certificate is directly related to is self-signed and acts as a trust anchor.
\end{explainbox}
\end{questioncard}

\begin{questioncard}{49}{tfcol}
\qtypebadge{tfcol}{True / False} \hfill \textcolor{darkslate!70}{\small Topic: Cryptography and PKI}

GCM means that provides authenticated encryption.

\begin{optlist}
\item True
\item False
\end{optlist}
\begin{explainbox}{}
\textbf{Correct Answer: True}\\[0.3em]
\textbf{Concept:} This statement correctly describes GCM.
\end{explainbox}
\end{questioncard}

\begin{questioncard}{50}{tfcol}
\qtypebadge{tfcol}{True / False} \hfill \textcolor{darkslate!70}{\small Topic: Cryptography and PKI}

GCM is unrelated to provides authenticated encryption.

\begin{optlist}
\item True
\item False
\end{optlist}
\begin{explainbox}{}
\textbf{Correct Answer: False}\\[0.3em]
\textbf{Concept:} GCM is directly related to provides authenticated encryption.
\end{explainbox}
\end{questioncard}

\begin{questioncard}{51}{tfcol}
\qtypebadge{tfcol}{True / False} \hfill \textcolor{darkslate!70}{\small Topic: Cryptography and PKI}

ECB means that should be avoided because it reveals plaintext patterns.

\begin{optlist}
\item True
\item False
\end{optlist}
\begin{explainbox}{}
\textbf{Correct Answer: True}\\[0.3em]
\textbf{Concept:} This statement correctly describes ECB.
\end{explainbox}
\end{questioncard}

\begin{questioncard}{52}{tfcol}
\qtypebadge{tfcol}{True / False} \hfill \textcolor{darkslate!70}{\small Topic: Cryptography and PKI}

ECB is unrelated to should be avoided because it reveals plaintext patterns.

\begin{optlist}
\item True
\item False
\end{optlist}
\begin{explainbox}{}
\textbf{Correct Answer: False}\\[0.3em]
\textbf{Concept:} ECB is directly related to should be avoided because it reveals plaintext patterns.
\end{explainbox}
\end{questioncard}

\section{Cryptography Basics}
\begin{questioncard}{1}{tfcol}
\qtypebadge{tfcol}{True / False} \hfill \textcolor{darkslate!70}{\small Topic: Cryptography Basics}

Symmetric encryption is generally faster than asymmetric encryption.

\begin{optlist}
\item True
\item False
\end{optlist}
\begin{explainbox}{}
\textbf{Correct Answer: True}\\[0.3em]
\textbf{Concept:} Symmetric algorithms are faster and suitable for bulk data encryption.
\end{explainbox}
\end{questioncard}

\begin{questioncard}{2}{tfcol}
\qtypebadge{tfcol}{True / False} \hfill \textcolor{darkslate!70}{\small Topic: Cryptography Basics}

Asymmetric encryption is ideal for encrypting large amounts of data directly.

\begin{optlist}
\item True
\item False
\end{optlist}
\begin{explainbox}{}
\textbf{Correct Answer: False}\\[0.3em]
\textbf{Concept:} Asymmetric encryption is slower; it is often used for key exchange and signatures, while symmetric encryption handles bulk data.
\end{explainbox}
\end{questioncard}

\section{Hash Algorithms}
\begin{questioncard}{1}{tfcol}
\qtypebadge{tfcol}{True / False} \hfill \textcolor{darkslate!70}{\small Topic: Hash Algorithms}

A hash function can be reversed to recover the original input.

\begin{optlist}
\item True
\item False
\end{optlist}
\begin{explainbox}{}
\textbf{Correct Answer: False}\\[0.3em]
\textbf{Concept:} Hash functions are one-way; the original input cannot be derived from the digest.
\end{explainbox}
\end{questioncard}

\section{PKI}
\begin{questioncard}{1}{tfcol}
\qtypebadge{tfcol}{True / False} \hfill \textcolor{darkslate!70}{\small Topic: PKI}

A digital certificate binds a public key to an identity.

\begin{optlist}
\item True
\item False
\end{optlist}
\begin{explainbox}{}
\textbf{Correct Answer: True}\\[0.3em]
\textbf{Concept:} Certificates contain a public key and identity information, signed by a trusted CA.
\end{explainbox}
\end{questioncard}

\begin{questioncard}{2}{tfcol}
\qtypebadge{tfcol}{True / False} \hfill \textcolor{darkslate!70}{\small Topic: PKI}

A CRL lists certificates that are still valid.

\begin{optlist}
\item True
\item False
\end{optlist}
\begin{explainbox}{}
\textbf{Correct Answer: False}\\[0.3em]
\textbf{Concept:} A CRL lists certificates that have been revoked and should no longer be trusted.
\end{explainbox}
\end{questioncard}

\section{VPN}
\begin{questioncard}{1}{tfcol}
\qtypebadge{tfcol}{True / False} \hfill \textcolor{darkslate!70}{\small Topic: VPN}

A VPN allows secure communication over an untrusted public network.

\begin{optlist}
\item True
\item False
\end{optlist}
\begin{explainbox}{}
\textbf{Correct Answer: True}\\[0.3em]
\textbf{Concept:} VPNs create encrypted tunnels over public networks.
\end{explainbox}
\end{questioncard}

\begin{questioncard}{2}{tfcol}
\qtypebadge{tfcol}{True / False} \hfill \textcolor{darkslate!70}{\small Topic: VPN}

GRE can tunnel multicast traffic.

\begin{optlist}
\item True
\item False
\end{optlist}
\begin{explainbox}{}
\textbf{Correct Answer: True}\\[0.3em]
\textbf{Concept:} GRE supports multicast tunneling, which is useful for routing protocols.
\end{explainbox}
\end{questioncard}

\begin{questioncard}{3}{tfcol}
\qtypebadge{tfcol}{True / False} \hfill \textcolor{darkslate!70}{\small Topic: VPN}

IPsec transport mode encrypts the entire original IP packet.

\begin{optlist}
\item True
\item False
\end{optlist}
\begin{explainbox}{}
\textbf{Correct Answer: False}\\[0.3em]
\textbf{Concept:} Transport mode encrypts only the payload; tunnel mode encrypts the entire packet.
\end{explainbox}
\end{questioncard}

\begin{questioncard}{4}{tfcol}
\qtypebadge{tfcol}{True / False} \hfill \textcolor{darkslate!70}{\small Topic: VPN}

IKE is responsible for negotiating IPsec keys and security associations.

\begin{optlist}
\item True
\item False
\end{optlist}
\begin{explainbox}{}
\textbf{Correct Answer: True}\\[0.3em]
\textbf{Concept:} IKE automates key exchange and SA negotiation.
\end{explainbox}
\end{questioncard}

\begin{questioncard}{5}{tfcol}
\qtypebadge{tfcol}{True / False} \hfill \textcolor{darkslate!70}{\small Topic: VPN}

L2TP provides encryption by itself.

\begin{optlist}
\item True
\item False
\end{optlist}
\begin{explainbox}{}
\textbf{Correct Answer: False}\\[0.3em]
\textbf{Concept:} L2TP does not encrypt traffic; it is usually combined with IPsec.
\end{explainbox}
\end{questioncard}

\begin{questioncard}{6}{tfcol}
\qtypebadge{tfcol}{True / False} \hfill \textcolor{darkslate!70}{\small Topic: VPN}

SSL VPN can be accessed through a standard web browser.

\begin{optlist}
\item True
\item False
\end{optlist}
\begin{explainbox}{}
\textbf{Correct Answer: True}\\[0.3em]
\textbf{Concept:} SSL VPN supports browser-based access via HTTPS.
\end{explainbox}
\end{questioncard}

\begin{questioncard}{7}{tfcol}
\qtypebadge{tfcol}{True / False} \hfill \textcolor{darkslate!70}{\small Topic: VPN}

GRE provides encryption by default.

\begin{optlist}
\item True
\item False
\end{optlist}
\begin{explainbox}{}
\textbf{Correct Answer: False}\\[0.3em]
\textbf{Concept:} GRE encapsulates but does not encrypt traffic.
\end{explainbox}
\end{questioncard}

\begin{questioncard}{8}{tfcol}
\qtypebadge{tfcol}{True / False} \hfill \textcolor{darkslate!70}{\small Topic: VPN}

SSL VPN commonly uses TCP port 443.

\begin{optlist}
\item True
\item False
\end{optlist}
\begin{explainbox}{}
\textbf{Correct Answer: True}\\[0.3em]
\textbf{Concept:} SSL VPN uses HTTPS/TCP 443.
\end{explainbox}
\end{questioncard}

\begin{questioncard}{9}{tfcol}
\qtypebadge{tfcol}{True / False} \hfill \textcolor{darkslate!70}{\small Topic: VPN}

IPsec tunnel mode encrypts the entire original IP packet.

\begin{optlist}
\item True
\item False
\end{optlist}
\begin{explainbox}{}
\textbf{Correct Answer: True}\\[0.3em]
\textbf{Concept:} Tunnel mode encapsulates the whole original packet.
\end{explainbox}
\end{questioncard}

\begin{questioncard}{10}{tfcol}
\qtypebadge{tfcol}{True / False} \hfill \textcolor{darkslate!70}{\small Topic: VPN}

VPN means that provides secure communication over public networks.

\begin{optlist}
\item True
\item False
\end{optlist}
\begin{explainbox}{}
\textbf{Correct Answer: True}\\[0.3em]
\textbf{Concept:} This statement correctly describes VPN.
\end{explainbox}
\end{questioncard}

\begin{questioncard}{11}{tfcol}
\qtypebadge{tfcol}{True / False} \hfill \textcolor{darkslate!70}{\small Topic: VPN}

VPN is unrelated to provides secure communication over public networks.

\begin{optlist}
\item True
\item False
\end{optlist}
\begin{explainbox}{}
\textbf{Correct Answer: False}\\[0.3em]
\textbf{Concept:} VPN is directly related to provides secure communication over public networks.
\end{explainbox}
\end{questioncard}

\begin{questioncard}{12}{tfcol}
\qtypebadge{tfcol}{True / False} \hfill \textcolor{darkslate!70}{\small Topic: VPN}

Site-to-site VPN means that connects entire networks at different locations.

\begin{optlist}
\item True
\item False
\end{optlist}
\begin{explainbox}{}
\textbf{Correct Answer: True}\\[0.3em]
\textbf{Concept:} This statement correctly describes Site-to-site VPN.
\end{explainbox}
\end{questioncard}

\begin{questioncard}{13}{tfcol}
\qtypebadge{tfcol}{True / False} \hfill \textcolor{darkslate!70}{\small Topic: VPN}

Site-to-site VPN is unrelated to connects entire networks at different locations.

\begin{optlist}
\item True
\item False
\end{optlist}
\begin{explainbox}{}
\textbf{Correct Answer: False}\\[0.3em]
\textbf{Concept:} Site-to-site VPN is directly related to connects entire networks at different locations.
\end{explainbox}
\end{questioncard}

\begin{questioncard}{14}{tfcol}
\qtypebadge{tfcol}{True / False} \hfill \textcolor{darkslate!70}{\small Topic: VPN}

Remote-access VPN means that connects individual users to a corporate network.

\begin{optlist}
\item True
\item False
\end{optlist}
\begin{explainbox}{}
\textbf{Correct Answer: True}\\[0.3em]
\textbf{Concept:} This statement correctly describes Remote-access VPN.
\end{explainbox}
\end{questioncard}

\begin{questioncard}{15}{tfcol}
\qtypebadge{tfcol}{True / False} \hfill \textcolor{darkslate!70}{\small Topic: VPN}

Remote-access VPN is unrelated to connects individual users to a corporate network.

\begin{optlist}
\item True
\item False
\end{optlist}
\begin{explainbox}{}
\textbf{Correct Answer: False}\\[0.3em]
\textbf{Concept:} Remote-access VPN is directly related to connects individual users to a corporate network.
\end{explainbox}
\end{questioncard}

\begin{questioncard}{16}{tfcol}
\qtypebadge{tfcol}{True / False} \hfill \textcolor{darkslate!70}{\small Topic: VPN}

GRE means that is IP protocol 47 and lacks built-in encryption.

\begin{optlist}
\item True
\item False
\end{optlist}
\begin{explainbox}{}
\textbf{Correct Answer: True}\\[0.3em]
\textbf{Concept:} This statement correctly describes GRE.
\end{explainbox}
\end{questioncard}

\begin{questioncard}{17}{tfcol}
\qtypebadge{tfcol}{True / False} \hfill \textcolor{darkslate!70}{\small Topic: VPN}

GRE is unrelated to is IP protocol 47 and lacks built-in encryption.

\begin{optlist}
\item True
\item False
\end{optlist}
\begin{explainbox}{}
\textbf{Correct Answer: False}\\[0.3em]
\textbf{Concept:} GRE is directly related to is IP protocol 47 and lacks built-in encryption.
\end{explainbox}
\end{questioncard}

\begin{questioncard}{18}{tfcol}
\qtypebadge{tfcol}{True / False} \hfill \textcolor{darkslate!70}{\small Topic: VPN}

GRE means that can tunnel multicast and non-IP protocols.

\begin{optlist}
\item True
\item False
\end{optlist}
\begin{explainbox}{}
\textbf{Correct Answer: True}\\[0.3em]
\textbf{Concept:} This statement correctly describes GRE.
\end{explainbox}
\end{questioncard}

\begin{questioncard}{19}{tfcol}
\qtypebadge{tfcol}{True / False} \hfill \textcolor{darkslate!70}{\small Topic: VPN}

GRE is unrelated to can tunnel multicast and non-IP protocols.

\begin{optlist}
\item True
\item False
\end{optlist}
\begin{explainbox}{}
\textbf{Correct Answer: False}\\[0.3em]
\textbf{Concept:} GRE is directly related to can tunnel multicast and non-IP protocols.
\end{explainbox}
\end{questioncard}

\begin{questioncard}{20}{tfcol}
\qtypebadge{tfcol}{True / False} \hfill \textcolor{darkslate!70}{\small Topic: VPN}

IPsec means that provides confidentiality, integrity, and authentication.

\begin{optlist}
\item True
\item False
\end{optlist}
\begin{explainbox}{}
\textbf{Correct Answer: True}\\[0.3em]
\textbf{Concept:} This statement correctly describes IPsec.
\end{explainbox}
\end{questioncard}

\begin{questioncard}{21}{tfcol}
\qtypebadge{tfcol}{True / False} \hfill \textcolor{darkslate!70}{\small Topic: VPN}

IPsec is unrelated to provides confidentiality, integrity, and authentication.

\begin{optlist}
\item True
\item False
\end{optlist}
\begin{explainbox}{}
\textbf{Correct Answer: False}\\[0.3em]
\textbf{Concept:} IPsec is directly related to provides confidentiality, integrity, and authentication.
\end{explainbox}
\end{questioncard}

\begin{questioncard}{22}{tfcol}
\qtypebadge{tfcol}{True / False} \hfill \textcolor{darkslate!70}{\small Topic: VPN}

ESP means that provides encryption and authentication.

\begin{optlist}
\item True
\item False
\end{optlist}
\begin{explainbox}{}
\textbf{Correct Answer: True}\\[0.3em]
\textbf{Concept:} This statement correctly describes ESP.
\end{explainbox}
\end{questioncard}

\begin{questioncard}{23}{tfcol}
\qtypebadge{tfcol}{True / False} \hfill \textcolor{darkslate!70}{\small Topic: VPN}

ESP is unrelated to provides encryption and authentication.

\begin{optlist}
\item True
\item False
\end{optlist}
\begin{explainbox}{}
\textbf{Correct Answer: False}\\[0.3em]
\textbf{Concept:} ESP is directly related to provides encryption and authentication.
\end{explainbox}
\end{questioncard}

\begin{questioncard}{24}{tfcol}
\qtypebadge{tfcol}{True / False} \hfill \textcolor{darkslate!70}{\small Topic: VPN}

AH means that provides authentication without encryption.

\begin{optlist}
\item True
\item False
\end{optlist}
\begin{explainbox}{}
\textbf{Correct Answer: True}\\[0.3em]
\textbf{Concept:} This statement correctly describes AH.
\end{explainbox}
\end{questioncard}

\begin{questioncard}{25}{tfcol}
\qtypebadge{tfcol}{True / False} \hfill \textcolor{darkslate!70}{\small Topic: VPN}

AH is unrelated to provides authentication without encryption.

\begin{optlist}
\item True
\item False
\end{optlist}
\begin{explainbox}{}
\textbf{Correct Answer: False}\\[0.3em]
\textbf{Concept:} AH is directly related to provides authentication without encryption.
\end{explainbox}
\end{questioncard}

\begin{questioncard}{26}{tfcol}
\qtypebadge{tfcol}{True / False} \hfill \textcolor{darkslate!70}{\small Topic: VPN}

IKE means that negotiates IPsec keys and security associations.

\begin{optlist}
\item True
\item False
\end{optlist}
\begin{explainbox}{}
\textbf{Correct Answer: True}\\[0.3em]
\textbf{Concept:} This statement correctly describes IKE.
\end{explainbox}
\end{questioncard}

\begin{questioncard}{27}{tfcol}
\qtypebadge{tfcol}{True / False} \hfill \textcolor{darkslate!70}{\small Topic: VPN}

IKE is unrelated to negotiates IPsec keys and security associations.

\begin{optlist}
\item True
\item False
\end{optlist}
\begin{explainbox}{}
\textbf{Correct Answer: False}\\[0.3em]
\textbf{Concept:} IKE is directly related to negotiates IPsec keys and security associations.
\end{explainbox}
\end{questioncard}

\begin{questioncard}{28}{tfcol}
\qtypebadge{tfcol}{True / False} \hfill \textcolor{darkslate!70}{\small Topic: VPN}

Tunnel mode means that encrypts the entire original IP packet.

\begin{optlist}
\item True
\item False
\end{optlist}
\begin{explainbox}{}
\textbf{Correct Answer: True}\\[0.3em]
\textbf{Concept:} This statement correctly describes Tunnel mode.
\end{explainbox}
\end{questioncard}

\begin{questioncard}{29}{tfcol}
\qtypebadge{tfcol}{True / False} \hfill \textcolor{darkslate!70}{\small Topic: VPN}

Tunnel mode is unrelated to encrypts the entire original IP packet.

\begin{optlist}
\item True
\item False
\end{optlist}
\begin{explainbox}{}
\textbf{Correct Answer: False}\\[0.3em]
\textbf{Concept:} Tunnel mode is directly related to encrypts the entire original IP packet.
\end{explainbox}
\end{questioncard}

\begin{questioncard}{30}{tfcol}
\qtypebadge{tfcol}{True / False} \hfill \textcolor{darkslate!70}{\small Topic: VPN}

Transport mode means that encrypts only the IP payload.

\begin{optlist}
\item True
\item False
\end{optlist}
\begin{explainbox}{}
\textbf{Correct Answer: True}\\[0.3em]
\textbf{Concept:} This statement correctly describes Transport mode.
\end{explainbox}
\end{questioncard}

\begin{questioncard}{31}{tfcol}
\qtypebadge{tfcol}{True / False} \hfill \textcolor{darkslate!70}{\small Topic: VPN}

Transport mode is unrelated to encrypts only the IP payload.

\begin{optlist}
\item True
\item False
\end{optlist}
\begin{explainbox}{}
\textbf{Correct Answer: False}\\[0.3em]
\textbf{Concept:} Transport mode is directly related to encrypts only the IP payload.
\end{explainbox}
\end{questioncard}

\begin{questioncard}{32}{tfcol}
\qtypebadge{tfcol}{True / False} \hfill \textcolor{darkslate!70}{\small Topic: VPN}

L2TP means that uses UDP port 1701 and lacks encryption.

\begin{optlist}
\item True
\item False
\end{optlist}
\begin{explainbox}{}
\textbf{Correct Answer: True}\\[0.3em]
\textbf{Concept:} This statement correctly describes L2TP.
\end{explainbox}
\end{questioncard}

\begin{questioncard}{33}{tfcol}
\qtypebadge{tfcol}{True / False} \hfill \textcolor{darkslate!70}{\small Topic: VPN}

L2TP is unrelated to uses UDP port 1701 and lacks encryption.

\begin{optlist}
\item True
\item False
\end{optlist}
\begin{explainbox}{}
\textbf{Correct Answer: False}\\[0.3em]
\textbf{Concept:} L2TP is directly related to uses UDP port 1701 and lacks encryption.
\end{explainbox}
\end{questioncard}

\begin{questioncard}{34}{tfcol}
\qtypebadge{tfcol}{True / False} \hfill \textcolor{darkslate!70}{\small Topic: VPN}

L2TP over IPsec means that combines L2TP tunneling with IPsec encryption.

\begin{optlist}
\item True
\item False
\end{optlist}
\begin{explainbox}{}
\textbf{Correct Answer: True}\\[0.3em]
\textbf{Concept:} This statement correctly describes L2TP over IPsec.
\end{explainbox}
\end{questioncard}

\begin{questioncard}{35}{tfcol}
\qtypebadge{tfcol}{True / False} \hfill \textcolor{darkslate!70}{\small Topic: VPN}

L2TP over IPsec is unrelated to combines L2TP tunneling with IPsec encryption.

\begin{optlist}
\item True
\item False
\end{optlist}
\begin{explainbox}{}
\textbf{Correct Answer: False}\\[0.3em]
\textbf{Concept:} L2TP over IPsec is directly related to combines L2TP tunneling with IPsec encryption.
\end{explainbox}
\end{questioncard}

\begin{questioncard}{36}{tfcol}
\qtypebadge{tfcol}{True / False} \hfill \textcolor{darkslate!70}{\small Topic: VPN}

SSL VPN means that commonly uses TCP port 443.

\begin{optlist}
\item True
\item False
\end{optlist}
\begin{explainbox}{}
\textbf{Correct Answer: True}\\[0.3em]
\textbf{Concept:} This statement correctly describes SSL VPN.
\end{explainbox}
\end{questioncard}

\begin{questioncard}{37}{tfcol}
\qtypebadge{tfcol}{True / False} \hfill \textcolor{darkslate!70}{\small Topic: VPN}

SSL VPN is unrelated to commonly uses TCP port 443.

\begin{optlist}
\item True
\item False
\end{optlist}
\begin{explainbox}{}
\textbf{Correct Answer: False}\\[0.3em]
\textbf{Concept:} SSL VPN is directly related to commonly uses TCP port 443.
\end{explainbox}
\end{questioncard}

\begin{questioncard}{38}{tfcol}
\qtypebadge{tfcol}{True / False} \hfill \textcolor{darkslate!70}{\small Topic: VPN}

Web proxy mode means that provides browser-based access to web apps.

\begin{optlist}
\item True
\item False
\end{optlist}
\begin{explainbox}{}
\textbf{Correct Answer: True}\\[0.3em]
\textbf{Concept:} This statement correctly describes Web proxy mode.
\end{explainbox}
\end{questioncard}

\begin{questioncard}{39}{tfcol}
\qtypebadge{tfcol}{True / False} \hfill \textcolor{darkslate!70}{\small Topic: VPN}

Web proxy mode is unrelated to provides browser-based access to web apps.

\begin{optlist}
\item True
\item False
\end{optlist}
\begin{explainbox}{}
\textbf{Correct Answer: False}\\[0.3em]
\textbf{Concept:} Web proxy mode is directly related to provides browser-based access to web apps.
\end{explainbox}
\end{questioncard}

\begin{questioncard}{40}{tfcol}
\qtypebadge{tfcol}{True / False} \hfill \textcolor{darkslate!70}{\small Topic: VPN}

Network extension mode means that assigns a virtual IP to the remote host.

\begin{optlist}
\item True
\item False
\end{optlist}
\begin{explainbox}{}
\textbf{Correct Answer: True}\\[0.3em]
\textbf{Concept:} This statement correctly describes Network extension mode.
\end{explainbox}
\end{questioncard}

\begin{questioncard}{41}{tfcol}
\qtypebadge{tfcol}{True / False} \hfill \textcolor{darkslate!70}{\small Topic: VPN}

Network extension mode is unrelated to assigns a virtual IP to the remote host.

\begin{optlist}
\item True
\item False
\end{optlist}
\begin{explainbox}{}
\textbf{Correct Answer: False}\\[0.3em]
\textbf{Concept:} Network extension mode is directly related to assigns a virtual IP to the remote host.
\end{explainbox}
\end{questioncard}

\begin{questioncard}{42}{tfcol}
\qtypebadge{tfcol}{True / False} \hfill \textcolor{darkslate!70}{\small Topic: VPN}

IKEv2 means that is faster and more reliable than IKEv1.

\begin{optlist}
\item True
\item False
\end{optlist}
\begin{explainbox}{}
\textbf{Correct Answer: True}\\[0.3em]
\textbf{Concept:} This statement correctly describes IKEv2.
\end{explainbox}
\end{questioncard}

\begin{questioncard}{43}{tfcol}
\qtypebadge{tfcol}{True / False} \hfill \textcolor{darkslate!70}{\small Topic: VPN}

IKEv2 is unrelated to is faster and more reliable than IKEv1.

\begin{optlist}
\item True
\item False
\end{optlist}
\begin{explainbox}{}
\textbf{Correct Answer: False}\\[0.3em]
\textbf{Concept:} IKEv2 is directly related to is faster and more reliable than IKEv1.
\end{explainbox}
\end{questioncard}

\section{GRE VPN}
\begin{questioncard}{1}{tfcol}
\qtypebadge{tfcol}{True / False} \hfill \textcolor{darkslate!70}{\small Topic: GRE VPN}

GRE provides strong encryption by default.

\begin{optlist}
\item True
\item False
\end{optlist}
\begin{explainbox}{}
\textbf{Correct Answer: False}\\[0.3em]
\textbf{Concept:} GRE only encapsulates traffic; encryption must be added separately, often with IPsec.
\end{explainbox}
\end{questioncard}

\section{IPsec VPN}
\begin{questioncard}{1}{tfcol}
\qtypebadge{tfcol}{True / False} \hfill \textcolor{darkslate!70}{\small Topic: IPsec VPN}

ESP provides both encryption and authentication.

\begin{optlist}
\item True
\item False
\end{optlist}
\begin{explainbox}{}
\textbf{Correct Answer: True}\\[0.3em]
\textbf{Concept:} ESP can encrypt payloads and authenticate packets.
\end{explainbox}
\end{questioncard}

\begin{questioncard}{2}{tfcol}
\qtypebadge{tfcol}{True / False} \hfill \textcolor{darkslate!70}{\small Topic: IPsec VPN}

AH provides payload encryption.

\begin{optlist}
\item True
\item False
\end{optlist}
\begin{explainbox}{}
\textbf{Correct Answer: False}\\[0.3em]
\textbf{Concept:} AH provides authentication and integrity but does not encrypt the payload.
\end{explainbox}
\end{questioncard}

\section{L2TP VPN}
\begin{questioncard}{1}{tfcol}
\qtypebadge{tfcol}{True / False} \hfill \textcolor{darkslate!70}{\small Topic: L2TP VPN}

L2TP alone does not provide encryption.

\begin{optlist}
\item True
\item False
\end{optlist}
\begin{explainbox}{}
\textbf{Correct Answer: True}\\[0.3em]
\textbf{Concept:} L2TP provides tunneling but relies on IPsec or other mechanisms for encryption.
\end{explainbox}
\end{questioncard}

\section{SSL VPN}
\begin{questioncard}{1}{tfcol}
\qtypebadge{tfcol}{True / False} \hfill \textcolor{darkslate!70}{\small Topic: SSL VPN}

SSL VPN typically uses TCP port 443.

\begin{optlist}
\item True
\item False
\end{optlist}
\begin{explainbox}{}
\textbf{Correct Answer: True}\\[0.3em]
\textbf{Concept:} SSL/TLS VPNs commonly use HTTPS/TCP 443, which is usually allowed through firewalls.
\end{explainbox}
\end{questioncard}

\end{document}
